\chapter{Computational tools and analyses}
\label{ch:computational}
% [draft: All three of these papers used the Solexa 1G, released in 2006 and rebranded in 2007 to the Illumina Genome Analyzer I after its acquisition, to sequence the DNA libraries generated. allowing focus to shift towards determining, and eventually modeling, how the genome functioned and self-organized.]
% % !!! SUPER IMPORTANT PAPER !!! https://academic.oup.com/nar/article-abstract/5/9/3157/2380868?redirectedFrom=fulltext&login=true

% iRegulon\sjcitep[]{janky2014iregulon} provides a computational alternative for identifying master regulators when experimental TF binding data is unavailable. First, a genome-wide assessment of motif enrichment near TSSs is generated by scoring occurrence of \num{10000} motif PWMs aggregated from various databases using ClusterBuster\sjcitep[]{frith2003cluster}, generating for each PWM a ranking of the genes it may regulate. A set of genes, derived from e.g. differential gene expression or co-expression data, can then be checked for overrepresentation in these PWM-gene rankings using an area under the rank curve approach. This returns a set of motifs that may regulate the gene set, which can then be mapped to TFs using orthology or similarity if annotation is absent. Applying this approach to validated gene sets from 115 ENCODE ChIP-seq datasets returned 94 correct TF hits, and was able to identify 307 direct p53 target genes using ENCODE ChIP-seq data.

\section{Single-cell data pre-processing}
% The newly available human genome, in conjunction with large amounts of data generated by microarray and sequencing-based quantitative assays, necessitated the development of new bioinformatics tools.
Once data is sequenced, and raw \verb|.fastq| data is available, the first step in data analysis is alignment to a reference genome. A number of algorithms exist for this task, and the optimal choice depends on several factors. Short-read aligners such as \verb|bowtie2| are optimized for reads < \SI{200}{\bp}, while long-read aligners like \verb|minimap2| are better at handling reads longer than \SI{1000}{\bp}. The type of sequencing technology used also plays a role: sequencing by synthesis (Illumina) generates more indels due to incorporation errors, and these need to be explicitly corrected for by \verb|bwa-mem|, whereas \verb|minimap2| is better suited for handling the higher single-base error rates found in Nanopore/PacBio reads.

Whereas genomic reads, such as generated by ChIP-seq, can be directly aligned to the genome, mRNA transcripts overlap discontinuously with the genome due to splicing, and are thus more challenging to align. Aligning short-read RNA-seq data to the genome requires splice site annotations, which required expressed sequence tags (ESTs) or full-length cDNA sequencing to generate. TopHat ended the need for ESTs by introducing computational splice site identification\sjcitep[.]{trapnell2009tophat} The algorithm identifies exons by looking for splice site motifs near the flanks of the reads that only partially map ("initially unmapped") to the reference genome with \verb|bowtie|. Multiple reads are then aggregated to statistically validate a splice site. The resulting read alignments can then be passed to Cufflinks\sjcitep[,]{trapnellTranscriptAssemblyQuantification2010} which assembles them into possible transcript isoforms and can provide a statistical estimate of their relative quantities without long-read cDNA sequencing data.

% Newer \verb|HISAT2|\sjcitep[]{pertea2016transcript} are used to handle spliced reads.

Like many classes of algorithms, a trade-off between memory and compute can also be made, for example \verb|STAR|\sjcitep[]{dobin2013star} requires large amounts of memory due to its use of an uncompressed genome index, but is faster, while \verb|HISAT2| is more memory-efficient, but takes longer.

% STAR's high mapping speed and accuracy were crucial for analyzing the large ENCODE transcriptome (Djebali et al., 2012) dataset (>80 billion Illumina reads)

The diversity in sequencing technologies and alignment algorithms led to the use of several independent alignment formats, each tied to specific sequencing technologies or alignment algorithms. Analyses on different datasets (e.g. variant calling, genotype analysis) required cross parsing or conversion. The introduction of the sequence alignment map (\verb|SAM|) format as part of the 1000 Genomes Project in 2009\sjcitep[,]{liSequenceAlignmentMapFormat2009} provided a unified standard for storing aligned sequencing data which has since become the de facto standard. Each read is stored in a single line containing 11 mandatory fields including: read name, reference name (chromosome or contig), position, mapping quality, the actual sequence, several other alignment quality parameters as well as optional tags.

\begin{lstlisting}
@HD     VN:1.6  SO:coordinate
@SQ     SN:chr1 LN:248956422
@SQ     SN:chr2 LN:242193529

read001  99  chr1  1000  60  50M  =  1150  200  ACGTACGTACGTACGTACGTACGTACGTACGTACGTACGTACGTACG  IIIIIIIIIIIIIIIIIIIIIIIIIIIIIIIIIIIIIIIIIIIIII  NM:i:0  MD:Z:50
read001  147 chr1  1150 60  50M  =  1000 -200 ACGTACGTACGTACGTACGTACGTACGTACGTACGTACGTACGTACG  IIIIIIIIIIIIIIIIIIIIIIIIIIIIIIIIIIIIIIIIIIIIII  NM:i:0  MD:Z:50
\end{lstlisting}

An indexed binary representation (\verb|BAM|) saves storage space and speed up searching and processing with random access. This also allows users to quickly inspect desired genomic regions using visualisation tools like \verb|igv|\sjcitep[.]{thorvaldsdottirIntegrativeGenomicsViewer2013}. Alignment quality itself is also a useful metric for assessing an experiment, where low genomic or transcriptomic alignment rates are indicative of low quality. When this occurs independently of per-base sequencing quality (i.e. sequencing run issues), either sample quality (degradation, viability) or library preparation (low input DNA, low reverse transcription efficiency, over-amplification or over-fragmentation) are usually at fault. Additionally, metrics such as reads mapped to the mitochondrial genome can indicate sample damage. These cases require optimization searches of better extraction and library prep parameters

After alignment, the next step is to vectorize the data by counting aligned reads in genomic regions. This generates a cell by feature count matrix, unlocking linear algebra (\autoref{tab:feature_count_matrix}). The region sets used for counting differ by assay used. scRNA-seq requires gene annotations, whereas for ATAC-seq several strategies exist: ideally, previously validated sample-specific peak regions are used to limit data sparsity. These peak sets can be called from previous local data with peak callers like MACS\sjcitep[.]{zhang2008macs} MACS uses a modeling-based approach: reads are shifted by half the learned fragment length to center to true binding sites (ChIP-seq bias). A sliding window scans the genome and uses a Poisson background (scaled to control when available) to call regions whose enrichment is statistically significant. It also removes PCR duplicates beyond what random sampling predicts. MACS outperformed contemporary tools (ChIPSeq Peak Finder, FindPeaks, QuEST) by calling more true peaks, showing higher motif enrichment and better summit-to-motif resolution while keeping lower false positives.  MACS2 is a faster, multithreaded version with better QC metrics and broader assay support (ChIP-seq, ATAC-seq, DNase-seq) compatible with paired-end reads\sjcitep[.]{gaspar2020macs2} When sample-specific peak sets are not available, public data or genomic bins can be used. Particularly, the ENCODE SCREEN regions is a set of \num{926535} human and \num{339815} mouse candidate cis-regulatory elements (cCREs) identified based on ATAC-seq, DNase hypersensitivity, CHiP-seq, histone modifications (H3K4me3, H3K27ac), and CTCF binding aggregated from over \num{9000} efforts\sjcitep[.]{encode2020expanded} These regions represent potential enhancers, promoters, and insulators. Finally, the selected region set should be filtered for problematic regions to remove false peaks\sjcitep[.]{amemiya2019encode}

\begin{table}[ht]
	\centering
	\begin{tabular}{lcccc}
		\hline
		                    & \textbf{barcode\_1} & \textbf{barcode\_2} & $\cdots$ & \textbf{barcode\_N} \\ \hline
		\textbf{feature\_1} & 3                   & 2                   & $\cdots$ & 13                  \\
		\textbf{feature\_2} & 2                   & 3                   & $\cdots$ & 1                   \\
		\textbf{feature\_3} & 1                   & 14                  & $\cdots$ & 18                  \\
		$\vdots$            & $\vdots$            & $\vdots$            & $\ddots$ & $\vdots$            \\
		\textbf{feature\_M} & 25                  & 0                   & $\cdots$ & 0                   \\ \hline
	\end{tabular}
	\caption{Example feature count matrix across different cells.}
	\label{tab:feature_count_matrix}
\end{table}

Raw count matrices contain all counts in all barcodes sequenced in the dataset. Filtering true cells from background noise barcodes is a dataset-defining step for any single-cell dataset. The barcode-rank plot (also referred to as "knee" or "elbow plot") visualises the distribution of reads across all sequenced barcodes in the sequencing dataset (\autoref{fig:kneeplot}). It is obtained by counting the number of reads per barcode, and ranking each barcode in a log-log figure - similarly in concept to a Zipf or Lorenz curve.

\begin{figure}[h]
	\centering
	\resizebox{0.8\textwidth}{!}{%% Creator: Matplotlib, PGF backend
%%
%% To include the figure in your LaTeX document, write
%%   \input{<filename>.pgf}
%%
%% Make sure the required packages are loaded in your preamble
%%   \usepackage{pgf}
%%
%% Also ensure that all the required font packages are loaded; for instance,
%% the lmodern package is sometimes necessary when using math font.
%%   \usepackage{lmodern}
%%
%% Figures using additional raster images can only be included by \input if
%% they are in the same directory as the main LaTeX file. For loading figures
%% from other directories you can use the `import` package
%%   \usepackage{import}
%%
%% and then include the figures with
%%   \import{<path to file>}{<filename>.pgf}
%%
%% Matplotlib used the following preamble
%%   \def\mathdefault#1{#1}
%%   \everymath=\expandafter{\the\everymath\displaystyle}
%%   \IfFileExists{scrextend.sty}{
%%     \usepackage[fontsize=10.000000pt]{scrextend}
%%   }{
%%     \renewcommand{\normalsize}{\fontsize{10.000000}{12.000000}\selectfont}
%%     \normalsize
%%   }
%%   
%%   \ifdefined\pdftexversion\else  % non-pdftex case.
%%     \usepackage{fontspec}
%%     \setmainfont{DejaVuSans.ttf}[Path=\detokenize{/Users/florianderop/.mydefaultenv/lib/python3.13/site-packages/matplotlib/mpl-data/fonts/ttf/}]
%%     \setsansfont{DejaVuSans.ttf}[Path=\detokenize{/Users/florianderop/.mydefaultenv/lib/python3.13/site-packages/matplotlib/mpl-data/fonts/ttf/}]
%%     \setmonofont{DejaVuSansMono.ttf}[Path=\detokenize{/Users/florianderop/.mydefaultenv/lib/python3.13/site-packages/matplotlib/mpl-data/fonts/ttf/}]
%%   \fi
%%   \makeatletter\@ifpackageloaded{underscore}{}{\usepackage[strings]{underscore}}\makeatother
%%
\begingroup%
\makeatletter%
\begin{pgfpicture}%
\pgfpathrectangle{\pgfpointorigin}{\pgfqpoint{8.000000in}{4.000000in}}%
\pgfusepath{use as bounding box, clip}%
\begin{pgfscope}%
\pgfsetbuttcap%
\pgfsetmiterjoin%
\definecolor{currentfill}{rgb}{1.000000,1.000000,1.000000}%
\pgfsetfillcolor{currentfill}%
\pgfsetlinewidth{0.000000pt}%
\definecolor{currentstroke}{rgb}{1.000000,1.000000,1.000000}%
\pgfsetstrokecolor{currentstroke}%
\pgfsetdash{}{0pt}%
\pgfpathmoveto{\pgfqpoint{0.000000in}{0.000000in}}%
\pgfpathlineto{\pgfqpoint{8.000000in}{0.000000in}}%
\pgfpathlineto{\pgfqpoint{8.000000in}{4.000000in}}%
\pgfpathlineto{\pgfqpoint{0.000000in}{4.000000in}}%
\pgfpathlineto{\pgfqpoint{0.000000in}{0.000000in}}%
\pgfpathclose%
\pgfusepath{fill}%
\end{pgfscope}%
\begin{pgfscope}%
\pgfsetbuttcap%
\pgfsetmiterjoin%
\definecolor{currentfill}{rgb}{1.000000,1.000000,1.000000}%
\pgfsetfillcolor{currentfill}%
\pgfsetlinewidth{0.000000pt}%
\definecolor{currentstroke}{rgb}{0.000000,0.000000,0.000000}%
\pgfsetstrokecolor{currentstroke}%
\pgfsetstrokeopacity{0.000000}%
\pgfsetdash{}{0pt}%
\pgfpathmoveto{\pgfqpoint{1.000000in}{0.440000in}}%
\pgfpathlineto{\pgfqpoint{7.200000in}{0.440000in}}%
\pgfpathlineto{\pgfqpoint{7.200000in}{3.520000in}}%
\pgfpathlineto{\pgfqpoint{1.000000in}{3.520000in}}%
\pgfpathlineto{\pgfqpoint{1.000000in}{0.440000in}}%
\pgfpathclose%
\pgfusepath{fill}%
\end{pgfscope}%
\begin{pgfscope}%
\pgfpathrectangle{\pgfqpoint{1.000000in}{0.440000in}}{\pgfqpoint{6.200000in}{3.080000in}}%
\pgfusepath{clip}%
\pgfsetbuttcap%
\pgfsetroundjoin%
\pgfsetlinewidth{1.003750pt}%
\definecolor{currentstroke}{rgb}{0.690196,0.690196,0.690196}%
\pgfsetstrokecolor{currentstroke}%
\pgfsetstrokeopacity{0.600000}%
\pgfsetdash{{1.000000pt}{1.650000pt}}{0.000000pt}%
\pgfpathmoveto{\pgfqpoint{1.281818in}{0.440000in}}%
\pgfpathlineto{\pgfqpoint{1.281818in}{3.520000in}}%
\pgfusepath{stroke}%
\end{pgfscope}%
\begin{pgfscope}%
\pgfsetbuttcap%
\pgfsetroundjoin%
\definecolor{currentfill}{rgb}{0.000000,0.000000,0.000000}%
\pgfsetfillcolor{currentfill}%
\pgfsetlinewidth{0.803000pt}%
\definecolor{currentstroke}{rgb}{0.000000,0.000000,0.000000}%
\pgfsetstrokecolor{currentstroke}%
\pgfsetdash{}{0pt}%
\pgfsys@defobject{currentmarker}{\pgfqpoint{0.000000in}{-0.048611in}}{\pgfqpoint{0.000000in}{0.000000in}}{%
\pgfpathmoveto{\pgfqpoint{0.000000in}{0.000000in}}%
\pgfpathlineto{\pgfqpoint{0.000000in}{-0.048611in}}%
\pgfusepath{stroke,fill}%
}%
\begin{pgfscope}%
\pgfsys@transformshift{1.281818in}{0.440000in}%
\pgfsys@useobject{currentmarker}{}%
\end{pgfscope}%
\end{pgfscope}%
\begin{pgfscope}%
\definecolor{textcolor}{rgb}{0.000000,0.000000,0.000000}%
\pgfsetstrokecolor{textcolor}%
\pgfsetfillcolor{textcolor}%
\pgftext[x=1.281818in,y=0.342778in,,top]{\color{textcolor}{\rmfamily\fontsize{10.000000}{12.000000}\selectfont\catcode`\^=\active\def^{\ifmmode\sp\else\^{}\fi}\catcode`\%=\active\def%{\%}$\mathdefault{10^{0}}$}}%
\end{pgfscope}%
\begin{pgfscope}%
\pgfpathrectangle{\pgfqpoint{1.000000in}{0.440000in}}{\pgfqpoint{6.200000in}{3.080000in}}%
\pgfusepath{clip}%
\pgfsetbuttcap%
\pgfsetroundjoin%
\pgfsetlinewidth{1.003750pt}%
\definecolor{currentstroke}{rgb}{0.690196,0.690196,0.690196}%
\pgfsetstrokecolor{currentstroke}%
\pgfsetstrokeopacity{0.600000}%
\pgfsetdash{{1.000000pt}{1.650000pt}}{0.000000pt}%
\pgfpathmoveto{\pgfqpoint{2.236634in}{0.440000in}}%
\pgfpathlineto{\pgfqpoint{2.236634in}{3.520000in}}%
\pgfusepath{stroke}%
\end{pgfscope}%
\begin{pgfscope}%
\pgfsetbuttcap%
\pgfsetroundjoin%
\definecolor{currentfill}{rgb}{0.000000,0.000000,0.000000}%
\pgfsetfillcolor{currentfill}%
\pgfsetlinewidth{0.803000pt}%
\definecolor{currentstroke}{rgb}{0.000000,0.000000,0.000000}%
\pgfsetstrokecolor{currentstroke}%
\pgfsetdash{}{0pt}%
\pgfsys@defobject{currentmarker}{\pgfqpoint{0.000000in}{-0.048611in}}{\pgfqpoint{0.000000in}{0.000000in}}{%
\pgfpathmoveto{\pgfqpoint{0.000000in}{0.000000in}}%
\pgfpathlineto{\pgfqpoint{0.000000in}{-0.048611in}}%
\pgfusepath{stroke,fill}%
}%
\begin{pgfscope}%
\pgfsys@transformshift{2.236634in}{0.440000in}%
\pgfsys@useobject{currentmarker}{}%
\end{pgfscope}%
\end{pgfscope}%
\begin{pgfscope}%
\definecolor{textcolor}{rgb}{0.000000,0.000000,0.000000}%
\pgfsetstrokecolor{textcolor}%
\pgfsetfillcolor{textcolor}%
\pgftext[x=2.236634in,y=0.342778in,,top]{\color{textcolor}{\rmfamily\fontsize{10.000000}{12.000000}\selectfont\catcode`\^=\active\def^{\ifmmode\sp\else\^{}\fi}\catcode`\%=\active\def%{\%}$\mathdefault{10^{1}}$}}%
\end{pgfscope}%
\begin{pgfscope}%
\pgfpathrectangle{\pgfqpoint{1.000000in}{0.440000in}}{\pgfqpoint{6.200000in}{3.080000in}}%
\pgfusepath{clip}%
\pgfsetbuttcap%
\pgfsetroundjoin%
\pgfsetlinewidth{1.003750pt}%
\definecolor{currentstroke}{rgb}{0.690196,0.690196,0.690196}%
\pgfsetstrokecolor{currentstroke}%
\pgfsetstrokeopacity{0.600000}%
\pgfsetdash{{1.000000pt}{1.650000pt}}{0.000000pt}%
\pgfpathmoveto{\pgfqpoint{3.191450in}{0.440000in}}%
\pgfpathlineto{\pgfqpoint{3.191450in}{3.520000in}}%
\pgfusepath{stroke}%
\end{pgfscope}%
\begin{pgfscope}%
\pgfsetbuttcap%
\pgfsetroundjoin%
\definecolor{currentfill}{rgb}{0.000000,0.000000,0.000000}%
\pgfsetfillcolor{currentfill}%
\pgfsetlinewidth{0.803000pt}%
\definecolor{currentstroke}{rgb}{0.000000,0.000000,0.000000}%
\pgfsetstrokecolor{currentstroke}%
\pgfsetdash{}{0pt}%
\pgfsys@defobject{currentmarker}{\pgfqpoint{0.000000in}{-0.048611in}}{\pgfqpoint{0.000000in}{0.000000in}}{%
\pgfpathmoveto{\pgfqpoint{0.000000in}{0.000000in}}%
\pgfpathlineto{\pgfqpoint{0.000000in}{-0.048611in}}%
\pgfusepath{stroke,fill}%
}%
\begin{pgfscope}%
\pgfsys@transformshift{3.191450in}{0.440000in}%
\pgfsys@useobject{currentmarker}{}%
\end{pgfscope}%
\end{pgfscope}%
\begin{pgfscope}%
\definecolor{textcolor}{rgb}{0.000000,0.000000,0.000000}%
\pgfsetstrokecolor{textcolor}%
\pgfsetfillcolor{textcolor}%
\pgftext[x=3.191450in,y=0.342778in,,top]{\color{textcolor}{\rmfamily\fontsize{10.000000}{12.000000}\selectfont\catcode`\^=\active\def^{\ifmmode\sp\else\^{}\fi}\catcode`\%=\active\def%{\%}$\mathdefault{10^{2}}$}}%
\end{pgfscope}%
\begin{pgfscope}%
\pgfpathrectangle{\pgfqpoint{1.000000in}{0.440000in}}{\pgfqpoint{6.200000in}{3.080000in}}%
\pgfusepath{clip}%
\pgfsetbuttcap%
\pgfsetroundjoin%
\pgfsetlinewidth{1.003750pt}%
\definecolor{currentstroke}{rgb}{0.690196,0.690196,0.690196}%
\pgfsetstrokecolor{currentstroke}%
\pgfsetstrokeopacity{0.600000}%
\pgfsetdash{{1.000000pt}{1.650000pt}}{0.000000pt}%
\pgfpathmoveto{\pgfqpoint{4.146266in}{0.440000in}}%
\pgfpathlineto{\pgfqpoint{4.146266in}{3.520000in}}%
\pgfusepath{stroke}%
\end{pgfscope}%
\begin{pgfscope}%
\pgfsetbuttcap%
\pgfsetroundjoin%
\definecolor{currentfill}{rgb}{0.000000,0.000000,0.000000}%
\pgfsetfillcolor{currentfill}%
\pgfsetlinewidth{0.803000pt}%
\definecolor{currentstroke}{rgb}{0.000000,0.000000,0.000000}%
\pgfsetstrokecolor{currentstroke}%
\pgfsetdash{}{0pt}%
\pgfsys@defobject{currentmarker}{\pgfqpoint{0.000000in}{-0.048611in}}{\pgfqpoint{0.000000in}{0.000000in}}{%
\pgfpathmoveto{\pgfqpoint{0.000000in}{0.000000in}}%
\pgfpathlineto{\pgfqpoint{0.000000in}{-0.048611in}}%
\pgfusepath{stroke,fill}%
}%
\begin{pgfscope}%
\pgfsys@transformshift{4.146266in}{0.440000in}%
\pgfsys@useobject{currentmarker}{}%
\end{pgfscope}%
\end{pgfscope}%
\begin{pgfscope}%
\definecolor{textcolor}{rgb}{0.000000,0.000000,0.000000}%
\pgfsetstrokecolor{textcolor}%
\pgfsetfillcolor{textcolor}%
\pgftext[x=4.146266in,y=0.342778in,,top]{\color{textcolor}{\rmfamily\fontsize{10.000000}{12.000000}\selectfont\catcode`\^=\active\def^{\ifmmode\sp\else\^{}\fi}\catcode`\%=\active\def%{\%}$\mathdefault{10^{3}}$}}%
\end{pgfscope}%
\begin{pgfscope}%
\pgfpathrectangle{\pgfqpoint{1.000000in}{0.440000in}}{\pgfqpoint{6.200000in}{3.080000in}}%
\pgfusepath{clip}%
\pgfsetbuttcap%
\pgfsetroundjoin%
\pgfsetlinewidth{1.003750pt}%
\definecolor{currentstroke}{rgb}{0.690196,0.690196,0.690196}%
\pgfsetstrokecolor{currentstroke}%
\pgfsetstrokeopacity{0.600000}%
\pgfsetdash{{1.000000pt}{1.650000pt}}{0.000000pt}%
\pgfpathmoveto{\pgfqpoint{5.101082in}{0.440000in}}%
\pgfpathlineto{\pgfqpoint{5.101082in}{3.520000in}}%
\pgfusepath{stroke}%
\end{pgfscope}%
\begin{pgfscope}%
\pgfsetbuttcap%
\pgfsetroundjoin%
\definecolor{currentfill}{rgb}{0.000000,0.000000,0.000000}%
\pgfsetfillcolor{currentfill}%
\pgfsetlinewidth{0.803000pt}%
\definecolor{currentstroke}{rgb}{0.000000,0.000000,0.000000}%
\pgfsetstrokecolor{currentstroke}%
\pgfsetdash{}{0pt}%
\pgfsys@defobject{currentmarker}{\pgfqpoint{0.000000in}{-0.048611in}}{\pgfqpoint{0.000000in}{0.000000in}}{%
\pgfpathmoveto{\pgfqpoint{0.000000in}{0.000000in}}%
\pgfpathlineto{\pgfqpoint{0.000000in}{-0.048611in}}%
\pgfusepath{stroke,fill}%
}%
\begin{pgfscope}%
\pgfsys@transformshift{5.101082in}{0.440000in}%
\pgfsys@useobject{currentmarker}{}%
\end{pgfscope}%
\end{pgfscope}%
\begin{pgfscope}%
\definecolor{textcolor}{rgb}{0.000000,0.000000,0.000000}%
\pgfsetstrokecolor{textcolor}%
\pgfsetfillcolor{textcolor}%
\pgftext[x=5.101082in,y=0.342778in,,top]{\color{textcolor}{\rmfamily\fontsize{10.000000}{12.000000}\selectfont\catcode`\^=\active\def^{\ifmmode\sp\else\^{}\fi}\catcode`\%=\active\def%{\%}$\mathdefault{10^{4}}$}}%
\end{pgfscope}%
\begin{pgfscope}%
\pgfpathrectangle{\pgfqpoint{1.000000in}{0.440000in}}{\pgfqpoint{6.200000in}{3.080000in}}%
\pgfusepath{clip}%
\pgfsetbuttcap%
\pgfsetroundjoin%
\pgfsetlinewidth{1.003750pt}%
\definecolor{currentstroke}{rgb}{0.690196,0.690196,0.690196}%
\pgfsetstrokecolor{currentstroke}%
\pgfsetstrokeopacity{0.600000}%
\pgfsetdash{{1.000000pt}{1.650000pt}}{0.000000pt}%
\pgfpathmoveto{\pgfqpoint{6.055898in}{0.440000in}}%
\pgfpathlineto{\pgfqpoint{6.055898in}{3.520000in}}%
\pgfusepath{stroke}%
\end{pgfscope}%
\begin{pgfscope}%
\pgfsetbuttcap%
\pgfsetroundjoin%
\definecolor{currentfill}{rgb}{0.000000,0.000000,0.000000}%
\pgfsetfillcolor{currentfill}%
\pgfsetlinewidth{0.803000pt}%
\definecolor{currentstroke}{rgb}{0.000000,0.000000,0.000000}%
\pgfsetstrokecolor{currentstroke}%
\pgfsetdash{}{0pt}%
\pgfsys@defobject{currentmarker}{\pgfqpoint{0.000000in}{-0.048611in}}{\pgfqpoint{0.000000in}{0.000000in}}{%
\pgfpathmoveto{\pgfqpoint{0.000000in}{0.000000in}}%
\pgfpathlineto{\pgfqpoint{0.000000in}{-0.048611in}}%
\pgfusepath{stroke,fill}%
}%
\begin{pgfscope}%
\pgfsys@transformshift{6.055898in}{0.440000in}%
\pgfsys@useobject{currentmarker}{}%
\end{pgfscope}%
\end{pgfscope}%
\begin{pgfscope}%
\definecolor{textcolor}{rgb}{0.000000,0.000000,0.000000}%
\pgfsetstrokecolor{textcolor}%
\pgfsetfillcolor{textcolor}%
\pgftext[x=6.055898in,y=0.342778in,,top]{\color{textcolor}{\rmfamily\fontsize{10.000000}{12.000000}\selectfont\catcode`\^=\active\def^{\ifmmode\sp\else\^{}\fi}\catcode`\%=\active\def%{\%}$\mathdefault{10^{5}}$}}%
\end{pgfscope}%
\begin{pgfscope}%
\pgfpathrectangle{\pgfqpoint{1.000000in}{0.440000in}}{\pgfqpoint{6.200000in}{3.080000in}}%
\pgfusepath{clip}%
\pgfsetbuttcap%
\pgfsetroundjoin%
\pgfsetlinewidth{1.003750pt}%
\definecolor{currentstroke}{rgb}{0.690196,0.690196,0.690196}%
\pgfsetstrokecolor{currentstroke}%
\pgfsetstrokeopacity{0.600000}%
\pgfsetdash{{1.000000pt}{1.650000pt}}{0.000000pt}%
\pgfpathmoveto{\pgfqpoint{7.010714in}{0.440000in}}%
\pgfpathlineto{\pgfqpoint{7.010714in}{3.520000in}}%
\pgfusepath{stroke}%
\end{pgfscope}%
\begin{pgfscope}%
\pgfsetbuttcap%
\pgfsetroundjoin%
\definecolor{currentfill}{rgb}{0.000000,0.000000,0.000000}%
\pgfsetfillcolor{currentfill}%
\pgfsetlinewidth{0.803000pt}%
\definecolor{currentstroke}{rgb}{0.000000,0.000000,0.000000}%
\pgfsetstrokecolor{currentstroke}%
\pgfsetdash{}{0pt}%
\pgfsys@defobject{currentmarker}{\pgfqpoint{0.000000in}{-0.048611in}}{\pgfqpoint{0.000000in}{0.000000in}}{%
\pgfpathmoveto{\pgfqpoint{0.000000in}{0.000000in}}%
\pgfpathlineto{\pgfqpoint{0.000000in}{-0.048611in}}%
\pgfusepath{stroke,fill}%
}%
\begin{pgfscope}%
\pgfsys@transformshift{7.010714in}{0.440000in}%
\pgfsys@useobject{currentmarker}{}%
\end{pgfscope}%
\end{pgfscope}%
\begin{pgfscope}%
\definecolor{textcolor}{rgb}{0.000000,0.000000,0.000000}%
\pgfsetstrokecolor{textcolor}%
\pgfsetfillcolor{textcolor}%
\pgftext[x=7.010714in,y=0.342778in,,top]{\color{textcolor}{\rmfamily\fontsize{10.000000}{12.000000}\selectfont\catcode`\^=\active\def^{\ifmmode\sp\else\^{}\fi}\catcode`\%=\active\def%{\%}$\mathdefault{10^{6}}$}}%
\end{pgfscope}%
\begin{pgfscope}%
\pgfpathrectangle{\pgfqpoint{1.000000in}{0.440000in}}{\pgfqpoint{6.200000in}{3.080000in}}%
\pgfusepath{clip}%
\pgfsetbuttcap%
\pgfsetroundjoin%
\pgfsetlinewidth{1.003750pt}%
\definecolor{currentstroke}{rgb}{0.690196,0.690196,0.690196}%
\pgfsetstrokecolor{currentstroke}%
\pgfsetstrokeopacity{0.600000}%
\pgfsetdash{{1.000000pt}{1.650000pt}}{0.000000pt}%
\pgfpathmoveto{\pgfqpoint{1.069993in}{0.440000in}}%
\pgfpathlineto{\pgfqpoint{1.069993in}{3.520000in}}%
\pgfusepath{stroke}%
\end{pgfscope}%
\begin{pgfscope}%
\pgfsetbuttcap%
\pgfsetroundjoin%
\definecolor{currentfill}{rgb}{0.000000,0.000000,0.000000}%
\pgfsetfillcolor{currentfill}%
\pgfsetlinewidth{0.602250pt}%
\definecolor{currentstroke}{rgb}{0.000000,0.000000,0.000000}%
\pgfsetstrokecolor{currentstroke}%
\pgfsetdash{}{0pt}%
\pgfsys@defobject{currentmarker}{\pgfqpoint{0.000000in}{-0.027778in}}{\pgfqpoint{0.000000in}{0.000000in}}{%
\pgfpathmoveto{\pgfqpoint{0.000000in}{0.000000in}}%
\pgfpathlineto{\pgfqpoint{0.000000in}{-0.027778in}}%
\pgfusepath{stroke,fill}%
}%
\begin{pgfscope}%
\pgfsys@transformshift{1.069993in}{0.440000in}%
\pgfsys@useobject{currentmarker}{}%
\end{pgfscope}%
\end{pgfscope}%
\begin{pgfscope}%
\pgfpathrectangle{\pgfqpoint{1.000000in}{0.440000in}}{\pgfqpoint{6.200000in}{3.080000in}}%
\pgfusepath{clip}%
\pgfsetbuttcap%
\pgfsetroundjoin%
\pgfsetlinewidth{1.003750pt}%
\definecolor{currentstroke}{rgb}{0.690196,0.690196,0.690196}%
\pgfsetstrokecolor{currentstroke}%
\pgfsetstrokeopacity{0.600000}%
\pgfsetdash{{1.000000pt}{1.650000pt}}{0.000000pt}%
\pgfpathmoveto{\pgfqpoint{1.133915in}{0.440000in}}%
\pgfpathlineto{\pgfqpoint{1.133915in}{3.520000in}}%
\pgfusepath{stroke}%
\end{pgfscope}%
\begin{pgfscope}%
\pgfsetbuttcap%
\pgfsetroundjoin%
\definecolor{currentfill}{rgb}{0.000000,0.000000,0.000000}%
\pgfsetfillcolor{currentfill}%
\pgfsetlinewidth{0.602250pt}%
\definecolor{currentstroke}{rgb}{0.000000,0.000000,0.000000}%
\pgfsetstrokecolor{currentstroke}%
\pgfsetdash{}{0pt}%
\pgfsys@defobject{currentmarker}{\pgfqpoint{0.000000in}{-0.027778in}}{\pgfqpoint{0.000000in}{0.000000in}}{%
\pgfpathmoveto{\pgfqpoint{0.000000in}{0.000000in}}%
\pgfpathlineto{\pgfqpoint{0.000000in}{-0.027778in}}%
\pgfusepath{stroke,fill}%
}%
\begin{pgfscope}%
\pgfsys@transformshift{1.133915in}{0.440000in}%
\pgfsys@useobject{currentmarker}{}%
\end{pgfscope}%
\end{pgfscope}%
\begin{pgfscope}%
\pgfpathrectangle{\pgfqpoint{1.000000in}{0.440000in}}{\pgfqpoint{6.200000in}{3.080000in}}%
\pgfusepath{clip}%
\pgfsetbuttcap%
\pgfsetroundjoin%
\pgfsetlinewidth{1.003750pt}%
\definecolor{currentstroke}{rgb}{0.690196,0.690196,0.690196}%
\pgfsetstrokecolor{currentstroke}%
\pgfsetstrokeopacity{0.600000}%
\pgfsetdash{{1.000000pt}{1.650000pt}}{0.000000pt}%
\pgfpathmoveto{\pgfqpoint{1.189287in}{0.440000in}}%
\pgfpathlineto{\pgfqpoint{1.189287in}{3.520000in}}%
\pgfusepath{stroke}%
\end{pgfscope}%
\begin{pgfscope}%
\pgfsetbuttcap%
\pgfsetroundjoin%
\definecolor{currentfill}{rgb}{0.000000,0.000000,0.000000}%
\pgfsetfillcolor{currentfill}%
\pgfsetlinewidth{0.602250pt}%
\definecolor{currentstroke}{rgb}{0.000000,0.000000,0.000000}%
\pgfsetstrokecolor{currentstroke}%
\pgfsetdash{}{0pt}%
\pgfsys@defobject{currentmarker}{\pgfqpoint{0.000000in}{-0.027778in}}{\pgfqpoint{0.000000in}{0.000000in}}{%
\pgfpathmoveto{\pgfqpoint{0.000000in}{0.000000in}}%
\pgfpathlineto{\pgfqpoint{0.000000in}{-0.027778in}}%
\pgfusepath{stroke,fill}%
}%
\begin{pgfscope}%
\pgfsys@transformshift{1.189287in}{0.440000in}%
\pgfsys@useobject{currentmarker}{}%
\end{pgfscope}%
\end{pgfscope}%
\begin{pgfscope}%
\pgfpathrectangle{\pgfqpoint{1.000000in}{0.440000in}}{\pgfqpoint{6.200000in}{3.080000in}}%
\pgfusepath{clip}%
\pgfsetbuttcap%
\pgfsetroundjoin%
\pgfsetlinewidth{1.003750pt}%
\definecolor{currentstroke}{rgb}{0.690196,0.690196,0.690196}%
\pgfsetstrokecolor{currentstroke}%
\pgfsetstrokeopacity{0.600000}%
\pgfsetdash{{1.000000pt}{1.650000pt}}{0.000000pt}%
\pgfpathmoveto{\pgfqpoint{1.238128in}{0.440000in}}%
\pgfpathlineto{\pgfqpoint{1.238128in}{3.520000in}}%
\pgfusepath{stroke}%
\end{pgfscope}%
\begin{pgfscope}%
\pgfsetbuttcap%
\pgfsetroundjoin%
\definecolor{currentfill}{rgb}{0.000000,0.000000,0.000000}%
\pgfsetfillcolor{currentfill}%
\pgfsetlinewidth{0.602250pt}%
\definecolor{currentstroke}{rgb}{0.000000,0.000000,0.000000}%
\pgfsetstrokecolor{currentstroke}%
\pgfsetdash{}{0pt}%
\pgfsys@defobject{currentmarker}{\pgfqpoint{0.000000in}{-0.027778in}}{\pgfqpoint{0.000000in}{0.000000in}}{%
\pgfpathmoveto{\pgfqpoint{0.000000in}{0.000000in}}%
\pgfpathlineto{\pgfqpoint{0.000000in}{-0.027778in}}%
\pgfusepath{stroke,fill}%
}%
\begin{pgfscope}%
\pgfsys@transformshift{1.238128in}{0.440000in}%
\pgfsys@useobject{currentmarker}{}%
\end{pgfscope}%
\end{pgfscope}%
\begin{pgfscope}%
\pgfpathrectangle{\pgfqpoint{1.000000in}{0.440000in}}{\pgfqpoint{6.200000in}{3.080000in}}%
\pgfusepath{clip}%
\pgfsetbuttcap%
\pgfsetroundjoin%
\pgfsetlinewidth{1.003750pt}%
\definecolor{currentstroke}{rgb}{0.690196,0.690196,0.690196}%
\pgfsetstrokecolor{currentstroke}%
\pgfsetstrokeopacity{0.600000}%
\pgfsetdash{{1.000000pt}{1.650000pt}}{0.000000pt}%
\pgfpathmoveto{\pgfqpoint{1.569246in}{0.440000in}}%
\pgfpathlineto{\pgfqpoint{1.569246in}{3.520000in}}%
\pgfusepath{stroke}%
\end{pgfscope}%
\begin{pgfscope}%
\pgfsetbuttcap%
\pgfsetroundjoin%
\definecolor{currentfill}{rgb}{0.000000,0.000000,0.000000}%
\pgfsetfillcolor{currentfill}%
\pgfsetlinewidth{0.602250pt}%
\definecolor{currentstroke}{rgb}{0.000000,0.000000,0.000000}%
\pgfsetstrokecolor{currentstroke}%
\pgfsetdash{}{0pt}%
\pgfsys@defobject{currentmarker}{\pgfqpoint{0.000000in}{-0.027778in}}{\pgfqpoint{0.000000in}{0.000000in}}{%
\pgfpathmoveto{\pgfqpoint{0.000000in}{0.000000in}}%
\pgfpathlineto{\pgfqpoint{0.000000in}{-0.027778in}}%
\pgfusepath{stroke,fill}%
}%
\begin{pgfscope}%
\pgfsys@transformshift{1.569246in}{0.440000in}%
\pgfsys@useobject{currentmarker}{}%
\end{pgfscope}%
\end{pgfscope}%
\begin{pgfscope}%
\pgfpathrectangle{\pgfqpoint{1.000000in}{0.440000in}}{\pgfqpoint{6.200000in}{3.080000in}}%
\pgfusepath{clip}%
\pgfsetbuttcap%
\pgfsetroundjoin%
\pgfsetlinewidth{1.003750pt}%
\definecolor{currentstroke}{rgb}{0.690196,0.690196,0.690196}%
\pgfsetstrokecolor{currentstroke}%
\pgfsetstrokeopacity{0.600000}%
\pgfsetdash{{1.000000pt}{1.650000pt}}{0.000000pt}%
\pgfpathmoveto{\pgfqpoint{1.737381in}{0.440000in}}%
\pgfpathlineto{\pgfqpoint{1.737381in}{3.520000in}}%
\pgfusepath{stroke}%
\end{pgfscope}%
\begin{pgfscope}%
\pgfsetbuttcap%
\pgfsetroundjoin%
\definecolor{currentfill}{rgb}{0.000000,0.000000,0.000000}%
\pgfsetfillcolor{currentfill}%
\pgfsetlinewidth{0.602250pt}%
\definecolor{currentstroke}{rgb}{0.000000,0.000000,0.000000}%
\pgfsetstrokecolor{currentstroke}%
\pgfsetdash{}{0pt}%
\pgfsys@defobject{currentmarker}{\pgfqpoint{0.000000in}{-0.027778in}}{\pgfqpoint{0.000000in}{0.000000in}}{%
\pgfpathmoveto{\pgfqpoint{0.000000in}{0.000000in}}%
\pgfpathlineto{\pgfqpoint{0.000000in}{-0.027778in}}%
\pgfusepath{stroke,fill}%
}%
\begin{pgfscope}%
\pgfsys@transformshift{1.737381in}{0.440000in}%
\pgfsys@useobject{currentmarker}{}%
\end{pgfscope}%
\end{pgfscope}%
\begin{pgfscope}%
\pgfpathrectangle{\pgfqpoint{1.000000in}{0.440000in}}{\pgfqpoint{6.200000in}{3.080000in}}%
\pgfusepath{clip}%
\pgfsetbuttcap%
\pgfsetroundjoin%
\pgfsetlinewidth{1.003750pt}%
\definecolor{currentstroke}{rgb}{0.690196,0.690196,0.690196}%
\pgfsetstrokecolor{currentstroke}%
\pgfsetstrokeopacity{0.600000}%
\pgfsetdash{{1.000000pt}{1.650000pt}}{0.000000pt}%
\pgfpathmoveto{\pgfqpoint{1.856675in}{0.440000in}}%
\pgfpathlineto{\pgfqpoint{1.856675in}{3.520000in}}%
\pgfusepath{stroke}%
\end{pgfscope}%
\begin{pgfscope}%
\pgfsetbuttcap%
\pgfsetroundjoin%
\definecolor{currentfill}{rgb}{0.000000,0.000000,0.000000}%
\pgfsetfillcolor{currentfill}%
\pgfsetlinewidth{0.602250pt}%
\definecolor{currentstroke}{rgb}{0.000000,0.000000,0.000000}%
\pgfsetstrokecolor{currentstroke}%
\pgfsetdash{}{0pt}%
\pgfsys@defobject{currentmarker}{\pgfqpoint{0.000000in}{-0.027778in}}{\pgfqpoint{0.000000in}{0.000000in}}{%
\pgfpathmoveto{\pgfqpoint{0.000000in}{0.000000in}}%
\pgfpathlineto{\pgfqpoint{0.000000in}{-0.027778in}}%
\pgfusepath{stroke,fill}%
}%
\begin{pgfscope}%
\pgfsys@transformshift{1.856675in}{0.440000in}%
\pgfsys@useobject{currentmarker}{}%
\end{pgfscope}%
\end{pgfscope}%
\begin{pgfscope}%
\pgfpathrectangle{\pgfqpoint{1.000000in}{0.440000in}}{\pgfqpoint{6.200000in}{3.080000in}}%
\pgfusepath{clip}%
\pgfsetbuttcap%
\pgfsetroundjoin%
\pgfsetlinewidth{1.003750pt}%
\definecolor{currentstroke}{rgb}{0.690196,0.690196,0.690196}%
\pgfsetstrokecolor{currentstroke}%
\pgfsetstrokeopacity{0.600000}%
\pgfsetdash{{1.000000pt}{1.650000pt}}{0.000000pt}%
\pgfpathmoveto{\pgfqpoint{1.949206in}{0.440000in}}%
\pgfpathlineto{\pgfqpoint{1.949206in}{3.520000in}}%
\pgfusepath{stroke}%
\end{pgfscope}%
\begin{pgfscope}%
\pgfsetbuttcap%
\pgfsetroundjoin%
\definecolor{currentfill}{rgb}{0.000000,0.000000,0.000000}%
\pgfsetfillcolor{currentfill}%
\pgfsetlinewidth{0.602250pt}%
\definecolor{currentstroke}{rgb}{0.000000,0.000000,0.000000}%
\pgfsetstrokecolor{currentstroke}%
\pgfsetdash{}{0pt}%
\pgfsys@defobject{currentmarker}{\pgfqpoint{0.000000in}{-0.027778in}}{\pgfqpoint{0.000000in}{0.000000in}}{%
\pgfpathmoveto{\pgfqpoint{0.000000in}{0.000000in}}%
\pgfpathlineto{\pgfqpoint{0.000000in}{-0.027778in}}%
\pgfusepath{stroke,fill}%
}%
\begin{pgfscope}%
\pgfsys@transformshift{1.949206in}{0.440000in}%
\pgfsys@useobject{currentmarker}{}%
\end{pgfscope}%
\end{pgfscope}%
\begin{pgfscope}%
\pgfpathrectangle{\pgfqpoint{1.000000in}{0.440000in}}{\pgfqpoint{6.200000in}{3.080000in}}%
\pgfusepath{clip}%
\pgfsetbuttcap%
\pgfsetroundjoin%
\pgfsetlinewidth{1.003750pt}%
\definecolor{currentstroke}{rgb}{0.690196,0.690196,0.690196}%
\pgfsetstrokecolor{currentstroke}%
\pgfsetstrokeopacity{0.600000}%
\pgfsetdash{{1.000000pt}{1.650000pt}}{0.000000pt}%
\pgfpathmoveto{\pgfqpoint{2.024809in}{0.440000in}}%
\pgfpathlineto{\pgfqpoint{2.024809in}{3.520000in}}%
\pgfusepath{stroke}%
\end{pgfscope}%
\begin{pgfscope}%
\pgfsetbuttcap%
\pgfsetroundjoin%
\definecolor{currentfill}{rgb}{0.000000,0.000000,0.000000}%
\pgfsetfillcolor{currentfill}%
\pgfsetlinewidth{0.602250pt}%
\definecolor{currentstroke}{rgb}{0.000000,0.000000,0.000000}%
\pgfsetstrokecolor{currentstroke}%
\pgfsetdash{}{0pt}%
\pgfsys@defobject{currentmarker}{\pgfqpoint{0.000000in}{-0.027778in}}{\pgfqpoint{0.000000in}{0.000000in}}{%
\pgfpathmoveto{\pgfqpoint{0.000000in}{0.000000in}}%
\pgfpathlineto{\pgfqpoint{0.000000in}{-0.027778in}}%
\pgfusepath{stroke,fill}%
}%
\begin{pgfscope}%
\pgfsys@transformshift{2.024809in}{0.440000in}%
\pgfsys@useobject{currentmarker}{}%
\end{pgfscope}%
\end{pgfscope}%
\begin{pgfscope}%
\pgfpathrectangle{\pgfqpoint{1.000000in}{0.440000in}}{\pgfqpoint{6.200000in}{3.080000in}}%
\pgfusepath{clip}%
\pgfsetbuttcap%
\pgfsetroundjoin%
\pgfsetlinewidth{1.003750pt}%
\definecolor{currentstroke}{rgb}{0.690196,0.690196,0.690196}%
\pgfsetstrokecolor{currentstroke}%
\pgfsetstrokeopacity{0.600000}%
\pgfsetdash{{1.000000pt}{1.650000pt}}{0.000000pt}%
\pgfpathmoveto{\pgfqpoint{2.088731in}{0.440000in}}%
\pgfpathlineto{\pgfqpoint{2.088731in}{3.520000in}}%
\pgfusepath{stroke}%
\end{pgfscope}%
\begin{pgfscope}%
\pgfsetbuttcap%
\pgfsetroundjoin%
\definecolor{currentfill}{rgb}{0.000000,0.000000,0.000000}%
\pgfsetfillcolor{currentfill}%
\pgfsetlinewidth{0.602250pt}%
\definecolor{currentstroke}{rgb}{0.000000,0.000000,0.000000}%
\pgfsetstrokecolor{currentstroke}%
\pgfsetdash{}{0pt}%
\pgfsys@defobject{currentmarker}{\pgfqpoint{0.000000in}{-0.027778in}}{\pgfqpoint{0.000000in}{0.000000in}}{%
\pgfpathmoveto{\pgfqpoint{0.000000in}{0.000000in}}%
\pgfpathlineto{\pgfqpoint{0.000000in}{-0.027778in}}%
\pgfusepath{stroke,fill}%
}%
\begin{pgfscope}%
\pgfsys@transformshift{2.088731in}{0.440000in}%
\pgfsys@useobject{currentmarker}{}%
\end{pgfscope}%
\end{pgfscope}%
\begin{pgfscope}%
\pgfpathrectangle{\pgfqpoint{1.000000in}{0.440000in}}{\pgfqpoint{6.200000in}{3.080000in}}%
\pgfusepath{clip}%
\pgfsetbuttcap%
\pgfsetroundjoin%
\pgfsetlinewidth{1.003750pt}%
\definecolor{currentstroke}{rgb}{0.690196,0.690196,0.690196}%
\pgfsetstrokecolor{currentstroke}%
\pgfsetstrokeopacity{0.600000}%
\pgfsetdash{{1.000000pt}{1.650000pt}}{0.000000pt}%
\pgfpathmoveto{\pgfqpoint{2.144103in}{0.440000in}}%
\pgfpathlineto{\pgfqpoint{2.144103in}{3.520000in}}%
\pgfusepath{stroke}%
\end{pgfscope}%
\begin{pgfscope}%
\pgfsetbuttcap%
\pgfsetroundjoin%
\definecolor{currentfill}{rgb}{0.000000,0.000000,0.000000}%
\pgfsetfillcolor{currentfill}%
\pgfsetlinewidth{0.602250pt}%
\definecolor{currentstroke}{rgb}{0.000000,0.000000,0.000000}%
\pgfsetstrokecolor{currentstroke}%
\pgfsetdash{}{0pt}%
\pgfsys@defobject{currentmarker}{\pgfqpoint{0.000000in}{-0.027778in}}{\pgfqpoint{0.000000in}{0.000000in}}{%
\pgfpathmoveto{\pgfqpoint{0.000000in}{0.000000in}}%
\pgfpathlineto{\pgfqpoint{0.000000in}{-0.027778in}}%
\pgfusepath{stroke,fill}%
}%
\begin{pgfscope}%
\pgfsys@transformshift{2.144103in}{0.440000in}%
\pgfsys@useobject{currentmarker}{}%
\end{pgfscope}%
\end{pgfscope}%
\begin{pgfscope}%
\pgfpathrectangle{\pgfqpoint{1.000000in}{0.440000in}}{\pgfqpoint{6.200000in}{3.080000in}}%
\pgfusepath{clip}%
\pgfsetbuttcap%
\pgfsetroundjoin%
\pgfsetlinewidth{1.003750pt}%
\definecolor{currentstroke}{rgb}{0.690196,0.690196,0.690196}%
\pgfsetstrokecolor{currentstroke}%
\pgfsetstrokeopacity{0.600000}%
\pgfsetdash{{1.000000pt}{1.650000pt}}{0.000000pt}%
\pgfpathmoveto{\pgfqpoint{2.192944in}{0.440000in}}%
\pgfpathlineto{\pgfqpoint{2.192944in}{3.520000in}}%
\pgfusepath{stroke}%
\end{pgfscope}%
\begin{pgfscope}%
\pgfsetbuttcap%
\pgfsetroundjoin%
\definecolor{currentfill}{rgb}{0.000000,0.000000,0.000000}%
\pgfsetfillcolor{currentfill}%
\pgfsetlinewidth{0.602250pt}%
\definecolor{currentstroke}{rgb}{0.000000,0.000000,0.000000}%
\pgfsetstrokecolor{currentstroke}%
\pgfsetdash{}{0pt}%
\pgfsys@defobject{currentmarker}{\pgfqpoint{0.000000in}{-0.027778in}}{\pgfqpoint{0.000000in}{0.000000in}}{%
\pgfpathmoveto{\pgfqpoint{0.000000in}{0.000000in}}%
\pgfpathlineto{\pgfqpoint{0.000000in}{-0.027778in}}%
\pgfusepath{stroke,fill}%
}%
\begin{pgfscope}%
\pgfsys@transformshift{2.192944in}{0.440000in}%
\pgfsys@useobject{currentmarker}{}%
\end{pgfscope}%
\end{pgfscope}%
\begin{pgfscope}%
\pgfpathrectangle{\pgfqpoint{1.000000in}{0.440000in}}{\pgfqpoint{6.200000in}{3.080000in}}%
\pgfusepath{clip}%
\pgfsetbuttcap%
\pgfsetroundjoin%
\pgfsetlinewidth{1.003750pt}%
\definecolor{currentstroke}{rgb}{0.690196,0.690196,0.690196}%
\pgfsetstrokecolor{currentstroke}%
\pgfsetstrokeopacity{0.600000}%
\pgfsetdash{{1.000000pt}{1.650000pt}}{0.000000pt}%
\pgfpathmoveto{\pgfqpoint{2.524062in}{0.440000in}}%
\pgfpathlineto{\pgfqpoint{2.524062in}{3.520000in}}%
\pgfusepath{stroke}%
\end{pgfscope}%
\begin{pgfscope}%
\pgfsetbuttcap%
\pgfsetroundjoin%
\definecolor{currentfill}{rgb}{0.000000,0.000000,0.000000}%
\pgfsetfillcolor{currentfill}%
\pgfsetlinewidth{0.602250pt}%
\definecolor{currentstroke}{rgb}{0.000000,0.000000,0.000000}%
\pgfsetstrokecolor{currentstroke}%
\pgfsetdash{}{0pt}%
\pgfsys@defobject{currentmarker}{\pgfqpoint{0.000000in}{-0.027778in}}{\pgfqpoint{0.000000in}{0.000000in}}{%
\pgfpathmoveto{\pgfqpoint{0.000000in}{0.000000in}}%
\pgfpathlineto{\pgfqpoint{0.000000in}{-0.027778in}}%
\pgfusepath{stroke,fill}%
}%
\begin{pgfscope}%
\pgfsys@transformshift{2.524062in}{0.440000in}%
\pgfsys@useobject{currentmarker}{}%
\end{pgfscope}%
\end{pgfscope}%
\begin{pgfscope}%
\pgfpathrectangle{\pgfqpoint{1.000000in}{0.440000in}}{\pgfqpoint{6.200000in}{3.080000in}}%
\pgfusepath{clip}%
\pgfsetbuttcap%
\pgfsetroundjoin%
\pgfsetlinewidth{1.003750pt}%
\definecolor{currentstroke}{rgb}{0.690196,0.690196,0.690196}%
\pgfsetstrokecolor{currentstroke}%
\pgfsetstrokeopacity{0.600000}%
\pgfsetdash{{1.000000pt}{1.650000pt}}{0.000000pt}%
\pgfpathmoveto{\pgfqpoint{2.692197in}{0.440000in}}%
\pgfpathlineto{\pgfqpoint{2.692197in}{3.520000in}}%
\pgfusepath{stroke}%
\end{pgfscope}%
\begin{pgfscope}%
\pgfsetbuttcap%
\pgfsetroundjoin%
\definecolor{currentfill}{rgb}{0.000000,0.000000,0.000000}%
\pgfsetfillcolor{currentfill}%
\pgfsetlinewidth{0.602250pt}%
\definecolor{currentstroke}{rgb}{0.000000,0.000000,0.000000}%
\pgfsetstrokecolor{currentstroke}%
\pgfsetdash{}{0pt}%
\pgfsys@defobject{currentmarker}{\pgfqpoint{0.000000in}{-0.027778in}}{\pgfqpoint{0.000000in}{0.000000in}}{%
\pgfpathmoveto{\pgfqpoint{0.000000in}{0.000000in}}%
\pgfpathlineto{\pgfqpoint{0.000000in}{-0.027778in}}%
\pgfusepath{stroke,fill}%
}%
\begin{pgfscope}%
\pgfsys@transformshift{2.692197in}{0.440000in}%
\pgfsys@useobject{currentmarker}{}%
\end{pgfscope}%
\end{pgfscope}%
\begin{pgfscope}%
\pgfpathrectangle{\pgfqpoint{1.000000in}{0.440000in}}{\pgfqpoint{6.200000in}{3.080000in}}%
\pgfusepath{clip}%
\pgfsetbuttcap%
\pgfsetroundjoin%
\pgfsetlinewidth{1.003750pt}%
\definecolor{currentstroke}{rgb}{0.690196,0.690196,0.690196}%
\pgfsetstrokecolor{currentstroke}%
\pgfsetstrokeopacity{0.600000}%
\pgfsetdash{{1.000000pt}{1.650000pt}}{0.000000pt}%
\pgfpathmoveto{\pgfqpoint{2.811491in}{0.440000in}}%
\pgfpathlineto{\pgfqpoint{2.811491in}{3.520000in}}%
\pgfusepath{stroke}%
\end{pgfscope}%
\begin{pgfscope}%
\pgfsetbuttcap%
\pgfsetroundjoin%
\definecolor{currentfill}{rgb}{0.000000,0.000000,0.000000}%
\pgfsetfillcolor{currentfill}%
\pgfsetlinewidth{0.602250pt}%
\definecolor{currentstroke}{rgb}{0.000000,0.000000,0.000000}%
\pgfsetstrokecolor{currentstroke}%
\pgfsetdash{}{0pt}%
\pgfsys@defobject{currentmarker}{\pgfqpoint{0.000000in}{-0.027778in}}{\pgfqpoint{0.000000in}{0.000000in}}{%
\pgfpathmoveto{\pgfqpoint{0.000000in}{0.000000in}}%
\pgfpathlineto{\pgfqpoint{0.000000in}{-0.027778in}}%
\pgfusepath{stroke,fill}%
}%
\begin{pgfscope}%
\pgfsys@transformshift{2.811491in}{0.440000in}%
\pgfsys@useobject{currentmarker}{}%
\end{pgfscope}%
\end{pgfscope}%
\begin{pgfscope}%
\pgfpathrectangle{\pgfqpoint{1.000000in}{0.440000in}}{\pgfqpoint{6.200000in}{3.080000in}}%
\pgfusepath{clip}%
\pgfsetbuttcap%
\pgfsetroundjoin%
\pgfsetlinewidth{1.003750pt}%
\definecolor{currentstroke}{rgb}{0.690196,0.690196,0.690196}%
\pgfsetstrokecolor{currentstroke}%
\pgfsetstrokeopacity{0.600000}%
\pgfsetdash{{1.000000pt}{1.650000pt}}{0.000000pt}%
\pgfpathmoveto{\pgfqpoint{2.904022in}{0.440000in}}%
\pgfpathlineto{\pgfqpoint{2.904022in}{3.520000in}}%
\pgfusepath{stroke}%
\end{pgfscope}%
\begin{pgfscope}%
\pgfsetbuttcap%
\pgfsetroundjoin%
\definecolor{currentfill}{rgb}{0.000000,0.000000,0.000000}%
\pgfsetfillcolor{currentfill}%
\pgfsetlinewidth{0.602250pt}%
\definecolor{currentstroke}{rgb}{0.000000,0.000000,0.000000}%
\pgfsetstrokecolor{currentstroke}%
\pgfsetdash{}{0pt}%
\pgfsys@defobject{currentmarker}{\pgfqpoint{0.000000in}{-0.027778in}}{\pgfqpoint{0.000000in}{0.000000in}}{%
\pgfpathmoveto{\pgfqpoint{0.000000in}{0.000000in}}%
\pgfpathlineto{\pgfqpoint{0.000000in}{-0.027778in}}%
\pgfusepath{stroke,fill}%
}%
\begin{pgfscope}%
\pgfsys@transformshift{2.904022in}{0.440000in}%
\pgfsys@useobject{currentmarker}{}%
\end{pgfscope}%
\end{pgfscope}%
\begin{pgfscope}%
\pgfpathrectangle{\pgfqpoint{1.000000in}{0.440000in}}{\pgfqpoint{6.200000in}{3.080000in}}%
\pgfusepath{clip}%
\pgfsetbuttcap%
\pgfsetroundjoin%
\pgfsetlinewidth{1.003750pt}%
\definecolor{currentstroke}{rgb}{0.690196,0.690196,0.690196}%
\pgfsetstrokecolor{currentstroke}%
\pgfsetstrokeopacity{0.600000}%
\pgfsetdash{{1.000000pt}{1.650000pt}}{0.000000pt}%
\pgfpathmoveto{\pgfqpoint{2.979625in}{0.440000in}}%
\pgfpathlineto{\pgfqpoint{2.979625in}{3.520000in}}%
\pgfusepath{stroke}%
\end{pgfscope}%
\begin{pgfscope}%
\pgfsetbuttcap%
\pgfsetroundjoin%
\definecolor{currentfill}{rgb}{0.000000,0.000000,0.000000}%
\pgfsetfillcolor{currentfill}%
\pgfsetlinewidth{0.602250pt}%
\definecolor{currentstroke}{rgb}{0.000000,0.000000,0.000000}%
\pgfsetstrokecolor{currentstroke}%
\pgfsetdash{}{0pt}%
\pgfsys@defobject{currentmarker}{\pgfqpoint{0.000000in}{-0.027778in}}{\pgfqpoint{0.000000in}{0.000000in}}{%
\pgfpathmoveto{\pgfqpoint{0.000000in}{0.000000in}}%
\pgfpathlineto{\pgfqpoint{0.000000in}{-0.027778in}}%
\pgfusepath{stroke,fill}%
}%
\begin{pgfscope}%
\pgfsys@transformshift{2.979625in}{0.440000in}%
\pgfsys@useobject{currentmarker}{}%
\end{pgfscope}%
\end{pgfscope}%
\begin{pgfscope}%
\pgfpathrectangle{\pgfqpoint{1.000000in}{0.440000in}}{\pgfqpoint{6.200000in}{3.080000in}}%
\pgfusepath{clip}%
\pgfsetbuttcap%
\pgfsetroundjoin%
\pgfsetlinewidth{1.003750pt}%
\definecolor{currentstroke}{rgb}{0.690196,0.690196,0.690196}%
\pgfsetstrokecolor{currentstroke}%
\pgfsetstrokeopacity{0.600000}%
\pgfsetdash{{1.000000pt}{1.650000pt}}{0.000000pt}%
\pgfpathmoveto{\pgfqpoint{3.043547in}{0.440000in}}%
\pgfpathlineto{\pgfqpoint{3.043547in}{3.520000in}}%
\pgfusepath{stroke}%
\end{pgfscope}%
\begin{pgfscope}%
\pgfsetbuttcap%
\pgfsetroundjoin%
\definecolor{currentfill}{rgb}{0.000000,0.000000,0.000000}%
\pgfsetfillcolor{currentfill}%
\pgfsetlinewidth{0.602250pt}%
\definecolor{currentstroke}{rgb}{0.000000,0.000000,0.000000}%
\pgfsetstrokecolor{currentstroke}%
\pgfsetdash{}{0pt}%
\pgfsys@defobject{currentmarker}{\pgfqpoint{0.000000in}{-0.027778in}}{\pgfqpoint{0.000000in}{0.000000in}}{%
\pgfpathmoveto{\pgfqpoint{0.000000in}{0.000000in}}%
\pgfpathlineto{\pgfqpoint{0.000000in}{-0.027778in}}%
\pgfusepath{stroke,fill}%
}%
\begin{pgfscope}%
\pgfsys@transformshift{3.043547in}{0.440000in}%
\pgfsys@useobject{currentmarker}{}%
\end{pgfscope}%
\end{pgfscope}%
\begin{pgfscope}%
\pgfpathrectangle{\pgfqpoint{1.000000in}{0.440000in}}{\pgfqpoint{6.200000in}{3.080000in}}%
\pgfusepath{clip}%
\pgfsetbuttcap%
\pgfsetroundjoin%
\pgfsetlinewidth{1.003750pt}%
\definecolor{currentstroke}{rgb}{0.690196,0.690196,0.690196}%
\pgfsetstrokecolor{currentstroke}%
\pgfsetstrokeopacity{0.600000}%
\pgfsetdash{{1.000000pt}{1.650000pt}}{0.000000pt}%
\pgfpathmoveto{\pgfqpoint{3.098919in}{0.440000in}}%
\pgfpathlineto{\pgfqpoint{3.098919in}{3.520000in}}%
\pgfusepath{stroke}%
\end{pgfscope}%
\begin{pgfscope}%
\pgfsetbuttcap%
\pgfsetroundjoin%
\definecolor{currentfill}{rgb}{0.000000,0.000000,0.000000}%
\pgfsetfillcolor{currentfill}%
\pgfsetlinewidth{0.602250pt}%
\definecolor{currentstroke}{rgb}{0.000000,0.000000,0.000000}%
\pgfsetstrokecolor{currentstroke}%
\pgfsetdash{}{0pt}%
\pgfsys@defobject{currentmarker}{\pgfqpoint{0.000000in}{-0.027778in}}{\pgfqpoint{0.000000in}{0.000000in}}{%
\pgfpathmoveto{\pgfqpoint{0.000000in}{0.000000in}}%
\pgfpathlineto{\pgfqpoint{0.000000in}{-0.027778in}}%
\pgfusepath{stroke,fill}%
}%
\begin{pgfscope}%
\pgfsys@transformshift{3.098919in}{0.440000in}%
\pgfsys@useobject{currentmarker}{}%
\end{pgfscope}%
\end{pgfscope}%
\begin{pgfscope}%
\pgfpathrectangle{\pgfqpoint{1.000000in}{0.440000in}}{\pgfqpoint{6.200000in}{3.080000in}}%
\pgfusepath{clip}%
\pgfsetbuttcap%
\pgfsetroundjoin%
\pgfsetlinewidth{1.003750pt}%
\definecolor{currentstroke}{rgb}{0.690196,0.690196,0.690196}%
\pgfsetstrokecolor{currentstroke}%
\pgfsetstrokeopacity{0.600000}%
\pgfsetdash{{1.000000pt}{1.650000pt}}{0.000000pt}%
\pgfpathmoveto{\pgfqpoint{3.147760in}{0.440000in}}%
\pgfpathlineto{\pgfqpoint{3.147760in}{3.520000in}}%
\pgfusepath{stroke}%
\end{pgfscope}%
\begin{pgfscope}%
\pgfsetbuttcap%
\pgfsetroundjoin%
\definecolor{currentfill}{rgb}{0.000000,0.000000,0.000000}%
\pgfsetfillcolor{currentfill}%
\pgfsetlinewidth{0.602250pt}%
\definecolor{currentstroke}{rgb}{0.000000,0.000000,0.000000}%
\pgfsetstrokecolor{currentstroke}%
\pgfsetdash{}{0pt}%
\pgfsys@defobject{currentmarker}{\pgfqpoint{0.000000in}{-0.027778in}}{\pgfqpoint{0.000000in}{0.000000in}}{%
\pgfpathmoveto{\pgfqpoint{0.000000in}{0.000000in}}%
\pgfpathlineto{\pgfqpoint{0.000000in}{-0.027778in}}%
\pgfusepath{stroke,fill}%
}%
\begin{pgfscope}%
\pgfsys@transformshift{3.147760in}{0.440000in}%
\pgfsys@useobject{currentmarker}{}%
\end{pgfscope}%
\end{pgfscope}%
\begin{pgfscope}%
\pgfpathrectangle{\pgfqpoint{1.000000in}{0.440000in}}{\pgfqpoint{6.200000in}{3.080000in}}%
\pgfusepath{clip}%
\pgfsetbuttcap%
\pgfsetroundjoin%
\pgfsetlinewidth{1.003750pt}%
\definecolor{currentstroke}{rgb}{0.690196,0.690196,0.690196}%
\pgfsetstrokecolor{currentstroke}%
\pgfsetstrokeopacity{0.600000}%
\pgfsetdash{{1.000000pt}{1.650000pt}}{0.000000pt}%
\pgfpathmoveto{\pgfqpoint{3.478878in}{0.440000in}}%
\pgfpathlineto{\pgfqpoint{3.478878in}{3.520000in}}%
\pgfusepath{stroke}%
\end{pgfscope}%
\begin{pgfscope}%
\pgfsetbuttcap%
\pgfsetroundjoin%
\definecolor{currentfill}{rgb}{0.000000,0.000000,0.000000}%
\pgfsetfillcolor{currentfill}%
\pgfsetlinewidth{0.602250pt}%
\definecolor{currentstroke}{rgb}{0.000000,0.000000,0.000000}%
\pgfsetstrokecolor{currentstroke}%
\pgfsetdash{}{0pt}%
\pgfsys@defobject{currentmarker}{\pgfqpoint{0.000000in}{-0.027778in}}{\pgfqpoint{0.000000in}{0.000000in}}{%
\pgfpathmoveto{\pgfqpoint{0.000000in}{0.000000in}}%
\pgfpathlineto{\pgfqpoint{0.000000in}{-0.027778in}}%
\pgfusepath{stroke,fill}%
}%
\begin{pgfscope}%
\pgfsys@transformshift{3.478878in}{0.440000in}%
\pgfsys@useobject{currentmarker}{}%
\end{pgfscope}%
\end{pgfscope}%
\begin{pgfscope}%
\pgfpathrectangle{\pgfqpoint{1.000000in}{0.440000in}}{\pgfqpoint{6.200000in}{3.080000in}}%
\pgfusepath{clip}%
\pgfsetbuttcap%
\pgfsetroundjoin%
\pgfsetlinewidth{1.003750pt}%
\definecolor{currentstroke}{rgb}{0.690196,0.690196,0.690196}%
\pgfsetstrokecolor{currentstroke}%
\pgfsetstrokeopacity{0.600000}%
\pgfsetdash{{1.000000pt}{1.650000pt}}{0.000000pt}%
\pgfpathmoveto{\pgfqpoint{3.647013in}{0.440000in}}%
\pgfpathlineto{\pgfqpoint{3.647013in}{3.520000in}}%
\pgfusepath{stroke}%
\end{pgfscope}%
\begin{pgfscope}%
\pgfsetbuttcap%
\pgfsetroundjoin%
\definecolor{currentfill}{rgb}{0.000000,0.000000,0.000000}%
\pgfsetfillcolor{currentfill}%
\pgfsetlinewidth{0.602250pt}%
\definecolor{currentstroke}{rgb}{0.000000,0.000000,0.000000}%
\pgfsetstrokecolor{currentstroke}%
\pgfsetdash{}{0pt}%
\pgfsys@defobject{currentmarker}{\pgfqpoint{0.000000in}{-0.027778in}}{\pgfqpoint{0.000000in}{0.000000in}}{%
\pgfpathmoveto{\pgfqpoint{0.000000in}{0.000000in}}%
\pgfpathlineto{\pgfqpoint{0.000000in}{-0.027778in}}%
\pgfusepath{stroke,fill}%
}%
\begin{pgfscope}%
\pgfsys@transformshift{3.647013in}{0.440000in}%
\pgfsys@useobject{currentmarker}{}%
\end{pgfscope}%
\end{pgfscope}%
\begin{pgfscope}%
\pgfpathrectangle{\pgfqpoint{1.000000in}{0.440000in}}{\pgfqpoint{6.200000in}{3.080000in}}%
\pgfusepath{clip}%
\pgfsetbuttcap%
\pgfsetroundjoin%
\pgfsetlinewidth{1.003750pt}%
\definecolor{currentstroke}{rgb}{0.690196,0.690196,0.690196}%
\pgfsetstrokecolor{currentstroke}%
\pgfsetstrokeopacity{0.600000}%
\pgfsetdash{{1.000000pt}{1.650000pt}}{0.000000pt}%
\pgfpathmoveto{\pgfqpoint{3.766306in}{0.440000in}}%
\pgfpathlineto{\pgfqpoint{3.766306in}{3.520000in}}%
\pgfusepath{stroke}%
\end{pgfscope}%
\begin{pgfscope}%
\pgfsetbuttcap%
\pgfsetroundjoin%
\definecolor{currentfill}{rgb}{0.000000,0.000000,0.000000}%
\pgfsetfillcolor{currentfill}%
\pgfsetlinewidth{0.602250pt}%
\definecolor{currentstroke}{rgb}{0.000000,0.000000,0.000000}%
\pgfsetstrokecolor{currentstroke}%
\pgfsetdash{}{0pt}%
\pgfsys@defobject{currentmarker}{\pgfqpoint{0.000000in}{-0.027778in}}{\pgfqpoint{0.000000in}{0.000000in}}{%
\pgfpathmoveto{\pgfqpoint{0.000000in}{0.000000in}}%
\pgfpathlineto{\pgfqpoint{0.000000in}{-0.027778in}}%
\pgfusepath{stroke,fill}%
}%
\begin{pgfscope}%
\pgfsys@transformshift{3.766306in}{0.440000in}%
\pgfsys@useobject{currentmarker}{}%
\end{pgfscope}%
\end{pgfscope}%
\begin{pgfscope}%
\pgfpathrectangle{\pgfqpoint{1.000000in}{0.440000in}}{\pgfqpoint{6.200000in}{3.080000in}}%
\pgfusepath{clip}%
\pgfsetbuttcap%
\pgfsetroundjoin%
\pgfsetlinewidth{1.003750pt}%
\definecolor{currentstroke}{rgb}{0.690196,0.690196,0.690196}%
\pgfsetstrokecolor{currentstroke}%
\pgfsetstrokeopacity{0.600000}%
\pgfsetdash{{1.000000pt}{1.650000pt}}{0.000000pt}%
\pgfpathmoveto{\pgfqpoint{3.858838in}{0.440000in}}%
\pgfpathlineto{\pgfqpoint{3.858838in}{3.520000in}}%
\pgfusepath{stroke}%
\end{pgfscope}%
\begin{pgfscope}%
\pgfsetbuttcap%
\pgfsetroundjoin%
\definecolor{currentfill}{rgb}{0.000000,0.000000,0.000000}%
\pgfsetfillcolor{currentfill}%
\pgfsetlinewidth{0.602250pt}%
\definecolor{currentstroke}{rgb}{0.000000,0.000000,0.000000}%
\pgfsetstrokecolor{currentstroke}%
\pgfsetdash{}{0pt}%
\pgfsys@defobject{currentmarker}{\pgfqpoint{0.000000in}{-0.027778in}}{\pgfqpoint{0.000000in}{0.000000in}}{%
\pgfpathmoveto{\pgfqpoint{0.000000in}{0.000000in}}%
\pgfpathlineto{\pgfqpoint{0.000000in}{-0.027778in}}%
\pgfusepath{stroke,fill}%
}%
\begin{pgfscope}%
\pgfsys@transformshift{3.858838in}{0.440000in}%
\pgfsys@useobject{currentmarker}{}%
\end{pgfscope}%
\end{pgfscope}%
\begin{pgfscope}%
\pgfpathrectangle{\pgfqpoint{1.000000in}{0.440000in}}{\pgfqpoint{6.200000in}{3.080000in}}%
\pgfusepath{clip}%
\pgfsetbuttcap%
\pgfsetroundjoin%
\pgfsetlinewidth{1.003750pt}%
\definecolor{currentstroke}{rgb}{0.690196,0.690196,0.690196}%
\pgfsetstrokecolor{currentstroke}%
\pgfsetstrokeopacity{0.600000}%
\pgfsetdash{{1.000000pt}{1.650000pt}}{0.000000pt}%
\pgfpathmoveto{\pgfqpoint{3.934441in}{0.440000in}}%
\pgfpathlineto{\pgfqpoint{3.934441in}{3.520000in}}%
\pgfusepath{stroke}%
\end{pgfscope}%
\begin{pgfscope}%
\pgfsetbuttcap%
\pgfsetroundjoin%
\definecolor{currentfill}{rgb}{0.000000,0.000000,0.000000}%
\pgfsetfillcolor{currentfill}%
\pgfsetlinewidth{0.602250pt}%
\definecolor{currentstroke}{rgb}{0.000000,0.000000,0.000000}%
\pgfsetstrokecolor{currentstroke}%
\pgfsetdash{}{0pt}%
\pgfsys@defobject{currentmarker}{\pgfqpoint{0.000000in}{-0.027778in}}{\pgfqpoint{0.000000in}{0.000000in}}{%
\pgfpathmoveto{\pgfqpoint{0.000000in}{0.000000in}}%
\pgfpathlineto{\pgfqpoint{0.000000in}{-0.027778in}}%
\pgfusepath{stroke,fill}%
}%
\begin{pgfscope}%
\pgfsys@transformshift{3.934441in}{0.440000in}%
\pgfsys@useobject{currentmarker}{}%
\end{pgfscope}%
\end{pgfscope}%
\begin{pgfscope}%
\pgfpathrectangle{\pgfqpoint{1.000000in}{0.440000in}}{\pgfqpoint{6.200000in}{3.080000in}}%
\pgfusepath{clip}%
\pgfsetbuttcap%
\pgfsetroundjoin%
\pgfsetlinewidth{1.003750pt}%
\definecolor{currentstroke}{rgb}{0.690196,0.690196,0.690196}%
\pgfsetstrokecolor{currentstroke}%
\pgfsetstrokeopacity{0.600000}%
\pgfsetdash{{1.000000pt}{1.650000pt}}{0.000000pt}%
\pgfpathmoveto{\pgfqpoint{3.998363in}{0.440000in}}%
\pgfpathlineto{\pgfqpoint{3.998363in}{3.520000in}}%
\pgfusepath{stroke}%
\end{pgfscope}%
\begin{pgfscope}%
\pgfsetbuttcap%
\pgfsetroundjoin%
\definecolor{currentfill}{rgb}{0.000000,0.000000,0.000000}%
\pgfsetfillcolor{currentfill}%
\pgfsetlinewidth{0.602250pt}%
\definecolor{currentstroke}{rgb}{0.000000,0.000000,0.000000}%
\pgfsetstrokecolor{currentstroke}%
\pgfsetdash{}{0pt}%
\pgfsys@defobject{currentmarker}{\pgfqpoint{0.000000in}{-0.027778in}}{\pgfqpoint{0.000000in}{0.000000in}}{%
\pgfpathmoveto{\pgfqpoint{0.000000in}{0.000000in}}%
\pgfpathlineto{\pgfqpoint{0.000000in}{-0.027778in}}%
\pgfusepath{stroke,fill}%
}%
\begin{pgfscope}%
\pgfsys@transformshift{3.998363in}{0.440000in}%
\pgfsys@useobject{currentmarker}{}%
\end{pgfscope}%
\end{pgfscope}%
\begin{pgfscope}%
\pgfpathrectangle{\pgfqpoint{1.000000in}{0.440000in}}{\pgfqpoint{6.200000in}{3.080000in}}%
\pgfusepath{clip}%
\pgfsetbuttcap%
\pgfsetroundjoin%
\pgfsetlinewidth{1.003750pt}%
\definecolor{currentstroke}{rgb}{0.690196,0.690196,0.690196}%
\pgfsetstrokecolor{currentstroke}%
\pgfsetstrokeopacity{0.600000}%
\pgfsetdash{{1.000000pt}{1.650000pt}}{0.000000pt}%
\pgfpathmoveto{\pgfqpoint{4.053735in}{0.440000in}}%
\pgfpathlineto{\pgfqpoint{4.053735in}{3.520000in}}%
\pgfusepath{stroke}%
\end{pgfscope}%
\begin{pgfscope}%
\pgfsetbuttcap%
\pgfsetroundjoin%
\definecolor{currentfill}{rgb}{0.000000,0.000000,0.000000}%
\pgfsetfillcolor{currentfill}%
\pgfsetlinewidth{0.602250pt}%
\definecolor{currentstroke}{rgb}{0.000000,0.000000,0.000000}%
\pgfsetstrokecolor{currentstroke}%
\pgfsetdash{}{0pt}%
\pgfsys@defobject{currentmarker}{\pgfqpoint{0.000000in}{-0.027778in}}{\pgfqpoint{0.000000in}{0.000000in}}{%
\pgfpathmoveto{\pgfqpoint{0.000000in}{0.000000in}}%
\pgfpathlineto{\pgfqpoint{0.000000in}{-0.027778in}}%
\pgfusepath{stroke,fill}%
}%
\begin{pgfscope}%
\pgfsys@transformshift{4.053735in}{0.440000in}%
\pgfsys@useobject{currentmarker}{}%
\end{pgfscope}%
\end{pgfscope}%
\begin{pgfscope}%
\pgfpathrectangle{\pgfqpoint{1.000000in}{0.440000in}}{\pgfqpoint{6.200000in}{3.080000in}}%
\pgfusepath{clip}%
\pgfsetbuttcap%
\pgfsetroundjoin%
\pgfsetlinewidth{1.003750pt}%
\definecolor{currentstroke}{rgb}{0.690196,0.690196,0.690196}%
\pgfsetstrokecolor{currentstroke}%
\pgfsetstrokeopacity{0.600000}%
\pgfsetdash{{1.000000pt}{1.650000pt}}{0.000000pt}%
\pgfpathmoveto{\pgfqpoint{4.102576in}{0.440000in}}%
\pgfpathlineto{\pgfqpoint{4.102576in}{3.520000in}}%
\pgfusepath{stroke}%
\end{pgfscope}%
\begin{pgfscope}%
\pgfsetbuttcap%
\pgfsetroundjoin%
\definecolor{currentfill}{rgb}{0.000000,0.000000,0.000000}%
\pgfsetfillcolor{currentfill}%
\pgfsetlinewidth{0.602250pt}%
\definecolor{currentstroke}{rgb}{0.000000,0.000000,0.000000}%
\pgfsetstrokecolor{currentstroke}%
\pgfsetdash{}{0pt}%
\pgfsys@defobject{currentmarker}{\pgfqpoint{0.000000in}{-0.027778in}}{\pgfqpoint{0.000000in}{0.000000in}}{%
\pgfpathmoveto{\pgfqpoint{0.000000in}{0.000000in}}%
\pgfpathlineto{\pgfqpoint{0.000000in}{-0.027778in}}%
\pgfusepath{stroke,fill}%
}%
\begin{pgfscope}%
\pgfsys@transformshift{4.102576in}{0.440000in}%
\pgfsys@useobject{currentmarker}{}%
\end{pgfscope}%
\end{pgfscope}%
\begin{pgfscope}%
\pgfpathrectangle{\pgfqpoint{1.000000in}{0.440000in}}{\pgfqpoint{6.200000in}{3.080000in}}%
\pgfusepath{clip}%
\pgfsetbuttcap%
\pgfsetroundjoin%
\pgfsetlinewidth{1.003750pt}%
\definecolor{currentstroke}{rgb}{0.690196,0.690196,0.690196}%
\pgfsetstrokecolor{currentstroke}%
\pgfsetstrokeopacity{0.600000}%
\pgfsetdash{{1.000000pt}{1.650000pt}}{0.000000pt}%
\pgfpathmoveto{\pgfqpoint{4.433694in}{0.440000in}}%
\pgfpathlineto{\pgfqpoint{4.433694in}{3.520000in}}%
\pgfusepath{stroke}%
\end{pgfscope}%
\begin{pgfscope}%
\pgfsetbuttcap%
\pgfsetroundjoin%
\definecolor{currentfill}{rgb}{0.000000,0.000000,0.000000}%
\pgfsetfillcolor{currentfill}%
\pgfsetlinewidth{0.602250pt}%
\definecolor{currentstroke}{rgb}{0.000000,0.000000,0.000000}%
\pgfsetstrokecolor{currentstroke}%
\pgfsetdash{}{0pt}%
\pgfsys@defobject{currentmarker}{\pgfqpoint{0.000000in}{-0.027778in}}{\pgfqpoint{0.000000in}{0.000000in}}{%
\pgfpathmoveto{\pgfqpoint{0.000000in}{0.000000in}}%
\pgfpathlineto{\pgfqpoint{0.000000in}{-0.027778in}}%
\pgfusepath{stroke,fill}%
}%
\begin{pgfscope}%
\pgfsys@transformshift{4.433694in}{0.440000in}%
\pgfsys@useobject{currentmarker}{}%
\end{pgfscope}%
\end{pgfscope}%
\begin{pgfscope}%
\pgfpathrectangle{\pgfqpoint{1.000000in}{0.440000in}}{\pgfqpoint{6.200000in}{3.080000in}}%
\pgfusepath{clip}%
\pgfsetbuttcap%
\pgfsetroundjoin%
\pgfsetlinewidth{1.003750pt}%
\definecolor{currentstroke}{rgb}{0.690196,0.690196,0.690196}%
\pgfsetstrokecolor{currentstroke}%
\pgfsetstrokeopacity{0.600000}%
\pgfsetdash{{1.000000pt}{1.650000pt}}{0.000000pt}%
\pgfpathmoveto{\pgfqpoint{4.601829in}{0.440000in}}%
\pgfpathlineto{\pgfqpoint{4.601829in}{3.520000in}}%
\pgfusepath{stroke}%
\end{pgfscope}%
\begin{pgfscope}%
\pgfsetbuttcap%
\pgfsetroundjoin%
\definecolor{currentfill}{rgb}{0.000000,0.000000,0.000000}%
\pgfsetfillcolor{currentfill}%
\pgfsetlinewidth{0.602250pt}%
\definecolor{currentstroke}{rgb}{0.000000,0.000000,0.000000}%
\pgfsetstrokecolor{currentstroke}%
\pgfsetdash{}{0pt}%
\pgfsys@defobject{currentmarker}{\pgfqpoint{0.000000in}{-0.027778in}}{\pgfqpoint{0.000000in}{0.000000in}}{%
\pgfpathmoveto{\pgfqpoint{0.000000in}{0.000000in}}%
\pgfpathlineto{\pgfqpoint{0.000000in}{-0.027778in}}%
\pgfusepath{stroke,fill}%
}%
\begin{pgfscope}%
\pgfsys@transformshift{4.601829in}{0.440000in}%
\pgfsys@useobject{currentmarker}{}%
\end{pgfscope}%
\end{pgfscope}%
\begin{pgfscope}%
\pgfpathrectangle{\pgfqpoint{1.000000in}{0.440000in}}{\pgfqpoint{6.200000in}{3.080000in}}%
\pgfusepath{clip}%
\pgfsetbuttcap%
\pgfsetroundjoin%
\pgfsetlinewidth{1.003750pt}%
\definecolor{currentstroke}{rgb}{0.690196,0.690196,0.690196}%
\pgfsetstrokecolor{currentstroke}%
\pgfsetstrokeopacity{0.600000}%
\pgfsetdash{{1.000000pt}{1.650000pt}}{0.000000pt}%
\pgfpathmoveto{\pgfqpoint{4.721122in}{0.440000in}}%
\pgfpathlineto{\pgfqpoint{4.721122in}{3.520000in}}%
\pgfusepath{stroke}%
\end{pgfscope}%
\begin{pgfscope}%
\pgfsetbuttcap%
\pgfsetroundjoin%
\definecolor{currentfill}{rgb}{0.000000,0.000000,0.000000}%
\pgfsetfillcolor{currentfill}%
\pgfsetlinewidth{0.602250pt}%
\definecolor{currentstroke}{rgb}{0.000000,0.000000,0.000000}%
\pgfsetstrokecolor{currentstroke}%
\pgfsetdash{}{0pt}%
\pgfsys@defobject{currentmarker}{\pgfqpoint{0.000000in}{-0.027778in}}{\pgfqpoint{0.000000in}{0.000000in}}{%
\pgfpathmoveto{\pgfqpoint{0.000000in}{0.000000in}}%
\pgfpathlineto{\pgfqpoint{0.000000in}{-0.027778in}}%
\pgfusepath{stroke,fill}%
}%
\begin{pgfscope}%
\pgfsys@transformshift{4.721122in}{0.440000in}%
\pgfsys@useobject{currentmarker}{}%
\end{pgfscope}%
\end{pgfscope}%
\begin{pgfscope}%
\pgfpathrectangle{\pgfqpoint{1.000000in}{0.440000in}}{\pgfqpoint{6.200000in}{3.080000in}}%
\pgfusepath{clip}%
\pgfsetbuttcap%
\pgfsetroundjoin%
\pgfsetlinewidth{1.003750pt}%
\definecolor{currentstroke}{rgb}{0.690196,0.690196,0.690196}%
\pgfsetstrokecolor{currentstroke}%
\pgfsetstrokeopacity{0.600000}%
\pgfsetdash{{1.000000pt}{1.650000pt}}{0.000000pt}%
\pgfpathmoveto{\pgfqpoint{4.813654in}{0.440000in}}%
\pgfpathlineto{\pgfqpoint{4.813654in}{3.520000in}}%
\pgfusepath{stroke}%
\end{pgfscope}%
\begin{pgfscope}%
\pgfsetbuttcap%
\pgfsetroundjoin%
\definecolor{currentfill}{rgb}{0.000000,0.000000,0.000000}%
\pgfsetfillcolor{currentfill}%
\pgfsetlinewidth{0.602250pt}%
\definecolor{currentstroke}{rgb}{0.000000,0.000000,0.000000}%
\pgfsetstrokecolor{currentstroke}%
\pgfsetdash{}{0pt}%
\pgfsys@defobject{currentmarker}{\pgfqpoint{0.000000in}{-0.027778in}}{\pgfqpoint{0.000000in}{0.000000in}}{%
\pgfpathmoveto{\pgfqpoint{0.000000in}{0.000000in}}%
\pgfpathlineto{\pgfqpoint{0.000000in}{-0.027778in}}%
\pgfusepath{stroke,fill}%
}%
\begin{pgfscope}%
\pgfsys@transformshift{4.813654in}{0.440000in}%
\pgfsys@useobject{currentmarker}{}%
\end{pgfscope}%
\end{pgfscope}%
\begin{pgfscope}%
\pgfpathrectangle{\pgfqpoint{1.000000in}{0.440000in}}{\pgfqpoint{6.200000in}{3.080000in}}%
\pgfusepath{clip}%
\pgfsetbuttcap%
\pgfsetroundjoin%
\pgfsetlinewidth{1.003750pt}%
\definecolor{currentstroke}{rgb}{0.690196,0.690196,0.690196}%
\pgfsetstrokecolor{currentstroke}%
\pgfsetstrokeopacity{0.600000}%
\pgfsetdash{{1.000000pt}{1.650000pt}}{0.000000pt}%
\pgfpathmoveto{\pgfqpoint{4.889257in}{0.440000in}}%
\pgfpathlineto{\pgfqpoint{4.889257in}{3.520000in}}%
\pgfusepath{stroke}%
\end{pgfscope}%
\begin{pgfscope}%
\pgfsetbuttcap%
\pgfsetroundjoin%
\definecolor{currentfill}{rgb}{0.000000,0.000000,0.000000}%
\pgfsetfillcolor{currentfill}%
\pgfsetlinewidth{0.602250pt}%
\definecolor{currentstroke}{rgb}{0.000000,0.000000,0.000000}%
\pgfsetstrokecolor{currentstroke}%
\pgfsetdash{}{0pt}%
\pgfsys@defobject{currentmarker}{\pgfqpoint{0.000000in}{-0.027778in}}{\pgfqpoint{0.000000in}{0.000000in}}{%
\pgfpathmoveto{\pgfqpoint{0.000000in}{0.000000in}}%
\pgfpathlineto{\pgfqpoint{0.000000in}{-0.027778in}}%
\pgfusepath{stroke,fill}%
}%
\begin{pgfscope}%
\pgfsys@transformshift{4.889257in}{0.440000in}%
\pgfsys@useobject{currentmarker}{}%
\end{pgfscope}%
\end{pgfscope}%
\begin{pgfscope}%
\pgfpathrectangle{\pgfqpoint{1.000000in}{0.440000in}}{\pgfqpoint{6.200000in}{3.080000in}}%
\pgfusepath{clip}%
\pgfsetbuttcap%
\pgfsetroundjoin%
\pgfsetlinewidth{1.003750pt}%
\definecolor{currentstroke}{rgb}{0.690196,0.690196,0.690196}%
\pgfsetstrokecolor{currentstroke}%
\pgfsetstrokeopacity{0.600000}%
\pgfsetdash{{1.000000pt}{1.650000pt}}{0.000000pt}%
\pgfpathmoveto{\pgfqpoint{4.953179in}{0.440000in}}%
\pgfpathlineto{\pgfqpoint{4.953179in}{3.520000in}}%
\pgfusepath{stroke}%
\end{pgfscope}%
\begin{pgfscope}%
\pgfsetbuttcap%
\pgfsetroundjoin%
\definecolor{currentfill}{rgb}{0.000000,0.000000,0.000000}%
\pgfsetfillcolor{currentfill}%
\pgfsetlinewidth{0.602250pt}%
\definecolor{currentstroke}{rgb}{0.000000,0.000000,0.000000}%
\pgfsetstrokecolor{currentstroke}%
\pgfsetdash{}{0pt}%
\pgfsys@defobject{currentmarker}{\pgfqpoint{0.000000in}{-0.027778in}}{\pgfqpoint{0.000000in}{0.000000in}}{%
\pgfpathmoveto{\pgfqpoint{0.000000in}{0.000000in}}%
\pgfpathlineto{\pgfqpoint{0.000000in}{-0.027778in}}%
\pgfusepath{stroke,fill}%
}%
\begin{pgfscope}%
\pgfsys@transformshift{4.953179in}{0.440000in}%
\pgfsys@useobject{currentmarker}{}%
\end{pgfscope}%
\end{pgfscope}%
\begin{pgfscope}%
\pgfpathrectangle{\pgfqpoint{1.000000in}{0.440000in}}{\pgfqpoint{6.200000in}{3.080000in}}%
\pgfusepath{clip}%
\pgfsetbuttcap%
\pgfsetroundjoin%
\pgfsetlinewidth{1.003750pt}%
\definecolor{currentstroke}{rgb}{0.690196,0.690196,0.690196}%
\pgfsetstrokecolor{currentstroke}%
\pgfsetstrokeopacity{0.600000}%
\pgfsetdash{{1.000000pt}{1.650000pt}}{0.000000pt}%
\pgfpathmoveto{\pgfqpoint{5.008551in}{0.440000in}}%
\pgfpathlineto{\pgfqpoint{5.008551in}{3.520000in}}%
\pgfusepath{stroke}%
\end{pgfscope}%
\begin{pgfscope}%
\pgfsetbuttcap%
\pgfsetroundjoin%
\definecolor{currentfill}{rgb}{0.000000,0.000000,0.000000}%
\pgfsetfillcolor{currentfill}%
\pgfsetlinewidth{0.602250pt}%
\definecolor{currentstroke}{rgb}{0.000000,0.000000,0.000000}%
\pgfsetstrokecolor{currentstroke}%
\pgfsetdash{}{0pt}%
\pgfsys@defobject{currentmarker}{\pgfqpoint{0.000000in}{-0.027778in}}{\pgfqpoint{0.000000in}{0.000000in}}{%
\pgfpathmoveto{\pgfqpoint{0.000000in}{0.000000in}}%
\pgfpathlineto{\pgfqpoint{0.000000in}{-0.027778in}}%
\pgfusepath{stroke,fill}%
}%
\begin{pgfscope}%
\pgfsys@transformshift{5.008551in}{0.440000in}%
\pgfsys@useobject{currentmarker}{}%
\end{pgfscope}%
\end{pgfscope}%
\begin{pgfscope}%
\pgfpathrectangle{\pgfqpoint{1.000000in}{0.440000in}}{\pgfqpoint{6.200000in}{3.080000in}}%
\pgfusepath{clip}%
\pgfsetbuttcap%
\pgfsetroundjoin%
\pgfsetlinewidth{1.003750pt}%
\definecolor{currentstroke}{rgb}{0.690196,0.690196,0.690196}%
\pgfsetstrokecolor{currentstroke}%
\pgfsetstrokeopacity{0.600000}%
\pgfsetdash{{1.000000pt}{1.650000pt}}{0.000000pt}%
\pgfpathmoveto{\pgfqpoint{5.057392in}{0.440000in}}%
\pgfpathlineto{\pgfqpoint{5.057392in}{3.520000in}}%
\pgfusepath{stroke}%
\end{pgfscope}%
\begin{pgfscope}%
\pgfsetbuttcap%
\pgfsetroundjoin%
\definecolor{currentfill}{rgb}{0.000000,0.000000,0.000000}%
\pgfsetfillcolor{currentfill}%
\pgfsetlinewidth{0.602250pt}%
\definecolor{currentstroke}{rgb}{0.000000,0.000000,0.000000}%
\pgfsetstrokecolor{currentstroke}%
\pgfsetdash{}{0pt}%
\pgfsys@defobject{currentmarker}{\pgfqpoint{0.000000in}{-0.027778in}}{\pgfqpoint{0.000000in}{0.000000in}}{%
\pgfpathmoveto{\pgfqpoint{0.000000in}{0.000000in}}%
\pgfpathlineto{\pgfqpoint{0.000000in}{-0.027778in}}%
\pgfusepath{stroke,fill}%
}%
\begin{pgfscope}%
\pgfsys@transformshift{5.057392in}{0.440000in}%
\pgfsys@useobject{currentmarker}{}%
\end{pgfscope}%
\end{pgfscope}%
\begin{pgfscope}%
\pgfpathrectangle{\pgfqpoint{1.000000in}{0.440000in}}{\pgfqpoint{6.200000in}{3.080000in}}%
\pgfusepath{clip}%
\pgfsetbuttcap%
\pgfsetroundjoin%
\pgfsetlinewidth{1.003750pt}%
\definecolor{currentstroke}{rgb}{0.690196,0.690196,0.690196}%
\pgfsetstrokecolor{currentstroke}%
\pgfsetstrokeopacity{0.600000}%
\pgfsetdash{{1.000000pt}{1.650000pt}}{0.000000pt}%
\pgfpathmoveto{\pgfqpoint{5.388510in}{0.440000in}}%
\pgfpathlineto{\pgfqpoint{5.388510in}{3.520000in}}%
\pgfusepath{stroke}%
\end{pgfscope}%
\begin{pgfscope}%
\pgfsetbuttcap%
\pgfsetroundjoin%
\definecolor{currentfill}{rgb}{0.000000,0.000000,0.000000}%
\pgfsetfillcolor{currentfill}%
\pgfsetlinewidth{0.602250pt}%
\definecolor{currentstroke}{rgb}{0.000000,0.000000,0.000000}%
\pgfsetstrokecolor{currentstroke}%
\pgfsetdash{}{0pt}%
\pgfsys@defobject{currentmarker}{\pgfqpoint{0.000000in}{-0.027778in}}{\pgfqpoint{0.000000in}{0.000000in}}{%
\pgfpathmoveto{\pgfqpoint{0.000000in}{0.000000in}}%
\pgfpathlineto{\pgfqpoint{0.000000in}{-0.027778in}}%
\pgfusepath{stroke,fill}%
}%
\begin{pgfscope}%
\pgfsys@transformshift{5.388510in}{0.440000in}%
\pgfsys@useobject{currentmarker}{}%
\end{pgfscope}%
\end{pgfscope}%
\begin{pgfscope}%
\pgfpathrectangle{\pgfqpoint{1.000000in}{0.440000in}}{\pgfqpoint{6.200000in}{3.080000in}}%
\pgfusepath{clip}%
\pgfsetbuttcap%
\pgfsetroundjoin%
\pgfsetlinewidth{1.003750pt}%
\definecolor{currentstroke}{rgb}{0.690196,0.690196,0.690196}%
\pgfsetstrokecolor{currentstroke}%
\pgfsetstrokeopacity{0.600000}%
\pgfsetdash{{1.000000pt}{1.650000pt}}{0.000000pt}%
\pgfpathmoveto{\pgfqpoint{5.556645in}{0.440000in}}%
\pgfpathlineto{\pgfqpoint{5.556645in}{3.520000in}}%
\pgfusepath{stroke}%
\end{pgfscope}%
\begin{pgfscope}%
\pgfsetbuttcap%
\pgfsetroundjoin%
\definecolor{currentfill}{rgb}{0.000000,0.000000,0.000000}%
\pgfsetfillcolor{currentfill}%
\pgfsetlinewidth{0.602250pt}%
\definecolor{currentstroke}{rgb}{0.000000,0.000000,0.000000}%
\pgfsetstrokecolor{currentstroke}%
\pgfsetdash{}{0pt}%
\pgfsys@defobject{currentmarker}{\pgfqpoint{0.000000in}{-0.027778in}}{\pgfqpoint{0.000000in}{0.000000in}}{%
\pgfpathmoveto{\pgfqpoint{0.000000in}{0.000000in}}%
\pgfpathlineto{\pgfqpoint{0.000000in}{-0.027778in}}%
\pgfusepath{stroke,fill}%
}%
\begin{pgfscope}%
\pgfsys@transformshift{5.556645in}{0.440000in}%
\pgfsys@useobject{currentmarker}{}%
\end{pgfscope}%
\end{pgfscope}%
\begin{pgfscope}%
\pgfpathrectangle{\pgfqpoint{1.000000in}{0.440000in}}{\pgfqpoint{6.200000in}{3.080000in}}%
\pgfusepath{clip}%
\pgfsetbuttcap%
\pgfsetroundjoin%
\pgfsetlinewidth{1.003750pt}%
\definecolor{currentstroke}{rgb}{0.690196,0.690196,0.690196}%
\pgfsetstrokecolor{currentstroke}%
\pgfsetstrokeopacity{0.600000}%
\pgfsetdash{{1.000000pt}{1.650000pt}}{0.000000pt}%
\pgfpathmoveto{\pgfqpoint{5.675938in}{0.440000in}}%
\pgfpathlineto{\pgfqpoint{5.675938in}{3.520000in}}%
\pgfusepath{stroke}%
\end{pgfscope}%
\begin{pgfscope}%
\pgfsetbuttcap%
\pgfsetroundjoin%
\definecolor{currentfill}{rgb}{0.000000,0.000000,0.000000}%
\pgfsetfillcolor{currentfill}%
\pgfsetlinewidth{0.602250pt}%
\definecolor{currentstroke}{rgb}{0.000000,0.000000,0.000000}%
\pgfsetstrokecolor{currentstroke}%
\pgfsetdash{}{0pt}%
\pgfsys@defobject{currentmarker}{\pgfqpoint{0.000000in}{-0.027778in}}{\pgfqpoint{0.000000in}{0.000000in}}{%
\pgfpathmoveto{\pgfqpoint{0.000000in}{0.000000in}}%
\pgfpathlineto{\pgfqpoint{0.000000in}{-0.027778in}}%
\pgfusepath{stroke,fill}%
}%
\begin{pgfscope}%
\pgfsys@transformshift{5.675938in}{0.440000in}%
\pgfsys@useobject{currentmarker}{}%
\end{pgfscope}%
\end{pgfscope}%
\begin{pgfscope}%
\pgfpathrectangle{\pgfqpoint{1.000000in}{0.440000in}}{\pgfqpoint{6.200000in}{3.080000in}}%
\pgfusepath{clip}%
\pgfsetbuttcap%
\pgfsetroundjoin%
\pgfsetlinewidth{1.003750pt}%
\definecolor{currentstroke}{rgb}{0.690196,0.690196,0.690196}%
\pgfsetstrokecolor{currentstroke}%
\pgfsetstrokeopacity{0.600000}%
\pgfsetdash{{1.000000pt}{1.650000pt}}{0.000000pt}%
\pgfpathmoveto{\pgfqpoint{5.768469in}{0.440000in}}%
\pgfpathlineto{\pgfqpoint{5.768469in}{3.520000in}}%
\pgfusepath{stroke}%
\end{pgfscope}%
\begin{pgfscope}%
\pgfsetbuttcap%
\pgfsetroundjoin%
\definecolor{currentfill}{rgb}{0.000000,0.000000,0.000000}%
\pgfsetfillcolor{currentfill}%
\pgfsetlinewidth{0.602250pt}%
\definecolor{currentstroke}{rgb}{0.000000,0.000000,0.000000}%
\pgfsetstrokecolor{currentstroke}%
\pgfsetdash{}{0pt}%
\pgfsys@defobject{currentmarker}{\pgfqpoint{0.000000in}{-0.027778in}}{\pgfqpoint{0.000000in}{0.000000in}}{%
\pgfpathmoveto{\pgfqpoint{0.000000in}{0.000000in}}%
\pgfpathlineto{\pgfqpoint{0.000000in}{-0.027778in}}%
\pgfusepath{stroke,fill}%
}%
\begin{pgfscope}%
\pgfsys@transformshift{5.768469in}{0.440000in}%
\pgfsys@useobject{currentmarker}{}%
\end{pgfscope}%
\end{pgfscope}%
\begin{pgfscope}%
\pgfpathrectangle{\pgfqpoint{1.000000in}{0.440000in}}{\pgfqpoint{6.200000in}{3.080000in}}%
\pgfusepath{clip}%
\pgfsetbuttcap%
\pgfsetroundjoin%
\pgfsetlinewidth{1.003750pt}%
\definecolor{currentstroke}{rgb}{0.690196,0.690196,0.690196}%
\pgfsetstrokecolor{currentstroke}%
\pgfsetstrokeopacity{0.600000}%
\pgfsetdash{{1.000000pt}{1.650000pt}}{0.000000pt}%
\pgfpathmoveto{\pgfqpoint{5.844073in}{0.440000in}}%
\pgfpathlineto{\pgfqpoint{5.844073in}{3.520000in}}%
\pgfusepath{stroke}%
\end{pgfscope}%
\begin{pgfscope}%
\pgfsetbuttcap%
\pgfsetroundjoin%
\definecolor{currentfill}{rgb}{0.000000,0.000000,0.000000}%
\pgfsetfillcolor{currentfill}%
\pgfsetlinewidth{0.602250pt}%
\definecolor{currentstroke}{rgb}{0.000000,0.000000,0.000000}%
\pgfsetstrokecolor{currentstroke}%
\pgfsetdash{}{0pt}%
\pgfsys@defobject{currentmarker}{\pgfqpoint{0.000000in}{-0.027778in}}{\pgfqpoint{0.000000in}{0.000000in}}{%
\pgfpathmoveto{\pgfqpoint{0.000000in}{0.000000in}}%
\pgfpathlineto{\pgfqpoint{0.000000in}{-0.027778in}}%
\pgfusepath{stroke,fill}%
}%
\begin{pgfscope}%
\pgfsys@transformshift{5.844073in}{0.440000in}%
\pgfsys@useobject{currentmarker}{}%
\end{pgfscope}%
\end{pgfscope}%
\begin{pgfscope}%
\pgfpathrectangle{\pgfqpoint{1.000000in}{0.440000in}}{\pgfqpoint{6.200000in}{3.080000in}}%
\pgfusepath{clip}%
\pgfsetbuttcap%
\pgfsetroundjoin%
\pgfsetlinewidth{1.003750pt}%
\definecolor{currentstroke}{rgb}{0.690196,0.690196,0.690196}%
\pgfsetstrokecolor{currentstroke}%
\pgfsetstrokeopacity{0.600000}%
\pgfsetdash{{1.000000pt}{1.650000pt}}{0.000000pt}%
\pgfpathmoveto{\pgfqpoint{5.907995in}{0.440000in}}%
\pgfpathlineto{\pgfqpoint{5.907995in}{3.520000in}}%
\pgfusepath{stroke}%
\end{pgfscope}%
\begin{pgfscope}%
\pgfsetbuttcap%
\pgfsetroundjoin%
\definecolor{currentfill}{rgb}{0.000000,0.000000,0.000000}%
\pgfsetfillcolor{currentfill}%
\pgfsetlinewidth{0.602250pt}%
\definecolor{currentstroke}{rgb}{0.000000,0.000000,0.000000}%
\pgfsetstrokecolor{currentstroke}%
\pgfsetdash{}{0pt}%
\pgfsys@defobject{currentmarker}{\pgfqpoint{0.000000in}{-0.027778in}}{\pgfqpoint{0.000000in}{0.000000in}}{%
\pgfpathmoveto{\pgfqpoint{0.000000in}{0.000000in}}%
\pgfpathlineto{\pgfqpoint{0.000000in}{-0.027778in}}%
\pgfusepath{stroke,fill}%
}%
\begin{pgfscope}%
\pgfsys@transformshift{5.907995in}{0.440000in}%
\pgfsys@useobject{currentmarker}{}%
\end{pgfscope}%
\end{pgfscope}%
\begin{pgfscope}%
\pgfpathrectangle{\pgfqpoint{1.000000in}{0.440000in}}{\pgfqpoint{6.200000in}{3.080000in}}%
\pgfusepath{clip}%
\pgfsetbuttcap%
\pgfsetroundjoin%
\pgfsetlinewidth{1.003750pt}%
\definecolor{currentstroke}{rgb}{0.690196,0.690196,0.690196}%
\pgfsetstrokecolor{currentstroke}%
\pgfsetstrokeopacity{0.600000}%
\pgfsetdash{{1.000000pt}{1.650000pt}}{0.000000pt}%
\pgfpathmoveto{\pgfqpoint{5.963366in}{0.440000in}}%
\pgfpathlineto{\pgfqpoint{5.963366in}{3.520000in}}%
\pgfusepath{stroke}%
\end{pgfscope}%
\begin{pgfscope}%
\pgfsetbuttcap%
\pgfsetroundjoin%
\definecolor{currentfill}{rgb}{0.000000,0.000000,0.000000}%
\pgfsetfillcolor{currentfill}%
\pgfsetlinewidth{0.602250pt}%
\definecolor{currentstroke}{rgb}{0.000000,0.000000,0.000000}%
\pgfsetstrokecolor{currentstroke}%
\pgfsetdash{}{0pt}%
\pgfsys@defobject{currentmarker}{\pgfqpoint{0.000000in}{-0.027778in}}{\pgfqpoint{0.000000in}{0.000000in}}{%
\pgfpathmoveto{\pgfqpoint{0.000000in}{0.000000in}}%
\pgfpathlineto{\pgfqpoint{0.000000in}{-0.027778in}}%
\pgfusepath{stroke,fill}%
}%
\begin{pgfscope}%
\pgfsys@transformshift{5.963366in}{0.440000in}%
\pgfsys@useobject{currentmarker}{}%
\end{pgfscope}%
\end{pgfscope}%
\begin{pgfscope}%
\pgfpathrectangle{\pgfqpoint{1.000000in}{0.440000in}}{\pgfqpoint{6.200000in}{3.080000in}}%
\pgfusepath{clip}%
\pgfsetbuttcap%
\pgfsetroundjoin%
\pgfsetlinewidth{1.003750pt}%
\definecolor{currentstroke}{rgb}{0.690196,0.690196,0.690196}%
\pgfsetstrokecolor{currentstroke}%
\pgfsetstrokeopacity{0.600000}%
\pgfsetdash{{1.000000pt}{1.650000pt}}{0.000000pt}%
\pgfpathmoveto{\pgfqpoint{6.012208in}{0.440000in}}%
\pgfpathlineto{\pgfqpoint{6.012208in}{3.520000in}}%
\pgfusepath{stroke}%
\end{pgfscope}%
\begin{pgfscope}%
\pgfsetbuttcap%
\pgfsetroundjoin%
\definecolor{currentfill}{rgb}{0.000000,0.000000,0.000000}%
\pgfsetfillcolor{currentfill}%
\pgfsetlinewidth{0.602250pt}%
\definecolor{currentstroke}{rgb}{0.000000,0.000000,0.000000}%
\pgfsetstrokecolor{currentstroke}%
\pgfsetdash{}{0pt}%
\pgfsys@defobject{currentmarker}{\pgfqpoint{0.000000in}{-0.027778in}}{\pgfqpoint{0.000000in}{0.000000in}}{%
\pgfpathmoveto{\pgfqpoint{0.000000in}{0.000000in}}%
\pgfpathlineto{\pgfqpoint{0.000000in}{-0.027778in}}%
\pgfusepath{stroke,fill}%
}%
\begin{pgfscope}%
\pgfsys@transformshift{6.012208in}{0.440000in}%
\pgfsys@useobject{currentmarker}{}%
\end{pgfscope}%
\end{pgfscope}%
\begin{pgfscope}%
\pgfpathrectangle{\pgfqpoint{1.000000in}{0.440000in}}{\pgfqpoint{6.200000in}{3.080000in}}%
\pgfusepath{clip}%
\pgfsetbuttcap%
\pgfsetroundjoin%
\pgfsetlinewidth{1.003750pt}%
\definecolor{currentstroke}{rgb}{0.690196,0.690196,0.690196}%
\pgfsetstrokecolor{currentstroke}%
\pgfsetstrokeopacity{0.600000}%
\pgfsetdash{{1.000000pt}{1.650000pt}}{0.000000pt}%
\pgfpathmoveto{\pgfqpoint{6.343326in}{0.440000in}}%
\pgfpathlineto{\pgfqpoint{6.343326in}{3.520000in}}%
\pgfusepath{stroke}%
\end{pgfscope}%
\begin{pgfscope}%
\pgfsetbuttcap%
\pgfsetroundjoin%
\definecolor{currentfill}{rgb}{0.000000,0.000000,0.000000}%
\pgfsetfillcolor{currentfill}%
\pgfsetlinewidth{0.602250pt}%
\definecolor{currentstroke}{rgb}{0.000000,0.000000,0.000000}%
\pgfsetstrokecolor{currentstroke}%
\pgfsetdash{}{0pt}%
\pgfsys@defobject{currentmarker}{\pgfqpoint{0.000000in}{-0.027778in}}{\pgfqpoint{0.000000in}{0.000000in}}{%
\pgfpathmoveto{\pgfqpoint{0.000000in}{0.000000in}}%
\pgfpathlineto{\pgfqpoint{0.000000in}{-0.027778in}}%
\pgfusepath{stroke,fill}%
}%
\begin{pgfscope}%
\pgfsys@transformshift{6.343326in}{0.440000in}%
\pgfsys@useobject{currentmarker}{}%
\end{pgfscope}%
\end{pgfscope}%
\begin{pgfscope}%
\pgfpathrectangle{\pgfqpoint{1.000000in}{0.440000in}}{\pgfqpoint{6.200000in}{3.080000in}}%
\pgfusepath{clip}%
\pgfsetbuttcap%
\pgfsetroundjoin%
\pgfsetlinewidth{1.003750pt}%
\definecolor{currentstroke}{rgb}{0.690196,0.690196,0.690196}%
\pgfsetstrokecolor{currentstroke}%
\pgfsetstrokeopacity{0.600000}%
\pgfsetdash{{1.000000pt}{1.650000pt}}{0.000000pt}%
\pgfpathmoveto{\pgfqpoint{6.511461in}{0.440000in}}%
\pgfpathlineto{\pgfqpoint{6.511461in}{3.520000in}}%
\pgfusepath{stroke}%
\end{pgfscope}%
\begin{pgfscope}%
\pgfsetbuttcap%
\pgfsetroundjoin%
\definecolor{currentfill}{rgb}{0.000000,0.000000,0.000000}%
\pgfsetfillcolor{currentfill}%
\pgfsetlinewidth{0.602250pt}%
\definecolor{currentstroke}{rgb}{0.000000,0.000000,0.000000}%
\pgfsetstrokecolor{currentstroke}%
\pgfsetdash{}{0pt}%
\pgfsys@defobject{currentmarker}{\pgfqpoint{0.000000in}{-0.027778in}}{\pgfqpoint{0.000000in}{0.000000in}}{%
\pgfpathmoveto{\pgfqpoint{0.000000in}{0.000000in}}%
\pgfpathlineto{\pgfqpoint{0.000000in}{-0.027778in}}%
\pgfusepath{stroke,fill}%
}%
\begin{pgfscope}%
\pgfsys@transformshift{6.511461in}{0.440000in}%
\pgfsys@useobject{currentmarker}{}%
\end{pgfscope}%
\end{pgfscope}%
\begin{pgfscope}%
\pgfpathrectangle{\pgfqpoint{1.000000in}{0.440000in}}{\pgfqpoint{6.200000in}{3.080000in}}%
\pgfusepath{clip}%
\pgfsetbuttcap%
\pgfsetroundjoin%
\pgfsetlinewidth{1.003750pt}%
\definecolor{currentstroke}{rgb}{0.690196,0.690196,0.690196}%
\pgfsetstrokecolor{currentstroke}%
\pgfsetstrokeopacity{0.600000}%
\pgfsetdash{{1.000000pt}{1.650000pt}}{0.000000pt}%
\pgfpathmoveto{\pgfqpoint{6.630754in}{0.440000in}}%
\pgfpathlineto{\pgfqpoint{6.630754in}{3.520000in}}%
\pgfusepath{stroke}%
\end{pgfscope}%
\begin{pgfscope}%
\pgfsetbuttcap%
\pgfsetroundjoin%
\definecolor{currentfill}{rgb}{0.000000,0.000000,0.000000}%
\pgfsetfillcolor{currentfill}%
\pgfsetlinewidth{0.602250pt}%
\definecolor{currentstroke}{rgb}{0.000000,0.000000,0.000000}%
\pgfsetstrokecolor{currentstroke}%
\pgfsetdash{}{0pt}%
\pgfsys@defobject{currentmarker}{\pgfqpoint{0.000000in}{-0.027778in}}{\pgfqpoint{0.000000in}{0.000000in}}{%
\pgfpathmoveto{\pgfqpoint{0.000000in}{0.000000in}}%
\pgfpathlineto{\pgfqpoint{0.000000in}{-0.027778in}}%
\pgfusepath{stroke,fill}%
}%
\begin{pgfscope}%
\pgfsys@transformshift{6.630754in}{0.440000in}%
\pgfsys@useobject{currentmarker}{}%
\end{pgfscope}%
\end{pgfscope}%
\begin{pgfscope}%
\pgfpathrectangle{\pgfqpoint{1.000000in}{0.440000in}}{\pgfqpoint{6.200000in}{3.080000in}}%
\pgfusepath{clip}%
\pgfsetbuttcap%
\pgfsetroundjoin%
\pgfsetlinewidth{1.003750pt}%
\definecolor{currentstroke}{rgb}{0.690196,0.690196,0.690196}%
\pgfsetstrokecolor{currentstroke}%
\pgfsetstrokeopacity{0.600000}%
\pgfsetdash{{1.000000pt}{1.650000pt}}{0.000000pt}%
\pgfpathmoveto{\pgfqpoint{6.723285in}{0.440000in}}%
\pgfpathlineto{\pgfqpoint{6.723285in}{3.520000in}}%
\pgfusepath{stroke}%
\end{pgfscope}%
\begin{pgfscope}%
\pgfsetbuttcap%
\pgfsetroundjoin%
\definecolor{currentfill}{rgb}{0.000000,0.000000,0.000000}%
\pgfsetfillcolor{currentfill}%
\pgfsetlinewidth{0.602250pt}%
\definecolor{currentstroke}{rgb}{0.000000,0.000000,0.000000}%
\pgfsetstrokecolor{currentstroke}%
\pgfsetdash{}{0pt}%
\pgfsys@defobject{currentmarker}{\pgfqpoint{0.000000in}{-0.027778in}}{\pgfqpoint{0.000000in}{0.000000in}}{%
\pgfpathmoveto{\pgfqpoint{0.000000in}{0.000000in}}%
\pgfpathlineto{\pgfqpoint{0.000000in}{-0.027778in}}%
\pgfusepath{stroke,fill}%
}%
\begin{pgfscope}%
\pgfsys@transformshift{6.723285in}{0.440000in}%
\pgfsys@useobject{currentmarker}{}%
\end{pgfscope}%
\end{pgfscope}%
\begin{pgfscope}%
\pgfpathrectangle{\pgfqpoint{1.000000in}{0.440000in}}{\pgfqpoint{6.200000in}{3.080000in}}%
\pgfusepath{clip}%
\pgfsetbuttcap%
\pgfsetroundjoin%
\pgfsetlinewidth{1.003750pt}%
\definecolor{currentstroke}{rgb}{0.690196,0.690196,0.690196}%
\pgfsetstrokecolor{currentstroke}%
\pgfsetstrokeopacity{0.600000}%
\pgfsetdash{{1.000000pt}{1.650000pt}}{0.000000pt}%
\pgfpathmoveto{\pgfqpoint{6.798889in}{0.440000in}}%
\pgfpathlineto{\pgfqpoint{6.798889in}{3.520000in}}%
\pgfusepath{stroke}%
\end{pgfscope}%
\begin{pgfscope}%
\pgfsetbuttcap%
\pgfsetroundjoin%
\definecolor{currentfill}{rgb}{0.000000,0.000000,0.000000}%
\pgfsetfillcolor{currentfill}%
\pgfsetlinewidth{0.602250pt}%
\definecolor{currentstroke}{rgb}{0.000000,0.000000,0.000000}%
\pgfsetstrokecolor{currentstroke}%
\pgfsetdash{}{0pt}%
\pgfsys@defobject{currentmarker}{\pgfqpoint{0.000000in}{-0.027778in}}{\pgfqpoint{0.000000in}{0.000000in}}{%
\pgfpathmoveto{\pgfqpoint{0.000000in}{0.000000in}}%
\pgfpathlineto{\pgfqpoint{0.000000in}{-0.027778in}}%
\pgfusepath{stroke,fill}%
}%
\begin{pgfscope}%
\pgfsys@transformshift{6.798889in}{0.440000in}%
\pgfsys@useobject{currentmarker}{}%
\end{pgfscope}%
\end{pgfscope}%
\begin{pgfscope}%
\pgfpathrectangle{\pgfqpoint{1.000000in}{0.440000in}}{\pgfqpoint{6.200000in}{3.080000in}}%
\pgfusepath{clip}%
\pgfsetbuttcap%
\pgfsetroundjoin%
\pgfsetlinewidth{1.003750pt}%
\definecolor{currentstroke}{rgb}{0.690196,0.690196,0.690196}%
\pgfsetstrokecolor{currentstroke}%
\pgfsetstrokeopacity{0.600000}%
\pgfsetdash{{1.000000pt}{1.650000pt}}{0.000000pt}%
\pgfpathmoveto{\pgfqpoint{6.862811in}{0.440000in}}%
\pgfpathlineto{\pgfqpoint{6.862811in}{3.520000in}}%
\pgfusepath{stroke}%
\end{pgfscope}%
\begin{pgfscope}%
\pgfsetbuttcap%
\pgfsetroundjoin%
\definecolor{currentfill}{rgb}{0.000000,0.000000,0.000000}%
\pgfsetfillcolor{currentfill}%
\pgfsetlinewidth{0.602250pt}%
\definecolor{currentstroke}{rgb}{0.000000,0.000000,0.000000}%
\pgfsetstrokecolor{currentstroke}%
\pgfsetdash{}{0pt}%
\pgfsys@defobject{currentmarker}{\pgfqpoint{0.000000in}{-0.027778in}}{\pgfqpoint{0.000000in}{0.000000in}}{%
\pgfpathmoveto{\pgfqpoint{0.000000in}{0.000000in}}%
\pgfpathlineto{\pgfqpoint{0.000000in}{-0.027778in}}%
\pgfusepath{stroke,fill}%
}%
\begin{pgfscope}%
\pgfsys@transformshift{6.862811in}{0.440000in}%
\pgfsys@useobject{currentmarker}{}%
\end{pgfscope}%
\end{pgfscope}%
\begin{pgfscope}%
\pgfpathrectangle{\pgfqpoint{1.000000in}{0.440000in}}{\pgfqpoint{6.200000in}{3.080000in}}%
\pgfusepath{clip}%
\pgfsetbuttcap%
\pgfsetroundjoin%
\pgfsetlinewidth{1.003750pt}%
\definecolor{currentstroke}{rgb}{0.690196,0.690196,0.690196}%
\pgfsetstrokecolor{currentstroke}%
\pgfsetstrokeopacity{0.600000}%
\pgfsetdash{{1.000000pt}{1.650000pt}}{0.000000pt}%
\pgfpathmoveto{\pgfqpoint{6.918182in}{0.440000in}}%
\pgfpathlineto{\pgfqpoint{6.918182in}{3.520000in}}%
\pgfusepath{stroke}%
\end{pgfscope}%
\begin{pgfscope}%
\pgfsetbuttcap%
\pgfsetroundjoin%
\definecolor{currentfill}{rgb}{0.000000,0.000000,0.000000}%
\pgfsetfillcolor{currentfill}%
\pgfsetlinewidth{0.602250pt}%
\definecolor{currentstroke}{rgb}{0.000000,0.000000,0.000000}%
\pgfsetstrokecolor{currentstroke}%
\pgfsetdash{}{0pt}%
\pgfsys@defobject{currentmarker}{\pgfqpoint{0.000000in}{-0.027778in}}{\pgfqpoint{0.000000in}{0.000000in}}{%
\pgfpathmoveto{\pgfqpoint{0.000000in}{0.000000in}}%
\pgfpathlineto{\pgfqpoint{0.000000in}{-0.027778in}}%
\pgfusepath{stroke,fill}%
}%
\begin{pgfscope}%
\pgfsys@transformshift{6.918182in}{0.440000in}%
\pgfsys@useobject{currentmarker}{}%
\end{pgfscope}%
\end{pgfscope}%
\begin{pgfscope}%
\pgfpathrectangle{\pgfqpoint{1.000000in}{0.440000in}}{\pgfqpoint{6.200000in}{3.080000in}}%
\pgfusepath{clip}%
\pgfsetbuttcap%
\pgfsetroundjoin%
\pgfsetlinewidth{1.003750pt}%
\definecolor{currentstroke}{rgb}{0.690196,0.690196,0.690196}%
\pgfsetstrokecolor{currentstroke}%
\pgfsetstrokeopacity{0.600000}%
\pgfsetdash{{1.000000pt}{1.650000pt}}{0.000000pt}%
\pgfpathmoveto{\pgfqpoint{6.967024in}{0.440000in}}%
\pgfpathlineto{\pgfqpoint{6.967024in}{3.520000in}}%
\pgfusepath{stroke}%
\end{pgfscope}%
\begin{pgfscope}%
\pgfsetbuttcap%
\pgfsetroundjoin%
\definecolor{currentfill}{rgb}{0.000000,0.000000,0.000000}%
\pgfsetfillcolor{currentfill}%
\pgfsetlinewidth{0.602250pt}%
\definecolor{currentstroke}{rgb}{0.000000,0.000000,0.000000}%
\pgfsetstrokecolor{currentstroke}%
\pgfsetdash{}{0pt}%
\pgfsys@defobject{currentmarker}{\pgfqpoint{0.000000in}{-0.027778in}}{\pgfqpoint{0.000000in}{0.000000in}}{%
\pgfpathmoveto{\pgfqpoint{0.000000in}{0.000000in}}%
\pgfpathlineto{\pgfqpoint{0.000000in}{-0.027778in}}%
\pgfusepath{stroke,fill}%
}%
\begin{pgfscope}%
\pgfsys@transformshift{6.967024in}{0.440000in}%
\pgfsys@useobject{currentmarker}{}%
\end{pgfscope}%
\end{pgfscope}%
\begin{pgfscope}%
\definecolor{textcolor}{rgb}{0.000000,0.000000,0.000000}%
\pgfsetstrokecolor{textcolor}%
\pgfsetfillcolor{textcolor}%
\pgftext[x=4.100000in,y=0.152809in,,top]{\color{textcolor}{\rmfamily\fontsize{12.000000}{14.400000}\selectfont\catcode`\^=\active\def^{\ifmmode\sp\else\^{}\fi}\catcode`\%=\active\def%{\%}Rank}}%
\end{pgfscope}%
\begin{pgfscope}%
\pgfpathrectangle{\pgfqpoint{1.000000in}{0.440000in}}{\pgfqpoint{6.200000in}{3.080000in}}%
\pgfusepath{clip}%
\pgfsetbuttcap%
\pgfsetroundjoin%
\pgfsetlinewidth{1.003750pt}%
\definecolor{currentstroke}{rgb}{0.690196,0.690196,0.690196}%
\pgfsetstrokecolor{currentstroke}%
\pgfsetstrokeopacity{0.600000}%
\pgfsetdash{{1.000000pt}{1.650000pt}}{0.000000pt}%
\pgfpathmoveto{\pgfqpoint{1.000000in}{0.915579in}}%
\pgfpathlineto{\pgfqpoint{7.200000in}{0.915579in}}%
\pgfusepath{stroke}%
\end{pgfscope}%
\begin{pgfscope}%
\pgfsetbuttcap%
\pgfsetroundjoin%
\definecolor{currentfill}{rgb}{0.000000,0.000000,0.000000}%
\pgfsetfillcolor{currentfill}%
\pgfsetlinewidth{0.803000pt}%
\definecolor{currentstroke}{rgb}{0.000000,0.000000,0.000000}%
\pgfsetstrokecolor{currentstroke}%
\pgfsetdash{}{0pt}%
\pgfsys@defobject{currentmarker}{\pgfqpoint{-0.048611in}{0.000000in}}{\pgfqpoint{-0.000000in}{0.000000in}}{%
\pgfpathmoveto{\pgfqpoint{-0.000000in}{0.000000in}}%
\pgfpathlineto{\pgfqpoint{-0.048611in}{0.000000in}}%
\pgfusepath{stroke,fill}%
}%
\begin{pgfscope}%
\pgfsys@transformshift{1.000000in}{0.915579in}%
\pgfsys@useobject{currentmarker}{}%
\end{pgfscope}%
\end{pgfscope}%
\begin{pgfscope}%
\definecolor{textcolor}{rgb}{0.000000,0.000000,0.000000}%
\pgfsetstrokecolor{textcolor}%
\pgfsetfillcolor{textcolor}%
\pgftext[x=0.701581in, y=0.862818in, left, base]{\color{textcolor}{\rmfamily\fontsize{10.000000}{12.000000}\selectfont\catcode`\^=\active\def^{\ifmmode\sp\else\^{}\fi}\catcode`\%=\active\def%{\%}$\mathdefault{10^{1}}$}}%
\end{pgfscope}%
\begin{pgfscope}%
\pgfpathrectangle{\pgfqpoint{1.000000in}{0.440000in}}{\pgfqpoint{6.200000in}{3.080000in}}%
\pgfusepath{clip}%
\pgfsetbuttcap%
\pgfsetroundjoin%
\pgfsetlinewidth{1.003750pt}%
\definecolor{currentstroke}{rgb}{0.690196,0.690196,0.690196}%
\pgfsetstrokecolor{currentstroke}%
\pgfsetstrokeopacity{0.600000}%
\pgfsetdash{{1.000000pt}{1.650000pt}}{0.000000pt}%
\pgfpathmoveto{\pgfqpoint{1.000000in}{1.557371in}}%
\pgfpathlineto{\pgfqpoint{7.200000in}{1.557371in}}%
\pgfusepath{stroke}%
\end{pgfscope}%
\begin{pgfscope}%
\pgfsetbuttcap%
\pgfsetroundjoin%
\definecolor{currentfill}{rgb}{0.000000,0.000000,0.000000}%
\pgfsetfillcolor{currentfill}%
\pgfsetlinewidth{0.803000pt}%
\definecolor{currentstroke}{rgb}{0.000000,0.000000,0.000000}%
\pgfsetstrokecolor{currentstroke}%
\pgfsetdash{}{0pt}%
\pgfsys@defobject{currentmarker}{\pgfqpoint{-0.048611in}{0.000000in}}{\pgfqpoint{-0.000000in}{0.000000in}}{%
\pgfpathmoveto{\pgfqpoint{-0.000000in}{0.000000in}}%
\pgfpathlineto{\pgfqpoint{-0.048611in}{0.000000in}}%
\pgfusepath{stroke,fill}%
}%
\begin{pgfscope}%
\pgfsys@transformshift{1.000000in}{1.557371in}%
\pgfsys@useobject{currentmarker}{}%
\end{pgfscope}%
\end{pgfscope}%
\begin{pgfscope}%
\definecolor{textcolor}{rgb}{0.000000,0.000000,0.000000}%
\pgfsetstrokecolor{textcolor}%
\pgfsetfillcolor{textcolor}%
\pgftext[x=0.701581in, y=1.504610in, left, base]{\color{textcolor}{\rmfamily\fontsize{10.000000}{12.000000}\selectfont\catcode`\^=\active\def^{\ifmmode\sp\else\^{}\fi}\catcode`\%=\active\def%{\%}$\mathdefault{10^{2}}$}}%
\end{pgfscope}%
\begin{pgfscope}%
\pgfpathrectangle{\pgfqpoint{1.000000in}{0.440000in}}{\pgfqpoint{6.200000in}{3.080000in}}%
\pgfusepath{clip}%
\pgfsetbuttcap%
\pgfsetroundjoin%
\pgfsetlinewidth{1.003750pt}%
\definecolor{currentstroke}{rgb}{0.690196,0.690196,0.690196}%
\pgfsetstrokecolor{currentstroke}%
\pgfsetstrokeopacity{0.600000}%
\pgfsetdash{{1.000000pt}{1.650000pt}}{0.000000pt}%
\pgfpathmoveto{\pgfqpoint{1.000000in}{2.199163in}}%
\pgfpathlineto{\pgfqpoint{7.200000in}{2.199163in}}%
\pgfusepath{stroke}%
\end{pgfscope}%
\begin{pgfscope}%
\pgfsetbuttcap%
\pgfsetroundjoin%
\definecolor{currentfill}{rgb}{0.000000,0.000000,0.000000}%
\pgfsetfillcolor{currentfill}%
\pgfsetlinewidth{0.803000pt}%
\definecolor{currentstroke}{rgb}{0.000000,0.000000,0.000000}%
\pgfsetstrokecolor{currentstroke}%
\pgfsetdash{}{0pt}%
\pgfsys@defobject{currentmarker}{\pgfqpoint{-0.048611in}{0.000000in}}{\pgfqpoint{-0.000000in}{0.000000in}}{%
\pgfpathmoveto{\pgfqpoint{-0.000000in}{0.000000in}}%
\pgfpathlineto{\pgfqpoint{-0.048611in}{0.000000in}}%
\pgfusepath{stroke,fill}%
}%
\begin{pgfscope}%
\pgfsys@transformshift{1.000000in}{2.199163in}%
\pgfsys@useobject{currentmarker}{}%
\end{pgfscope}%
\end{pgfscope}%
\begin{pgfscope}%
\definecolor{textcolor}{rgb}{0.000000,0.000000,0.000000}%
\pgfsetstrokecolor{textcolor}%
\pgfsetfillcolor{textcolor}%
\pgftext[x=0.701581in, y=2.146402in, left, base]{\color{textcolor}{\rmfamily\fontsize{10.000000}{12.000000}\selectfont\catcode`\^=\active\def^{\ifmmode\sp\else\^{}\fi}\catcode`\%=\active\def%{\%}$\mathdefault{10^{3}}$}}%
\end{pgfscope}%
\begin{pgfscope}%
\pgfpathrectangle{\pgfqpoint{1.000000in}{0.440000in}}{\pgfqpoint{6.200000in}{3.080000in}}%
\pgfusepath{clip}%
\pgfsetbuttcap%
\pgfsetroundjoin%
\pgfsetlinewidth{1.003750pt}%
\definecolor{currentstroke}{rgb}{0.690196,0.690196,0.690196}%
\pgfsetstrokecolor{currentstroke}%
\pgfsetstrokeopacity{0.600000}%
\pgfsetdash{{1.000000pt}{1.650000pt}}{0.000000pt}%
\pgfpathmoveto{\pgfqpoint{1.000000in}{2.840956in}}%
\pgfpathlineto{\pgfqpoint{7.200000in}{2.840956in}}%
\pgfusepath{stroke}%
\end{pgfscope}%
\begin{pgfscope}%
\pgfsetbuttcap%
\pgfsetroundjoin%
\definecolor{currentfill}{rgb}{0.000000,0.000000,0.000000}%
\pgfsetfillcolor{currentfill}%
\pgfsetlinewidth{0.803000pt}%
\definecolor{currentstroke}{rgb}{0.000000,0.000000,0.000000}%
\pgfsetstrokecolor{currentstroke}%
\pgfsetdash{}{0pt}%
\pgfsys@defobject{currentmarker}{\pgfqpoint{-0.048611in}{0.000000in}}{\pgfqpoint{-0.000000in}{0.000000in}}{%
\pgfpathmoveto{\pgfqpoint{-0.000000in}{0.000000in}}%
\pgfpathlineto{\pgfqpoint{-0.048611in}{0.000000in}}%
\pgfusepath{stroke,fill}%
}%
\begin{pgfscope}%
\pgfsys@transformshift{1.000000in}{2.840956in}%
\pgfsys@useobject{currentmarker}{}%
\end{pgfscope}%
\end{pgfscope}%
\begin{pgfscope}%
\definecolor{textcolor}{rgb}{0.000000,0.000000,0.000000}%
\pgfsetstrokecolor{textcolor}%
\pgfsetfillcolor{textcolor}%
\pgftext[x=0.701581in, y=2.788194in, left, base]{\color{textcolor}{\rmfamily\fontsize{10.000000}{12.000000}\selectfont\catcode`\^=\active\def^{\ifmmode\sp\else\^{}\fi}\catcode`\%=\active\def%{\%}$\mathdefault{10^{4}}$}}%
\end{pgfscope}%
\begin{pgfscope}%
\pgfpathrectangle{\pgfqpoint{1.000000in}{0.440000in}}{\pgfqpoint{6.200000in}{3.080000in}}%
\pgfusepath{clip}%
\pgfsetbuttcap%
\pgfsetroundjoin%
\pgfsetlinewidth{1.003750pt}%
\definecolor{currentstroke}{rgb}{0.690196,0.690196,0.690196}%
\pgfsetstrokecolor{currentstroke}%
\pgfsetstrokeopacity{0.600000}%
\pgfsetdash{{1.000000pt}{1.650000pt}}{0.000000pt}%
\pgfpathmoveto{\pgfqpoint{1.000000in}{3.482748in}}%
\pgfpathlineto{\pgfqpoint{7.200000in}{3.482748in}}%
\pgfusepath{stroke}%
\end{pgfscope}%
\begin{pgfscope}%
\pgfsetbuttcap%
\pgfsetroundjoin%
\definecolor{currentfill}{rgb}{0.000000,0.000000,0.000000}%
\pgfsetfillcolor{currentfill}%
\pgfsetlinewidth{0.803000pt}%
\definecolor{currentstroke}{rgb}{0.000000,0.000000,0.000000}%
\pgfsetstrokecolor{currentstroke}%
\pgfsetdash{}{0pt}%
\pgfsys@defobject{currentmarker}{\pgfqpoint{-0.048611in}{0.000000in}}{\pgfqpoint{-0.000000in}{0.000000in}}{%
\pgfpathmoveto{\pgfqpoint{-0.000000in}{0.000000in}}%
\pgfpathlineto{\pgfqpoint{-0.048611in}{0.000000in}}%
\pgfusepath{stroke,fill}%
}%
\begin{pgfscope}%
\pgfsys@transformshift{1.000000in}{3.482748in}%
\pgfsys@useobject{currentmarker}{}%
\end{pgfscope}%
\end{pgfscope}%
\begin{pgfscope}%
\definecolor{textcolor}{rgb}{0.000000,0.000000,0.000000}%
\pgfsetstrokecolor{textcolor}%
\pgfsetfillcolor{textcolor}%
\pgftext[x=0.701581in, y=3.429986in, left, base]{\color{textcolor}{\rmfamily\fontsize{10.000000}{12.000000}\selectfont\catcode`\^=\active\def^{\ifmmode\sp\else\^{}\fi}\catcode`\%=\active\def%{\%}$\mathdefault{10^{5}}$}}%
\end{pgfscope}%
\begin{pgfscope}%
\pgfpathrectangle{\pgfqpoint{1.000000in}{0.440000in}}{\pgfqpoint{6.200000in}{3.080000in}}%
\pgfusepath{clip}%
\pgfsetbuttcap%
\pgfsetroundjoin%
\pgfsetlinewidth{1.003750pt}%
\definecolor{currentstroke}{rgb}{0.690196,0.690196,0.690196}%
\pgfsetstrokecolor{currentstroke}%
\pgfsetstrokeopacity{0.600000}%
\pgfsetdash{{1.000000pt}{1.650000pt}}{0.000000pt}%
\pgfpathmoveto{\pgfqpoint{1.000000in}{0.466986in}}%
\pgfpathlineto{\pgfqpoint{7.200000in}{0.466986in}}%
\pgfusepath{stroke}%
\end{pgfscope}%
\begin{pgfscope}%
\pgfsetbuttcap%
\pgfsetroundjoin%
\definecolor{currentfill}{rgb}{0.000000,0.000000,0.000000}%
\pgfsetfillcolor{currentfill}%
\pgfsetlinewidth{0.602250pt}%
\definecolor{currentstroke}{rgb}{0.000000,0.000000,0.000000}%
\pgfsetstrokecolor{currentstroke}%
\pgfsetdash{}{0pt}%
\pgfsys@defobject{currentmarker}{\pgfqpoint{-0.027778in}{0.000000in}}{\pgfqpoint{-0.000000in}{0.000000in}}{%
\pgfpathmoveto{\pgfqpoint{-0.000000in}{0.000000in}}%
\pgfpathlineto{\pgfqpoint{-0.027778in}{0.000000in}}%
\pgfusepath{stroke,fill}%
}%
\begin{pgfscope}%
\pgfsys@transformshift{1.000000in}{0.466986in}%
\pgfsys@useobject{currentmarker}{}%
\end{pgfscope}%
\end{pgfscope}%
\begin{pgfscope}%
\pgfpathrectangle{\pgfqpoint{1.000000in}{0.440000in}}{\pgfqpoint{6.200000in}{3.080000in}}%
\pgfusepath{clip}%
\pgfsetbuttcap%
\pgfsetroundjoin%
\pgfsetlinewidth{1.003750pt}%
\definecolor{currentstroke}{rgb}{0.690196,0.690196,0.690196}%
\pgfsetstrokecolor{currentstroke}%
\pgfsetstrokeopacity{0.600000}%
\pgfsetdash{{1.000000pt}{1.650000pt}}{0.000000pt}%
\pgfpathmoveto{\pgfqpoint{1.000000in}{0.580000in}}%
\pgfpathlineto{\pgfqpoint{7.200000in}{0.580000in}}%
\pgfusepath{stroke}%
\end{pgfscope}%
\begin{pgfscope}%
\pgfsetbuttcap%
\pgfsetroundjoin%
\definecolor{currentfill}{rgb}{0.000000,0.000000,0.000000}%
\pgfsetfillcolor{currentfill}%
\pgfsetlinewidth{0.602250pt}%
\definecolor{currentstroke}{rgb}{0.000000,0.000000,0.000000}%
\pgfsetstrokecolor{currentstroke}%
\pgfsetdash{}{0pt}%
\pgfsys@defobject{currentmarker}{\pgfqpoint{-0.027778in}{0.000000in}}{\pgfqpoint{-0.000000in}{0.000000in}}{%
\pgfpathmoveto{\pgfqpoint{-0.000000in}{0.000000in}}%
\pgfpathlineto{\pgfqpoint{-0.027778in}{0.000000in}}%
\pgfusepath{stroke,fill}%
}%
\begin{pgfscope}%
\pgfsys@transformshift{1.000000in}{0.580000in}%
\pgfsys@useobject{currentmarker}{}%
\end{pgfscope}%
\end{pgfscope}%
\begin{pgfscope}%
\pgfpathrectangle{\pgfqpoint{1.000000in}{0.440000in}}{\pgfqpoint{6.200000in}{3.080000in}}%
\pgfusepath{clip}%
\pgfsetbuttcap%
\pgfsetroundjoin%
\pgfsetlinewidth{1.003750pt}%
\definecolor{currentstroke}{rgb}{0.690196,0.690196,0.690196}%
\pgfsetstrokecolor{currentstroke}%
\pgfsetstrokeopacity{0.600000}%
\pgfsetdash{{1.000000pt}{1.650000pt}}{0.000000pt}%
\pgfpathmoveto{\pgfqpoint{1.000000in}{0.660185in}}%
\pgfpathlineto{\pgfqpoint{7.200000in}{0.660185in}}%
\pgfusepath{stroke}%
\end{pgfscope}%
\begin{pgfscope}%
\pgfsetbuttcap%
\pgfsetroundjoin%
\definecolor{currentfill}{rgb}{0.000000,0.000000,0.000000}%
\pgfsetfillcolor{currentfill}%
\pgfsetlinewidth{0.602250pt}%
\definecolor{currentstroke}{rgb}{0.000000,0.000000,0.000000}%
\pgfsetstrokecolor{currentstroke}%
\pgfsetdash{}{0pt}%
\pgfsys@defobject{currentmarker}{\pgfqpoint{-0.027778in}{0.000000in}}{\pgfqpoint{-0.000000in}{0.000000in}}{%
\pgfpathmoveto{\pgfqpoint{-0.000000in}{0.000000in}}%
\pgfpathlineto{\pgfqpoint{-0.027778in}{0.000000in}}%
\pgfusepath{stroke,fill}%
}%
\begin{pgfscope}%
\pgfsys@transformshift{1.000000in}{0.660185in}%
\pgfsys@useobject{currentmarker}{}%
\end{pgfscope}%
\end{pgfscope}%
\begin{pgfscope}%
\pgfpathrectangle{\pgfqpoint{1.000000in}{0.440000in}}{\pgfqpoint{6.200000in}{3.080000in}}%
\pgfusepath{clip}%
\pgfsetbuttcap%
\pgfsetroundjoin%
\pgfsetlinewidth{1.003750pt}%
\definecolor{currentstroke}{rgb}{0.690196,0.690196,0.690196}%
\pgfsetstrokecolor{currentstroke}%
\pgfsetstrokeopacity{0.600000}%
\pgfsetdash{{1.000000pt}{1.650000pt}}{0.000000pt}%
\pgfpathmoveto{\pgfqpoint{1.000000in}{0.722381in}}%
\pgfpathlineto{\pgfqpoint{7.200000in}{0.722381in}}%
\pgfusepath{stroke}%
\end{pgfscope}%
\begin{pgfscope}%
\pgfsetbuttcap%
\pgfsetroundjoin%
\definecolor{currentfill}{rgb}{0.000000,0.000000,0.000000}%
\pgfsetfillcolor{currentfill}%
\pgfsetlinewidth{0.602250pt}%
\definecolor{currentstroke}{rgb}{0.000000,0.000000,0.000000}%
\pgfsetstrokecolor{currentstroke}%
\pgfsetdash{}{0pt}%
\pgfsys@defobject{currentmarker}{\pgfqpoint{-0.027778in}{0.000000in}}{\pgfqpoint{-0.000000in}{0.000000in}}{%
\pgfpathmoveto{\pgfqpoint{-0.000000in}{0.000000in}}%
\pgfpathlineto{\pgfqpoint{-0.027778in}{0.000000in}}%
\pgfusepath{stroke,fill}%
}%
\begin{pgfscope}%
\pgfsys@transformshift{1.000000in}{0.722381in}%
\pgfsys@useobject{currentmarker}{}%
\end{pgfscope}%
\end{pgfscope}%
\begin{pgfscope}%
\pgfpathrectangle{\pgfqpoint{1.000000in}{0.440000in}}{\pgfqpoint{6.200000in}{3.080000in}}%
\pgfusepath{clip}%
\pgfsetbuttcap%
\pgfsetroundjoin%
\pgfsetlinewidth{1.003750pt}%
\definecolor{currentstroke}{rgb}{0.690196,0.690196,0.690196}%
\pgfsetstrokecolor{currentstroke}%
\pgfsetstrokeopacity{0.600000}%
\pgfsetdash{{1.000000pt}{1.650000pt}}{0.000000pt}%
\pgfpathmoveto{\pgfqpoint{1.000000in}{0.773199in}}%
\pgfpathlineto{\pgfqpoint{7.200000in}{0.773199in}}%
\pgfusepath{stroke}%
\end{pgfscope}%
\begin{pgfscope}%
\pgfsetbuttcap%
\pgfsetroundjoin%
\definecolor{currentfill}{rgb}{0.000000,0.000000,0.000000}%
\pgfsetfillcolor{currentfill}%
\pgfsetlinewidth{0.602250pt}%
\definecolor{currentstroke}{rgb}{0.000000,0.000000,0.000000}%
\pgfsetstrokecolor{currentstroke}%
\pgfsetdash{}{0pt}%
\pgfsys@defobject{currentmarker}{\pgfqpoint{-0.027778in}{0.000000in}}{\pgfqpoint{-0.000000in}{0.000000in}}{%
\pgfpathmoveto{\pgfqpoint{-0.000000in}{0.000000in}}%
\pgfpathlineto{\pgfqpoint{-0.027778in}{0.000000in}}%
\pgfusepath{stroke,fill}%
}%
\begin{pgfscope}%
\pgfsys@transformshift{1.000000in}{0.773199in}%
\pgfsys@useobject{currentmarker}{}%
\end{pgfscope}%
\end{pgfscope}%
\begin{pgfscope}%
\pgfpathrectangle{\pgfqpoint{1.000000in}{0.440000in}}{\pgfqpoint{6.200000in}{3.080000in}}%
\pgfusepath{clip}%
\pgfsetbuttcap%
\pgfsetroundjoin%
\pgfsetlinewidth{1.003750pt}%
\definecolor{currentstroke}{rgb}{0.690196,0.690196,0.690196}%
\pgfsetstrokecolor{currentstroke}%
\pgfsetstrokeopacity{0.600000}%
\pgfsetdash{{1.000000pt}{1.650000pt}}{0.000000pt}%
\pgfpathmoveto{\pgfqpoint{1.000000in}{0.816165in}}%
\pgfpathlineto{\pgfqpoint{7.200000in}{0.816165in}}%
\pgfusepath{stroke}%
\end{pgfscope}%
\begin{pgfscope}%
\pgfsetbuttcap%
\pgfsetroundjoin%
\definecolor{currentfill}{rgb}{0.000000,0.000000,0.000000}%
\pgfsetfillcolor{currentfill}%
\pgfsetlinewidth{0.602250pt}%
\definecolor{currentstroke}{rgb}{0.000000,0.000000,0.000000}%
\pgfsetstrokecolor{currentstroke}%
\pgfsetdash{}{0pt}%
\pgfsys@defobject{currentmarker}{\pgfqpoint{-0.027778in}{0.000000in}}{\pgfqpoint{-0.000000in}{0.000000in}}{%
\pgfpathmoveto{\pgfqpoint{-0.000000in}{0.000000in}}%
\pgfpathlineto{\pgfqpoint{-0.027778in}{0.000000in}}%
\pgfusepath{stroke,fill}%
}%
\begin{pgfscope}%
\pgfsys@transformshift{1.000000in}{0.816165in}%
\pgfsys@useobject{currentmarker}{}%
\end{pgfscope}%
\end{pgfscope}%
\begin{pgfscope}%
\pgfpathrectangle{\pgfqpoint{1.000000in}{0.440000in}}{\pgfqpoint{6.200000in}{3.080000in}}%
\pgfusepath{clip}%
\pgfsetbuttcap%
\pgfsetroundjoin%
\pgfsetlinewidth{1.003750pt}%
\definecolor{currentstroke}{rgb}{0.690196,0.690196,0.690196}%
\pgfsetstrokecolor{currentstroke}%
\pgfsetstrokeopacity{0.600000}%
\pgfsetdash{{1.000000pt}{1.650000pt}}{0.000000pt}%
\pgfpathmoveto{\pgfqpoint{1.000000in}{0.853383in}}%
\pgfpathlineto{\pgfqpoint{7.200000in}{0.853383in}}%
\pgfusepath{stroke}%
\end{pgfscope}%
\begin{pgfscope}%
\pgfsetbuttcap%
\pgfsetroundjoin%
\definecolor{currentfill}{rgb}{0.000000,0.000000,0.000000}%
\pgfsetfillcolor{currentfill}%
\pgfsetlinewidth{0.602250pt}%
\definecolor{currentstroke}{rgb}{0.000000,0.000000,0.000000}%
\pgfsetstrokecolor{currentstroke}%
\pgfsetdash{}{0pt}%
\pgfsys@defobject{currentmarker}{\pgfqpoint{-0.027778in}{0.000000in}}{\pgfqpoint{-0.000000in}{0.000000in}}{%
\pgfpathmoveto{\pgfqpoint{-0.000000in}{0.000000in}}%
\pgfpathlineto{\pgfqpoint{-0.027778in}{0.000000in}}%
\pgfusepath{stroke,fill}%
}%
\begin{pgfscope}%
\pgfsys@transformshift{1.000000in}{0.853383in}%
\pgfsys@useobject{currentmarker}{}%
\end{pgfscope}%
\end{pgfscope}%
\begin{pgfscope}%
\pgfpathrectangle{\pgfqpoint{1.000000in}{0.440000in}}{\pgfqpoint{6.200000in}{3.080000in}}%
\pgfusepath{clip}%
\pgfsetbuttcap%
\pgfsetroundjoin%
\pgfsetlinewidth{1.003750pt}%
\definecolor{currentstroke}{rgb}{0.690196,0.690196,0.690196}%
\pgfsetstrokecolor{currentstroke}%
\pgfsetstrokeopacity{0.600000}%
\pgfsetdash{{1.000000pt}{1.650000pt}}{0.000000pt}%
\pgfpathmoveto{\pgfqpoint{1.000000in}{0.886213in}}%
\pgfpathlineto{\pgfqpoint{7.200000in}{0.886213in}}%
\pgfusepath{stroke}%
\end{pgfscope}%
\begin{pgfscope}%
\pgfsetbuttcap%
\pgfsetroundjoin%
\definecolor{currentfill}{rgb}{0.000000,0.000000,0.000000}%
\pgfsetfillcolor{currentfill}%
\pgfsetlinewidth{0.602250pt}%
\definecolor{currentstroke}{rgb}{0.000000,0.000000,0.000000}%
\pgfsetstrokecolor{currentstroke}%
\pgfsetdash{}{0pt}%
\pgfsys@defobject{currentmarker}{\pgfqpoint{-0.027778in}{0.000000in}}{\pgfqpoint{-0.000000in}{0.000000in}}{%
\pgfpathmoveto{\pgfqpoint{-0.000000in}{0.000000in}}%
\pgfpathlineto{\pgfqpoint{-0.027778in}{0.000000in}}%
\pgfusepath{stroke,fill}%
}%
\begin{pgfscope}%
\pgfsys@transformshift{1.000000in}{0.886213in}%
\pgfsys@useobject{currentmarker}{}%
\end{pgfscope}%
\end{pgfscope}%
\begin{pgfscope}%
\pgfpathrectangle{\pgfqpoint{1.000000in}{0.440000in}}{\pgfqpoint{6.200000in}{3.080000in}}%
\pgfusepath{clip}%
\pgfsetbuttcap%
\pgfsetroundjoin%
\pgfsetlinewidth{1.003750pt}%
\definecolor{currentstroke}{rgb}{0.690196,0.690196,0.690196}%
\pgfsetstrokecolor{currentstroke}%
\pgfsetstrokeopacity{0.600000}%
\pgfsetdash{{1.000000pt}{1.650000pt}}{0.000000pt}%
\pgfpathmoveto{\pgfqpoint{1.000000in}{1.108778in}}%
\pgfpathlineto{\pgfqpoint{7.200000in}{1.108778in}}%
\pgfusepath{stroke}%
\end{pgfscope}%
\begin{pgfscope}%
\pgfsetbuttcap%
\pgfsetroundjoin%
\definecolor{currentfill}{rgb}{0.000000,0.000000,0.000000}%
\pgfsetfillcolor{currentfill}%
\pgfsetlinewidth{0.602250pt}%
\definecolor{currentstroke}{rgb}{0.000000,0.000000,0.000000}%
\pgfsetstrokecolor{currentstroke}%
\pgfsetdash{}{0pt}%
\pgfsys@defobject{currentmarker}{\pgfqpoint{-0.027778in}{0.000000in}}{\pgfqpoint{-0.000000in}{0.000000in}}{%
\pgfpathmoveto{\pgfqpoint{-0.000000in}{0.000000in}}%
\pgfpathlineto{\pgfqpoint{-0.027778in}{0.000000in}}%
\pgfusepath{stroke,fill}%
}%
\begin{pgfscope}%
\pgfsys@transformshift{1.000000in}{1.108778in}%
\pgfsys@useobject{currentmarker}{}%
\end{pgfscope}%
\end{pgfscope}%
\begin{pgfscope}%
\pgfpathrectangle{\pgfqpoint{1.000000in}{0.440000in}}{\pgfqpoint{6.200000in}{3.080000in}}%
\pgfusepath{clip}%
\pgfsetbuttcap%
\pgfsetroundjoin%
\pgfsetlinewidth{1.003750pt}%
\definecolor{currentstroke}{rgb}{0.690196,0.690196,0.690196}%
\pgfsetstrokecolor{currentstroke}%
\pgfsetstrokeopacity{0.600000}%
\pgfsetdash{{1.000000pt}{1.650000pt}}{0.000000pt}%
\pgfpathmoveto{\pgfqpoint{1.000000in}{1.221792in}}%
\pgfpathlineto{\pgfqpoint{7.200000in}{1.221792in}}%
\pgfusepath{stroke}%
\end{pgfscope}%
\begin{pgfscope}%
\pgfsetbuttcap%
\pgfsetroundjoin%
\definecolor{currentfill}{rgb}{0.000000,0.000000,0.000000}%
\pgfsetfillcolor{currentfill}%
\pgfsetlinewidth{0.602250pt}%
\definecolor{currentstroke}{rgb}{0.000000,0.000000,0.000000}%
\pgfsetstrokecolor{currentstroke}%
\pgfsetdash{}{0pt}%
\pgfsys@defobject{currentmarker}{\pgfqpoint{-0.027778in}{0.000000in}}{\pgfqpoint{-0.000000in}{0.000000in}}{%
\pgfpathmoveto{\pgfqpoint{-0.000000in}{0.000000in}}%
\pgfpathlineto{\pgfqpoint{-0.027778in}{0.000000in}}%
\pgfusepath{stroke,fill}%
}%
\begin{pgfscope}%
\pgfsys@transformshift{1.000000in}{1.221792in}%
\pgfsys@useobject{currentmarker}{}%
\end{pgfscope}%
\end{pgfscope}%
\begin{pgfscope}%
\pgfpathrectangle{\pgfqpoint{1.000000in}{0.440000in}}{\pgfqpoint{6.200000in}{3.080000in}}%
\pgfusepath{clip}%
\pgfsetbuttcap%
\pgfsetroundjoin%
\pgfsetlinewidth{1.003750pt}%
\definecolor{currentstroke}{rgb}{0.690196,0.690196,0.690196}%
\pgfsetstrokecolor{currentstroke}%
\pgfsetstrokeopacity{0.600000}%
\pgfsetdash{{1.000000pt}{1.650000pt}}{0.000000pt}%
\pgfpathmoveto{\pgfqpoint{1.000000in}{1.301977in}}%
\pgfpathlineto{\pgfqpoint{7.200000in}{1.301977in}}%
\pgfusepath{stroke}%
\end{pgfscope}%
\begin{pgfscope}%
\pgfsetbuttcap%
\pgfsetroundjoin%
\definecolor{currentfill}{rgb}{0.000000,0.000000,0.000000}%
\pgfsetfillcolor{currentfill}%
\pgfsetlinewidth{0.602250pt}%
\definecolor{currentstroke}{rgb}{0.000000,0.000000,0.000000}%
\pgfsetstrokecolor{currentstroke}%
\pgfsetdash{}{0pt}%
\pgfsys@defobject{currentmarker}{\pgfqpoint{-0.027778in}{0.000000in}}{\pgfqpoint{-0.000000in}{0.000000in}}{%
\pgfpathmoveto{\pgfqpoint{-0.000000in}{0.000000in}}%
\pgfpathlineto{\pgfqpoint{-0.027778in}{0.000000in}}%
\pgfusepath{stroke,fill}%
}%
\begin{pgfscope}%
\pgfsys@transformshift{1.000000in}{1.301977in}%
\pgfsys@useobject{currentmarker}{}%
\end{pgfscope}%
\end{pgfscope}%
\begin{pgfscope}%
\pgfpathrectangle{\pgfqpoint{1.000000in}{0.440000in}}{\pgfqpoint{6.200000in}{3.080000in}}%
\pgfusepath{clip}%
\pgfsetbuttcap%
\pgfsetroundjoin%
\pgfsetlinewidth{1.003750pt}%
\definecolor{currentstroke}{rgb}{0.690196,0.690196,0.690196}%
\pgfsetstrokecolor{currentstroke}%
\pgfsetstrokeopacity{0.600000}%
\pgfsetdash{{1.000000pt}{1.650000pt}}{0.000000pt}%
\pgfpathmoveto{\pgfqpoint{1.000000in}{1.364173in}}%
\pgfpathlineto{\pgfqpoint{7.200000in}{1.364173in}}%
\pgfusepath{stroke}%
\end{pgfscope}%
\begin{pgfscope}%
\pgfsetbuttcap%
\pgfsetroundjoin%
\definecolor{currentfill}{rgb}{0.000000,0.000000,0.000000}%
\pgfsetfillcolor{currentfill}%
\pgfsetlinewidth{0.602250pt}%
\definecolor{currentstroke}{rgb}{0.000000,0.000000,0.000000}%
\pgfsetstrokecolor{currentstroke}%
\pgfsetdash{}{0pt}%
\pgfsys@defobject{currentmarker}{\pgfqpoint{-0.027778in}{0.000000in}}{\pgfqpoint{-0.000000in}{0.000000in}}{%
\pgfpathmoveto{\pgfqpoint{-0.000000in}{0.000000in}}%
\pgfpathlineto{\pgfqpoint{-0.027778in}{0.000000in}}%
\pgfusepath{stroke,fill}%
}%
\begin{pgfscope}%
\pgfsys@transformshift{1.000000in}{1.364173in}%
\pgfsys@useobject{currentmarker}{}%
\end{pgfscope}%
\end{pgfscope}%
\begin{pgfscope}%
\pgfpathrectangle{\pgfqpoint{1.000000in}{0.440000in}}{\pgfqpoint{6.200000in}{3.080000in}}%
\pgfusepath{clip}%
\pgfsetbuttcap%
\pgfsetroundjoin%
\pgfsetlinewidth{1.003750pt}%
\definecolor{currentstroke}{rgb}{0.690196,0.690196,0.690196}%
\pgfsetstrokecolor{currentstroke}%
\pgfsetstrokeopacity{0.600000}%
\pgfsetdash{{1.000000pt}{1.650000pt}}{0.000000pt}%
\pgfpathmoveto{\pgfqpoint{1.000000in}{1.414991in}}%
\pgfpathlineto{\pgfqpoint{7.200000in}{1.414991in}}%
\pgfusepath{stroke}%
\end{pgfscope}%
\begin{pgfscope}%
\pgfsetbuttcap%
\pgfsetroundjoin%
\definecolor{currentfill}{rgb}{0.000000,0.000000,0.000000}%
\pgfsetfillcolor{currentfill}%
\pgfsetlinewidth{0.602250pt}%
\definecolor{currentstroke}{rgb}{0.000000,0.000000,0.000000}%
\pgfsetstrokecolor{currentstroke}%
\pgfsetdash{}{0pt}%
\pgfsys@defobject{currentmarker}{\pgfqpoint{-0.027778in}{0.000000in}}{\pgfqpoint{-0.000000in}{0.000000in}}{%
\pgfpathmoveto{\pgfqpoint{-0.000000in}{0.000000in}}%
\pgfpathlineto{\pgfqpoint{-0.027778in}{0.000000in}}%
\pgfusepath{stroke,fill}%
}%
\begin{pgfscope}%
\pgfsys@transformshift{1.000000in}{1.414991in}%
\pgfsys@useobject{currentmarker}{}%
\end{pgfscope}%
\end{pgfscope}%
\begin{pgfscope}%
\pgfpathrectangle{\pgfqpoint{1.000000in}{0.440000in}}{\pgfqpoint{6.200000in}{3.080000in}}%
\pgfusepath{clip}%
\pgfsetbuttcap%
\pgfsetroundjoin%
\pgfsetlinewidth{1.003750pt}%
\definecolor{currentstroke}{rgb}{0.690196,0.690196,0.690196}%
\pgfsetstrokecolor{currentstroke}%
\pgfsetstrokeopacity{0.600000}%
\pgfsetdash{{1.000000pt}{1.650000pt}}{0.000000pt}%
\pgfpathmoveto{\pgfqpoint{1.000000in}{1.457957in}}%
\pgfpathlineto{\pgfqpoint{7.200000in}{1.457957in}}%
\pgfusepath{stroke}%
\end{pgfscope}%
\begin{pgfscope}%
\pgfsetbuttcap%
\pgfsetroundjoin%
\definecolor{currentfill}{rgb}{0.000000,0.000000,0.000000}%
\pgfsetfillcolor{currentfill}%
\pgfsetlinewidth{0.602250pt}%
\definecolor{currentstroke}{rgb}{0.000000,0.000000,0.000000}%
\pgfsetstrokecolor{currentstroke}%
\pgfsetdash{}{0pt}%
\pgfsys@defobject{currentmarker}{\pgfqpoint{-0.027778in}{0.000000in}}{\pgfqpoint{-0.000000in}{0.000000in}}{%
\pgfpathmoveto{\pgfqpoint{-0.000000in}{0.000000in}}%
\pgfpathlineto{\pgfqpoint{-0.027778in}{0.000000in}}%
\pgfusepath{stroke,fill}%
}%
\begin{pgfscope}%
\pgfsys@transformshift{1.000000in}{1.457957in}%
\pgfsys@useobject{currentmarker}{}%
\end{pgfscope}%
\end{pgfscope}%
\begin{pgfscope}%
\pgfpathrectangle{\pgfqpoint{1.000000in}{0.440000in}}{\pgfqpoint{6.200000in}{3.080000in}}%
\pgfusepath{clip}%
\pgfsetbuttcap%
\pgfsetroundjoin%
\pgfsetlinewidth{1.003750pt}%
\definecolor{currentstroke}{rgb}{0.690196,0.690196,0.690196}%
\pgfsetstrokecolor{currentstroke}%
\pgfsetstrokeopacity{0.600000}%
\pgfsetdash{{1.000000pt}{1.650000pt}}{0.000000pt}%
\pgfpathmoveto{\pgfqpoint{1.000000in}{1.495175in}}%
\pgfpathlineto{\pgfqpoint{7.200000in}{1.495175in}}%
\pgfusepath{stroke}%
\end{pgfscope}%
\begin{pgfscope}%
\pgfsetbuttcap%
\pgfsetroundjoin%
\definecolor{currentfill}{rgb}{0.000000,0.000000,0.000000}%
\pgfsetfillcolor{currentfill}%
\pgfsetlinewidth{0.602250pt}%
\definecolor{currentstroke}{rgb}{0.000000,0.000000,0.000000}%
\pgfsetstrokecolor{currentstroke}%
\pgfsetdash{}{0pt}%
\pgfsys@defobject{currentmarker}{\pgfqpoint{-0.027778in}{0.000000in}}{\pgfqpoint{-0.000000in}{0.000000in}}{%
\pgfpathmoveto{\pgfqpoint{-0.000000in}{0.000000in}}%
\pgfpathlineto{\pgfqpoint{-0.027778in}{0.000000in}}%
\pgfusepath{stroke,fill}%
}%
\begin{pgfscope}%
\pgfsys@transformshift{1.000000in}{1.495175in}%
\pgfsys@useobject{currentmarker}{}%
\end{pgfscope}%
\end{pgfscope}%
\begin{pgfscope}%
\pgfpathrectangle{\pgfqpoint{1.000000in}{0.440000in}}{\pgfqpoint{6.200000in}{3.080000in}}%
\pgfusepath{clip}%
\pgfsetbuttcap%
\pgfsetroundjoin%
\pgfsetlinewidth{1.003750pt}%
\definecolor{currentstroke}{rgb}{0.690196,0.690196,0.690196}%
\pgfsetstrokecolor{currentstroke}%
\pgfsetstrokeopacity{0.600000}%
\pgfsetdash{{1.000000pt}{1.650000pt}}{0.000000pt}%
\pgfpathmoveto{\pgfqpoint{1.000000in}{1.528005in}}%
\pgfpathlineto{\pgfqpoint{7.200000in}{1.528005in}}%
\pgfusepath{stroke}%
\end{pgfscope}%
\begin{pgfscope}%
\pgfsetbuttcap%
\pgfsetroundjoin%
\definecolor{currentfill}{rgb}{0.000000,0.000000,0.000000}%
\pgfsetfillcolor{currentfill}%
\pgfsetlinewidth{0.602250pt}%
\definecolor{currentstroke}{rgb}{0.000000,0.000000,0.000000}%
\pgfsetstrokecolor{currentstroke}%
\pgfsetdash{}{0pt}%
\pgfsys@defobject{currentmarker}{\pgfqpoint{-0.027778in}{0.000000in}}{\pgfqpoint{-0.000000in}{0.000000in}}{%
\pgfpathmoveto{\pgfqpoint{-0.000000in}{0.000000in}}%
\pgfpathlineto{\pgfqpoint{-0.027778in}{0.000000in}}%
\pgfusepath{stroke,fill}%
}%
\begin{pgfscope}%
\pgfsys@transformshift{1.000000in}{1.528005in}%
\pgfsys@useobject{currentmarker}{}%
\end{pgfscope}%
\end{pgfscope}%
\begin{pgfscope}%
\pgfpathrectangle{\pgfqpoint{1.000000in}{0.440000in}}{\pgfqpoint{6.200000in}{3.080000in}}%
\pgfusepath{clip}%
\pgfsetbuttcap%
\pgfsetroundjoin%
\pgfsetlinewidth{1.003750pt}%
\definecolor{currentstroke}{rgb}{0.690196,0.690196,0.690196}%
\pgfsetstrokecolor{currentstroke}%
\pgfsetstrokeopacity{0.600000}%
\pgfsetdash{{1.000000pt}{1.650000pt}}{0.000000pt}%
\pgfpathmoveto{\pgfqpoint{1.000000in}{1.750570in}}%
\pgfpathlineto{\pgfqpoint{7.200000in}{1.750570in}}%
\pgfusepath{stroke}%
\end{pgfscope}%
\begin{pgfscope}%
\pgfsetbuttcap%
\pgfsetroundjoin%
\definecolor{currentfill}{rgb}{0.000000,0.000000,0.000000}%
\pgfsetfillcolor{currentfill}%
\pgfsetlinewidth{0.602250pt}%
\definecolor{currentstroke}{rgb}{0.000000,0.000000,0.000000}%
\pgfsetstrokecolor{currentstroke}%
\pgfsetdash{}{0pt}%
\pgfsys@defobject{currentmarker}{\pgfqpoint{-0.027778in}{0.000000in}}{\pgfqpoint{-0.000000in}{0.000000in}}{%
\pgfpathmoveto{\pgfqpoint{-0.000000in}{0.000000in}}%
\pgfpathlineto{\pgfqpoint{-0.027778in}{0.000000in}}%
\pgfusepath{stroke,fill}%
}%
\begin{pgfscope}%
\pgfsys@transformshift{1.000000in}{1.750570in}%
\pgfsys@useobject{currentmarker}{}%
\end{pgfscope}%
\end{pgfscope}%
\begin{pgfscope}%
\pgfpathrectangle{\pgfqpoint{1.000000in}{0.440000in}}{\pgfqpoint{6.200000in}{3.080000in}}%
\pgfusepath{clip}%
\pgfsetbuttcap%
\pgfsetroundjoin%
\pgfsetlinewidth{1.003750pt}%
\definecolor{currentstroke}{rgb}{0.690196,0.690196,0.690196}%
\pgfsetstrokecolor{currentstroke}%
\pgfsetstrokeopacity{0.600000}%
\pgfsetdash{{1.000000pt}{1.650000pt}}{0.000000pt}%
\pgfpathmoveto{\pgfqpoint{1.000000in}{1.863584in}}%
\pgfpathlineto{\pgfqpoint{7.200000in}{1.863584in}}%
\pgfusepath{stroke}%
\end{pgfscope}%
\begin{pgfscope}%
\pgfsetbuttcap%
\pgfsetroundjoin%
\definecolor{currentfill}{rgb}{0.000000,0.000000,0.000000}%
\pgfsetfillcolor{currentfill}%
\pgfsetlinewidth{0.602250pt}%
\definecolor{currentstroke}{rgb}{0.000000,0.000000,0.000000}%
\pgfsetstrokecolor{currentstroke}%
\pgfsetdash{}{0pt}%
\pgfsys@defobject{currentmarker}{\pgfqpoint{-0.027778in}{0.000000in}}{\pgfqpoint{-0.000000in}{0.000000in}}{%
\pgfpathmoveto{\pgfqpoint{-0.000000in}{0.000000in}}%
\pgfpathlineto{\pgfqpoint{-0.027778in}{0.000000in}}%
\pgfusepath{stroke,fill}%
}%
\begin{pgfscope}%
\pgfsys@transformshift{1.000000in}{1.863584in}%
\pgfsys@useobject{currentmarker}{}%
\end{pgfscope}%
\end{pgfscope}%
\begin{pgfscope}%
\pgfpathrectangle{\pgfqpoint{1.000000in}{0.440000in}}{\pgfqpoint{6.200000in}{3.080000in}}%
\pgfusepath{clip}%
\pgfsetbuttcap%
\pgfsetroundjoin%
\pgfsetlinewidth{1.003750pt}%
\definecolor{currentstroke}{rgb}{0.690196,0.690196,0.690196}%
\pgfsetstrokecolor{currentstroke}%
\pgfsetstrokeopacity{0.600000}%
\pgfsetdash{{1.000000pt}{1.650000pt}}{0.000000pt}%
\pgfpathmoveto{\pgfqpoint{1.000000in}{1.943769in}}%
\pgfpathlineto{\pgfqpoint{7.200000in}{1.943769in}}%
\pgfusepath{stroke}%
\end{pgfscope}%
\begin{pgfscope}%
\pgfsetbuttcap%
\pgfsetroundjoin%
\definecolor{currentfill}{rgb}{0.000000,0.000000,0.000000}%
\pgfsetfillcolor{currentfill}%
\pgfsetlinewidth{0.602250pt}%
\definecolor{currentstroke}{rgb}{0.000000,0.000000,0.000000}%
\pgfsetstrokecolor{currentstroke}%
\pgfsetdash{}{0pt}%
\pgfsys@defobject{currentmarker}{\pgfqpoint{-0.027778in}{0.000000in}}{\pgfqpoint{-0.000000in}{0.000000in}}{%
\pgfpathmoveto{\pgfqpoint{-0.000000in}{0.000000in}}%
\pgfpathlineto{\pgfqpoint{-0.027778in}{0.000000in}}%
\pgfusepath{stroke,fill}%
}%
\begin{pgfscope}%
\pgfsys@transformshift{1.000000in}{1.943769in}%
\pgfsys@useobject{currentmarker}{}%
\end{pgfscope}%
\end{pgfscope}%
\begin{pgfscope}%
\pgfpathrectangle{\pgfqpoint{1.000000in}{0.440000in}}{\pgfqpoint{6.200000in}{3.080000in}}%
\pgfusepath{clip}%
\pgfsetbuttcap%
\pgfsetroundjoin%
\pgfsetlinewidth{1.003750pt}%
\definecolor{currentstroke}{rgb}{0.690196,0.690196,0.690196}%
\pgfsetstrokecolor{currentstroke}%
\pgfsetstrokeopacity{0.600000}%
\pgfsetdash{{1.000000pt}{1.650000pt}}{0.000000pt}%
\pgfpathmoveto{\pgfqpoint{1.000000in}{2.005965in}}%
\pgfpathlineto{\pgfqpoint{7.200000in}{2.005965in}}%
\pgfusepath{stroke}%
\end{pgfscope}%
\begin{pgfscope}%
\pgfsetbuttcap%
\pgfsetroundjoin%
\definecolor{currentfill}{rgb}{0.000000,0.000000,0.000000}%
\pgfsetfillcolor{currentfill}%
\pgfsetlinewidth{0.602250pt}%
\definecolor{currentstroke}{rgb}{0.000000,0.000000,0.000000}%
\pgfsetstrokecolor{currentstroke}%
\pgfsetdash{}{0pt}%
\pgfsys@defobject{currentmarker}{\pgfqpoint{-0.027778in}{0.000000in}}{\pgfqpoint{-0.000000in}{0.000000in}}{%
\pgfpathmoveto{\pgfqpoint{-0.000000in}{0.000000in}}%
\pgfpathlineto{\pgfqpoint{-0.027778in}{0.000000in}}%
\pgfusepath{stroke,fill}%
}%
\begin{pgfscope}%
\pgfsys@transformshift{1.000000in}{2.005965in}%
\pgfsys@useobject{currentmarker}{}%
\end{pgfscope}%
\end{pgfscope}%
\begin{pgfscope}%
\pgfpathrectangle{\pgfqpoint{1.000000in}{0.440000in}}{\pgfqpoint{6.200000in}{3.080000in}}%
\pgfusepath{clip}%
\pgfsetbuttcap%
\pgfsetroundjoin%
\pgfsetlinewidth{1.003750pt}%
\definecolor{currentstroke}{rgb}{0.690196,0.690196,0.690196}%
\pgfsetstrokecolor{currentstroke}%
\pgfsetstrokeopacity{0.600000}%
\pgfsetdash{{1.000000pt}{1.650000pt}}{0.000000pt}%
\pgfpathmoveto{\pgfqpoint{1.000000in}{2.056783in}}%
\pgfpathlineto{\pgfqpoint{7.200000in}{2.056783in}}%
\pgfusepath{stroke}%
\end{pgfscope}%
\begin{pgfscope}%
\pgfsetbuttcap%
\pgfsetroundjoin%
\definecolor{currentfill}{rgb}{0.000000,0.000000,0.000000}%
\pgfsetfillcolor{currentfill}%
\pgfsetlinewidth{0.602250pt}%
\definecolor{currentstroke}{rgb}{0.000000,0.000000,0.000000}%
\pgfsetstrokecolor{currentstroke}%
\pgfsetdash{}{0pt}%
\pgfsys@defobject{currentmarker}{\pgfqpoint{-0.027778in}{0.000000in}}{\pgfqpoint{-0.000000in}{0.000000in}}{%
\pgfpathmoveto{\pgfqpoint{-0.000000in}{0.000000in}}%
\pgfpathlineto{\pgfqpoint{-0.027778in}{0.000000in}}%
\pgfusepath{stroke,fill}%
}%
\begin{pgfscope}%
\pgfsys@transformshift{1.000000in}{2.056783in}%
\pgfsys@useobject{currentmarker}{}%
\end{pgfscope}%
\end{pgfscope}%
\begin{pgfscope}%
\pgfpathrectangle{\pgfqpoint{1.000000in}{0.440000in}}{\pgfqpoint{6.200000in}{3.080000in}}%
\pgfusepath{clip}%
\pgfsetbuttcap%
\pgfsetroundjoin%
\pgfsetlinewidth{1.003750pt}%
\definecolor{currentstroke}{rgb}{0.690196,0.690196,0.690196}%
\pgfsetstrokecolor{currentstroke}%
\pgfsetstrokeopacity{0.600000}%
\pgfsetdash{{1.000000pt}{1.650000pt}}{0.000000pt}%
\pgfpathmoveto{\pgfqpoint{1.000000in}{2.099749in}}%
\pgfpathlineto{\pgfqpoint{7.200000in}{2.099749in}}%
\pgfusepath{stroke}%
\end{pgfscope}%
\begin{pgfscope}%
\pgfsetbuttcap%
\pgfsetroundjoin%
\definecolor{currentfill}{rgb}{0.000000,0.000000,0.000000}%
\pgfsetfillcolor{currentfill}%
\pgfsetlinewidth{0.602250pt}%
\definecolor{currentstroke}{rgb}{0.000000,0.000000,0.000000}%
\pgfsetstrokecolor{currentstroke}%
\pgfsetdash{}{0pt}%
\pgfsys@defobject{currentmarker}{\pgfqpoint{-0.027778in}{0.000000in}}{\pgfqpoint{-0.000000in}{0.000000in}}{%
\pgfpathmoveto{\pgfqpoint{-0.000000in}{0.000000in}}%
\pgfpathlineto{\pgfqpoint{-0.027778in}{0.000000in}}%
\pgfusepath{stroke,fill}%
}%
\begin{pgfscope}%
\pgfsys@transformshift{1.000000in}{2.099749in}%
\pgfsys@useobject{currentmarker}{}%
\end{pgfscope}%
\end{pgfscope}%
\begin{pgfscope}%
\pgfpathrectangle{\pgfqpoint{1.000000in}{0.440000in}}{\pgfqpoint{6.200000in}{3.080000in}}%
\pgfusepath{clip}%
\pgfsetbuttcap%
\pgfsetroundjoin%
\pgfsetlinewidth{1.003750pt}%
\definecolor{currentstroke}{rgb}{0.690196,0.690196,0.690196}%
\pgfsetstrokecolor{currentstroke}%
\pgfsetstrokeopacity{0.600000}%
\pgfsetdash{{1.000000pt}{1.650000pt}}{0.000000pt}%
\pgfpathmoveto{\pgfqpoint{1.000000in}{2.136967in}}%
\pgfpathlineto{\pgfqpoint{7.200000in}{2.136967in}}%
\pgfusepath{stroke}%
\end{pgfscope}%
\begin{pgfscope}%
\pgfsetbuttcap%
\pgfsetroundjoin%
\definecolor{currentfill}{rgb}{0.000000,0.000000,0.000000}%
\pgfsetfillcolor{currentfill}%
\pgfsetlinewidth{0.602250pt}%
\definecolor{currentstroke}{rgb}{0.000000,0.000000,0.000000}%
\pgfsetstrokecolor{currentstroke}%
\pgfsetdash{}{0pt}%
\pgfsys@defobject{currentmarker}{\pgfqpoint{-0.027778in}{0.000000in}}{\pgfqpoint{-0.000000in}{0.000000in}}{%
\pgfpathmoveto{\pgfqpoint{-0.000000in}{0.000000in}}%
\pgfpathlineto{\pgfqpoint{-0.027778in}{0.000000in}}%
\pgfusepath{stroke,fill}%
}%
\begin{pgfscope}%
\pgfsys@transformshift{1.000000in}{2.136967in}%
\pgfsys@useobject{currentmarker}{}%
\end{pgfscope}%
\end{pgfscope}%
\begin{pgfscope}%
\pgfpathrectangle{\pgfqpoint{1.000000in}{0.440000in}}{\pgfqpoint{6.200000in}{3.080000in}}%
\pgfusepath{clip}%
\pgfsetbuttcap%
\pgfsetroundjoin%
\pgfsetlinewidth{1.003750pt}%
\definecolor{currentstroke}{rgb}{0.690196,0.690196,0.690196}%
\pgfsetstrokecolor{currentstroke}%
\pgfsetstrokeopacity{0.600000}%
\pgfsetdash{{1.000000pt}{1.650000pt}}{0.000000pt}%
\pgfpathmoveto{\pgfqpoint{1.000000in}{2.169797in}}%
\pgfpathlineto{\pgfqpoint{7.200000in}{2.169797in}}%
\pgfusepath{stroke}%
\end{pgfscope}%
\begin{pgfscope}%
\pgfsetbuttcap%
\pgfsetroundjoin%
\definecolor{currentfill}{rgb}{0.000000,0.000000,0.000000}%
\pgfsetfillcolor{currentfill}%
\pgfsetlinewidth{0.602250pt}%
\definecolor{currentstroke}{rgb}{0.000000,0.000000,0.000000}%
\pgfsetstrokecolor{currentstroke}%
\pgfsetdash{}{0pt}%
\pgfsys@defobject{currentmarker}{\pgfqpoint{-0.027778in}{0.000000in}}{\pgfqpoint{-0.000000in}{0.000000in}}{%
\pgfpathmoveto{\pgfqpoint{-0.000000in}{0.000000in}}%
\pgfpathlineto{\pgfqpoint{-0.027778in}{0.000000in}}%
\pgfusepath{stroke,fill}%
}%
\begin{pgfscope}%
\pgfsys@transformshift{1.000000in}{2.169797in}%
\pgfsys@useobject{currentmarker}{}%
\end{pgfscope}%
\end{pgfscope}%
\begin{pgfscope}%
\pgfpathrectangle{\pgfqpoint{1.000000in}{0.440000in}}{\pgfqpoint{6.200000in}{3.080000in}}%
\pgfusepath{clip}%
\pgfsetbuttcap%
\pgfsetroundjoin%
\pgfsetlinewidth{1.003750pt}%
\definecolor{currentstroke}{rgb}{0.690196,0.690196,0.690196}%
\pgfsetstrokecolor{currentstroke}%
\pgfsetstrokeopacity{0.600000}%
\pgfsetdash{{1.000000pt}{1.650000pt}}{0.000000pt}%
\pgfpathmoveto{\pgfqpoint{1.000000in}{2.392362in}}%
\pgfpathlineto{\pgfqpoint{7.200000in}{2.392362in}}%
\pgfusepath{stroke}%
\end{pgfscope}%
\begin{pgfscope}%
\pgfsetbuttcap%
\pgfsetroundjoin%
\definecolor{currentfill}{rgb}{0.000000,0.000000,0.000000}%
\pgfsetfillcolor{currentfill}%
\pgfsetlinewidth{0.602250pt}%
\definecolor{currentstroke}{rgb}{0.000000,0.000000,0.000000}%
\pgfsetstrokecolor{currentstroke}%
\pgfsetdash{}{0pt}%
\pgfsys@defobject{currentmarker}{\pgfqpoint{-0.027778in}{0.000000in}}{\pgfqpoint{-0.000000in}{0.000000in}}{%
\pgfpathmoveto{\pgfqpoint{-0.000000in}{0.000000in}}%
\pgfpathlineto{\pgfqpoint{-0.027778in}{0.000000in}}%
\pgfusepath{stroke,fill}%
}%
\begin{pgfscope}%
\pgfsys@transformshift{1.000000in}{2.392362in}%
\pgfsys@useobject{currentmarker}{}%
\end{pgfscope}%
\end{pgfscope}%
\begin{pgfscope}%
\pgfpathrectangle{\pgfqpoint{1.000000in}{0.440000in}}{\pgfqpoint{6.200000in}{3.080000in}}%
\pgfusepath{clip}%
\pgfsetbuttcap%
\pgfsetroundjoin%
\pgfsetlinewidth{1.003750pt}%
\definecolor{currentstroke}{rgb}{0.690196,0.690196,0.690196}%
\pgfsetstrokecolor{currentstroke}%
\pgfsetstrokeopacity{0.600000}%
\pgfsetdash{{1.000000pt}{1.650000pt}}{0.000000pt}%
\pgfpathmoveto{\pgfqpoint{1.000000in}{2.505376in}}%
\pgfpathlineto{\pgfqpoint{7.200000in}{2.505376in}}%
\pgfusepath{stroke}%
\end{pgfscope}%
\begin{pgfscope}%
\pgfsetbuttcap%
\pgfsetroundjoin%
\definecolor{currentfill}{rgb}{0.000000,0.000000,0.000000}%
\pgfsetfillcolor{currentfill}%
\pgfsetlinewidth{0.602250pt}%
\definecolor{currentstroke}{rgb}{0.000000,0.000000,0.000000}%
\pgfsetstrokecolor{currentstroke}%
\pgfsetdash{}{0pt}%
\pgfsys@defobject{currentmarker}{\pgfqpoint{-0.027778in}{0.000000in}}{\pgfqpoint{-0.000000in}{0.000000in}}{%
\pgfpathmoveto{\pgfqpoint{-0.000000in}{0.000000in}}%
\pgfpathlineto{\pgfqpoint{-0.027778in}{0.000000in}}%
\pgfusepath{stroke,fill}%
}%
\begin{pgfscope}%
\pgfsys@transformshift{1.000000in}{2.505376in}%
\pgfsys@useobject{currentmarker}{}%
\end{pgfscope}%
\end{pgfscope}%
\begin{pgfscope}%
\pgfpathrectangle{\pgfqpoint{1.000000in}{0.440000in}}{\pgfqpoint{6.200000in}{3.080000in}}%
\pgfusepath{clip}%
\pgfsetbuttcap%
\pgfsetroundjoin%
\pgfsetlinewidth{1.003750pt}%
\definecolor{currentstroke}{rgb}{0.690196,0.690196,0.690196}%
\pgfsetstrokecolor{currentstroke}%
\pgfsetstrokeopacity{0.600000}%
\pgfsetdash{{1.000000pt}{1.650000pt}}{0.000000pt}%
\pgfpathmoveto{\pgfqpoint{1.000000in}{2.585561in}}%
\pgfpathlineto{\pgfqpoint{7.200000in}{2.585561in}}%
\pgfusepath{stroke}%
\end{pgfscope}%
\begin{pgfscope}%
\pgfsetbuttcap%
\pgfsetroundjoin%
\definecolor{currentfill}{rgb}{0.000000,0.000000,0.000000}%
\pgfsetfillcolor{currentfill}%
\pgfsetlinewidth{0.602250pt}%
\definecolor{currentstroke}{rgb}{0.000000,0.000000,0.000000}%
\pgfsetstrokecolor{currentstroke}%
\pgfsetdash{}{0pt}%
\pgfsys@defobject{currentmarker}{\pgfqpoint{-0.027778in}{0.000000in}}{\pgfqpoint{-0.000000in}{0.000000in}}{%
\pgfpathmoveto{\pgfqpoint{-0.000000in}{0.000000in}}%
\pgfpathlineto{\pgfqpoint{-0.027778in}{0.000000in}}%
\pgfusepath{stroke,fill}%
}%
\begin{pgfscope}%
\pgfsys@transformshift{1.000000in}{2.585561in}%
\pgfsys@useobject{currentmarker}{}%
\end{pgfscope}%
\end{pgfscope}%
\begin{pgfscope}%
\pgfpathrectangle{\pgfqpoint{1.000000in}{0.440000in}}{\pgfqpoint{6.200000in}{3.080000in}}%
\pgfusepath{clip}%
\pgfsetbuttcap%
\pgfsetroundjoin%
\pgfsetlinewidth{1.003750pt}%
\definecolor{currentstroke}{rgb}{0.690196,0.690196,0.690196}%
\pgfsetstrokecolor{currentstroke}%
\pgfsetstrokeopacity{0.600000}%
\pgfsetdash{{1.000000pt}{1.650000pt}}{0.000000pt}%
\pgfpathmoveto{\pgfqpoint{1.000000in}{2.647757in}}%
\pgfpathlineto{\pgfqpoint{7.200000in}{2.647757in}}%
\pgfusepath{stroke}%
\end{pgfscope}%
\begin{pgfscope}%
\pgfsetbuttcap%
\pgfsetroundjoin%
\definecolor{currentfill}{rgb}{0.000000,0.000000,0.000000}%
\pgfsetfillcolor{currentfill}%
\pgfsetlinewidth{0.602250pt}%
\definecolor{currentstroke}{rgb}{0.000000,0.000000,0.000000}%
\pgfsetstrokecolor{currentstroke}%
\pgfsetdash{}{0pt}%
\pgfsys@defobject{currentmarker}{\pgfqpoint{-0.027778in}{0.000000in}}{\pgfqpoint{-0.000000in}{0.000000in}}{%
\pgfpathmoveto{\pgfqpoint{-0.000000in}{0.000000in}}%
\pgfpathlineto{\pgfqpoint{-0.027778in}{0.000000in}}%
\pgfusepath{stroke,fill}%
}%
\begin{pgfscope}%
\pgfsys@transformshift{1.000000in}{2.647757in}%
\pgfsys@useobject{currentmarker}{}%
\end{pgfscope}%
\end{pgfscope}%
\begin{pgfscope}%
\pgfpathrectangle{\pgfqpoint{1.000000in}{0.440000in}}{\pgfqpoint{6.200000in}{3.080000in}}%
\pgfusepath{clip}%
\pgfsetbuttcap%
\pgfsetroundjoin%
\pgfsetlinewidth{1.003750pt}%
\definecolor{currentstroke}{rgb}{0.690196,0.690196,0.690196}%
\pgfsetstrokecolor{currentstroke}%
\pgfsetstrokeopacity{0.600000}%
\pgfsetdash{{1.000000pt}{1.650000pt}}{0.000000pt}%
\pgfpathmoveto{\pgfqpoint{1.000000in}{2.698575in}}%
\pgfpathlineto{\pgfqpoint{7.200000in}{2.698575in}}%
\pgfusepath{stroke}%
\end{pgfscope}%
\begin{pgfscope}%
\pgfsetbuttcap%
\pgfsetroundjoin%
\definecolor{currentfill}{rgb}{0.000000,0.000000,0.000000}%
\pgfsetfillcolor{currentfill}%
\pgfsetlinewidth{0.602250pt}%
\definecolor{currentstroke}{rgb}{0.000000,0.000000,0.000000}%
\pgfsetstrokecolor{currentstroke}%
\pgfsetdash{}{0pt}%
\pgfsys@defobject{currentmarker}{\pgfqpoint{-0.027778in}{0.000000in}}{\pgfqpoint{-0.000000in}{0.000000in}}{%
\pgfpathmoveto{\pgfqpoint{-0.000000in}{0.000000in}}%
\pgfpathlineto{\pgfqpoint{-0.027778in}{0.000000in}}%
\pgfusepath{stroke,fill}%
}%
\begin{pgfscope}%
\pgfsys@transformshift{1.000000in}{2.698575in}%
\pgfsys@useobject{currentmarker}{}%
\end{pgfscope}%
\end{pgfscope}%
\begin{pgfscope}%
\pgfpathrectangle{\pgfqpoint{1.000000in}{0.440000in}}{\pgfqpoint{6.200000in}{3.080000in}}%
\pgfusepath{clip}%
\pgfsetbuttcap%
\pgfsetroundjoin%
\pgfsetlinewidth{1.003750pt}%
\definecolor{currentstroke}{rgb}{0.690196,0.690196,0.690196}%
\pgfsetstrokecolor{currentstroke}%
\pgfsetstrokeopacity{0.600000}%
\pgfsetdash{{1.000000pt}{1.650000pt}}{0.000000pt}%
\pgfpathmoveto{\pgfqpoint{1.000000in}{2.741541in}}%
\pgfpathlineto{\pgfqpoint{7.200000in}{2.741541in}}%
\pgfusepath{stroke}%
\end{pgfscope}%
\begin{pgfscope}%
\pgfsetbuttcap%
\pgfsetroundjoin%
\definecolor{currentfill}{rgb}{0.000000,0.000000,0.000000}%
\pgfsetfillcolor{currentfill}%
\pgfsetlinewidth{0.602250pt}%
\definecolor{currentstroke}{rgb}{0.000000,0.000000,0.000000}%
\pgfsetstrokecolor{currentstroke}%
\pgfsetdash{}{0pt}%
\pgfsys@defobject{currentmarker}{\pgfqpoint{-0.027778in}{0.000000in}}{\pgfqpoint{-0.000000in}{0.000000in}}{%
\pgfpathmoveto{\pgfqpoint{-0.000000in}{0.000000in}}%
\pgfpathlineto{\pgfqpoint{-0.027778in}{0.000000in}}%
\pgfusepath{stroke,fill}%
}%
\begin{pgfscope}%
\pgfsys@transformshift{1.000000in}{2.741541in}%
\pgfsys@useobject{currentmarker}{}%
\end{pgfscope}%
\end{pgfscope}%
\begin{pgfscope}%
\pgfpathrectangle{\pgfqpoint{1.000000in}{0.440000in}}{\pgfqpoint{6.200000in}{3.080000in}}%
\pgfusepath{clip}%
\pgfsetbuttcap%
\pgfsetroundjoin%
\pgfsetlinewidth{1.003750pt}%
\definecolor{currentstroke}{rgb}{0.690196,0.690196,0.690196}%
\pgfsetstrokecolor{currentstroke}%
\pgfsetstrokeopacity{0.600000}%
\pgfsetdash{{1.000000pt}{1.650000pt}}{0.000000pt}%
\pgfpathmoveto{\pgfqpoint{1.000000in}{2.778759in}}%
\pgfpathlineto{\pgfqpoint{7.200000in}{2.778759in}}%
\pgfusepath{stroke}%
\end{pgfscope}%
\begin{pgfscope}%
\pgfsetbuttcap%
\pgfsetroundjoin%
\definecolor{currentfill}{rgb}{0.000000,0.000000,0.000000}%
\pgfsetfillcolor{currentfill}%
\pgfsetlinewidth{0.602250pt}%
\definecolor{currentstroke}{rgb}{0.000000,0.000000,0.000000}%
\pgfsetstrokecolor{currentstroke}%
\pgfsetdash{}{0pt}%
\pgfsys@defobject{currentmarker}{\pgfqpoint{-0.027778in}{0.000000in}}{\pgfqpoint{-0.000000in}{0.000000in}}{%
\pgfpathmoveto{\pgfqpoint{-0.000000in}{0.000000in}}%
\pgfpathlineto{\pgfqpoint{-0.027778in}{0.000000in}}%
\pgfusepath{stroke,fill}%
}%
\begin{pgfscope}%
\pgfsys@transformshift{1.000000in}{2.778759in}%
\pgfsys@useobject{currentmarker}{}%
\end{pgfscope}%
\end{pgfscope}%
\begin{pgfscope}%
\pgfpathrectangle{\pgfqpoint{1.000000in}{0.440000in}}{\pgfqpoint{6.200000in}{3.080000in}}%
\pgfusepath{clip}%
\pgfsetbuttcap%
\pgfsetroundjoin%
\pgfsetlinewidth{1.003750pt}%
\definecolor{currentstroke}{rgb}{0.690196,0.690196,0.690196}%
\pgfsetstrokecolor{currentstroke}%
\pgfsetstrokeopacity{0.600000}%
\pgfsetdash{{1.000000pt}{1.650000pt}}{0.000000pt}%
\pgfpathmoveto{\pgfqpoint{1.000000in}{2.811589in}}%
\pgfpathlineto{\pgfqpoint{7.200000in}{2.811589in}}%
\pgfusepath{stroke}%
\end{pgfscope}%
\begin{pgfscope}%
\pgfsetbuttcap%
\pgfsetroundjoin%
\definecolor{currentfill}{rgb}{0.000000,0.000000,0.000000}%
\pgfsetfillcolor{currentfill}%
\pgfsetlinewidth{0.602250pt}%
\definecolor{currentstroke}{rgb}{0.000000,0.000000,0.000000}%
\pgfsetstrokecolor{currentstroke}%
\pgfsetdash{}{0pt}%
\pgfsys@defobject{currentmarker}{\pgfqpoint{-0.027778in}{0.000000in}}{\pgfqpoint{-0.000000in}{0.000000in}}{%
\pgfpathmoveto{\pgfqpoint{-0.000000in}{0.000000in}}%
\pgfpathlineto{\pgfqpoint{-0.027778in}{0.000000in}}%
\pgfusepath{stroke,fill}%
}%
\begin{pgfscope}%
\pgfsys@transformshift{1.000000in}{2.811589in}%
\pgfsys@useobject{currentmarker}{}%
\end{pgfscope}%
\end{pgfscope}%
\begin{pgfscope}%
\pgfpathrectangle{\pgfqpoint{1.000000in}{0.440000in}}{\pgfqpoint{6.200000in}{3.080000in}}%
\pgfusepath{clip}%
\pgfsetbuttcap%
\pgfsetroundjoin%
\pgfsetlinewidth{1.003750pt}%
\definecolor{currentstroke}{rgb}{0.690196,0.690196,0.690196}%
\pgfsetstrokecolor{currentstroke}%
\pgfsetstrokeopacity{0.600000}%
\pgfsetdash{{1.000000pt}{1.650000pt}}{0.000000pt}%
\pgfpathmoveto{\pgfqpoint{1.000000in}{3.034154in}}%
\pgfpathlineto{\pgfqpoint{7.200000in}{3.034154in}}%
\pgfusepath{stroke}%
\end{pgfscope}%
\begin{pgfscope}%
\pgfsetbuttcap%
\pgfsetroundjoin%
\definecolor{currentfill}{rgb}{0.000000,0.000000,0.000000}%
\pgfsetfillcolor{currentfill}%
\pgfsetlinewidth{0.602250pt}%
\definecolor{currentstroke}{rgb}{0.000000,0.000000,0.000000}%
\pgfsetstrokecolor{currentstroke}%
\pgfsetdash{}{0pt}%
\pgfsys@defobject{currentmarker}{\pgfqpoint{-0.027778in}{0.000000in}}{\pgfqpoint{-0.000000in}{0.000000in}}{%
\pgfpathmoveto{\pgfqpoint{-0.000000in}{0.000000in}}%
\pgfpathlineto{\pgfqpoint{-0.027778in}{0.000000in}}%
\pgfusepath{stroke,fill}%
}%
\begin{pgfscope}%
\pgfsys@transformshift{1.000000in}{3.034154in}%
\pgfsys@useobject{currentmarker}{}%
\end{pgfscope}%
\end{pgfscope}%
\begin{pgfscope}%
\pgfpathrectangle{\pgfqpoint{1.000000in}{0.440000in}}{\pgfqpoint{6.200000in}{3.080000in}}%
\pgfusepath{clip}%
\pgfsetbuttcap%
\pgfsetroundjoin%
\pgfsetlinewidth{1.003750pt}%
\definecolor{currentstroke}{rgb}{0.690196,0.690196,0.690196}%
\pgfsetstrokecolor{currentstroke}%
\pgfsetstrokeopacity{0.600000}%
\pgfsetdash{{1.000000pt}{1.650000pt}}{0.000000pt}%
\pgfpathmoveto{\pgfqpoint{1.000000in}{3.147168in}}%
\pgfpathlineto{\pgfqpoint{7.200000in}{3.147168in}}%
\pgfusepath{stroke}%
\end{pgfscope}%
\begin{pgfscope}%
\pgfsetbuttcap%
\pgfsetroundjoin%
\definecolor{currentfill}{rgb}{0.000000,0.000000,0.000000}%
\pgfsetfillcolor{currentfill}%
\pgfsetlinewidth{0.602250pt}%
\definecolor{currentstroke}{rgb}{0.000000,0.000000,0.000000}%
\pgfsetstrokecolor{currentstroke}%
\pgfsetdash{}{0pt}%
\pgfsys@defobject{currentmarker}{\pgfqpoint{-0.027778in}{0.000000in}}{\pgfqpoint{-0.000000in}{0.000000in}}{%
\pgfpathmoveto{\pgfqpoint{-0.000000in}{0.000000in}}%
\pgfpathlineto{\pgfqpoint{-0.027778in}{0.000000in}}%
\pgfusepath{stroke,fill}%
}%
\begin{pgfscope}%
\pgfsys@transformshift{1.000000in}{3.147168in}%
\pgfsys@useobject{currentmarker}{}%
\end{pgfscope}%
\end{pgfscope}%
\begin{pgfscope}%
\pgfpathrectangle{\pgfqpoint{1.000000in}{0.440000in}}{\pgfqpoint{6.200000in}{3.080000in}}%
\pgfusepath{clip}%
\pgfsetbuttcap%
\pgfsetroundjoin%
\pgfsetlinewidth{1.003750pt}%
\definecolor{currentstroke}{rgb}{0.690196,0.690196,0.690196}%
\pgfsetstrokecolor{currentstroke}%
\pgfsetstrokeopacity{0.600000}%
\pgfsetdash{{1.000000pt}{1.650000pt}}{0.000000pt}%
\pgfpathmoveto{\pgfqpoint{1.000000in}{3.227353in}}%
\pgfpathlineto{\pgfqpoint{7.200000in}{3.227353in}}%
\pgfusepath{stroke}%
\end{pgfscope}%
\begin{pgfscope}%
\pgfsetbuttcap%
\pgfsetroundjoin%
\definecolor{currentfill}{rgb}{0.000000,0.000000,0.000000}%
\pgfsetfillcolor{currentfill}%
\pgfsetlinewidth{0.602250pt}%
\definecolor{currentstroke}{rgb}{0.000000,0.000000,0.000000}%
\pgfsetstrokecolor{currentstroke}%
\pgfsetdash{}{0pt}%
\pgfsys@defobject{currentmarker}{\pgfqpoint{-0.027778in}{0.000000in}}{\pgfqpoint{-0.000000in}{0.000000in}}{%
\pgfpathmoveto{\pgfqpoint{-0.000000in}{0.000000in}}%
\pgfpathlineto{\pgfqpoint{-0.027778in}{0.000000in}}%
\pgfusepath{stroke,fill}%
}%
\begin{pgfscope}%
\pgfsys@transformshift{1.000000in}{3.227353in}%
\pgfsys@useobject{currentmarker}{}%
\end{pgfscope}%
\end{pgfscope}%
\begin{pgfscope}%
\pgfpathrectangle{\pgfqpoint{1.000000in}{0.440000in}}{\pgfqpoint{6.200000in}{3.080000in}}%
\pgfusepath{clip}%
\pgfsetbuttcap%
\pgfsetroundjoin%
\pgfsetlinewidth{1.003750pt}%
\definecolor{currentstroke}{rgb}{0.690196,0.690196,0.690196}%
\pgfsetstrokecolor{currentstroke}%
\pgfsetstrokeopacity{0.600000}%
\pgfsetdash{{1.000000pt}{1.650000pt}}{0.000000pt}%
\pgfpathmoveto{\pgfqpoint{1.000000in}{3.289549in}}%
\pgfpathlineto{\pgfqpoint{7.200000in}{3.289549in}}%
\pgfusepath{stroke}%
\end{pgfscope}%
\begin{pgfscope}%
\pgfsetbuttcap%
\pgfsetroundjoin%
\definecolor{currentfill}{rgb}{0.000000,0.000000,0.000000}%
\pgfsetfillcolor{currentfill}%
\pgfsetlinewidth{0.602250pt}%
\definecolor{currentstroke}{rgb}{0.000000,0.000000,0.000000}%
\pgfsetstrokecolor{currentstroke}%
\pgfsetdash{}{0pt}%
\pgfsys@defobject{currentmarker}{\pgfqpoint{-0.027778in}{0.000000in}}{\pgfqpoint{-0.000000in}{0.000000in}}{%
\pgfpathmoveto{\pgfqpoint{-0.000000in}{0.000000in}}%
\pgfpathlineto{\pgfqpoint{-0.027778in}{0.000000in}}%
\pgfusepath{stroke,fill}%
}%
\begin{pgfscope}%
\pgfsys@transformshift{1.000000in}{3.289549in}%
\pgfsys@useobject{currentmarker}{}%
\end{pgfscope}%
\end{pgfscope}%
\begin{pgfscope}%
\pgfpathrectangle{\pgfqpoint{1.000000in}{0.440000in}}{\pgfqpoint{6.200000in}{3.080000in}}%
\pgfusepath{clip}%
\pgfsetbuttcap%
\pgfsetroundjoin%
\pgfsetlinewidth{1.003750pt}%
\definecolor{currentstroke}{rgb}{0.690196,0.690196,0.690196}%
\pgfsetstrokecolor{currentstroke}%
\pgfsetstrokeopacity{0.600000}%
\pgfsetdash{{1.000000pt}{1.650000pt}}{0.000000pt}%
\pgfpathmoveto{\pgfqpoint{1.000000in}{3.340367in}}%
\pgfpathlineto{\pgfqpoint{7.200000in}{3.340367in}}%
\pgfusepath{stroke}%
\end{pgfscope}%
\begin{pgfscope}%
\pgfsetbuttcap%
\pgfsetroundjoin%
\definecolor{currentfill}{rgb}{0.000000,0.000000,0.000000}%
\pgfsetfillcolor{currentfill}%
\pgfsetlinewidth{0.602250pt}%
\definecolor{currentstroke}{rgb}{0.000000,0.000000,0.000000}%
\pgfsetstrokecolor{currentstroke}%
\pgfsetdash{}{0pt}%
\pgfsys@defobject{currentmarker}{\pgfqpoint{-0.027778in}{0.000000in}}{\pgfqpoint{-0.000000in}{0.000000in}}{%
\pgfpathmoveto{\pgfqpoint{-0.000000in}{0.000000in}}%
\pgfpathlineto{\pgfqpoint{-0.027778in}{0.000000in}}%
\pgfusepath{stroke,fill}%
}%
\begin{pgfscope}%
\pgfsys@transformshift{1.000000in}{3.340367in}%
\pgfsys@useobject{currentmarker}{}%
\end{pgfscope}%
\end{pgfscope}%
\begin{pgfscope}%
\pgfpathrectangle{\pgfqpoint{1.000000in}{0.440000in}}{\pgfqpoint{6.200000in}{3.080000in}}%
\pgfusepath{clip}%
\pgfsetbuttcap%
\pgfsetroundjoin%
\pgfsetlinewidth{1.003750pt}%
\definecolor{currentstroke}{rgb}{0.690196,0.690196,0.690196}%
\pgfsetstrokecolor{currentstroke}%
\pgfsetstrokeopacity{0.600000}%
\pgfsetdash{{1.000000pt}{1.650000pt}}{0.000000pt}%
\pgfpathmoveto{\pgfqpoint{1.000000in}{3.383333in}}%
\pgfpathlineto{\pgfqpoint{7.200000in}{3.383333in}}%
\pgfusepath{stroke}%
\end{pgfscope}%
\begin{pgfscope}%
\pgfsetbuttcap%
\pgfsetroundjoin%
\definecolor{currentfill}{rgb}{0.000000,0.000000,0.000000}%
\pgfsetfillcolor{currentfill}%
\pgfsetlinewidth{0.602250pt}%
\definecolor{currentstroke}{rgb}{0.000000,0.000000,0.000000}%
\pgfsetstrokecolor{currentstroke}%
\pgfsetdash{}{0pt}%
\pgfsys@defobject{currentmarker}{\pgfqpoint{-0.027778in}{0.000000in}}{\pgfqpoint{-0.000000in}{0.000000in}}{%
\pgfpathmoveto{\pgfqpoint{-0.000000in}{0.000000in}}%
\pgfpathlineto{\pgfqpoint{-0.027778in}{0.000000in}}%
\pgfusepath{stroke,fill}%
}%
\begin{pgfscope}%
\pgfsys@transformshift{1.000000in}{3.383333in}%
\pgfsys@useobject{currentmarker}{}%
\end{pgfscope}%
\end{pgfscope}%
\begin{pgfscope}%
\pgfpathrectangle{\pgfqpoint{1.000000in}{0.440000in}}{\pgfqpoint{6.200000in}{3.080000in}}%
\pgfusepath{clip}%
\pgfsetbuttcap%
\pgfsetroundjoin%
\pgfsetlinewidth{1.003750pt}%
\definecolor{currentstroke}{rgb}{0.690196,0.690196,0.690196}%
\pgfsetstrokecolor{currentstroke}%
\pgfsetstrokeopacity{0.600000}%
\pgfsetdash{{1.000000pt}{1.650000pt}}{0.000000pt}%
\pgfpathmoveto{\pgfqpoint{1.000000in}{3.420551in}}%
\pgfpathlineto{\pgfqpoint{7.200000in}{3.420551in}}%
\pgfusepath{stroke}%
\end{pgfscope}%
\begin{pgfscope}%
\pgfsetbuttcap%
\pgfsetroundjoin%
\definecolor{currentfill}{rgb}{0.000000,0.000000,0.000000}%
\pgfsetfillcolor{currentfill}%
\pgfsetlinewidth{0.602250pt}%
\definecolor{currentstroke}{rgb}{0.000000,0.000000,0.000000}%
\pgfsetstrokecolor{currentstroke}%
\pgfsetdash{}{0pt}%
\pgfsys@defobject{currentmarker}{\pgfqpoint{-0.027778in}{0.000000in}}{\pgfqpoint{-0.000000in}{0.000000in}}{%
\pgfpathmoveto{\pgfqpoint{-0.000000in}{0.000000in}}%
\pgfpathlineto{\pgfqpoint{-0.027778in}{0.000000in}}%
\pgfusepath{stroke,fill}%
}%
\begin{pgfscope}%
\pgfsys@transformshift{1.000000in}{3.420551in}%
\pgfsys@useobject{currentmarker}{}%
\end{pgfscope}%
\end{pgfscope}%
\begin{pgfscope}%
\pgfpathrectangle{\pgfqpoint{1.000000in}{0.440000in}}{\pgfqpoint{6.200000in}{3.080000in}}%
\pgfusepath{clip}%
\pgfsetbuttcap%
\pgfsetroundjoin%
\pgfsetlinewidth{1.003750pt}%
\definecolor{currentstroke}{rgb}{0.690196,0.690196,0.690196}%
\pgfsetstrokecolor{currentstroke}%
\pgfsetstrokeopacity{0.600000}%
\pgfsetdash{{1.000000pt}{1.650000pt}}{0.000000pt}%
\pgfpathmoveto{\pgfqpoint{1.000000in}{3.453381in}}%
\pgfpathlineto{\pgfqpoint{7.200000in}{3.453381in}}%
\pgfusepath{stroke}%
\end{pgfscope}%
\begin{pgfscope}%
\pgfsetbuttcap%
\pgfsetroundjoin%
\definecolor{currentfill}{rgb}{0.000000,0.000000,0.000000}%
\pgfsetfillcolor{currentfill}%
\pgfsetlinewidth{0.602250pt}%
\definecolor{currentstroke}{rgb}{0.000000,0.000000,0.000000}%
\pgfsetstrokecolor{currentstroke}%
\pgfsetdash{}{0pt}%
\pgfsys@defobject{currentmarker}{\pgfqpoint{-0.027778in}{0.000000in}}{\pgfqpoint{-0.000000in}{0.000000in}}{%
\pgfpathmoveto{\pgfqpoint{-0.000000in}{0.000000in}}%
\pgfpathlineto{\pgfqpoint{-0.027778in}{0.000000in}}%
\pgfusepath{stroke,fill}%
}%
\begin{pgfscope}%
\pgfsys@transformshift{1.000000in}{3.453381in}%
\pgfsys@useobject{currentmarker}{}%
\end{pgfscope}%
\end{pgfscope}%
\begin{pgfscope}%
\definecolor{textcolor}{rgb}{0.000000,0.000000,0.000000}%
\pgfsetstrokecolor{textcolor}%
\pgfsetfillcolor{textcolor}%
\pgftext[x=0.646026in,y=1.980000in,,bottom,rotate=90.000000]{\color{textcolor}{\rmfamily\fontsize{12.000000}{14.400000}\selectfont\catcode`\^=\active\def^{\ifmmode\sp\else\^{}\fi}\catcode`\%=\active\def%{\%}Number of reads per cell}}%
\end{pgfscope}%
\begin{pgfscope}%
\pgfpathrectangle{\pgfqpoint{1.000000in}{0.440000in}}{\pgfqpoint{6.200000in}{3.080000in}}%
\pgfusepath{clip}%
\pgfsetrectcap%
\pgfsetroundjoin%
\pgfsetlinewidth{1.204500pt}%
\definecolor{currentstroke}{rgb}{0.121569,0.466667,0.705882}%
\pgfsetstrokecolor{currentstroke}%
\pgfsetdash{}{0pt}%
\pgfpathmoveto{\pgfqpoint{0.990000in}{3.034154in}}%
\pgfpathlineto{\pgfqpoint{5.101040in}{3.034154in}}%
\pgfpathlineto{\pgfqpoint{5.101040in}{0.430000in}}%
\pgfpathlineto{\pgfqpoint{5.101040in}{0.430000in}}%
\pgfusepath{stroke}%
\end{pgfscope}%
\begin{pgfscope}%
\pgfpathrectangle{\pgfqpoint{1.000000in}{0.440000in}}{\pgfqpoint{6.200000in}{3.080000in}}%
\pgfusepath{clip}%
\pgfsetrectcap%
\pgfsetroundjoin%
\pgfsetlinewidth{1.204500pt}%
\definecolor{currentstroke}{rgb}{1.000000,0.498039,0.054902}%
\pgfsetstrokecolor{currentstroke}%
\pgfsetdash{}{0pt}%
\pgfpathmoveto{\pgfqpoint{0.990000in}{1.812766in}}%
\pgfpathlineto{\pgfqpoint{6.918182in}{1.812766in}}%
\pgfpathlineto{\pgfqpoint{6.918182in}{1.812766in}}%
\pgfusepath{stroke}%
\end{pgfscope}%
\begin{pgfscope}%
\pgfpathrectangle{\pgfqpoint{1.000000in}{0.440000in}}{\pgfqpoint{6.200000in}{3.080000in}}%
\pgfusepath{clip}%
\pgfsetrectcap%
\pgfsetroundjoin%
\pgfsetlinewidth{1.204500pt}%
\definecolor{currentstroke}{rgb}{0.172549,0.627451,0.172549}%
\pgfsetstrokecolor{currentstroke}%
\pgfsetdash{}{0pt}%
\pgfpathmoveto{\pgfqpoint{0.990000in}{3.171143in}}%
\pgfpathlineto{\pgfqpoint{1.281818in}{3.171137in}}%
\pgfpathlineto{\pgfqpoint{1.569246in}{3.139711in}}%
\pgfpathlineto{\pgfqpoint{1.737381in}{3.138400in}}%
\pgfpathlineto{\pgfqpoint{1.856675in}{3.112465in}}%
\pgfpathlineto{\pgfqpoint{1.949206in}{3.111696in}}%
\pgfpathlineto{\pgfqpoint{2.024809in}{3.108395in}}%
\pgfpathlineto{\pgfqpoint{2.088731in}{3.103787in}}%
\pgfpathlineto{\pgfqpoint{2.144103in}{3.096729in}}%
\pgfpathlineto{\pgfqpoint{2.192944in}{3.092511in}}%
\pgfpathlineto{\pgfqpoint{2.276156in}{3.078106in}}%
\pgfpathlineto{\pgfqpoint{2.376159in}{3.076566in}}%
\pgfpathlineto{\pgfqpoint{2.404769in}{3.072867in}}%
\pgfpathlineto{\pgfqpoint{2.456670in}{3.070161in}}%
\pgfpathlineto{\pgfqpoint{2.524062in}{3.066512in}}%
\pgfpathlineto{\pgfqpoint{2.599666in}{3.062651in}}%
\pgfpathlineto{\pgfqpoint{2.663588in}{3.059870in}}%
\pgfpathlineto{\pgfqpoint{2.692197in}{3.058660in}}%
\pgfpathlineto{\pgfqpoint{2.718959in}{3.057470in}}%
\pgfpathlineto{\pgfqpoint{2.744099in}{3.055941in}}%
\pgfpathlineto{\pgfqpoint{2.756119in}{3.053078in}}%
\pgfpathlineto{\pgfqpoint{2.779162in}{3.052021in}}%
\pgfpathlineto{\pgfqpoint{2.790221in}{3.050369in}}%
\pgfpathlineto{\pgfqpoint{2.841480in}{3.048045in}}%
\pgfpathlineto{\pgfqpoint{2.878364in}{3.046543in}}%
\pgfpathlineto{\pgfqpoint{2.928184in}{3.045140in}}%
\pgfpathlineto{\pgfqpoint{2.943544in}{3.044039in}}%
\pgfpathlineto{\pgfqpoint{2.965567in}{3.043595in}}%
\pgfpathlineto{\pgfqpoint{2.972656in}{3.042339in}}%
\pgfpathlineto{\pgfqpoint{2.979625in}{3.042258in}}%
\pgfpathlineto{\pgfqpoint{2.999857in}{3.038469in}}%
\pgfpathlineto{\pgfqpoint{3.025383in}{3.037079in}}%
\pgfpathlineto{\pgfqpoint{3.043547in}{3.035530in}}%
\pgfpathlineto{\pgfqpoint{3.072156in}{3.034043in}}%
\pgfpathlineto{\pgfqpoint{3.109158in}{3.031943in}}%
\pgfpathlineto{\pgfqpoint{3.156874in}{3.028949in}}%
\pgfpathlineto{\pgfqpoint{3.165792in}{3.027825in}}%
\pgfpathlineto{\pgfqpoint{3.238444in}{3.023139in}}%
\pgfpathlineto{\pgfqpoint{3.245784in}{3.020882in}}%
\pgfpathlineto{\pgfqpoint{3.277293in}{3.019872in}}%
\pgfpathlineto{\pgfqpoint{3.293816in}{3.018784in}}%
\pgfpathlineto{\pgfqpoint{3.309705in}{3.017337in}}%
\pgfpathlineto{\pgfqpoint{3.315895in}{3.016225in}}%
\pgfpathlineto{\pgfqpoint{3.336857in}{3.013116in}}%
\pgfpathlineto{\pgfqpoint{3.367796in}{3.011125in}}%
\pgfpathlineto{\pgfqpoint{3.381131in}{3.009379in}}%
\pgfpathlineto{\pgfqpoint{3.391498in}{3.008800in}}%
\pgfpathlineto{\pgfqpoint{3.399107in}{3.007361in}}%
\pgfpathlineto{\pgfqpoint{3.418740in}{3.005900in}}%
\pgfpathlineto{\pgfqpoint{3.430555in}{3.004834in}}%
\pgfpathlineto{\pgfqpoint{3.563032in}{2.994456in}}%
\pgfpathlineto{\pgfqpoint{3.606383in}{2.990099in}}%
\pgfpathlineto{\pgfqpoint{3.613937in}{2.989331in}}%
\pgfpathlineto{\pgfqpoint{3.642845in}{2.987755in}}%
\pgfpathlineto{\pgfqpoint{3.653867in}{2.986832in}}%
\pgfpathlineto{\pgfqpoint{3.823359in}{2.971313in}}%
\pgfpathlineto{\pgfqpoint{3.829635in}{2.970403in}}%
\pgfpathlineto{\pgfqpoint{3.839310in}{2.969051in}}%
\pgfpathlineto{\pgfqpoint{3.851306in}{2.967604in}}%
\pgfpathlineto{\pgfqpoint{3.870289in}{2.966203in}}%
\pgfpathlineto{\pgfqpoint{3.875898in}{2.965366in}}%
\pgfpathlineto{\pgfqpoint{3.889983in}{2.964293in}}%
\pgfpathlineto{\pgfqpoint{3.926768in}{2.960634in}}%
\pgfpathlineto{\pgfqpoint{3.932363in}{2.960090in}}%
\pgfpathlineto{\pgfqpoint{3.949374in}{2.958558in}}%
\pgfpathlineto{\pgfqpoint{3.954014in}{2.958101in}}%
\pgfpathlineto{\pgfqpoint{3.990590in}{2.954137in}}%
\pgfpathlineto{\pgfqpoint{4.073473in}{2.945480in}}%
\pgfpathlineto{\pgfqpoint{4.078874in}{2.944693in}}%
\pgfpathlineto{\pgfqpoint{4.216303in}{2.927356in}}%
\pgfpathlineto{\pgfqpoint{4.226337in}{2.926126in}}%
\pgfpathlineto{\pgfqpoint{4.246358in}{2.923858in}}%
\pgfpathlineto{\pgfqpoint{4.251537in}{2.922988in}}%
\pgfpathlineto{\pgfqpoint{4.272244in}{2.920610in}}%
\pgfpathlineto{\pgfqpoint{4.277414in}{2.919645in}}%
\pgfpathlineto{\pgfqpoint{4.304327in}{2.916666in}}%
\pgfpathlineto{\pgfqpoint{4.309115in}{2.916006in}}%
\pgfpathlineto{\pgfqpoint{4.332256in}{2.913138in}}%
\pgfpathlineto{\pgfqpoint{4.367520in}{2.908518in}}%
\pgfpathlineto{\pgfqpoint{4.380685in}{2.906234in}}%
\pgfpathlineto{\pgfqpoint{4.385837in}{2.905483in}}%
\pgfpathlineto{\pgfqpoint{4.405159in}{2.902571in}}%
\pgfpathlineto{\pgfqpoint{4.409797in}{2.901833in}}%
\pgfpathlineto{\pgfqpoint{4.418491in}{2.900733in}}%
\pgfpathlineto{\pgfqpoint{4.438845in}{2.897521in}}%
\pgfpathlineto{\pgfqpoint{4.441906in}{2.896792in}}%
\pgfpathlineto{\pgfqpoint{4.464647in}{2.893418in}}%
\pgfpathlineto{\pgfqpoint{4.466183in}{2.892932in}}%
\pgfpathlineto{\pgfqpoint{4.470000in}{2.891866in}}%
\pgfpathlineto{\pgfqpoint{4.480688in}{2.890353in}}%
\pgfpathlineto{\pgfqpoint{4.488210in}{2.889254in}}%
\pgfpathlineto{\pgfqpoint{4.492370in}{2.888338in}}%
\pgfpathlineto{\pgfqpoint{4.514790in}{2.884407in}}%
\pgfpathlineto{\pgfqpoint{4.516831in}{2.884097in}}%
\pgfpathlineto{\pgfqpoint{4.523730in}{2.882996in}}%
\pgfpathlineto{\pgfqpoint{4.527219in}{2.882228in}}%
\pgfpathlineto{\pgfqpoint{4.534112in}{2.881120in}}%
\pgfpathlineto{\pgfqpoint{4.536707in}{2.880830in}}%
\pgfpathlineto{\pgfqpoint{4.554281in}{2.877134in}}%
\pgfpathlineto{\pgfqpoint{4.556292in}{2.876546in}}%
\pgfpathlineto{\pgfqpoint{4.596823in}{2.868987in}}%
\pgfpathlineto{\pgfqpoint{4.602933in}{2.867951in}}%
\pgfpathlineto{\pgfqpoint{4.675817in}{2.853464in}}%
\pgfpathlineto{\pgfqpoint{4.677548in}{2.852904in}}%
\pgfpathlineto{\pgfqpoint{4.682925in}{2.851834in}}%
\pgfpathlineto{\pgfqpoint{4.687109in}{2.851082in}}%
\pgfpathlineto{\pgfqpoint{4.749663in}{2.835011in}}%
\pgfpathlineto{\pgfqpoint{4.751786in}{2.834213in}}%
\pgfpathlineto{\pgfqpoint{4.754090in}{2.833641in}}%
\pgfpathlineto{\pgfqpoint{4.823488in}{2.811744in}}%
\pgfpathlineto{\pgfqpoint{4.827839in}{2.809756in}}%
\pgfpathlineto{\pgfqpoint{4.831747in}{2.808725in}}%
\pgfpathlineto{\pgfqpoint{4.833964in}{2.807816in}}%
\pgfpathlineto{\pgfqpoint{4.867623in}{2.795624in}}%
\pgfpathlineto{\pgfqpoint{4.869585in}{2.794540in}}%
\pgfpathlineto{\pgfqpoint{4.922448in}{2.769334in}}%
\pgfpathlineto{\pgfqpoint{4.926195in}{2.766764in}}%
\pgfpathlineto{\pgfqpoint{4.939080in}{2.758157in}}%
\pgfpathlineto{\pgfqpoint{4.940670in}{2.756841in}}%
\pgfpathlineto{\pgfqpoint{4.953889in}{2.747450in}}%
\pgfpathlineto{\pgfqpoint{4.955424in}{2.746435in}}%
\pgfpathlineto{\pgfqpoint{4.958886in}{2.742891in}}%
\pgfpathlineto{\pgfqpoint{4.960635in}{2.741342in}}%
\pgfpathlineto{\pgfqpoint{4.964803in}{2.737166in}}%
\pgfpathlineto{\pgfqpoint{4.967330in}{2.734402in}}%
\pgfpathlineto{\pgfqpoint{4.970864in}{2.730577in}}%
\pgfpathlineto{\pgfqpoint{4.972903in}{2.727286in}}%
\pgfpathlineto{\pgfqpoint{4.973524in}{2.726572in}}%
\pgfpathlineto{\pgfqpoint{4.974875in}{2.724718in}}%
\pgfpathlineto{\pgfqpoint{4.977621in}{2.721142in}}%
\pgfpathlineto{\pgfqpoint{4.978569in}{2.720069in}}%
\pgfpathlineto{\pgfqpoint{4.988425in}{2.706859in}}%
\pgfpathlineto{\pgfqpoint{4.989186in}{2.704867in}}%
\pgfpathlineto{\pgfqpoint{4.995813in}{2.691184in}}%
\pgfpathlineto{\pgfqpoint{4.997626in}{2.688259in}}%
\pgfpathlineto{\pgfqpoint{4.999167in}{2.681970in}}%
\pgfpathlineto{\pgfqpoint{5.000332in}{2.679494in}}%
\pgfpathlineto{\pgfqpoint{5.012625in}{2.616713in}}%
\pgfpathlineto{\pgfqpoint{5.015490in}{2.590875in}}%
\pgfpathlineto{\pgfqpoint{5.016152in}{2.586396in}}%
\pgfpathlineto{\pgfqpoint{5.018840in}{2.564658in}}%
\pgfpathlineto{\pgfqpoint{5.019245in}{2.562093in}}%
\pgfpathlineto{\pgfqpoint{5.020959in}{2.553940in}}%
\pgfpathlineto{\pgfqpoint{5.021160in}{2.553236in}}%
\pgfpathlineto{\pgfqpoint{5.022716in}{2.545138in}}%
\pgfpathlineto{\pgfqpoint{5.024365in}{2.536466in}}%
\pgfpathlineto{\pgfqpoint{5.026108in}{2.525620in}}%
\pgfpathlineto{\pgfqpoint{5.028140in}{2.514875in}}%
\pgfpathlineto{\pgfqpoint{5.029966in}{2.505655in}}%
\pgfpathlineto{\pgfqpoint{5.032762in}{2.491079in}}%
\pgfpathlineto{\pgfqpoint{5.034324in}{2.484649in}}%
\pgfpathlineto{\pgfqpoint{5.035297in}{2.480210in}}%
\pgfpathlineto{\pgfqpoint{5.037091in}{2.472895in}}%
\pgfpathlineto{\pgfqpoint{5.037236in}{2.471954in}}%
\pgfpathlineto{\pgfqpoint{5.047743in}{2.431318in}}%
\pgfpathlineto{\pgfqpoint{5.050422in}{2.421826in}}%
\pgfpathlineto{\pgfqpoint{5.051919in}{2.415229in}}%
\pgfpathlineto{\pgfqpoint{5.052059in}{2.414844in}}%
\pgfpathlineto{\pgfqpoint{5.053503in}{2.409260in}}%
\pgfpathlineto{\pgfqpoint{5.053968in}{2.408209in}}%
\pgfpathlineto{\pgfqpoint{5.057714in}{2.394445in}}%
\pgfpathlineto{\pgfqpoint{5.059460in}{2.388715in}}%
\pgfpathlineto{\pgfqpoint{5.064654in}{2.369121in}}%
\pgfpathlineto{\pgfqpoint{5.066596in}{2.360819in}}%
\pgfpathlineto{\pgfqpoint{5.066911in}{2.359724in}}%
\pgfpathlineto{\pgfqpoint{5.075821in}{2.324890in}}%
\pgfpathlineto{\pgfqpoint{5.077580in}{2.313845in}}%
\pgfpathlineto{\pgfqpoint{5.077843in}{2.312920in}}%
\pgfpathlineto{\pgfqpoint{5.079463in}{2.303879in}}%
\pgfpathlineto{\pgfqpoint{5.084499in}{2.277811in}}%
\pgfpathlineto{\pgfqpoint{5.085835in}{2.270571in}}%
\pgfpathlineto{\pgfqpoint{5.088280in}{2.252755in}}%
\pgfpathlineto{\pgfqpoint{5.095866in}{2.183101in}}%
\pgfpathlineto{\pgfqpoint{5.095908in}{2.183101in}}%
\pgfpathlineto{\pgfqpoint{5.097626in}{2.156801in}}%
\pgfpathlineto{\pgfqpoint{5.098419in}{2.144190in}}%
\pgfpathlineto{\pgfqpoint{5.099753in}{2.115990in}}%
\pgfpathlineto{\pgfqpoint{5.102076in}{2.082926in}}%
\pgfpathlineto{\pgfqpoint{5.102159in}{2.082926in}}%
\pgfpathlineto{\pgfqpoint{5.103686in}{2.069939in}}%
\pgfpathlineto{\pgfqpoint{5.105413in}{2.056318in}}%
\pgfpathlineto{\pgfqpoint{5.105659in}{2.055386in}}%
\pgfpathlineto{\pgfqpoint{5.107460in}{2.047333in}}%
\pgfpathlineto{\pgfqpoint{5.107868in}{2.046371in}}%
\pgfpathlineto{\pgfqpoint{5.109375in}{2.041506in}}%
\pgfpathlineto{\pgfqpoint{5.109822in}{2.040523in}}%
\pgfpathlineto{\pgfqpoint{5.111200in}{2.037054in}}%
\pgfpathlineto{\pgfqpoint{5.111645in}{2.036055in}}%
\pgfpathlineto{\pgfqpoint{5.113097in}{2.031515in}}%
\pgfpathlineto{\pgfqpoint{5.113500in}{2.030496in}}%
\pgfpathlineto{\pgfqpoint{5.114786in}{2.027932in}}%
\pgfpathlineto{\pgfqpoint{5.115187in}{2.026899in}}%
\pgfpathlineto{\pgfqpoint{5.116467in}{2.024302in}}%
\pgfpathlineto{\pgfqpoint{5.117066in}{2.023256in}}%
\pgfpathlineto{\pgfqpoint{5.118301in}{2.021152in}}%
\pgfpathlineto{\pgfqpoint{5.118937in}{2.020094in}}%
\pgfpathlineto{\pgfqpoint{5.120167in}{2.018500in}}%
\pgfpathlineto{\pgfqpoint{5.120721in}{2.017432in}}%
\pgfpathlineto{\pgfqpoint{5.122063in}{2.014204in}}%
\pgfpathlineto{\pgfqpoint{5.123166in}{2.013119in}}%
\pgfpathlineto{\pgfqpoint{5.124422in}{2.010937in}}%
\pgfpathlineto{\pgfqpoint{5.125557in}{2.009840in}}%
\pgfpathlineto{\pgfqpoint{5.126806in}{2.007632in}}%
\pgfpathlineto{\pgfqpoint{5.127896in}{2.006522in}}%
\pgfpathlineto{\pgfqpoint{5.129138in}{2.004287in}}%
\pgfpathlineto{\pgfqpoint{5.130028in}{2.003726in}}%
\pgfpathlineto{\pgfqpoint{5.131187in}{2.002600in}}%
\pgfpathlineto{\pgfqpoint{5.131957in}{2.002035in}}%
\pgfpathlineto{\pgfqpoint{5.133149in}{2.000334in}}%
\pgfpathlineto{\pgfqpoint{5.134069in}{1.999764in}}%
\pgfpathlineto{\pgfqpoint{5.135293in}{1.998049in}}%
\pgfpathlineto{\pgfqpoint{5.136437in}{1.997475in}}%
\pgfpathlineto{\pgfqpoint{5.137653in}{1.995746in}}%
\pgfpathlineto{\pgfqpoint{5.138184in}{1.995167in}}%
\pgfpathlineto{\pgfqpoint{5.139396in}{1.993423in}}%
\pgfpathlineto{\pgfqpoint{5.140378in}{1.992839in}}%
\pgfpathlineto{\pgfqpoint{5.141546in}{1.991668in}}%
\pgfpathlineto{\pgfqpoint{5.142485in}{1.991081in}}%
\pgfpathlineto{\pgfqpoint{5.143646in}{1.989902in}}%
\pgfpathlineto{\pgfqpoint{5.144207in}{1.989311in}}%
\pgfpathlineto{\pgfqpoint{5.145402in}{1.987530in}}%
\pgfpathlineto{\pgfqpoint{5.146741in}{1.986934in}}%
\pgfpathlineto{\pgfqpoint{5.147928in}{1.985737in}}%
\pgfpathlineto{\pgfqpoint{5.148520in}{1.985137in}}%
\pgfpathlineto{\pgfqpoint{5.149702in}{1.983933in}}%
\pgfpathlineto{\pgfqpoint{5.150218in}{1.983329in}}%
\pgfpathlineto{\pgfqpoint{5.151395in}{1.982118in}}%
\pgfpathlineto{\pgfqpoint{5.152605in}{1.981510in}}%
\pgfpathlineto{\pgfqpoint{5.153775in}{1.980290in}}%
\pgfpathlineto{\pgfqpoint{5.155124in}{1.979678in}}%
\pgfpathlineto{\pgfqpoint{5.156287in}{1.978450in}}%
\pgfpathlineto{\pgfqpoint{5.157737in}{1.977834in}}%
\pgfpathlineto{\pgfqpoint{5.158893in}{1.976598in}}%
\pgfpathlineto{\pgfqpoint{5.160369in}{1.975978in}}%
\pgfpathlineto{\pgfqpoint{5.161553in}{1.974734in}}%
\pgfpathlineto{\pgfqpoint{5.163449in}{1.974109in}}%
\pgfpathlineto{\pgfqpoint{5.164624in}{1.972857in}}%
\pgfpathlineto{\pgfqpoint{5.166187in}{1.972228in}}%
\pgfpathlineto{\pgfqpoint{5.167355in}{1.970967in}}%
\pgfpathlineto{\pgfqpoint{5.169470in}{1.970334in}}%
\pgfpathlineto{\pgfqpoint{5.170629in}{1.969064in}}%
\pgfpathlineto{\pgfqpoint{5.172727in}{1.968427in}}%
\pgfpathlineto{\pgfqpoint{5.173842in}{1.967789in}}%
\pgfpathlineto{\pgfqpoint{5.175266in}{1.967149in}}%
\pgfpathlineto{\pgfqpoint{5.176409in}{1.966507in}}%
\pgfpathlineto{\pgfqpoint{5.177617in}{1.965864in}}%
\pgfpathlineto{\pgfqpoint{5.178788in}{1.964574in}}%
\pgfpathlineto{\pgfqpoint{5.181632in}{1.963926in}}%
\pgfpathlineto{\pgfqpoint{5.182791in}{1.962627in}}%
\pgfpathlineto{\pgfqpoint{5.185642in}{1.961975in}}%
\pgfpathlineto{\pgfqpoint{5.186756in}{1.961322in}}%
\pgfpathlineto{\pgfqpoint{5.188170in}{1.960666in}}%
\pgfpathlineto{\pgfqpoint{5.189311in}{1.960010in}}%
\pgfpathlineto{\pgfqpoint{5.190884in}{1.959352in}}%
\pgfpathlineto{\pgfqpoint{5.192018in}{1.958692in}}%
\pgfpathlineto{\pgfqpoint{5.193712in}{1.958031in}}%
\pgfpathlineto{\pgfqpoint{5.194839in}{1.957368in}}%
\pgfpathlineto{\pgfqpoint{5.196456in}{1.956703in}}%
\pgfpathlineto{\pgfqpoint{5.197575in}{1.956037in}}%
\pgfpathlineto{\pgfqpoint{5.200032in}{1.955370in}}%
\pgfpathlineto{\pgfqpoint{5.201174in}{1.954701in}}%
\pgfpathlineto{\pgfqpoint{5.203707in}{1.954030in}}%
\pgfpathlineto{\pgfqpoint{5.204838in}{1.953357in}}%
\pgfpathlineto{\pgfqpoint{5.207028in}{1.952683in}}%
\pgfpathlineto{\pgfqpoint{5.208151in}{1.952008in}}%
\pgfpathlineto{\pgfqpoint{5.210323in}{1.951330in}}%
\pgfpathlineto{\pgfqpoint{5.211437in}{1.950651in}}%
\pgfpathlineto{\pgfqpoint{5.213592in}{1.949971in}}%
\pgfpathlineto{\pgfqpoint{5.214729in}{1.949288in}}%
\pgfpathlineto{\pgfqpoint{5.217274in}{1.948604in}}%
\pgfpathlineto{\pgfqpoint{5.218401in}{1.947919in}}%
\pgfpathlineto{\pgfqpoint{5.221514in}{1.947231in}}%
\pgfpathlineto{\pgfqpoint{5.222629in}{1.946542in}}%
\pgfpathlineto{\pgfqpoint{5.224388in}{1.945851in}}%
\pgfpathlineto{\pgfqpoint{5.225527in}{1.945159in}}%
\pgfpathlineto{\pgfqpoint{5.228129in}{1.944465in}}%
\pgfpathlineto{\pgfqpoint{5.229257in}{1.943769in}}%
\pgfpathlineto{\pgfqpoint{5.231867in}{1.943071in}}%
\pgfpathlineto{\pgfqpoint{5.232985in}{1.942372in}}%
\pgfpathlineto{\pgfqpoint{5.235841in}{1.941670in}}%
\pgfpathlineto{\pgfqpoint{5.236978in}{1.940967in}}%
\pgfpathlineto{\pgfqpoint{5.240637in}{1.940263in}}%
\pgfpathlineto{\pgfqpoint{5.241761in}{1.939556in}}%
\pgfpathlineto{\pgfqpoint{5.245261in}{1.938848in}}%
\pgfpathlineto{\pgfqpoint{5.246401in}{1.938138in}}%
\pgfpathlineto{\pgfqpoint{5.249544in}{1.937426in}}%
\pgfpathlineto{\pgfqpoint{5.250673in}{1.936712in}}%
\pgfpathlineto{\pgfqpoint{5.255015in}{1.935996in}}%
\pgfpathlineto{\pgfqpoint{5.256130in}{1.935279in}}%
\pgfpathlineto{\pgfqpoint{5.259992in}{1.934560in}}%
\pgfpathlineto{\pgfqpoint{5.261121in}{1.933839in}}%
\pgfpathlineto{\pgfqpoint{5.265161in}{1.933115in}}%
\pgfpathlineto{\pgfqpoint{5.266276in}{1.932391in}}%
\pgfpathlineto{\pgfqpoint{5.269935in}{1.931664in}}%
\pgfpathlineto{\pgfqpoint{5.271065in}{1.930935in}}%
\pgfpathlineto{\pgfqpoint{5.275172in}{1.930204in}}%
\pgfpathlineto{\pgfqpoint{5.276288in}{1.929472in}}%
\pgfpathlineto{\pgfqpoint{5.280775in}{1.928737in}}%
\pgfpathlineto{\pgfqpoint{5.281903in}{1.928001in}}%
\pgfpathlineto{\pgfqpoint{5.286569in}{1.927263in}}%
\pgfpathlineto{\pgfqpoint{5.287681in}{1.926522in}}%
\pgfpathlineto{\pgfqpoint{5.291130in}{1.925780in}}%
\pgfpathlineto{\pgfqpoint{5.292256in}{1.925036in}}%
\pgfpathlineto{\pgfqpoint{5.296962in}{1.924290in}}%
\pgfpathlineto{\pgfqpoint{5.298099in}{1.923541in}}%
\pgfpathlineto{\pgfqpoint{5.303045in}{1.922791in}}%
\pgfpathlineto{\pgfqpoint{5.304165in}{1.922039in}}%
\pgfpathlineto{\pgfqpoint{5.308362in}{1.921284in}}%
\pgfpathlineto{\pgfqpoint{5.309492in}{1.920528in}}%
\pgfpathlineto{\pgfqpoint{5.314107in}{1.919770in}}%
\pgfpathlineto{\pgfqpoint{5.315222in}{1.919009in}}%
\pgfpathlineto{\pgfqpoint{5.319799in}{1.918247in}}%
\pgfpathlineto{\pgfqpoint{5.320923in}{1.917482in}}%
\pgfpathlineto{\pgfqpoint{5.325968in}{1.916715in}}%
\pgfpathlineto{\pgfqpoint{5.327100in}{1.915946in}}%
\pgfpathlineto{\pgfqpoint{5.333967in}{1.915175in}}%
\pgfpathlineto{\pgfqpoint{5.335100in}{1.914402in}}%
\pgfpathlineto{\pgfqpoint{5.340745in}{1.913627in}}%
\pgfpathlineto{\pgfqpoint{5.341861in}{1.912849in}}%
\pgfpathlineto{\pgfqpoint{5.347713in}{1.912070in}}%
\pgfpathlineto{\pgfqpoint{5.348832in}{1.911288in}}%
\pgfpathlineto{\pgfqpoint{5.354745in}{1.910504in}}%
\pgfpathlineto{\pgfqpoint{5.355868in}{1.909717in}}%
\pgfpathlineto{\pgfqpoint{5.362367in}{1.908929in}}%
\pgfpathlineto{\pgfqpoint{5.363491in}{1.908138in}}%
\pgfpathlineto{\pgfqpoint{5.369938in}{1.907345in}}%
\pgfpathlineto{\pgfqpoint{5.371064in}{1.906550in}}%
\pgfpathlineto{\pgfqpoint{5.378224in}{1.905752in}}%
\pgfpathlineto{\pgfqpoint{5.379349in}{1.904953in}}%
\pgfpathlineto{\pgfqpoint{5.385868in}{1.904151in}}%
\pgfpathlineto{\pgfqpoint{5.386994in}{1.903346in}}%
\pgfpathlineto{\pgfqpoint{5.393579in}{1.902539in}}%
\pgfpathlineto{\pgfqpoint{5.394704in}{1.901730in}}%
\pgfpathlineto{\pgfqpoint{5.400848in}{1.900919in}}%
\pgfpathlineto{\pgfqpoint{5.401973in}{1.900105in}}%
\pgfpathlineto{\pgfqpoint{5.408762in}{1.899289in}}%
\pgfpathlineto{\pgfqpoint{5.409886in}{1.898470in}}%
\pgfpathlineto{\pgfqpoint{5.416644in}{1.897649in}}%
\pgfpathlineto{\pgfqpoint{5.417766in}{1.896826in}}%
\pgfpathlineto{\pgfqpoint{5.424778in}{1.896000in}}%
\pgfpathlineto{\pgfqpoint{5.425897in}{1.895172in}}%
\pgfpathlineto{\pgfqpoint{5.433463in}{1.894341in}}%
\pgfpathlineto{\pgfqpoint{5.434578in}{1.893508in}}%
\pgfpathlineto{\pgfqpoint{5.442279in}{1.892672in}}%
\pgfpathlineto{\pgfqpoint{5.443407in}{1.891834in}}%
\pgfpathlineto{\pgfqpoint{5.451162in}{1.890993in}}%
\pgfpathlineto{\pgfqpoint{5.452284in}{1.890150in}}%
\pgfpathlineto{\pgfqpoint{5.460260in}{1.889304in}}%
\pgfpathlineto{\pgfqpoint{5.461375in}{1.888455in}}%
\pgfpathlineto{\pgfqpoint{5.469725in}{1.887604in}}%
\pgfpathlineto{\pgfqpoint{5.470849in}{1.886750in}}%
\pgfpathlineto{\pgfqpoint{5.478512in}{1.885894in}}%
\pgfpathlineto{\pgfqpoint{5.479629in}{1.885035in}}%
\pgfpathlineto{\pgfqpoint{5.487885in}{1.884174in}}%
\pgfpathlineto{\pgfqpoint{5.489009in}{1.883309in}}%
\pgfpathlineto{\pgfqpoint{5.497321in}{1.882442in}}%
\pgfpathlineto{\pgfqpoint{5.498436in}{1.881573in}}%
\pgfpathlineto{\pgfqpoint{5.507482in}{1.880700in}}%
\pgfpathlineto{\pgfqpoint{5.508601in}{1.879825in}}%
\pgfpathlineto{\pgfqpoint{5.516959in}{1.878947in}}%
\pgfpathlineto{\pgfqpoint{5.518083in}{1.878067in}}%
\pgfpathlineto{\pgfqpoint{5.527131in}{1.877183in}}%
\pgfpathlineto{\pgfqpoint{5.528243in}{1.876297in}}%
\pgfpathlineto{\pgfqpoint{5.537725in}{1.875408in}}%
\pgfpathlineto{\pgfqpoint{5.538838in}{1.874516in}}%
\pgfpathlineto{\pgfqpoint{5.548211in}{1.873621in}}%
\pgfpathlineto{\pgfqpoint{5.549324in}{1.872723in}}%
\pgfpathlineto{\pgfqpoint{5.558699in}{1.871823in}}%
\pgfpathlineto{\pgfqpoint{5.559812in}{1.870919in}}%
\pgfpathlineto{\pgfqpoint{5.569157in}{1.870013in}}%
\pgfpathlineto{\pgfqpoint{5.570268in}{1.869104in}}%
\pgfpathlineto{\pgfqpoint{5.579854in}{1.868191in}}%
\pgfpathlineto{\pgfqpoint{5.580977in}{1.867276in}}%
\pgfpathlineto{\pgfqpoint{5.590983in}{1.866357in}}%
\pgfpathlineto{\pgfqpoint{5.592101in}{1.865436in}}%
\pgfpathlineto{\pgfqpoint{5.602106in}{1.864512in}}%
\pgfpathlineto{\pgfqpoint{5.603219in}{1.863584in}}%
\pgfpathlineto{\pgfqpoint{5.613324in}{1.862653in}}%
\pgfpathlineto{\pgfqpoint{5.614444in}{1.861720in}}%
\pgfpathlineto{\pgfqpoint{5.624047in}{1.860783in}}%
\pgfpathlineto{\pgfqpoint{5.625162in}{1.859843in}}%
\pgfpathlineto{\pgfqpoint{5.635461in}{1.858899in}}%
\pgfpathlineto{\pgfqpoint{5.636580in}{1.857953in}}%
\pgfpathlineto{\pgfqpoint{5.647003in}{1.857003in}}%
\pgfpathlineto{\pgfqpoint{5.648124in}{1.856050in}}%
\pgfpathlineto{\pgfqpoint{5.657940in}{1.855094in}}%
\pgfpathlineto{\pgfqpoint{5.659054in}{1.854135in}}%
\pgfpathlineto{\pgfqpoint{5.669334in}{1.853172in}}%
\pgfpathlineto{\pgfqpoint{5.670449in}{1.852206in}}%
\pgfpathlineto{\pgfqpoint{5.680803in}{1.851236in}}%
\pgfpathlineto{\pgfqpoint{5.681918in}{1.850264in}}%
\pgfpathlineto{\pgfqpoint{5.692899in}{1.849287in}}%
\pgfpathlineto{\pgfqpoint{5.694012in}{1.848308in}}%
\pgfpathlineto{\pgfqpoint{5.704633in}{1.847324in}}%
\pgfpathlineto{\pgfqpoint{5.705754in}{1.846338in}}%
\pgfpathlineto{\pgfqpoint{5.716628in}{1.845348in}}%
\pgfpathlineto{\pgfqpoint{5.717744in}{1.844354in}}%
\pgfpathlineto{\pgfqpoint{5.728696in}{1.843357in}}%
\pgfpathlineto{\pgfqpoint{5.729808in}{1.842356in}}%
\pgfpathlineto{\pgfqpoint{5.741167in}{1.841351in}}%
\pgfpathlineto{\pgfqpoint{5.742282in}{1.840343in}}%
\pgfpathlineto{\pgfqpoint{5.753326in}{1.839332in}}%
\pgfpathlineto{\pgfqpoint{5.754443in}{1.838316in}}%
\pgfpathlineto{\pgfqpoint{5.765815in}{1.837297in}}%
\pgfpathlineto{\pgfqpoint{5.766932in}{1.836274in}}%
\pgfpathlineto{\pgfqpoint{5.778150in}{1.835248in}}%
\pgfpathlineto{\pgfqpoint{5.779267in}{1.834217in}}%
\pgfpathlineto{\pgfqpoint{5.790962in}{1.833183in}}%
\pgfpathlineto{\pgfqpoint{5.792076in}{1.832145in}}%
\pgfpathlineto{\pgfqpoint{5.804068in}{1.831103in}}%
\pgfpathlineto{\pgfqpoint{5.805185in}{1.830057in}}%
\pgfpathlineto{\pgfqpoint{5.816425in}{1.829007in}}%
\pgfpathlineto{\pgfqpoint{5.817539in}{1.827953in}}%
\pgfpathlineto{\pgfqpoint{5.829049in}{1.826896in}}%
\pgfpathlineto{\pgfqpoint{5.830165in}{1.825834in}}%
\pgfpathlineto{\pgfqpoint{5.842175in}{1.824768in}}%
\pgfpathlineto{\pgfqpoint{5.843291in}{1.823698in}}%
\pgfpathlineto{\pgfqpoint{5.855208in}{1.822624in}}%
\pgfpathlineto{\pgfqpoint{5.856323in}{1.821546in}}%
\pgfpathlineto{\pgfqpoint{5.868535in}{1.820463in}}%
\pgfpathlineto{\pgfqpoint{5.869648in}{1.819377in}}%
\pgfpathlineto{\pgfqpoint{5.881334in}{1.818286in}}%
\pgfpathlineto{\pgfqpoint{5.882450in}{1.817191in}}%
\pgfpathlineto{\pgfqpoint{5.894882in}{1.816091in}}%
\pgfpathlineto{\pgfqpoint{5.895993in}{1.814987in}}%
\pgfpathlineto{\pgfqpoint{5.908066in}{1.813879in}}%
\pgfpathlineto{\pgfqpoint{5.909178in}{1.812766in}}%
\pgfpathlineto{\pgfqpoint{5.914898in}{1.812766in}}%
\pgfpathlineto{\pgfqpoint{5.916009in}{1.811649in}}%
\pgfpathlineto{\pgfqpoint{5.921762in}{1.811649in}}%
\pgfpathlineto{\pgfqpoint{5.922878in}{1.810527in}}%
\pgfpathlineto{\pgfqpoint{5.928452in}{1.810527in}}%
\pgfpathlineto{\pgfqpoint{5.929567in}{1.809401in}}%
\pgfpathlineto{\pgfqpoint{5.935303in}{1.809401in}}%
\pgfpathlineto{\pgfqpoint{5.936416in}{1.808270in}}%
\pgfpathlineto{\pgfqpoint{5.941594in}{1.808270in}}%
\pgfpathlineto{\pgfqpoint{5.942707in}{1.807135in}}%
\pgfpathlineto{\pgfqpoint{5.948201in}{1.807135in}}%
\pgfpathlineto{\pgfqpoint{5.949312in}{1.805995in}}%
\pgfpathlineto{\pgfqpoint{5.955047in}{1.805995in}}%
\pgfpathlineto{\pgfqpoint{5.956161in}{1.804850in}}%
\pgfpathlineto{\pgfqpoint{5.961913in}{1.804850in}}%
\pgfpathlineto{\pgfqpoint{5.963024in}{1.803701in}}%
\pgfpathlineto{\pgfqpoint{5.968809in}{1.803701in}}%
\pgfpathlineto{\pgfqpoint{5.969923in}{1.802547in}}%
\pgfpathlineto{\pgfqpoint{5.975644in}{1.802547in}}%
\pgfpathlineto{\pgfqpoint{5.976759in}{1.801388in}}%
\pgfpathlineto{\pgfqpoint{5.982575in}{1.801388in}}%
\pgfpathlineto{\pgfqpoint{5.983687in}{1.800224in}}%
\pgfpathlineto{\pgfqpoint{5.989524in}{1.800224in}}%
\pgfpathlineto{\pgfqpoint{5.990637in}{1.799056in}}%
\pgfpathlineto{\pgfqpoint{5.996430in}{1.799056in}}%
\pgfpathlineto{\pgfqpoint{5.997544in}{1.797882in}}%
\pgfpathlineto{\pgfqpoint{6.003463in}{1.797882in}}%
\pgfpathlineto{\pgfqpoint{6.004577in}{1.796703in}}%
\pgfpathlineto{\pgfqpoint{6.010411in}{1.796703in}}%
\pgfpathlineto{\pgfqpoint{6.011525in}{1.795520in}}%
\pgfpathlineto{\pgfqpoint{6.017259in}{1.795520in}}%
\pgfpathlineto{\pgfqpoint{6.018372in}{1.794331in}}%
\pgfpathlineto{\pgfqpoint{6.024165in}{1.794331in}}%
\pgfpathlineto{\pgfqpoint{6.025278in}{1.793138in}}%
\pgfpathlineto{\pgfqpoint{6.030866in}{1.793138in}}%
\pgfpathlineto{\pgfqpoint{6.031979in}{1.791939in}}%
\pgfpathlineto{\pgfqpoint{6.037863in}{1.791939in}}%
\pgfpathlineto{\pgfqpoint{6.038974in}{1.790735in}}%
\pgfpathlineto{\pgfqpoint{6.044995in}{1.790735in}}%
\pgfpathlineto{\pgfqpoint{6.046109in}{1.789525in}}%
\pgfpathlineto{\pgfqpoint{6.051952in}{1.789525in}}%
\pgfpathlineto{\pgfqpoint{6.053064in}{1.788311in}}%
\pgfpathlineto{\pgfqpoint{6.058794in}{1.788311in}}%
\pgfpathlineto{\pgfqpoint{6.059909in}{1.787091in}}%
\pgfpathlineto{\pgfqpoint{6.065428in}{1.787091in}}%
\pgfpathlineto{\pgfqpoint{6.066541in}{1.785866in}}%
\pgfpathlineto{\pgfqpoint{6.072397in}{1.785866in}}%
\pgfpathlineto{\pgfqpoint{6.073511in}{1.784635in}}%
\pgfpathlineto{\pgfqpoint{6.079449in}{1.784635in}}%
\pgfpathlineto{\pgfqpoint{6.080561in}{1.783399in}}%
\pgfpathlineto{\pgfqpoint{6.086326in}{1.783399in}}%
\pgfpathlineto{\pgfqpoint{6.087439in}{1.782158in}}%
\pgfpathlineto{\pgfqpoint{6.093379in}{1.782158in}}%
\pgfpathlineto{\pgfqpoint{6.094492in}{1.780911in}}%
\pgfpathlineto{\pgfqpoint{6.100493in}{1.780911in}}%
\pgfpathlineto{\pgfqpoint{6.101605in}{1.779658in}}%
\pgfpathlineto{\pgfqpoint{6.107373in}{1.779658in}}%
\pgfpathlineto{\pgfqpoint{6.108485in}{1.778400in}}%
\pgfpathlineto{\pgfqpoint{6.114354in}{1.778400in}}%
\pgfpathlineto{\pgfqpoint{6.115465in}{1.777136in}}%
\pgfpathlineto{\pgfqpoint{6.121120in}{1.777136in}}%
\pgfpathlineto{\pgfqpoint{6.122234in}{1.775866in}}%
\pgfpathlineto{\pgfqpoint{6.127892in}{1.775866in}}%
\pgfpathlineto{\pgfqpoint{6.129006in}{1.774590in}}%
\pgfpathlineto{\pgfqpoint{6.134768in}{1.774590in}}%
\pgfpathlineto{\pgfqpoint{6.135880in}{1.773309in}}%
\pgfpathlineto{\pgfqpoint{6.141612in}{1.773309in}}%
\pgfpathlineto{\pgfqpoint{6.142724in}{1.772021in}}%
\pgfpathlineto{\pgfqpoint{6.148591in}{1.772021in}}%
\pgfpathlineto{\pgfqpoint{6.149704in}{1.770728in}}%
\pgfpathlineto{\pgfqpoint{6.155608in}{1.770728in}}%
\pgfpathlineto{\pgfqpoint{6.156722in}{1.769428in}}%
\pgfpathlineto{\pgfqpoint{6.162489in}{1.769428in}}%
\pgfpathlineto{\pgfqpoint{6.163600in}{1.768123in}}%
\pgfpathlineto{\pgfqpoint{6.169393in}{1.768123in}}%
\pgfpathlineto{\pgfqpoint{6.170505in}{1.766811in}}%
\pgfpathlineto{\pgfqpoint{6.176091in}{1.766811in}}%
\pgfpathlineto{\pgfqpoint{6.177204in}{1.765493in}}%
\pgfpathlineto{\pgfqpoint{6.183122in}{1.765493in}}%
\pgfpathlineto{\pgfqpoint{6.184234in}{1.764169in}}%
\pgfpathlineto{\pgfqpoint{6.189994in}{1.764169in}}%
\pgfpathlineto{\pgfqpoint{6.191106in}{1.762839in}}%
\pgfpathlineto{\pgfqpoint{6.196984in}{1.762839in}}%
\pgfpathlineto{\pgfqpoint{6.198098in}{1.761502in}}%
\pgfpathlineto{\pgfqpoint{6.203728in}{1.761502in}}%
\pgfpathlineto{\pgfqpoint{6.204841in}{1.760159in}}%
\pgfpathlineto{\pgfqpoint{6.211097in}{1.760159in}}%
\pgfpathlineto{\pgfqpoint{6.212210in}{1.758809in}}%
\pgfpathlineto{\pgfqpoint{6.218015in}{1.758809in}}%
\pgfpathlineto{\pgfqpoint{6.219127in}{1.757453in}}%
\pgfpathlineto{\pgfqpoint{6.224803in}{1.757453in}}%
\pgfpathlineto{\pgfqpoint{6.225916in}{1.756090in}}%
\pgfpathlineto{\pgfqpoint{6.231612in}{1.756090in}}%
\pgfpathlineto{\pgfqpoint{6.232723in}{1.754720in}}%
\pgfpathlineto{\pgfqpoint{6.238453in}{1.754720in}}%
\pgfpathlineto{\pgfqpoint{6.239564in}{1.753344in}}%
\pgfpathlineto{\pgfqpoint{6.245576in}{1.753344in}}%
\pgfpathlineto{\pgfqpoint{6.246687in}{1.751960in}}%
\pgfpathlineto{\pgfqpoint{6.252285in}{1.751960in}}%
\pgfpathlineto{\pgfqpoint{6.253396in}{1.750570in}}%
\pgfpathlineto{\pgfqpoint{6.259182in}{1.750570in}}%
\pgfpathlineto{\pgfqpoint{6.260295in}{1.749173in}}%
\pgfpathlineto{\pgfqpoint{6.266505in}{1.749173in}}%
\pgfpathlineto{\pgfqpoint{6.267616in}{1.747769in}}%
\pgfpathlineto{\pgfqpoint{6.273264in}{1.747769in}}%
\pgfpathlineto{\pgfqpoint{6.274377in}{1.746358in}}%
\pgfpathlineto{\pgfqpoint{6.280096in}{1.746358in}}%
\pgfpathlineto{\pgfqpoint{6.281208in}{1.744939in}}%
\pgfpathlineto{\pgfqpoint{6.286851in}{1.744939in}}%
\pgfpathlineto{\pgfqpoint{6.287964in}{1.743513in}}%
\pgfpathlineto{\pgfqpoint{6.293703in}{1.743513in}}%
\pgfpathlineto{\pgfqpoint{6.294816in}{1.742080in}}%
\pgfpathlineto{\pgfqpoint{6.300646in}{1.742080in}}%
\pgfpathlineto{\pgfqpoint{6.301759in}{1.740640in}}%
\pgfpathlineto{\pgfqpoint{6.307506in}{1.740640in}}%
\pgfpathlineto{\pgfqpoint{6.308619in}{1.739192in}}%
\pgfpathlineto{\pgfqpoint{6.314242in}{1.739192in}}%
\pgfpathlineto{\pgfqpoint{6.315354in}{1.737736in}}%
\pgfpathlineto{\pgfqpoint{6.320959in}{1.737736in}}%
\pgfpathlineto{\pgfqpoint{6.322071in}{1.736273in}}%
\pgfpathlineto{\pgfqpoint{6.327989in}{1.736273in}}%
\pgfpathlineto{\pgfqpoint{6.329102in}{1.734802in}}%
\pgfpathlineto{\pgfqpoint{6.335069in}{1.734802in}}%
\pgfpathlineto{\pgfqpoint{6.336182in}{1.733324in}}%
\pgfpathlineto{\pgfqpoint{6.341770in}{1.733324in}}%
\pgfpathlineto{\pgfqpoint{6.342882in}{1.731837in}}%
\pgfpathlineto{\pgfqpoint{6.348528in}{1.731837in}}%
\pgfpathlineto{\pgfqpoint{6.349641in}{1.730343in}}%
\pgfpathlineto{\pgfqpoint{6.355555in}{1.730343in}}%
\pgfpathlineto{\pgfqpoint{6.356667in}{1.728840in}}%
\pgfpathlineto{\pgfqpoint{6.362379in}{1.728840in}}%
\pgfpathlineto{\pgfqpoint{6.363491in}{1.727329in}}%
\pgfpathlineto{\pgfqpoint{6.369109in}{1.727329in}}%
\pgfpathlineto{\pgfqpoint{6.370222in}{1.725810in}}%
\pgfpathlineto{\pgfqpoint{6.375771in}{1.725810in}}%
\pgfpathlineto{\pgfqpoint{6.376883in}{1.724283in}}%
\pgfpathlineto{\pgfqpoint{6.382415in}{1.724283in}}%
\pgfpathlineto{\pgfqpoint{6.383526in}{1.722748in}}%
\pgfpathlineto{\pgfqpoint{6.389132in}{1.722748in}}%
\pgfpathlineto{\pgfqpoint{6.390244in}{1.721203in}}%
\pgfpathlineto{\pgfqpoint{6.395828in}{1.721203in}}%
\pgfpathlineto{\pgfqpoint{6.396941in}{1.719651in}}%
\pgfpathlineto{\pgfqpoint{6.402628in}{1.719651in}}%
\pgfpathlineto{\pgfqpoint{6.403740in}{1.718089in}}%
\pgfpathlineto{\pgfqpoint{6.409328in}{1.718089in}}%
\pgfpathlineto{\pgfqpoint{6.410441in}{1.716519in}}%
\pgfpathlineto{\pgfqpoint{6.416015in}{1.716519in}}%
\pgfpathlineto{\pgfqpoint{6.417127in}{1.714940in}}%
\pgfpathlineto{\pgfqpoint{6.422653in}{1.714940in}}%
\pgfpathlineto{\pgfqpoint{6.423765in}{1.713351in}}%
\pgfpathlineto{\pgfqpoint{6.429103in}{1.713351in}}%
\pgfpathlineto{\pgfqpoint{6.430214in}{1.711754in}}%
\pgfpathlineto{\pgfqpoint{6.435784in}{1.711754in}}%
\pgfpathlineto{\pgfqpoint{6.436896in}{1.710148in}}%
\pgfpathlineto{\pgfqpoint{6.442394in}{1.710148in}}%
\pgfpathlineto{\pgfqpoint{6.443506in}{1.708532in}}%
\pgfpathlineto{\pgfqpoint{6.449121in}{1.708532in}}%
\pgfpathlineto{\pgfqpoint{6.450233in}{1.706906in}}%
\pgfpathlineto{\pgfqpoint{6.455756in}{1.706906in}}%
\pgfpathlineto{\pgfqpoint{6.456867in}{1.705272in}}%
\pgfpathlineto{\pgfqpoint{6.462324in}{1.705272in}}%
\pgfpathlineto{\pgfqpoint{6.463435in}{1.703627in}}%
\pgfpathlineto{\pgfqpoint{6.468912in}{1.703627in}}%
\pgfpathlineto{\pgfqpoint{6.470024in}{1.701973in}}%
\pgfpathlineto{\pgfqpoint{6.475369in}{1.701973in}}%
\pgfpathlineto{\pgfqpoint{6.476480in}{1.700309in}}%
\pgfpathlineto{\pgfqpoint{6.481818in}{1.700309in}}%
\pgfpathlineto{\pgfqpoint{6.482930in}{1.698635in}}%
\pgfpathlineto{\pgfqpoint{6.488178in}{1.698635in}}%
\pgfpathlineto{\pgfqpoint{6.489289in}{1.696951in}}%
\pgfpathlineto{\pgfqpoint{6.494556in}{1.696951in}}%
\pgfpathlineto{\pgfqpoint{6.495667in}{1.695257in}}%
\pgfpathlineto{\pgfqpoint{6.500955in}{1.695257in}}%
\pgfpathlineto{\pgfqpoint{6.502066in}{1.693552in}}%
\pgfpathlineto{\pgfqpoint{6.507251in}{1.693552in}}%
\pgfpathlineto{\pgfqpoint{6.508363in}{1.691837in}}%
\pgfpathlineto{\pgfqpoint{6.513775in}{1.691837in}}%
\pgfpathlineto{\pgfqpoint{6.514887in}{1.690111in}}%
\pgfpathlineto{\pgfqpoint{6.520118in}{1.690111in}}%
\pgfpathlineto{\pgfqpoint{6.521229in}{1.688374in}}%
\pgfpathlineto{\pgfqpoint{6.526332in}{1.688374in}}%
\pgfpathlineto{\pgfqpoint{6.527444in}{1.686627in}}%
\pgfpathlineto{\pgfqpoint{6.532715in}{1.686627in}}%
\pgfpathlineto{\pgfqpoint{6.533827in}{1.684868in}}%
\pgfpathlineto{\pgfqpoint{6.539015in}{1.684868in}}%
\pgfpathlineto{\pgfqpoint{6.540127in}{1.683098in}}%
\pgfpathlineto{\pgfqpoint{6.545111in}{1.683098in}}%
\pgfpathlineto{\pgfqpoint{6.546222in}{1.681317in}}%
\pgfpathlineto{\pgfqpoint{6.551300in}{1.681317in}}%
\pgfpathlineto{\pgfqpoint{6.552412in}{1.679525in}}%
\pgfpathlineto{\pgfqpoint{6.557542in}{1.679525in}}%
\pgfpathlineto{\pgfqpoint{6.558653in}{1.677721in}}%
\pgfpathlineto{\pgfqpoint{6.563771in}{1.677721in}}%
\pgfpathlineto{\pgfqpoint{6.564883in}{1.675905in}}%
\pgfpathlineto{\pgfqpoint{6.569817in}{1.675905in}}%
\pgfpathlineto{\pgfqpoint{6.570929in}{1.674077in}}%
\pgfpathlineto{\pgfqpoint{6.575809in}{1.674077in}}%
\pgfpathlineto{\pgfqpoint{6.576921in}{1.672237in}}%
\pgfpathlineto{\pgfqpoint{6.581771in}{1.672237in}}%
\pgfpathlineto{\pgfqpoint{6.582883in}{1.670385in}}%
\pgfpathlineto{\pgfqpoint{6.587824in}{1.670385in}}%
\pgfpathlineto{\pgfqpoint{6.588935in}{1.668521in}}%
\pgfpathlineto{\pgfqpoint{6.593807in}{1.668521in}}%
\pgfpathlineto{\pgfqpoint{6.594919in}{1.666644in}}%
\pgfpathlineto{\pgfqpoint{6.599686in}{1.666644in}}%
\pgfpathlineto{\pgfqpoint{6.600797in}{1.664754in}}%
\pgfpathlineto{\pgfqpoint{6.605619in}{1.664754in}}%
\pgfpathlineto{\pgfqpoint{6.606731in}{1.662852in}}%
\pgfpathlineto{\pgfqpoint{6.611501in}{1.662852in}}%
\pgfpathlineto{\pgfqpoint{6.612613in}{1.660936in}}%
\pgfpathlineto{\pgfqpoint{6.617347in}{1.660936in}}%
\pgfpathlineto{\pgfqpoint{6.618459in}{1.659007in}}%
\pgfpathlineto{\pgfqpoint{6.623139in}{1.659007in}}%
\pgfpathlineto{\pgfqpoint{6.624250in}{1.657065in}}%
\pgfpathlineto{\pgfqpoint{6.628906in}{1.657065in}}%
\pgfpathlineto{\pgfqpoint{6.630017in}{1.655109in}}%
\pgfpathlineto{\pgfqpoint{6.634784in}{1.655109in}}%
\pgfpathlineto{\pgfqpoint{6.635895in}{1.653139in}}%
\pgfpathlineto{\pgfqpoint{6.640443in}{1.653139in}}%
\pgfpathlineto{\pgfqpoint{6.641554in}{1.651155in}}%
\pgfpathlineto{\pgfqpoint{6.646128in}{1.651155in}}%
\pgfpathlineto{\pgfqpoint{6.647239in}{1.649157in}}%
\pgfpathlineto{\pgfqpoint{6.651669in}{1.649157in}}%
\pgfpathlineto{\pgfqpoint{6.652780in}{1.647145in}}%
\pgfpathlineto{\pgfqpoint{6.657244in}{1.647145in}}%
\pgfpathlineto{\pgfqpoint{6.658357in}{1.645118in}}%
\pgfpathlineto{\pgfqpoint{6.662757in}{1.645118in}}%
\pgfpathlineto{\pgfqpoint{6.663869in}{1.643076in}}%
\pgfpathlineto{\pgfqpoint{6.668314in}{1.643076in}}%
\pgfpathlineto{\pgfqpoint{6.669426in}{1.641019in}}%
\pgfpathlineto{\pgfqpoint{6.673785in}{1.641019in}}%
\pgfpathlineto{\pgfqpoint{6.674897in}{1.638946in}}%
\pgfpathlineto{\pgfqpoint{6.679166in}{1.638946in}}%
\pgfpathlineto{\pgfqpoint{6.680278in}{1.636858in}}%
\pgfpathlineto{\pgfqpoint{6.684599in}{1.636858in}}%
\pgfpathlineto{\pgfqpoint{6.685711in}{1.634755in}}%
\pgfpathlineto{\pgfqpoint{6.689832in}{1.634755in}}%
\pgfpathlineto{\pgfqpoint{6.690943in}{1.632635in}}%
\pgfpathlineto{\pgfqpoint{6.695255in}{1.632635in}}%
\pgfpathlineto{\pgfqpoint{6.696366in}{1.630499in}}%
\pgfpathlineto{\pgfqpoint{6.700581in}{1.630499in}}%
\pgfpathlineto{\pgfqpoint{6.701693in}{1.628347in}}%
\pgfpathlineto{\pgfqpoint{6.705749in}{1.628347in}}%
\pgfpathlineto{\pgfqpoint{6.706860in}{1.626178in}}%
\pgfpathlineto{\pgfqpoint{6.710810in}{1.626178in}}%
\pgfpathlineto{\pgfqpoint{6.711922in}{1.623992in}}%
\pgfpathlineto{\pgfqpoint{6.715897in}{1.623992in}}%
\pgfpathlineto{\pgfqpoint{6.717008in}{1.621788in}}%
\pgfpathlineto{\pgfqpoint{6.720900in}{1.621788in}}%
\pgfpathlineto{\pgfqpoint{6.722011in}{1.619568in}}%
\pgfpathlineto{\pgfqpoint{6.725909in}{1.619568in}}%
\pgfpathlineto{\pgfqpoint{6.727021in}{1.617329in}}%
\pgfpathlineto{\pgfqpoint{6.730804in}{1.617329in}}%
\pgfpathlineto{\pgfqpoint{6.731915in}{1.615072in}}%
\pgfpathlineto{\pgfqpoint{6.735704in}{1.615072in}}%
\pgfpathlineto{\pgfqpoint{6.736815in}{1.612797in}}%
\pgfpathlineto{\pgfqpoint{6.740557in}{1.612797in}}%
\pgfpathlineto{\pgfqpoint{6.741668in}{1.610502in}}%
\pgfpathlineto{\pgfqpoint{6.745327in}{1.610502in}}%
\pgfpathlineto{\pgfqpoint{6.746438in}{1.608189in}}%
\pgfpathlineto{\pgfqpoint{6.750061in}{1.608189in}}%
\pgfpathlineto{\pgfqpoint{6.751172in}{1.605857in}}%
\pgfpathlineto{\pgfqpoint{6.754633in}{1.605857in}}%
\pgfpathlineto{\pgfqpoint{6.755744in}{1.603505in}}%
\pgfpathlineto{\pgfqpoint{6.759222in}{1.603505in}}%
\pgfpathlineto{\pgfqpoint{6.760333in}{1.601133in}}%
\pgfpathlineto{\pgfqpoint{6.763854in}{1.601133in}}%
\pgfpathlineto{\pgfqpoint{6.764965in}{1.598740in}}%
\pgfpathlineto{\pgfqpoint{6.768396in}{1.598740in}}%
\pgfpathlineto{\pgfqpoint{6.769507in}{1.596327in}}%
\pgfpathlineto{\pgfqpoint{6.772821in}{1.596327in}}%
\pgfpathlineto{\pgfqpoint{6.773932in}{1.593893in}}%
\pgfpathlineto{\pgfqpoint{6.777190in}{1.593893in}}%
\pgfpathlineto{\pgfqpoint{6.778302in}{1.591437in}}%
\pgfpathlineto{\pgfqpoint{6.781511in}{1.591437in}}%
\pgfpathlineto{\pgfqpoint{6.782622in}{1.588959in}}%
\pgfpathlineto{\pgfqpoint{6.785824in}{1.588959in}}%
\pgfpathlineto{\pgfqpoint{6.786935in}{1.586459in}}%
\pgfpathlineto{\pgfqpoint{6.790178in}{1.586459in}}%
\pgfpathlineto{\pgfqpoint{6.791289in}{1.583937in}}%
\pgfpathlineto{\pgfqpoint{6.794165in}{1.583937in}}%
\pgfpathlineto{\pgfqpoint{6.795276in}{1.581391in}}%
\pgfpathlineto{\pgfqpoint{6.798220in}{1.581391in}}%
\pgfpathlineto{\pgfqpoint{6.799332in}{1.578823in}}%
\pgfpathlineto{\pgfqpoint{6.802266in}{1.578823in}}%
\pgfpathlineto{\pgfqpoint{6.803377in}{1.576230in}}%
\pgfpathlineto{\pgfqpoint{6.806257in}{1.576230in}}%
\pgfpathlineto{\pgfqpoint{6.807368in}{1.573613in}}%
\pgfpathlineto{\pgfqpoint{6.810105in}{1.573613in}}%
\pgfpathlineto{\pgfqpoint{6.811216in}{1.570971in}}%
\pgfpathlineto{\pgfqpoint{6.813874in}{1.570971in}}%
\pgfpathlineto{\pgfqpoint{6.814986in}{1.568303in}}%
\pgfpathlineto{\pgfqpoint{6.817681in}{1.568303in}}%
\pgfpathlineto{\pgfqpoint{6.818792in}{1.565610in}}%
\pgfpathlineto{\pgfqpoint{6.821307in}{1.565610in}}%
\pgfpathlineto{\pgfqpoint{6.822418in}{1.562891in}}%
\pgfpathlineto{\pgfqpoint{6.824968in}{1.562891in}}%
\pgfpathlineto{\pgfqpoint{6.826080in}{1.560145in}}%
\pgfpathlineto{\pgfqpoint{6.828513in}{1.560145in}}%
\pgfpathlineto{\pgfqpoint{6.829625in}{1.557371in}}%
\pgfpathlineto{\pgfqpoint{6.832060in}{1.557371in}}%
\pgfpathlineto{\pgfqpoint{6.833172in}{1.554570in}}%
\pgfpathlineto{\pgfqpoint{6.835515in}{1.554570in}}%
\pgfpathlineto{\pgfqpoint{6.836627in}{1.551740in}}%
\pgfpathlineto{\pgfqpoint{6.838840in}{1.551740in}}%
\pgfpathlineto{\pgfqpoint{6.839951in}{1.548882in}}%
\pgfpathlineto{\pgfqpoint{6.842017in}{1.548882in}}%
\pgfpathlineto{\pgfqpoint{6.843129in}{1.545993in}}%
\pgfpathlineto{\pgfqpoint{6.845230in}{1.545993in}}%
\pgfpathlineto{\pgfqpoint{6.846342in}{1.543075in}}%
\pgfpathlineto{\pgfqpoint{6.848304in}{1.543075in}}%
\pgfpathlineto{\pgfqpoint{6.849415in}{1.540125in}}%
\pgfpathlineto{\pgfqpoint{6.851410in}{1.540125in}}%
\pgfpathlineto{\pgfqpoint{6.852521in}{1.537144in}}%
\pgfpathlineto{\pgfqpoint{6.854397in}{1.537144in}}%
\pgfpathlineto{\pgfqpoint{6.855508in}{1.534131in}}%
\pgfpathlineto{\pgfqpoint{6.857359in}{1.534131in}}%
\pgfpathlineto{\pgfqpoint{6.858471in}{1.531085in}}%
\pgfpathlineto{\pgfqpoint{6.860186in}{1.531085in}}%
\pgfpathlineto{\pgfqpoint{6.861298in}{1.528005in}}%
\pgfpathlineto{\pgfqpoint{6.862966in}{1.528005in}}%
\pgfpathlineto{\pgfqpoint{6.864078in}{1.524890in}}%
\pgfpathlineto{\pgfqpoint{6.865624in}{1.524890in}}%
\pgfpathlineto{\pgfqpoint{6.866736in}{1.521741in}}%
\pgfpathlineto{\pgfqpoint{6.868221in}{1.521741in}}%
\pgfpathlineto{\pgfqpoint{6.869332in}{1.518555in}}%
\pgfpathlineto{\pgfqpoint{6.870806in}{1.518555in}}%
\pgfpathlineto{\pgfqpoint{6.871917in}{1.515333in}}%
\pgfpathlineto{\pgfqpoint{6.873184in}{1.515333in}}%
\pgfpathlineto{\pgfqpoint{6.874295in}{1.512073in}}%
\pgfpathlineto{\pgfqpoint{6.875529in}{1.512073in}}%
\pgfpathlineto{\pgfqpoint{6.876640in}{1.508774in}}%
\pgfpathlineto{\pgfqpoint{6.877885in}{1.508774in}}%
\pgfpathlineto{\pgfqpoint{6.878996in}{1.505436in}}%
\pgfpathlineto{\pgfqpoint{6.880163in}{1.505436in}}%
\pgfpathlineto{\pgfqpoint{6.881275in}{1.502058in}}%
\pgfpathlineto{\pgfqpoint{6.882393in}{1.502058in}}%
\pgfpathlineto{\pgfqpoint{6.883504in}{1.498638in}}%
\pgfpathlineto{\pgfqpoint{6.884417in}{1.498638in}}%
\pgfpathlineto{\pgfqpoint{6.885529in}{1.495175in}}%
\pgfpathlineto{\pgfqpoint{6.886410in}{1.495175in}}%
\pgfpathlineto{\pgfqpoint{6.887522in}{1.491669in}}%
\pgfpathlineto{\pgfqpoint{6.888357in}{1.491669in}}%
\pgfpathlineto{\pgfqpoint{6.889468in}{1.488119in}}%
\pgfpathlineto{\pgfqpoint{6.890214in}{1.488119in}}%
\pgfpathlineto{\pgfqpoint{6.891326in}{1.484522in}}%
\pgfpathlineto{\pgfqpoint{6.891985in}{1.484522in}}%
\pgfpathlineto{\pgfqpoint{6.893096in}{1.480879in}}%
\pgfpathlineto{\pgfqpoint{6.893732in}{1.480879in}}%
\pgfpathlineto{\pgfqpoint{6.894844in}{1.477187in}}%
\pgfpathlineto{\pgfqpoint{6.895338in}{1.477187in}}%
\pgfpathlineto{\pgfqpoint{6.896450in}{1.473445in}}%
\pgfpathlineto{\pgfqpoint{6.896904in}{1.473445in}}%
\pgfpathlineto{\pgfqpoint{6.898015in}{1.469653in}}%
\pgfpathlineto{\pgfqpoint{6.898434in}{1.469653in}}%
\pgfpathlineto{\pgfqpoint{6.899545in}{1.465809in}}%
\pgfpathlineto{\pgfqpoint{6.899840in}{1.465809in}}%
\pgfpathlineto{\pgfqpoint{6.900951in}{1.461910in}}%
\pgfpathlineto{\pgfqpoint{6.901195in}{1.461910in}}%
\pgfpathlineto{\pgfqpoint{6.902306in}{1.457957in}}%
\pgfpathlineto{\pgfqpoint{6.902439in}{1.457957in}}%
\pgfpathlineto{\pgfqpoint{6.903551in}{1.453946in}}%
\pgfpathlineto{\pgfqpoint{6.903640in}{1.453946in}}%
\pgfpathlineto{\pgfqpoint{6.904752in}{1.449877in}}%
\pgfpathlineto{\pgfqpoint{6.904801in}{1.449877in}}%
\pgfpathlineto{\pgfqpoint{6.905913in}{1.441556in}}%
\pgfpathlineto{\pgfqpoint{6.906927in}{1.441556in}}%
\pgfpathlineto{\pgfqpoint{6.908039in}{1.432979in}}%
\pgfpathlineto{\pgfqpoint{6.908821in}{1.432979in}}%
\pgfpathlineto{\pgfqpoint{6.909933in}{1.424130in}}%
\pgfpathlineto{\pgfqpoint{6.910453in}{1.424130in}}%
\pgfpathlineto{\pgfqpoint{6.911565in}{1.414991in}}%
\pgfpathlineto{\pgfqpoint{6.911924in}{1.414991in}}%
\pgfpathlineto{\pgfqpoint{6.913036in}{1.405541in}}%
\pgfpathlineto{\pgfqpoint{6.913154in}{1.405541in}}%
\pgfpathlineto{\pgfqpoint{6.914266in}{1.390738in}}%
\pgfpathlineto{\pgfqpoint{6.914685in}{1.390738in}}%
\pgfpathlineto{\pgfqpoint{6.915797in}{1.375105in}}%
\pgfpathlineto{\pgfqpoint{6.915833in}{1.375105in}}%
\pgfpathlineto{\pgfqpoint{6.916946in}{1.346926in}}%
\pgfpathlineto{\pgfqpoint{6.917107in}{1.346926in}}%
\pgfpathlineto{\pgfqpoint{6.918182in}{1.122377in}}%
\pgfpathlineto{\pgfqpoint{6.918182in}{1.122377in}}%
\pgfusepath{stroke}%
\end{pgfscope}%
\begin{pgfscope}%
\pgfpathrectangle{\pgfqpoint{1.000000in}{0.440000in}}{\pgfqpoint{6.200000in}{3.080000in}}%
\pgfusepath{clip}%
\pgfsetrectcap%
\pgfsetroundjoin%
\pgfsetlinewidth{1.204500pt}%
\definecolor{currentstroke}{rgb}{0.839216,0.152941,0.156863}%
\pgfsetstrokecolor{currentstroke}%
\pgfsetdash{}{0pt}%
\pgfpathmoveto{\pgfqpoint{0.990000in}{3.337222in}}%
\pgfpathlineto{\pgfqpoint{1.281818in}{3.337209in}}%
\pgfpathlineto{\pgfqpoint{1.569246in}{3.290528in}}%
\pgfpathlineto{\pgfqpoint{1.737381in}{3.274994in}}%
\pgfpathlineto{\pgfqpoint{1.856675in}{3.214949in}}%
\pgfpathlineto{\pgfqpoint{1.949206in}{3.201120in}}%
\pgfpathlineto{\pgfqpoint{2.024809in}{3.195739in}}%
\pgfpathlineto{\pgfqpoint{2.088731in}{3.187662in}}%
\pgfpathlineto{\pgfqpoint{2.144103in}{3.187293in}}%
\pgfpathlineto{\pgfqpoint{2.192944in}{3.178748in}}%
\pgfpathlineto{\pgfqpoint{2.236634in}{3.169109in}}%
\pgfpathlineto{\pgfqpoint{2.276156in}{3.157751in}}%
\pgfpathlineto{\pgfqpoint{2.312238in}{3.143474in}}%
\pgfpathlineto{\pgfqpoint{2.345429in}{3.139424in}}%
\pgfpathlineto{\pgfqpoint{2.376159in}{3.133731in}}%
\pgfpathlineto{\pgfqpoint{2.404769in}{3.133399in}}%
\pgfpathlineto{\pgfqpoint{2.431531in}{3.129823in}}%
\pgfpathlineto{\pgfqpoint{2.456670in}{3.122733in}}%
\pgfpathlineto{\pgfqpoint{2.480372in}{3.119520in}}%
\pgfpathlineto{\pgfqpoint{2.502792in}{3.109609in}}%
\pgfpathlineto{\pgfqpoint{2.524062in}{3.097840in}}%
\pgfpathlineto{\pgfqpoint{2.582018in}{3.084274in}}%
\pgfpathlineto{\pgfqpoint{2.599666in}{3.082534in}}%
\pgfpathlineto{\pgfqpoint{2.616594in}{3.078106in}}%
\pgfpathlineto{\pgfqpoint{2.632857in}{3.076889in}}%
\pgfpathlineto{\pgfqpoint{2.648507in}{3.074198in}}%
\pgfpathlineto{\pgfqpoint{2.663588in}{3.069204in}}%
\pgfpathlineto{\pgfqpoint{2.678139in}{3.066078in}}%
\pgfpathlineto{\pgfqpoint{2.692197in}{3.064919in}}%
\pgfpathlineto{\pgfqpoint{2.705794in}{3.058596in}}%
\pgfpathlineto{\pgfqpoint{2.718959in}{3.055837in}}%
\pgfpathlineto{\pgfqpoint{2.744099in}{3.054701in}}%
\pgfpathlineto{\pgfqpoint{2.756119in}{3.050592in}}%
\pgfpathlineto{\pgfqpoint{2.767801in}{3.050448in}}%
\pgfpathlineto{\pgfqpoint{2.779162in}{3.047554in}}%
\pgfpathlineto{\pgfqpoint{2.790221in}{3.041268in}}%
\pgfpathlineto{\pgfqpoint{2.800992in}{3.039086in}}%
\pgfpathlineto{\pgfqpoint{2.811491in}{3.038825in}}%
\pgfpathlineto{\pgfqpoint{2.821730in}{3.033708in}}%
\pgfpathlineto{\pgfqpoint{2.831722in}{3.033610in}}%
\pgfpathlineto{\pgfqpoint{2.841480in}{3.024887in}}%
\pgfpathlineto{\pgfqpoint{2.851013in}{3.022166in}}%
\pgfpathlineto{\pgfqpoint{2.860332in}{3.017293in}}%
\pgfpathlineto{\pgfqpoint{2.887094in}{3.014181in}}%
\pgfpathlineto{\pgfqpoint{2.895644in}{3.010261in}}%
\pgfpathlineto{\pgfqpoint{2.904022in}{3.009988in}}%
\pgfpathlineto{\pgfqpoint{2.912233in}{3.007515in}}%
\pgfpathlineto{\pgfqpoint{2.928184in}{3.006085in}}%
\pgfpathlineto{\pgfqpoint{2.951016in}{3.005499in}}%
\pgfpathlineto{\pgfqpoint{2.958355in}{3.004927in}}%
\pgfpathlineto{\pgfqpoint{2.965567in}{3.002517in}}%
\pgfpathlineto{\pgfqpoint{2.993222in}{2.997747in}}%
\pgfpathlineto{\pgfqpoint{3.019148in}{2.989069in}}%
\pgfpathlineto{\pgfqpoint{3.025383in}{2.987640in}}%
\pgfpathlineto{\pgfqpoint{3.031527in}{2.984893in}}%
\pgfpathlineto{\pgfqpoint{3.043547in}{2.983108in}}%
\pgfpathlineto{\pgfqpoint{3.055229in}{2.982773in}}%
\pgfpathlineto{\pgfqpoint{3.066590in}{2.981244in}}%
\pgfpathlineto{\pgfqpoint{3.083070in}{2.980146in}}%
\pgfpathlineto{\pgfqpoint{3.093703in}{2.977904in}}%
\pgfpathlineto{\pgfqpoint{3.098919in}{2.977204in}}%
\pgfpathlineto{\pgfqpoint{3.104070in}{2.972671in}}%
\pgfpathlineto{\pgfqpoint{3.114184in}{2.972358in}}%
\pgfpathlineto{\pgfqpoint{3.124058in}{2.967940in}}%
\pgfpathlineto{\pgfqpoint{3.128908in}{2.967179in}}%
\pgfpathlineto{\pgfqpoint{3.138441in}{2.962551in}}%
\pgfpathlineto{\pgfqpoint{3.147760in}{2.961883in}}%
\pgfpathlineto{\pgfqpoint{3.156874in}{2.957790in}}%
\pgfpathlineto{\pgfqpoint{3.161357in}{2.957680in}}%
\pgfpathlineto{\pgfqpoint{3.170180in}{2.954322in}}%
\pgfpathlineto{\pgfqpoint{3.187282in}{2.951468in}}%
\pgfpathlineto{\pgfqpoint{3.195576in}{2.949738in}}%
\pgfpathlineto{\pgfqpoint{3.219506in}{2.942649in}}%
\pgfpathlineto{\pgfqpoint{3.223364in}{2.942533in}}%
\pgfpathlineto{\pgfqpoint{3.230972in}{2.936446in}}%
\pgfpathlineto{\pgfqpoint{3.234725in}{2.934779in}}%
\pgfpathlineto{\pgfqpoint{3.238444in}{2.934361in}}%
\pgfpathlineto{\pgfqpoint{3.242130in}{2.932501in}}%
\pgfpathlineto{\pgfqpoint{3.256555in}{2.931817in}}%
\pgfpathlineto{\pgfqpoint{3.260084in}{2.929697in}}%
\pgfpathlineto{\pgfqpoint{3.303422in}{2.922759in}}%
\pgfpathlineto{\pgfqpoint{3.306576in}{2.922510in}}%
\pgfpathlineto{\pgfqpoint{3.315895in}{2.917599in}}%
\pgfpathlineto{\pgfqpoint{3.318955in}{2.914276in}}%
\pgfpathlineto{\pgfqpoint{3.328003in}{2.912514in}}%
\pgfpathlineto{\pgfqpoint{3.330975in}{2.909544in}}%
\pgfpathlineto{\pgfqpoint{3.348377in}{2.906807in}}%
\pgfpathlineto{\pgfqpoint{3.359585in}{2.903998in}}%
\pgfpathlineto{\pgfqpoint{3.362340in}{2.903664in}}%
\pgfpathlineto{\pgfqpoint{3.367796in}{2.900710in}}%
\pgfpathlineto{\pgfqpoint{3.378498in}{2.899899in}}%
\pgfpathlineto{\pgfqpoint{3.383747in}{2.898724in}}%
\pgfpathlineto{\pgfqpoint{3.396586in}{2.897225in}}%
\pgfpathlineto{\pgfqpoint{3.464105in}{2.884550in}}%
\pgfpathlineto{\pgfqpoint{3.474711in}{2.882204in}}%
\pgfpathlineto{\pgfqpoint{3.487090in}{2.880734in}}%
\pgfpathlineto{\pgfqpoint{3.497131in}{2.877965in}}%
\pgfpathlineto{\pgfqpoint{3.501080in}{2.874897in}}%
\pgfpathlineto{\pgfqpoint{3.520281in}{2.871646in}}%
\pgfpathlineto{\pgfqpoint{3.531389in}{2.869692in}}%
\pgfpathlineto{\pgfqpoint{3.535027in}{2.868558in}}%
\pgfpathlineto{\pgfqpoint{3.543983in}{2.866887in}}%
\pgfpathlineto{\pgfqpoint{3.552750in}{2.865359in}}%
\pgfpathlineto{\pgfqpoint{3.556206in}{2.864361in}}%
\pgfpathlineto{\pgfqpoint{3.559633in}{2.862665in}}%
\pgfpathlineto{\pgfqpoint{3.571409in}{2.858063in}}%
\pgfpathlineto{\pgfqpoint{3.581244in}{2.857223in}}%
\pgfpathlineto{\pgfqpoint{3.586075in}{2.854979in}}%
\pgfpathlineto{\pgfqpoint{3.590851in}{2.854103in}}%
\pgfpathlineto{\pgfqpoint{3.600240in}{2.852824in}}%
\pgfpathlineto{\pgfqpoint{3.606383in}{2.851458in}}%
\pgfpathlineto{\pgfqpoint{3.609421in}{2.849762in}}%
\pgfpathlineto{\pgfqpoint{3.615431in}{2.848490in}}%
\pgfpathlineto{\pgfqpoint{3.619882in}{2.847021in}}%
\pgfpathlineto{\pgfqpoint{3.641447in}{2.843867in}}%
\pgfpathlineto{\pgfqpoint{3.649768in}{2.841596in}}%
\pgfpathlineto{\pgfqpoint{3.660610in}{2.839866in}}%
\pgfpathlineto{\pgfqpoint{3.669869in}{2.838098in}}%
\pgfpathlineto{\pgfqpoint{3.677645in}{2.836828in}}%
\pgfpathlineto{\pgfqpoint{3.736213in}{2.825453in}}%
\pgfpathlineto{\pgfqpoint{3.739544in}{2.822640in}}%
\pgfpathlineto{\pgfqpoint{3.746126in}{2.821357in}}%
\pgfpathlineto{\pgfqpoint{3.770433in}{2.814975in}}%
\pgfpathlineto{\pgfqpoint{3.778564in}{2.813348in}}%
\pgfpathlineto{\pgfqpoint{3.799179in}{2.808631in}}%
\pgfpathlineto{\pgfqpoint{3.802992in}{2.807942in}}%
\pgfpathlineto{\pgfqpoint{3.811445in}{2.807062in}}%
\pgfpathlineto{\pgfqpoint{3.836694in}{2.798333in}}%
\pgfpathlineto{\pgfqpoint{3.838440in}{2.796998in}}%
\pgfpathlineto{\pgfqpoint{3.845351in}{2.795657in}}%
\pgfpathlineto{\pgfqpoint{3.852991in}{2.793782in}}%
\pgfpathlineto{\pgfqpoint{3.856342in}{2.793055in}}%
\pgfpathlineto{\pgfqpoint{3.858007in}{2.792259in}}%
\pgfpathlineto{\pgfqpoint{3.861318in}{2.791261in}}%
\pgfpathlineto{\pgfqpoint{3.866235in}{2.790093in}}%
\pgfpathlineto{\pgfqpoint{3.873503in}{2.788684in}}%
\pgfpathlineto{\pgfqpoint{3.887668in}{2.785948in}}%
\pgfpathlineto{\pgfqpoint{3.893811in}{2.784142in}}%
\pgfpathlineto{\pgfqpoint{3.895333in}{2.784040in}}%
\pgfpathlineto{\pgfqpoint{3.896849in}{2.782772in}}%
\pgfpathlineto{\pgfqpoint{3.907310in}{2.780150in}}%
\pgfpathlineto{\pgfqpoint{3.928875in}{2.774122in}}%
\pgfpathlineto{\pgfqpoint{3.932363in}{2.772986in}}%
\pgfpathlineto{\pgfqpoint{3.952694in}{2.767780in}}%
\pgfpathlineto{\pgfqpoint{3.959255in}{2.766145in}}%
\pgfpathlineto{\pgfqpoint{3.961851in}{2.765524in}}%
\pgfpathlineto{\pgfqpoint{3.967633in}{2.764022in}}%
\pgfpathlineto{\pgfqpoint{3.971443in}{2.762069in}}%
\pgfpathlineto{\pgfqpoint{3.973964in}{2.760845in}}%
\pgfpathlineto{\pgfqpoint{3.985121in}{2.757744in}}%
\pgfpathlineto{\pgfqpoint{3.989986in}{2.756653in}}%
\pgfpathlineto{\pgfqpoint{3.994195in}{2.755632in}}%
\pgfpathlineto{\pgfqpoint{3.996582in}{2.754912in}}%
\pgfpathlineto{\pgfqpoint{4.001902in}{2.753924in}}%
\pgfpathlineto{\pgfqpoint{4.100729in}{2.727537in}}%
\pgfpathlineto{\pgfqpoint{4.103036in}{2.726446in}}%
\pgfpathlineto{\pgfqpoint{4.113040in}{2.724464in}}%
\pgfpathlineto{\pgfqpoint{4.121049in}{2.721912in}}%
\pgfpathlineto{\pgfqpoint{4.150802in}{2.715079in}}%
\pgfpathlineto{\pgfqpoint{4.158523in}{2.712925in}}%
\pgfpathlineto{\pgfqpoint{4.166103in}{2.710977in}}%
\pgfpathlineto{\pgfqpoint{4.179712in}{2.707444in}}%
\pgfpathlineto{\pgfqpoint{4.181621in}{2.705548in}}%
\pgfpathlineto{\pgfqpoint{4.185411in}{2.704504in}}%
\pgfpathlineto{\pgfqpoint{4.195844in}{2.701900in}}%
\pgfpathlineto{\pgfqpoint{4.206020in}{2.699687in}}%
\pgfpathlineto{\pgfqpoint{4.207811in}{2.698900in}}%
\pgfpathlineto{\pgfqpoint{4.215953in}{2.696851in}}%
\pgfpathlineto{\pgfqpoint{4.221178in}{2.695304in}}%
\pgfpathlineto{\pgfqpoint{4.226337in}{2.694315in}}%
\pgfpathlineto{\pgfqpoint{4.229742in}{2.693228in}}%
\pgfpathlineto{\pgfqpoint{4.237468in}{2.691041in}}%
\pgfpathlineto{\pgfqpoint{4.241113in}{2.690085in}}%
\pgfpathlineto{\pgfqpoint{4.245706in}{2.688981in}}%
\pgfpathlineto{\pgfqpoint{4.255380in}{2.686615in}}%
\pgfpathlineto{\pgfqpoint{4.258871in}{2.685157in}}%
\pgfpathlineto{\pgfqpoint{4.260133in}{2.684131in}}%
\pgfpathlineto{\pgfqpoint{4.300057in}{2.672849in}}%
\pgfpathlineto{\pgfqpoint{4.303193in}{2.671828in}}%
\pgfpathlineto{\pgfqpoint{4.309954in}{2.670547in}}%
\pgfpathlineto{\pgfqpoint{4.333578in}{2.665100in}}%
\pgfpathlineto{\pgfqpoint{4.341940in}{2.662891in}}%
\pgfpathlineto{\pgfqpoint{4.347592in}{2.661833in}}%
\pgfpathlineto{\pgfqpoint{4.350896in}{2.660505in}}%
\pgfpathlineto{\pgfqpoint{4.356429in}{2.658206in}}%
\pgfpathlineto{\pgfqpoint{4.360406in}{2.656914in}}%
\pgfpathlineto{\pgfqpoint{4.376184in}{2.652401in}}%
\pgfpathlineto{\pgfqpoint{4.379506in}{2.651302in}}%
\pgfpathlineto{\pgfqpoint{4.382800in}{2.650199in}}%
\pgfpathlineto{\pgfqpoint{4.389082in}{2.648425in}}%
\pgfpathlineto{\pgfqpoint{4.392759in}{2.647366in}}%
\pgfpathlineto{\pgfqpoint{4.395042in}{2.646584in}}%
\pgfpathlineto{\pgfqpoint{4.400019in}{2.645518in}}%
\pgfpathlineto{\pgfqpoint{4.432240in}{2.636321in}}%
\pgfpathlineto{\pgfqpoint{4.454715in}{2.628905in}}%
\pgfpathlineto{\pgfqpoint{4.458443in}{2.627889in}}%
\pgfpathlineto{\pgfqpoint{4.462138in}{2.626869in}}%
\pgfpathlineto{\pgfqpoint{4.465608in}{2.625906in}}%
\pgfpathlineto{\pgfqpoint{4.469810in}{2.624880in}}%
\pgfpathlineto{\pgfqpoint{4.474158in}{2.623423in}}%
\pgfpathlineto{\pgfqpoint{4.480503in}{2.622326in}}%
\pgfpathlineto{\pgfqpoint{4.482167in}{2.621531in}}%
\pgfpathlineto{\pgfqpoint{4.610853in}{2.579930in}}%
\pgfpathlineto{\pgfqpoint{4.615693in}{2.578861in}}%
\pgfpathlineto{\pgfqpoint{4.623637in}{2.575847in}}%
\pgfpathlineto{\pgfqpoint{4.626643in}{2.574835in}}%
\pgfpathlineto{\pgfqpoint{4.628332in}{2.574110in}}%
\pgfpathlineto{\pgfqpoint{4.634254in}{2.571997in}}%
\pgfpathlineto{\pgfqpoint{4.637691in}{2.570970in}}%
\pgfpathlineto{\pgfqpoint{4.640093in}{2.569793in}}%
\pgfpathlineto{\pgfqpoint{4.643232in}{2.568685in}}%
\pgfpathlineto{\pgfqpoint{4.645228in}{2.567944in}}%
\pgfpathlineto{\pgfqpoint{4.648576in}{2.566902in}}%
\pgfpathlineto{\pgfqpoint{4.652754in}{2.565633in}}%
\pgfpathlineto{\pgfqpoint{4.656647in}{2.564583in}}%
\pgfpathlineto{\pgfqpoint{4.663613in}{2.561562in}}%
\pgfpathlineto{\pgfqpoint{4.667878in}{2.560496in}}%
\pgfpathlineto{\pgfqpoint{4.673497in}{2.558123in}}%
\pgfpathlineto{\pgfqpoint{4.678238in}{2.557122in}}%
\pgfpathlineto{\pgfqpoint{4.680073in}{2.556504in}}%
\pgfpathlineto{\pgfqpoint{4.684287in}{2.555419in}}%
\pgfpathlineto{\pgfqpoint{4.687896in}{2.554252in}}%
\pgfpathlineto{\pgfqpoint{4.691252in}{2.553158in}}%
\pgfpathlineto{\pgfqpoint{4.695354in}{2.551588in}}%
\pgfpathlineto{\pgfqpoint{4.698760in}{2.550484in}}%
\pgfpathlineto{\pgfqpoint{4.701595in}{2.549772in}}%
\pgfpathlineto{\pgfqpoint{4.704734in}{2.548740in}}%
\pgfpathlineto{\pgfqpoint{4.708064in}{2.547624in}}%
\pgfpathlineto{\pgfqpoint{4.710092in}{2.546584in}}%
\pgfpathlineto{\pgfqpoint{4.712745in}{2.545943in}}%
\pgfpathlineto{\pgfqpoint{4.715906in}{2.544896in}}%
\pgfpathlineto{\pgfqpoint{4.718522in}{2.544251in}}%
\pgfpathlineto{\pgfqpoint{4.721640in}{2.543198in}}%
\pgfpathlineto{\pgfqpoint{4.725351in}{2.541653in}}%
\pgfpathlineto{\pgfqpoint{4.728112in}{2.540590in}}%
\pgfpathlineto{\pgfqpoint{4.731260in}{2.539441in}}%
\pgfpathlineto{\pgfqpoint{4.734586in}{2.538371in}}%
\pgfpathlineto{\pgfqpoint{4.736887in}{2.537875in}}%
\pgfpathlineto{\pgfqpoint{4.739672in}{2.536798in}}%
\pgfpathlineto{\pgfqpoint{4.742242in}{2.536050in}}%
\pgfpathlineto{\pgfqpoint{4.745773in}{2.534966in}}%
\pgfpathlineto{\pgfqpoint{4.747723in}{2.533794in}}%
\pgfpathlineto{\pgfqpoint{4.750726in}{2.532701in}}%
\pgfpathlineto{\pgfqpoint{4.753707in}{2.531519in}}%
\pgfpathlineto{\pgfqpoint{4.756763in}{2.530417in}}%
\pgfpathlineto{\pgfqpoint{4.759418in}{2.529481in}}%
\pgfpathlineto{\pgfqpoint{4.762525in}{2.528371in}}%
\pgfpathlineto{\pgfqpoint{4.764678in}{2.527858in}}%
\pgfpathlineto{\pgfqpoint{4.767468in}{2.526827in}}%
\pgfpathlineto{\pgfqpoint{4.769595in}{2.526311in}}%
\pgfpathlineto{\pgfqpoint{4.772627in}{2.525274in}}%
\pgfpathlineto{\pgfqpoint{4.775365in}{2.524148in}}%
\pgfpathlineto{\pgfqpoint{4.777542in}{2.523103in}}%
\pgfpathlineto{\pgfqpoint{4.779888in}{2.522318in}}%
\pgfpathlineto{\pgfqpoint{4.783025in}{2.521266in}}%
\pgfpathlineto{\pgfqpoint{4.785251in}{2.520563in}}%
\pgfpathlineto{\pgfqpoint{4.788613in}{2.519506in}}%
\pgfpathlineto{\pgfqpoint{4.791247in}{2.518355in}}%
\pgfpathlineto{\pgfqpoint{4.794908in}{2.517289in}}%
\pgfpathlineto{\pgfqpoint{4.797761in}{2.515950in}}%
\pgfpathlineto{\pgfqpoint{4.800510in}{2.514875in}}%
\pgfpathlineto{\pgfqpoint{4.802474in}{2.514336in}}%
\pgfpathlineto{\pgfqpoint{4.805530in}{2.513254in}}%
\pgfpathlineto{\pgfqpoint{4.807807in}{2.512440in}}%
\pgfpathlineto{\pgfqpoint{4.811075in}{2.511351in}}%
\pgfpathlineto{\pgfqpoint{4.813405in}{2.510440in}}%
\pgfpathlineto{\pgfqpoint{4.817204in}{2.509343in}}%
\pgfpathlineto{\pgfqpoint{4.819255in}{2.508701in}}%
\pgfpathlineto{\pgfqpoint{4.822028in}{2.507689in}}%
\pgfpathlineto{\pgfqpoint{4.823488in}{2.507043in}}%
\pgfpathlineto{\pgfqpoint{4.825669in}{2.506211in}}%
\pgfpathlineto{\pgfqpoint{4.828719in}{2.505190in}}%
\pgfpathlineto{\pgfqpoint{4.830873in}{2.504352in}}%
\pgfpathlineto{\pgfqpoint{4.833332in}{2.503325in}}%
\pgfpathlineto{\pgfqpoint{4.835855in}{2.502011in}}%
\pgfpathlineto{\pgfqpoint{4.838598in}{2.501069in}}%
\pgfpathlineto{\pgfqpoint{4.840079in}{2.500313in}}%
\pgfpathlineto{\pgfqpoint{4.842639in}{2.499271in}}%
\pgfpathlineto{\pgfqpoint{4.844645in}{2.498510in}}%
\pgfpathlineto{\pgfqpoint{4.847559in}{2.497460in}}%
\pgfpathlineto{\pgfqpoint{4.849845in}{2.496311in}}%
\pgfpathlineto{\pgfqpoint{4.852497in}{2.495253in}}%
\pgfpathlineto{\pgfqpoint{4.854230in}{2.494771in}}%
\pgfpathlineto{\pgfqpoint{4.856779in}{2.493707in}}%
\pgfpathlineto{\pgfqpoint{4.858792in}{2.492834in}}%
\pgfpathlineto{\pgfqpoint{4.861240in}{2.491763in}}%
\pgfpathlineto{\pgfqpoint{4.863158in}{2.490981in}}%
\pgfpathlineto{\pgfqpoint{4.866530in}{2.489903in}}%
\pgfpathlineto{\pgfqpoint{4.868351in}{2.489313in}}%
\pgfpathlineto{\pgfqpoint{4.870598in}{2.488229in}}%
\pgfpathlineto{\pgfqpoint{4.872473in}{2.487536in}}%
\pgfpathlineto{\pgfqpoint{4.875628in}{2.486444in}}%
\pgfpathlineto{\pgfqpoint{4.877552in}{2.485648in}}%
\pgfpathlineto{\pgfqpoint{4.880739in}{2.484549in}}%
\pgfpathlineto{\pgfqpoint{4.882709in}{2.483646in}}%
\pgfpathlineto{\pgfqpoint{4.886622in}{2.482539in}}%
\pgfpathlineto{\pgfqpoint{4.888358in}{2.481933in}}%
\pgfpathlineto{\pgfqpoint{4.891600in}{2.480921in}}%
\pgfpathlineto{\pgfqpoint{4.893315in}{2.480312in}}%
\pgfpathlineto{\pgfqpoint{4.896519in}{2.479293in}}%
\pgfpathlineto{\pgfqpoint{4.898281in}{2.478578in}}%
\pgfpathlineto{\pgfqpoint{4.901246in}{2.477554in}}%
\pgfpathlineto{\pgfqpoint{4.903055in}{2.476731in}}%
\pgfpathlineto{\pgfqpoint{4.906715in}{2.475699in}}%
\pgfpathlineto{\pgfqpoint{4.908500in}{2.474871in}}%
\pgfpathlineto{\pgfqpoint{4.910869in}{2.473833in}}%
\pgfpathlineto{\pgfqpoint{4.912767in}{2.472895in}}%
\pgfpathlineto{\pgfqpoint{4.914916in}{2.471849in}}%
\pgfpathlineto{\pgfqpoint{4.916667in}{2.471115in}}%
\pgfpathlineto{\pgfqpoint{4.919568in}{2.470062in}}%
\pgfpathlineto{\pgfqpoint{4.921107in}{2.469640in}}%
\pgfpathlineto{\pgfqpoint{4.923786in}{2.468582in}}%
\pgfpathlineto{\pgfqpoint{4.925309in}{2.468157in}}%
\pgfpathlineto{\pgfqpoint{4.927836in}{2.467093in}}%
\pgfpathlineto{\pgfqpoint{4.929470in}{2.466453in}}%
\pgfpathlineto{\pgfqpoint{4.932968in}{2.465383in}}%
\pgfpathlineto{\pgfqpoint{4.934644in}{2.464631in}}%
\pgfpathlineto{\pgfqpoint{4.937545in}{2.463554in}}%
\pgfpathlineto{\pgfqpoint{4.939141in}{2.462905in}}%
\pgfpathlineto{\pgfqpoint{4.941402in}{2.461821in}}%
\pgfpathlineto{\pgfqpoint{4.943166in}{2.460842in}}%
\pgfpathlineto{\pgfqpoint{4.945587in}{2.459750in}}%
\pgfpathlineto{\pgfqpoint{4.947393in}{2.458763in}}%
\pgfpathlineto{\pgfqpoint{4.951041in}{2.457663in}}%
\pgfpathlineto{\pgfqpoint{4.952527in}{2.457221in}}%
\pgfpathlineto{\pgfqpoint{4.955365in}{2.456115in}}%
\pgfpathlineto{\pgfqpoint{4.956953in}{2.455449in}}%
\pgfpathlineto{\pgfqpoint{4.960053in}{2.454447in}}%
\pgfpathlineto{\pgfqpoint{4.961623in}{2.453777in}}%
\pgfpathlineto{\pgfqpoint{4.964400in}{2.452769in}}%
\pgfpathlineto{\pgfqpoint{4.965953in}{2.452095in}}%
\pgfpathlineto{\pgfqpoint{4.968644in}{2.451080in}}%
\pgfpathlineto{\pgfqpoint{4.970182in}{2.450402in}}%
\pgfpathlineto{\pgfqpoint{4.973467in}{2.449382in}}%
\pgfpathlineto{\pgfqpoint{4.975100in}{2.448586in}}%
\pgfpathlineto{\pgfqpoint{4.978625in}{2.447559in}}%
\pgfpathlineto{\pgfqpoint{4.980182in}{2.446872in}}%
\pgfpathlineto{\pgfqpoint{4.982562in}{2.445838in}}%
\pgfpathlineto{\pgfqpoint{4.983994in}{2.445378in}}%
\pgfpathlineto{\pgfqpoint{4.987281in}{2.444339in}}%
\pgfpathlineto{\pgfqpoint{4.988697in}{2.443876in}}%
\pgfpathlineto{\pgfqpoint{4.992055in}{2.442831in}}%
\pgfpathlineto{\pgfqpoint{4.993562in}{2.442133in}}%
\pgfpathlineto{\pgfqpoint{4.996561in}{2.441082in}}%
\pgfpathlineto{\pgfqpoint{4.998052in}{2.440379in}}%
\pgfpathlineto{\pgfqpoint{5.001493in}{2.439321in}}%
\pgfpathlineto{\pgfqpoint{5.002967in}{2.438732in}}%
\pgfpathlineto{\pgfqpoint{5.005429in}{2.437667in}}%
\pgfpathlineto{\pgfqpoint{5.006784in}{2.437312in}}%
\pgfpathlineto{\pgfqpoint{5.008654in}{2.436242in}}%
\pgfpathlineto{\pgfqpoint{5.010154in}{2.435527in}}%
\pgfpathlineto{\pgfqpoint{5.013190in}{2.434451in}}%
\pgfpathlineto{\pgfqpoint{5.014571in}{2.433971in}}%
\pgfpathlineto{\pgfqpoint{5.017574in}{2.432889in}}%
\pgfpathlineto{\pgfqpoint{5.018992in}{2.432285in}}%
\pgfpathlineto{\pgfqpoint{5.021863in}{2.431196in}}%
\pgfpathlineto{\pgfqpoint{5.023316in}{2.430589in}}%
\pgfpathlineto{\pgfqpoint{5.026108in}{2.429494in}}%
\pgfpathlineto{\pgfqpoint{5.027645in}{2.428639in}}%
\pgfpathlineto{\pgfqpoint{5.030998in}{2.427535in}}%
\pgfpathlineto{\pgfqpoint{5.032419in}{2.426920in}}%
\pgfpathlineto{\pgfqpoint{5.034859in}{2.425810in}}%
\pgfpathlineto{\pgfqpoint{5.036267in}{2.425191in}}%
\pgfpathlineto{\pgfqpoint{5.038733in}{2.424199in}}%
\pgfpathlineto{\pgfqpoint{5.040176in}{2.423452in}}%
\pgfpathlineto{\pgfqpoint{5.043238in}{2.422453in}}%
\pgfpathlineto{\pgfqpoint{5.044761in}{2.421576in}}%
\pgfpathlineto{\pgfqpoint{5.047554in}{2.420570in}}%
\pgfpathlineto{\pgfqpoint{5.048967in}{2.419939in}}%
\pgfpathlineto{\pgfqpoint{5.052106in}{2.418928in}}%
\pgfpathlineto{\pgfqpoint{5.053457in}{2.418420in}}%
\pgfpathlineto{\pgfqpoint{5.055822in}{2.417403in}}%
\pgfpathlineto{\pgfqpoint{5.057207in}{2.416765in}}%
\pgfpathlineto{\pgfqpoint{5.060239in}{2.415742in}}%
\pgfpathlineto{\pgfqpoint{5.061609in}{2.415101in}}%
\pgfpathlineto{\pgfqpoint{5.064291in}{2.414071in}}%
\pgfpathlineto{\pgfqpoint{5.065739in}{2.413297in}}%
\pgfpathlineto{\pgfqpoint{5.068529in}{2.412260in}}%
\pgfpathlineto{\pgfqpoint{5.069962in}{2.411481in}}%
\pgfpathlineto{\pgfqpoint{5.072680in}{2.410438in}}%
\pgfpathlineto{\pgfqpoint{5.074054in}{2.409784in}}%
\pgfpathlineto{\pgfqpoint{5.076613in}{2.408735in}}%
\pgfpathlineto{\pgfqpoint{5.077931in}{2.408209in}}%
\pgfpathlineto{\pgfqpoint{5.080030in}{2.407153in}}%
\pgfpathlineto{\pgfqpoint{5.081337in}{2.406624in}}%
\pgfpathlineto{\pgfqpoint{5.084024in}{2.405563in}}%
\pgfpathlineto{\pgfqpoint{5.085405in}{2.404897in}}%
\pgfpathlineto{\pgfqpoint{5.087381in}{2.403829in}}%
\pgfpathlineto{\pgfqpoint{5.088750in}{2.403160in}}%
\pgfpathlineto{\pgfqpoint{5.091560in}{2.402085in}}%
\pgfpathlineto{\pgfqpoint{5.092916in}{2.401412in}}%
\pgfpathlineto{\pgfqpoint{5.095908in}{2.400330in}}%
\pgfpathlineto{\pgfqpoint{5.097207in}{2.399788in}}%
\pgfpathlineto{\pgfqpoint{5.099961in}{2.398700in}}%
\pgfpathlineto{\pgfqpoint{5.101289in}{2.398018in}}%
\pgfpathlineto{\pgfqpoint{5.104180in}{2.396924in}}%
\pgfpathlineto{\pgfqpoint{5.105454in}{2.396512in}}%
\pgfpathlineto{\pgfqpoint{5.107215in}{2.395411in}}%
\pgfpathlineto{\pgfqpoint{5.108520in}{2.394859in}}%
\pgfpathlineto{\pgfqpoint{5.110916in}{2.393752in}}%
\pgfpathlineto{\pgfqpoint{5.112210in}{2.393197in}}%
\pgfpathlineto{\pgfqpoint{5.114344in}{2.392223in}}%
\pgfpathlineto{\pgfqpoint{5.115627in}{2.391664in}}%
\pgfpathlineto{\pgfqpoint{5.117385in}{2.390685in}}%
\pgfpathlineto{\pgfqpoint{5.118699in}{2.390123in}}%
\pgfpathlineto{\pgfqpoint{5.121590in}{2.389138in}}%
\pgfpathlineto{\pgfqpoint{5.122969in}{2.388291in}}%
\pgfpathlineto{\pgfqpoint{5.125518in}{2.387299in}}%
\pgfpathlineto{\pgfqpoint{5.126806in}{2.386731in}}%
\pgfpathlineto{\pgfqpoint{5.129177in}{2.385734in}}%
\pgfpathlineto{\pgfqpoint{5.130492in}{2.385019in}}%
\pgfpathlineto{\pgfqpoint{5.133110in}{2.384016in}}%
\pgfpathlineto{\pgfqpoint{5.134414in}{2.383441in}}%
\pgfpathlineto{\pgfqpoint{5.137007in}{2.382432in}}%
\pgfpathlineto{\pgfqpoint{5.138298in}{2.381854in}}%
\pgfpathlineto{\pgfqpoint{5.140793in}{2.380839in}}%
\pgfpathlineto{\pgfqpoint{5.142147in}{2.379966in}}%
\pgfpathlineto{\pgfqpoint{5.144618in}{2.378944in}}%
\pgfpathlineto{\pgfqpoint{5.145886in}{2.378359in}}%
\pgfpathlineto{\pgfqpoint{5.148298in}{2.377331in}}%
\pgfpathlineto{\pgfqpoint{5.149628in}{2.376594in}}%
\pgfpathlineto{\pgfqpoint{5.152129in}{2.375560in}}%
\pgfpathlineto{\pgfqpoint{5.153447in}{2.374819in}}%
\pgfpathlineto{\pgfqpoint{5.155779in}{2.373778in}}%
\pgfpathlineto{\pgfqpoint{5.157085in}{2.373032in}}%
\pgfpathlineto{\pgfqpoint{5.159433in}{2.371985in}}%
\pgfpathlineto{\pgfqpoint{5.160728in}{2.371234in}}%
\pgfpathlineto{\pgfqpoint{5.162949in}{2.370180in}}%
\pgfpathlineto{\pgfqpoint{5.164233in}{2.369576in}}%
\pgfpathlineto{\pgfqpoint{5.166895in}{2.368515in}}%
\pgfpathlineto{\pgfqpoint{5.168131in}{2.368059in}}%
\pgfpathlineto{\pgfqpoint{5.170874in}{2.366993in}}%
\pgfpathlineto{\pgfqpoint{5.172099in}{2.366534in}}%
\pgfpathlineto{\pgfqpoint{5.174364in}{2.365462in}}%
\pgfpathlineto{\pgfqpoint{5.175613in}{2.364847in}}%
\pgfpathlineto{\pgfqpoint{5.177824in}{2.363769in}}%
\pgfpathlineto{\pgfqpoint{5.179097in}{2.363150in}}%
\pgfpathlineto{\pgfqpoint{5.182075in}{2.362065in}}%
\pgfpathlineto{\pgfqpoint{5.183336in}{2.361443in}}%
\pgfpathlineto{\pgfqpoint{5.186013in}{2.360350in}}%
\pgfpathlineto{\pgfqpoint{5.187228in}{2.359881in}}%
\pgfpathlineto{\pgfqpoint{5.190015in}{2.358783in}}%
\pgfpathlineto{\pgfqpoint{5.191251in}{2.358311in}}%
\pgfpathlineto{\pgfqpoint{5.193281in}{2.357206in}}%
\pgfpathlineto{\pgfqpoint{5.194574in}{2.356415in}}%
\pgfpathlineto{\pgfqpoint{5.196423in}{2.355462in}}%
\pgfpathlineto{\pgfqpoint{5.197673in}{2.354825in}}%
\pgfpathlineto{\pgfqpoint{5.199934in}{2.353866in}}%
\pgfpathlineto{\pgfqpoint{5.201174in}{2.353226in}}%
\pgfpathlineto{\pgfqpoint{5.203771in}{2.352262in}}%
\pgfpathlineto{\pgfqpoint{5.205032in}{2.351617in}}%
\pgfpathlineto{\pgfqpoint{5.207606in}{2.350648in}}%
\pgfpathlineto{\pgfqpoint{5.208855in}{2.349999in}}%
\pgfpathlineto{\pgfqpoint{5.211627in}{2.349024in}}%
\pgfpathlineto{\pgfqpoint{5.212864in}{2.348372in}}%
\pgfpathlineto{\pgfqpoint{5.215453in}{2.347391in}}%
\pgfpathlineto{\pgfqpoint{5.216679in}{2.346900in}}%
\pgfpathlineto{\pgfqpoint{5.218932in}{2.345914in}}%
\pgfpathlineto{\pgfqpoint{5.220147in}{2.345419in}}%
\pgfpathlineto{\pgfqpoint{5.222258in}{2.344428in}}%
\pgfpathlineto{\pgfqpoint{5.223463in}{2.343931in}}%
\pgfpathlineto{\pgfqpoint{5.225250in}{2.342934in}}%
\pgfpathlineto{\pgfqpoint{5.226478in}{2.342268in}}%
\pgfpathlineto{\pgfqpoint{5.228800in}{2.341265in}}%
\pgfpathlineto{\pgfqpoint{5.230048in}{2.340595in}}%
\pgfpathlineto{\pgfqpoint{5.232532in}{2.339586in}}%
\pgfpathlineto{\pgfqpoint{5.233768in}{2.338912in}}%
\pgfpathlineto{\pgfqpoint{5.235931in}{2.337897in}}%
\pgfpathlineto{\pgfqpoint{5.237157in}{2.337218in}}%
\pgfpathlineto{\pgfqpoint{5.239569in}{2.336197in}}%
\pgfpathlineto{\pgfqpoint{5.240814in}{2.335515in}}%
\pgfpathlineto{\pgfqpoint{5.243058in}{2.334487in}}%
\pgfpathlineto{\pgfqpoint{5.244264in}{2.333972in}}%
\pgfpathlineto{\pgfqpoint{5.246547in}{2.332939in}}%
\pgfpathlineto{\pgfqpoint{5.247743in}{2.332422in}}%
\pgfpathlineto{\pgfqpoint{5.249979in}{2.331383in}}%
\pgfpathlineto{\pgfqpoint{5.251193in}{2.330688in}}%
\pgfpathlineto{\pgfqpoint{5.253209in}{2.329643in}}%
\pgfpathlineto{\pgfqpoint{5.254443in}{2.328944in}}%
\pgfpathlineto{\pgfqpoint{5.257127in}{2.327892in}}%
\pgfpathlineto{\pgfqpoint{5.258321in}{2.327365in}}%
\pgfpathlineto{\pgfqpoint{5.260811in}{2.326307in}}%
\pgfpathlineto{\pgfqpoint{5.262022in}{2.325600in}}%
\pgfpathlineto{\pgfqpoint{5.264350in}{2.324535in}}%
\pgfpathlineto{\pgfqpoint{5.265607in}{2.323645in}}%
\pgfpathlineto{\pgfqpoint{5.267998in}{2.322573in}}%
\pgfpathlineto{\pgfqpoint{5.269216in}{2.321856in}}%
\pgfpathlineto{\pgfqpoint{5.272082in}{2.320777in}}%
\pgfpathlineto{\pgfqpoint{5.273288in}{2.320236in}}%
\pgfpathlineto{\pgfqpoint{5.275772in}{2.319151in}}%
\pgfpathlineto{\pgfqpoint{5.276994in}{2.318425in}}%
\pgfpathlineto{\pgfqpoint{5.279645in}{2.317332in}}%
\pgfpathlineto{\pgfqpoint{5.280883in}{2.316419in}}%
\pgfpathlineto{\pgfqpoint{5.283482in}{2.315319in}}%
\pgfpathlineto{\pgfqpoint{5.284709in}{2.314583in}}%
\pgfpathlineto{\pgfqpoint{5.287495in}{2.313475in}}%
\pgfpathlineto{\pgfqpoint{5.288684in}{2.312920in}}%
\pgfpathlineto{\pgfqpoint{5.291156in}{2.311992in}}%
\pgfpathlineto{\pgfqpoint{5.292361in}{2.311247in}}%
\pgfpathlineto{\pgfqpoint{5.293900in}{2.310313in}}%
\pgfpathlineto{\pgfqpoint{5.295097in}{2.309751in}}%
\pgfpathlineto{\pgfqpoint{5.297014in}{2.308813in}}%
\pgfpathlineto{\pgfqpoint{5.298253in}{2.307871in}}%
\pgfpathlineto{\pgfqpoint{5.300695in}{2.306925in}}%
\pgfpathlineto{\pgfqpoint{5.301897in}{2.306167in}}%
\pgfpathlineto{\pgfqpoint{5.304521in}{2.305216in}}%
\pgfpathlineto{\pgfqpoint{5.305712in}{2.304644in}}%
\pgfpathlineto{\pgfqpoint{5.307682in}{2.303688in}}%
\pgfpathlineto{\pgfqpoint{5.308865in}{2.303112in}}%
\pgfpathlineto{\pgfqpoint{5.310995in}{2.302151in}}%
\pgfpathlineto{\pgfqpoint{5.312193in}{2.301379in}}%
\pgfpathlineto{\pgfqpoint{5.314504in}{2.300412in}}%
\pgfpathlineto{\pgfqpoint{5.315692in}{2.299830in}}%
\pgfpathlineto{\pgfqpoint{5.317591in}{2.298857in}}%
\pgfpathlineto{\pgfqpoint{5.318794in}{2.298076in}}%
\pgfpathlineto{\pgfqpoint{5.320898in}{2.297097in}}%
\pgfpathlineto{\pgfqpoint{5.322068in}{2.296508in}}%
\pgfpathlineto{\pgfqpoint{5.324229in}{2.295524in}}%
\pgfpathlineto{\pgfqpoint{5.325413in}{2.294931in}}%
\pgfpathlineto{\pgfqpoint{5.327772in}{2.293941in}}%
\pgfpathlineto{\pgfqpoint{5.328947in}{2.293345in}}%
\pgfpathlineto{\pgfqpoint{5.331048in}{2.292349in}}%
\pgfpathlineto{\pgfqpoint{5.332237in}{2.291750in}}%
\pgfpathlineto{\pgfqpoint{5.334038in}{2.290749in}}%
\pgfpathlineto{\pgfqpoint{5.335218in}{2.290146in}}%
\pgfpathlineto{\pgfqpoint{5.337499in}{2.289139in}}%
\pgfpathlineto{\pgfqpoint{5.338670in}{2.288533in}}%
\pgfpathlineto{\pgfqpoint{5.340559in}{2.287519in}}%
\pgfpathlineto{\pgfqpoint{5.341745in}{2.286910in}}%
\pgfpathlineto{\pgfqpoint{5.344059in}{2.285891in}}%
\pgfpathlineto{\pgfqpoint{5.345234in}{2.285277in}}%
\pgfpathlineto{\pgfqpoint{5.347324in}{2.284252in}}%
\pgfpathlineto{\pgfqpoint{5.348513in}{2.283429in}}%
\pgfpathlineto{\pgfqpoint{5.350767in}{2.282397in}}%
\pgfpathlineto{\pgfqpoint{5.351969in}{2.281569in}}%
\pgfpathlineto{\pgfqpoint{5.354047in}{2.280530in}}%
\pgfpathlineto{\pgfqpoint{5.355217in}{2.279905in}}%
\pgfpathlineto{\pgfqpoint{5.357748in}{2.278860in}}%
\pgfpathlineto{\pgfqpoint{5.358952in}{2.278021in}}%
\pgfpathlineto{\pgfqpoint{5.360840in}{2.276969in}}%
\pgfpathlineto{\pgfqpoint{5.362013in}{2.276336in}}%
\pgfpathlineto{\pgfqpoint{5.364547in}{2.275277in}}%
\pgfpathlineto{\pgfqpoint{5.365710in}{2.274640in}}%
\pgfpathlineto{\pgfqpoint{5.367851in}{2.273575in}}%
\pgfpathlineto{\pgfqpoint{5.369026in}{2.272934in}}%
\pgfpathlineto{\pgfqpoint{5.371409in}{2.271862in}}%
\pgfpathlineto{\pgfqpoint{5.372596in}{2.271002in}}%
\pgfpathlineto{\pgfqpoint{5.374895in}{2.269923in}}%
\pgfpathlineto{\pgfqpoint{5.376072in}{2.269273in}}%
\pgfpathlineto{\pgfqpoint{5.378543in}{2.268188in}}%
\pgfpathlineto{\pgfqpoint{5.379730in}{2.267316in}}%
\pgfpathlineto{\pgfqpoint{5.382159in}{2.266222in}}%
\pgfpathlineto{\pgfqpoint{5.383336in}{2.265564in}}%
\pgfpathlineto{\pgfqpoint{5.386140in}{2.264464in}}%
\pgfpathlineto{\pgfqpoint{5.387306in}{2.263802in}}%
\pgfpathlineto{\pgfqpoint{5.389318in}{2.262694in}}%
\pgfpathlineto{\pgfqpoint{5.390496in}{2.262028in}}%
\pgfpathlineto{\pgfqpoint{5.392061in}{2.261136in}}%
\pgfpathlineto{\pgfqpoint{5.393231in}{2.260466in}}%
\pgfpathlineto{\pgfqpoint{5.394970in}{2.259570in}}%
\pgfpathlineto{\pgfqpoint{5.396152in}{2.258671in}}%
\pgfpathlineto{\pgfqpoint{5.398345in}{2.257769in}}%
\pgfpathlineto{\pgfqpoint{5.399537in}{2.256864in}}%
\pgfpathlineto{\pgfqpoint{5.401853in}{2.255956in}}%
\pgfpathlineto{\pgfqpoint{5.403036in}{2.255045in}}%
\pgfpathlineto{\pgfqpoint{5.405172in}{2.254131in}}%
\pgfpathlineto{\pgfqpoint{5.406346in}{2.253444in}}%
\pgfpathlineto{\pgfqpoint{5.408386in}{2.252525in}}%
\pgfpathlineto{\pgfqpoint{5.409570in}{2.251603in}}%
\pgfpathlineto{\pgfqpoint{5.411968in}{2.250677in}}%
\pgfpathlineto{\pgfqpoint{5.413161in}{2.249749in}}%
\pgfpathlineto{\pgfqpoint{5.415693in}{2.248818in}}%
\pgfpathlineto{\pgfqpoint{5.416857in}{2.248117in}}%
\pgfpathlineto{\pgfqpoint{5.418461in}{2.247180in}}%
\pgfpathlineto{\pgfqpoint{5.419616in}{2.246475in}}%
\pgfpathlineto{\pgfqpoint{5.421421in}{2.245533in}}%
\pgfpathlineto{\pgfqpoint{5.422587in}{2.244824in}}%
\pgfpathlineto{\pgfqpoint{5.424379in}{2.243876in}}%
\pgfpathlineto{\pgfqpoint{5.425537in}{2.243163in}}%
\pgfpathlineto{\pgfqpoint{5.427485in}{2.242209in}}%
\pgfpathlineto{\pgfqpoint{5.428654in}{2.241492in}}%
\pgfpathlineto{\pgfqpoint{5.430363in}{2.240532in}}%
\pgfpathlineto{\pgfqpoint{5.431524in}{2.239810in}}%
\pgfpathlineto{\pgfqpoint{5.433760in}{2.238845in}}%
\pgfpathlineto{\pgfqpoint{5.434930in}{2.238119in}}%
\pgfpathlineto{\pgfqpoint{5.436724in}{2.237148in}}%
\pgfpathlineto{\pgfqpoint{5.437885in}{2.236417in}}%
\pgfpathlineto{\pgfqpoint{5.439960in}{2.235440in}}%
\pgfpathlineto{\pgfqpoint{5.441112in}{2.234705in}}%
\pgfpathlineto{\pgfqpoint{5.443280in}{2.233722in}}%
\pgfpathlineto{\pgfqpoint{5.444441in}{2.232982in}}%
\pgfpathlineto{\pgfqpoint{5.446231in}{2.231993in}}%
\pgfpathlineto{\pgfqpoint{5.447384in}{2.231249in}}%
\pgfpathlineto{\pgfqpoint{5.449429in}{2.230253in}}%
\pgfpathlineto{\pgfqpoint{5.450609in}{2.229254in}}%
\pgfpathlineto{\pgfqpoint{5.452905in}{2.228251in}}%
\pgfpathlineto{\pgfqpoint{5.454075in}{2.227245in}}%
\pgfpathlineto{\pgfqpoint{5.456493in}{2.226235in}}%
\pgfpathlineto{\pgfqpoint{5.457653in}{2.225475in}}%
\pgfpathlineto{\pgfqpoint{5.459789in}{2.224459in}}%
\pgfpathlineto{\pgfqpoint{5.460940in}{2.223694in}}%
\pgfpathlineto{\pgfqpoint{5.462746in}{2.222672in}}%
\pgfpathlineto{\pgfqpoint{5.463906in}{2.221902in}}%
\pgfpathlineto{\pgfqpoint{5.465889in}{2.220873in}}%
\pgfpathlineto{\pgfqpoint{5.467040in}{2.220098in}}%
\pgfpathlineto{\pgfqpoint{5.469077in}{2.219062in}}%
\pgfpathlineto{\pgfqpoint{5.470236in}{2.218282in}}%
\pgfpathlineto{\pgfqpoint{5.472562in}{2.217239in}}%
\pgfpathlineto{\pgfqpoint{5.473712in}{2.216454in}}%
\pgfpathlineto{\pgfqpoint{5.475750in}{2.215405in}}%
\pgfpathlineto{\pgfqpoint{5.476890in}{2.214878in}}%
\pgfpathlineto{\pgfqpoint{5.478662in}{2.213822in}}%
\pgfpathlineto{\pgfqpoint{5.479812in}{2.213028in}}%
\pgfpathlineto{\pgfqpoint{5.481621in}{2.211965in}}%
\pgfpathlineto{\pgfqpoint{5.482779in}{2.211165in}}%
\pgfpathlineto{\pgfqpoint{5.485019in}{2.210095in}}%
\pgfpathlineto{\pgfqpoint{5.486168in}{2.209290in}}%
\pgfpathlineto{\pgfqpoint{5.488504in}{2.208213in}}%
\pgfpathlineto{\pgfqpoint{5.489659in}{2.207402in}}%
\pgfpathlineto{\pgfqpoint{5.491620in}{2.206318in}}%
\pgfpathlineto{\pgfqpoint{5.492767in}{2.205502in}}%
\pgfpathlineto{\pgfqpoint{5.494809in}{2.204410in}}%
\pgfpathlineto{\pgfqpoint{5.495963in}{2.203588in}}%
\pgfpathlineto{\pgfqpoint{5.497879in}{2.202488in}}%
\pgfpathlineto{\pgfqpoint{5.499024in}{2.201661in}}%
\pgfpathlineto{\pgfqpoint{5.501210in}{2.200554in}}%
\pgfpathlineto{\pgfqpoint{5.502362in}{2.199720in}}%
\pgfpathlineto{\pgfqpoint{5.504185in}{2.198885in}}%
\pgfpathlineto{\pgfqpoint{5.505344in}{2.198046in}}%
\pgfpathlineto{\pgfqpoint{5.506906in}{2.197206in}}%
\pgfpathlineto{\pgfqpoint{5.508057in}{2.196362in}}%
\pgfpathlineto{\pgfqpoint{5.509624in}{2.195516in}}%
\pgfpathlineto{\pgfqpoint{5.510768in}{2.194668in}}%
\pgfpathlineto{\pgfqpoint{5.512494in}{2.193817in}}%
\pgfpathlineto{\pgfqpoint{5.513645in}{2.192963in}}%
\pgfpathlineto{\pgfqpoint{5.515512in}{2.192107in}}%
\pgfpathlineto{\pgfqpoint{5.516670in}{2.191248in}}%
\pgfpathlineto{\pgfqpoint{5.518341in}{2.190386in}}%
\pgfpathlineto{\pgfqpoint{5.519492in}{2.189522in}}%
\pgfpathlineto{\pgfqpoint{5.521362in}{2.188655in}}%
\pgfpathlineto{\pgfqpoint{5.522519in}{2.187785in}}%
\pgfpathlineto{\pgfqpoint{5.524525in}{2.186913in}}%
\pgfpathlineto{\pgfqpoint{5.525674in}{2.186038in}}%
\pgfpathlineto{\pgfqpoint{5.527383in}{2.185160in}}%
\pgfpathlineto{\pgfqpoint{5.528553in}{2.183985in}}%
\pgfpathlineto{\pgfqpoint{5.530310in}{2.183101in}}%
\pgfpathlineto{\pgfqpoint{5.531457in}{2.182214in}}%
\pgfpathlineto{\pgfqpoint{5.533581in}{2.181323in}}%
\pgfpathlineto{\pgfqpoint{5.534734in}{2.180431in}}%
\pgfpathlineto{\pgfqpoint{5.536421in}{2.179535in}}%
\pgfpathlineto{\pgfqpoint{5.537566in}{2.178636in}}%
\pgfpathlineto{\pgfqpoint{5.539789in}{2.177735in}}%
\pgfpathlineto{\pgfqpoint{5.540939in}{2.176830in}}%
\pgfpathlineto{\pgfqpoint{5.542815in}{2.175923in}}%
\pgfpathlineto{\pgfqpoint{5.543943in}{2.175316in}}%
\pgfpathlineto{\pgfqpoint{5.545280in}{2.174404in}}%
\pgfpathlineto{\pgfqpoint{5.546430in}{2.173488in}}%
\pgfpathlineto{\pgfqpoint{5.548451in}{2.172570in}}%
\pgfpathlineto{\pgfqpoint{5.549605in}{2.171649in}}%
\pgfpathlineto{\pgfqpoint{5.551303in}{2.170724in}}%
\pgfpathlineto{\pgfqpoint{5.552435in}{2.170106in}}%
\pgfpathlineto{\pgfqpoint{5.554191in}{2.169177in}}%
\pgfpathlineto{\pgfqpoint{5.555343in}{2.168244in}}%
\pgfpathlineto{\pgfqpoint{5.556990in}{2.167308in}}%
\pgfpathlineto{\pgfqpoint{5.558135in}{2.166369in}}%
\pgfpathlineto{\pgfqpoint{5.559990in}{2.165427in}}%
\pgfpathlineto{\pgfqpoint{5.561140in}{2.164481in}}%
\pgfpathlineto{\pgfqpoint{5.562914in}{2.163533in}}%
\pgfpathlineto{\pgfqpoint{5.564056in}{2.162581in}}%
\pgfpathlineto{\pgfqpoint{5.565777in}{2.161626in}}%
\pgfpathlineto{\pgfqpoint{5.566924in}{2.160668in}}%
\pgfpathlineto{\pgfqpoint{5.568888in}{2.159706in}}%
\pgfpathlineto{\pgfqpoint{5.570028in}{2.158741in}}%
\pgfpathlineto{\pgfqpoint{5.571711in}{2.157773in}}%
\pgfpathlineto{\pgfqpoint{5.572842in}{2.157125in}}%
\pgfpathlineto{\pgfqpoint{5.574593in}{2.156151in}}%
\pgfpathlineto{\pgfqpoint{5.575743in}{2.155174in}}%
\pgfpathlineto{\pgfqpoint{5.577298in}{2.154193in}}%
\pgfpathlineto{\pgfqpoint{5.578427in}{2.153537in}}%
\pgfpathlineto{\pgfqpoint{5.579933in}{2.152550in}}%
\pgfpathlineto{\pgfqpoint{5.581081in}{2.151560in}}%
\pgfpathlineto{\pgfqpoint{5.583018in}{2.150567in}}%
\pgfpathlineto{\pgfqpoint{5.584158in}{2.149569in}}%
\pgfpathlineto{\pgfqpoint{5.586132in}{2.148568in}}%
\pgfpathlineto{\pgfqpoint{5.587276in}{2.147564in}}%
\pgfpathlineto{\pgfqpoint{5.589121in}{2.146556in}}%
\pgfpathlineto{\pgfqpoint{5.590270in}{2.145544in}}%
\pgfpathlineto{\pgfqpoint{5.592012in}{2.144529in}}%
\pgfpathlineto{\pgfqpoint{5.593153in}{2.143510in}}%
\pgfpathlineto{\pgfqpoint{5.595148in}{2.142487in}}%
\pgfpathlineto{\pgfqpoint{5.596280in}{2.141803in}}%
\pgfpathlineto{\pgfqpoint{5.597747in}{2.140774in}}%
\pgfpathlineto{\pgfqpoint{5.598885in}{2.139741in}}%
\pgfpathlineto{\pgfqpoint{5.600716in}{2.138704in}}%
\pgfpathlineto{\pgfqpoint{5.601858in}{2.137663in}}%
\pgfpathlineto{\pgfqpoint{5.603590in}{2.136619in}}%
\pgfpathlineto{\pgfqpoint{5.604736in}{2.135570in}}%
\pgfpathlineto{\pgfqpoint{5.606909in}{2.134518in}}%
\pgfpathlineto{\pgfqpoint{5.608046in}{2.133461in}}%
\pgfpathlineto{\pgfqpoint{5.609946in}{2.132401in}}%
\pgfpathlineto{\pgfqpoint{5.611075in}{2.131692in}}%
\pgfpathlineto{\pgfqpoint{5.612636in}{2.130624in}}%
\pgfpathlineto{\pgfqpoint{5.613782in}{2.129553in}}%
\pgfpathlineto{\pgfqpoint{5.615944in}{2.128478in}}%
\pgfpathlineto{\pgfqpoint{5.617081in}{2.127398in}}%
\pgfpathlineto{\pgfqpoint{5.618952in}{2.126314in}}%
\pgfpathlineto{\pgfqpoint{5.620092in}{2.125226in}}%
\pgfpathlineto{\pgfqpoint{5.621939in}{2.124134in}}%
\pgfpathlineto{\pgfqpoint{5.623083in}{2.123037in}}%
\pgfpathlineto{\pgfqpoint{5.625138in}{2.121936in}}%
\pgfpathlineto{\pgfqpoint{5.626285in}{2.120831in}}%
\pgfpathlineto{\pgfqpoint{5.628209in}{2.119721in}}%
\pgfpathlineto{\pgfqpoint{5.629347in}{2.118607in}}%
\pgfpathlineto{\pgfqpoint{5.630817in}{2.117862in}}%
\pgfpathlineto{\pgfqpoint{5.631960in}{2.116740in}}%
\pgfpathlineto{\pgfqpoint{5.633353in}{2.115990in}}%
\pgfpathlineto{\pgfqpoint{5.634488in}{2.114861in}}%
\pgfpathlineto{\pgfqpoint{5.636055in}{2.114105in}}%
\pgfpathlineto{\pgfqpoint{5.637195in}{2.112968in}}%
\pgfpathlineto{\pgfqpoint{5.638468in}{2.112208in}}%
\pgfpathlineto{\pgfqpoint{5.639612in}{2.111063in}}%
\pgfpathlineto{\pgfqpoint{5.641148in}{2.110297in}}%
\pgfpathlineto{\pgfqpoint{5.642274in}{2.109530in}}%
\pgfpathlineto{\pgfqpoint{5.643487in}{2.108760in}}%
\pgfpathlineto{\pgfqpoint{5.644628in}{2.107601in}}%
\pgfpathlineto{\pgfqpoint{5.646057in}{2.106825in}}%
\pgfpathlineto{\pgfqpoint{5.647192in}{2.105658in}}%
\pgfpathlineto{\pgfqpoint{5.648744in}{2.104878in}}%
\pgfpathlineto{\pgfqpoint{5.649883in}{2.103702in}}%
\pgfpathlineto{\pgfqpoint{5.651415in}{2.102916in}}%
\pgfpathlineto{\pgfqpoint{5.652557in}{2.101732in}}%
\pgfpathlineto{\pgfqpoint{5.654221in}{2.100941in}}%
\pgfpathlineto{\pgfqpoint{5.655344in}{2.100147in}}%
\pgfpathlineto{\pgfqpoint{5.656487in}{2.099350in}}%
\pgfpathlineto{\pgfqpoint{5.657615in}{2.098552in}}%
\pgfpathlineto{\pgfqpoint{5.658643in}{2.097751in}}%
\pgfpathlineto{\pgfqpoint{5.659777in}{2.096545in}}%
\pgfpathlineto{\pgfqpoint{5.661229in}{2.095738in}}%
\pgfpathlineto{\pgfqpoint{5.662355in}{2.094929in}}%
\pgfpathlineto{\pgfqpoint{5.663265in}{2.094118in}}%
\pgfpathlineto{\pgfqpoint{5.664407in}{2.092896in}}%
\pgfpathlineto{\pgfqpoint{5.666003in}{2.092079in}}%
\pgfpathlineto{\pgfqpoint{5.667137in}{2.090848in}}%
\pgfpathlineto{\pgfqpoint{5.668733in}{2.090025in}}%
\pgfpathlineto{\pgfqpoint{5.669860in}{2.089199in}}%
\pgfpathlineto{\pgfqpoint{5.671121in}{2.088370in}}%
\pgfpathlineto{\pgfqpoint{5.672262in}{2.087123in}}%
\pgfpathlineto{\pgfqpoint{5.673808in}{2.086289in}}%
\pgfpathlineto{\pgfqpoint{5.674942in}{2.085032in}}%
\pgfpathlineto{\pgfqpoint{5.676394in}{2.084192in}}%
\pgfpathlineto{\pgfqpoint{5.677521in}{2.083348in}}%
\pgfpathlineto{\pgfqpoint{5.678656in}{2.082502in}}%
\pgfpathlineto{\pgfqpoint{5.679787in}{2.081654in}}%
\pgfpathlineto{\pgfqpoint{5.681110in}{2.080803in}}%
\pgfpathlineto{\pgfqpoint{5.682235in}{2.079949in}}%
\pgfpathlineto{\pgfqpoint{5.683478in}{2.079093in}}%
\pgfpathlineto{\pgfqpoint{5.684607in}{2.078234in}}%
\pgfpathlineto{\pgfqpoint{5.685712in}{2.077372in}}%
\pgfpathlineto{\pgfqpoint{5.686845in}{2.076075in}}%
\pgfpathlineto{\pgfqpoint{5.688899in}{2.075207in}}%
\pgfpathlineto{\pgfqpoint{5.690023in}{2.074335in}}%
\pgfpathlineto{\pgfqpoint{5.691434in}{2.073462in}}%
\pgfpathlineto{\pgfqpoint{5.692571in}{2.072146in}}%
\pgfpathlineto{\pgfqpoint{5.694121in}{2.071265in}}%
\pgfpathlineto{\pgfqpoint{5.695251in}{2.070382in}}%
\pgfpathlineto{\pgfqpoint{5.696604in}{2.069496in}}%
\pgfpathlineto{\pgfqpoint{5.697727in}{2.068607in}}%
\pgfpathlineto{\pgfqpoint{5.698955in}{2.067715in}}%
\pgfpathlineto{\pgfqpoint{5.700091in}{2.066371in}}%
\pgfpathlineto{\pgfqpoint{5.701770in}{2.065472in}}%
\pgfpathlineto{\pgfqpoint{5.702908in}{2.064118in}}%
\pgfpathlineto{\pgfqpoint{5.704353in}{2.063212in}}%
\pgfpathlineto{\pgfqpoint{5.705474in}{2.062302in}}%
\pgfpathlineto{\pgfqpoint{5.706506in}{2.061390in}}%
\pgfpathlineto{\pgfqpoint{5.707631in}{2.060475in}}%
\pgfpathlineto{\pgfqpoint{5.709088in}{2.059556in}}%
\pgfpathlineto{\pgfqpoint{5.710216in}{2.058635in}}%
\pgfpathlineto{\pgfqpoint{5.711407in}{2.057710in}}%
\pgfpathlineto{\pgfqpoint{5.712538in}{2.056318in}}%
\pgfpathlineto{\pgfqpoint{5.714215in}{2.055386in}}%
\pgfpathlineto{\pgfqpoint{5.715338in}{2.054450in}}%
\pgfpathlineto{\pgfqpoint{5.716449in}{2.053512in}}%
\pgfpathlineto{\pgfqpoint{5.717585in}{2.052098in}}%
\pgfpathlineto{\pgfqpoint{5.719204in}{2.051152in}}%
\pgfpathlineto{\pgfqpoint{5.720333in}{2.050202in}}%
\pgfpathlineto{\pgfqpoint{5.721449in}{2.049249in}}%
\pgfpathlineto{\pgfqpoint{5.722580in}{2.047814in}}%
\pgfpathlineto{\pgfqpoint{5.724365in}{2.046852in}}%
\pgfpathlineto{\pgfqpoint{5.725488in}{2.045888in}}%
\pgfpathlineto{\pgfqpoint{5.726793in}{2.044920in}}%
\pgfpathlineto{\pgfqpoint{5.727919in}{2.043949in}}%
\pgfpathlineto{\pgfqpoint{5.729115in}{2.042974in}}%
\pgfpathlineto{\pgfqpoint{5.730245in}{2.041996in}}%
\pgfpathlineto{\pgfqpoint{5.731344in}{2.041015in}}%
\pgfpathlineto{\pgfqpoint{5.732467in}{2.040030in}}%
\pgfpathlineto{\pgfqpoint{5.733650in}{2.039042in}}%
\pgfpathlineto{\pgfqpoint{5.734776in}{2.038050in}}%
\pgfpathlineto{\pgfqpoint{5.735890in}{2.037054in}}%
\pgfpathlineto{\pgfqpoint{5.737028in}{2.035555in}}%
\pgfpathlineto{\pgfqpoint{5.739063in}{2.034550in}}%
\pgfpathlineto{\pgfqpoint{5.740183in}{2.033542in}}%
\pgfpathlineto{\pgfqpoint{5.741610in}{2.032530in}}%
\pgfpathlineto{\pgfqpoint{5.742732in}{2.031515in}}%
\pgfpathlineto{\pgfqpoint{5.744106in}{2.030496in}}%
\pgfpathlineto{\pgfqpoint{5.745239in}{2.028960in}}%
\pgfpathlineto{\pgfqpoint{5.746868in}{2.027932in}}%
\pgfpathlineto{\pgfqpoint{5.747993in}{2.026899in}}%
\pgfpathlineto{\pgfqpoint{5.749237in}{2.025863in}}%
\pgfpathlineto{\pgfqpoint{5.750365in}{2.024823in}}%
\pgfpathlineto{\pgfqpoint{5.751758in}{2.023779in}}%
\pgfpathlineto{\pgfqpoint{5.752887in}{2.022206in}}%
\pgfpathlineto{\pgfqpoint{5.754949in}{2.021152in}}%
\pgfpathlineto{\pgfqpoint{5.756070in}{2.020094in}}%
\pgfpathlineto{\pgfqpoint{5.757324in}{2.019033in}}%
\pgfpathlineto{\pgfqpoint{5.758447in}{2.017967in}}%
\pgfpathlineto{\pgfqpoint{5.759643in}{2.016897in}}%
\pgfpathlineto{\pgfqpoint{5.760768in}{2.015823in}}%
\pgfpathlineto{\pgfqpoint{5.762270in}{2.014744in}}%
\pgfpathlineto{\pgfqpoint{5.763396in}{2.013662in}}%
\pgfpathlineto{\pgfqpoint{5.764779in}{2.012575in}}%
\pgfpathlineto{\pgfqpoint{5.765899in}{2.011484in}}%
\pgfpathlineto{\pgfqpoint{5.767273in}{2.010389in}}%
\pgfpathlineto{\pgfqpoint{5.768395in}{2.009290in}}%
\pgfpathlineto{\pgfqpoint{5.769885in}{2.008186in}}%
\pgfpathlineto{\pgfqpoint{5.771008in}{2.007078in}}%
\pgfpathlineto{\pgfqpoint{5.771617in}{2.006522in}}%
\pgfpathlineto{\pgfqpoint{5.772743in}{2.005407in}}%
\pgfpathlineto{\pgfqpoint{5.773825in}{2.004848in}}%
\pgfpathlineto{\pgfqpoint{5.774946in}{2.003726in}}%
\pgfpathlineto{\pgfqpoint{5.775900in}{2.003164in}}%
\pgfpathlineto{\pgfqpoint{5.777022in}{2.002035in}}%
\pgfpathlineto{\pgfqpoint{5.778093in}{2.001469in}}%
\pgfpathlineto{\pgfqpoint{5.779218in}{2.000334in}}%
\pgfpathlineto{\pgfqpoint{5.779969in}{1.999764in}}%
\pgfpathlineto{\pgfqpoint{5.781089in}{1.998622in}}%
\pgfpathlineto{\pgfqpoint{5.782069in}{1.998049in}}%
\pgfpathlineto{\pgfqpoint{5.783191in}{1.996900in}}%
\pgfpathlineto{\pgfqpoint{5.784167in}{1.996323in}}%
\pgfpathlineto{\pgfqpoint{5.785291in}{1.995167in}}%
\pgfpathlineto{\pgfqpoint{5.786309in}{1.994587in}}%
\pgfpathlineto{\pgfqpoint{5.787436in}{1.993423in}}%
\pgfpathlineto{\pgfqpoint{5.788472in}{1.992839in}}%
\pgfpathlineto{\pgfqpoint{5.789593in}{1.991668in}}%
\pgfpathlineto{\pgfqpoint{5.790561in}{1.991081in}}%
\pgfpathlineto{\pgfqpoint{5.791684in}{1.989902in}}%
\pgfpathlineto{\pgfqpoint{5.792554in}{1.989311in}}%
\pgfpathlineto{\pgfqpoint{5.793679in}{1.988125in}}%
\pgfpathlineto{\pgfqpoint{5.794575in}{1.987530in}}%
\pgfpathlineto{\pgfqpoint{5.795695in}{1.986336in}}%
\pgfpathlineto{\pgfqpoint{5.796820in}{1.985737in}}%
\pgfpathlineto{\pgfqpoint{5.797942in}{1.984536in}}%
\pgfpathlineto{\pgfqpoint{5.799237in}{1.983933in}}%
\pgfpathlineto{\pgfqpoint{5.800360in}{1.982724in}}%
\pgfpathlineto{\pgfqpoint{5.801480in}{1.982118in}}%
\pgfpathlineto{\pgfqpoint{5.802604in}{1.980900in}}%
\pgfpathlineto{\pgfqpoint{5.803520in}{1.980290in}}%
\pgfpathlineto{\pgfqpoint{5.804646in}{1.979065in}}%
\pgfpathlineto{\pgfqpoint{5.805837in}{1.978450in}}%
\pgfpathlineto{\pgfqpoint{5.806958in}{1.977217in}}%
\pgfpathlineto{\pgfqpoint{5.807969in}{1.976598in}}%
\pgfpathlineto{\pgfqpoint{5.809091in}{1.975356in}}%
\pgfpathlineto{\pgfqpoint{5.810397in}{1.974734in}}%
\pgfpathlineto{\pgfqpoint{5.811520in}{1.973484in}}%
\pgfpathlineto{\pgfqpoint{5.812648in}{1.972857in}}%
\pgfpathlineto{\pgfqpoint{5.813772in}{1.971598in}}%
\pgfpathlineto{\pgfqpoint{5.814878in}{1.970967in}}%
\pgfpathlineto{\pgfqpoint{5.816004in}{1.969700in}}%
\pgfpathlineto{\pgfqpoint{5.817134in}{1.969064in}}%
\pgfpathlineto{\pgfqpoint{5.818253in}{1.967789in}}%
\pgfpathlineto{\pgfqpoint{5.819568in}{1.967149in}}%
\pgfpathlineto{\pgfqpoint{5.820688in}{1.965864in}}%
\pgfpathlineto{\pgfqpoint{5.822126in}{1.965220in}}%
\pgfpathlineto{\pgfqpoint{5.823247in}{1.963926in}}%
\pgfpathlineto{\pgfqpoint{5.824212in}{1.963278in}}%
\pgfpathlineto{\pgfqpoint{5.825334in}{1.961975in}}%
\pgfpathlineto{\pgfqpoint{5.826518in}{1.961322in}}%
\pgfpathlineto{\pgfqpoint{5.827642in}{1.960010in}}%
\pgfpathlineto{\pgfqpoint{5.829120in}{1.959352in}}%
\pgfpathlineto{\pgfqpoint{5.830244in}{1.958031in}}%
\pgfpathlineto{\pgfqpoint{5.831706in}{1.957368in}}%
\pgfpathlineto{\pgfqpoint{5.832829in}{1.956037in}}%
\pgfpathlineto{\pgfqpoint{5.834148in}{1.955370in}}%
\pgfpathlineto{\pgfqpoint{5.835272in}{1.954030in}}%
\pgfpathlineto{\pgfqpoint{5.836822in}{1.953357in}}%
\pgfpathlineto{\pgfqpoint{5.837946in}{1.952008in}}%
\pgfpathlineto{\pgfqpoint{5.839521in}{1.951330in}}%
\pgfpathlineto{\pgfqpoint{5.840645in}{1.949971in}}%
\pgfpathlineto{\pgfqpoint{5.842043in}{1.949288in}}%
\pgfpathlineto{\pgfqpoint{5.843167in}{1.947919in}}%
\pgfpathlineto{\pgfqpoint{5.844846in}{1.947231in}}%
\pgfpathlineto{\pgfqpoint{5.845969in}{1.945851in}}%
\pgfpathlineto{\pgfqpoint{5.847439in}{1.945159in}}%
\pgfpathlineto{\pgfqpoint{5.848562in}{1.943769in}}%
\pgfpathlineto{\pgfqpoint{5.850431in}{1.943071in}}%
\pgfpathlineto{\pgfqpoint{5.851552in}{1.941670in}}%
\pgfpathlineto{\pgfqpoint{5.853394in}{1.940967in}}%
\pgfpathlineto{\pgfqpoint{5.854514in}{1.939556in}}%
\pgfpathlineto{\pgfqpoint{5.856129in}{1.938848in}}%
\pgfpathlineto{\pgfqpoint{5.857248in}{1.937426in}}%
\pgfpathlineto{\pgfqpoint{5.859032in}{1.936712in}}%
\pgfpathlineto{\pgfqpoint{5.860151in}{1.935279in}}%
\pgfpathlineto{\pgfqpoint{5.862054in}{1.934560in}}%
\pgfpathlineto{\pgfqpoint{5.863178in}{1.933115in}}%
\pgfpathlineto{\pgfqpoint{5.865127in}{1.932391in}}%
\pgfpathlineto{\pgfqpoint{5.866249in}{1.930935in}}%
\pgfpathlineto{\pgfqpoint{5.868470in}{1.930204in}}%
\pgfpathlineto{\pgfqpoint{5.869589in}{1.928737in}}%
\pgfpathlineto{\pgfqpoint{5.871599in}{1.928001in}}%
\pgfpathlineto{\pgfqpoint{5.872723in}{1.926522in}}%
\pgfpathlineto{\pgfqpoint{5.874904in}{1.925780in}}%
\pgfpathlineto{\pgfqpoint{5.876025in}{1.924290in}}%
\pgfpathlineto{\pgfqpoint{5.878227in}{1.923541in}}%
\pgfpathlineto{\pgfqpoint{5.879339in}{1.922791in}}%
\pgfpathlineto{\pgfqpoint{5.880556in}{1.922039in}}%
\pgfpathlineto{\pgfqpoint{5.881675in}{1.920528in}}%
\pgfpathlineto{\pgfqpoint{5.883909in}{1.919770in}}%
\pgfpathlineto{\pgfqpoint{5.885025in}{1.919009in}}%
\pgfpathlineto{\pgfqpoint{5.886363in}{1.918247in}}%
\pgfpathlineto{\pgfqpoint{5.887479in}{1.917482in}}%
\pgfpathlineto{\pgfqpoint{5.888778in}{1.916715in}}%
\pgfpathlineto{\pgfqpoint{5.889893in}{1.915946in}}%
\pgfpathlineto{\pgfqpoint{5.891228in}{1.915175in}}%
\pgfpathlineto{\pgfqpoint{5.892342in}{1.914402in}}%
\pgfpathlineto{\pgfqpoint{5.893927in}{1.913627in}}%
\pgfpathlineto{\pgfqpoint{5.895040in}{1.912849in}}%
\pgfpathlineto{\pgfqpoint{5.896535in}{1.912070in}}%
\pgfpathlineto{\pgfqpoint{5.897648in}{1.911288in}}%
\pgfpathlineto{\pgfqpoint{5.899043in}{1.910504in}}%
\pgfpathlineto{\pgfqpoint{5.900155in}{1.909717in}}%
\pgfpathlineto{\pgfqpoint{5.901661in}{1.908929in}}%
\pgfpathlineto{\pgfqpoint{5.902773in}{1.908138in}}%
\pgfpathlineto{\pgfqpoint{5.904234in}{1.907345in}}%
\pgfpathlineto{\pgfqpoint{5.905350in}{1.906550in}}%
\pgfpathlineto{\pgfqpoint{5.907111in}{1.905752in}}%
\pgfpathlineto{\pgfqpoint{5.908226in}{1.904953in}}%
\pgfpathlineto{\pgfqpoint{5.909868in}{1.904151in}}%
\pgfpathlineto{\pgfqpoint{5.910982in}{1.903346in}}%
\pgfpathlineto{\pgfqpoint{5.912789in}{1.902539in}}%
\pgfpathlineto{\pgfqpoint{5.913906in}{1.901730in}}%
\pgfpathlineto{\pgfqpoint{5.915532in}{1.900919in}}%
\pgfpathlineto{\pgfqpoint{5.916648in}{1.900105in}}%
\pgfpathlineto{\pgfqpoint{5.918667in}{1.899289in}}%
\pgfpathlineto{\pgfqpoint{5.919780in}{1.898470in}}%
\pgfpathlineto{\pgfqpoint{5.921653in}{1.897649in}}%
\pgfpathlineto{\pgfqpoint{5.922769in}{1.896826in}}%
\pgfpathlineto{\pgfqpoint{5.924646in}{1.896000in}}%
\pgfpathlineto{\pgfqpoint{5.925760in}{1.895172in}}%
\pgfpathlineto{\pgfqpoint{5.927707in}{1.894341in}}%
\pgfpathlineto{\pgfqpoint{5.928819in}{1.893508in}}%
\pgfpathlineto{\pgfqpoint{5.930819in}{1.892672in}}%
\pgfpathlineto{\pgfqpoint{5.931934in}{1.891834in}}%
\pgfpathlineto{\pgfqpoint{5.934009in}{1.890993in}}%
\pgfpathlineto{\pgfqpoint{5.935120in}{1.890150in}}%
\pgfpathlineto{\pgfqpoint{5.937223in}{1.889304in}}%
\pgfpathlineto{\pgfqpoint{5.938337in}{1.888455in}}%
\pgfpathlineto{\pgfqpoint{5.940544in}{1.887604in}}%
\pgfpathlineto{\pgfqpoint{5.941660in}{1.886750in}}%
\pgfpathlineto{\pgfqpoint{5.943817in}{1.885894in}}%
\pgfpathlineto{\pgfqpoint{5.944930in}{1.885035in}}%
\pgfpathlineto{\pgfqpoint{5.947366in}{1.884174in}}%
\pgfpathlineto{\pgfqpoint{5.948480in}{1.883309in}}%
\pgfpathlineto{\pgfqpoint{5.950998in}{1.882442in}}%
\pgfpathlineto{\pgfqpoint{5.952112in}{1.881573in}}%
\pgfpathlineto{\pgfqpoint{5.954730in}{1.880700in}}%
\pgfpathlineto{\pgfqpoint{5.955845in}{1.879825in}}%
\pgfpathlineto{\pgfqpoint{5.958712in}{1.878947in}}%
\pgfpathlineto{\pgfqpoint{5.959827in}{1.878067in}}%
\pgfpathlineto{\pgfqpoint{5.962625in}{1.877183in}}%
\pgfpathlineto{\pgfqpoint{5.963739in}{1.876297in}}%
\pgfpathlineto{\pgfqpoint{5.966496in}{1.875408in}}%
\pgfpathlineto{\pgfqpoint{5.967611in}{1.874516in}}%
\pgfpathlineto{\pgfqpoint{5.970540in}{1.873621in}}%
\pgfpathlineto{\pgfqpoint{5.971654in}{1.872723in}}%
\pgfpathlineto{\pgfqpoint{5.974424in}{1.871823in}}%
\pgfpathlineto{\pgfqpoint{5.975538in}{1.870919in}}%
\pgfpathlineto{\pgfqpoint{5.978327in}{1.870013in}}%
\pgfpathlineto{\pgfqpoint{5.979441in}{1.869104in}}%
\pgfpathlineto{\pgfqpoint{5.982303in}{1.868191in}}%
\pgfpathlineto{\pgfqpoint{5.983416in}{1.867276in}}%
\pgfpathlineto{\pgfqpoint{5.986579in}{1.866357in}}%
\pgfpathlineto{\pgfqpoint{5.987695in}{1.865436in}}%
\pgfpathlineto{\pgfqpoint{5.991127in}{1.864512in}}%
\pgfpathlineto{\pgfqpoint{5.992240in}{1.863584in}}%
\pgfpathlineto{\pgfqpoint{5.995520in}{1.862653in}}%
\pgfpathlineto{\pgfqpoint{5.996631in}{1.861720in}}%
\pgfpathlineto{\pgfqpoint{5.999909in}{1.860783in}}%
\pgfpathlineto{\pgfqpoint{6.001023in}{1.859843in}}%
\pgfpathlineto{\pgfqpoint{6.004769in}{1.858899in}}%
\pgfpathlineto{\pgfqpoint{6.005884in}{1.857953in}}%
\pgfpathlineto{\pgfqpoint{6.009578in}{1.857003in}}%
\pgfpathlineto{\pgfqpoint{6.010689in}{1.856050in}}%
\pgfpathlineto{\pgfqpoint{6.014656in}{1.855094in}}%
\pgfpathlineto{\pgfqpoint{6.015768in}{1.854135in}}%
\pgfpathlineto{\pgfqpoint{6.019714in}{1.853172in}}%
\pgfpathlineto{\pgfqpoint{6.020826in}{1.852206in}}%
\pgfpathlineto{\pgfqpoint{6.024536in}{1.851236in}}%
\pgfpathlineto{\pgfqpoint{6.025649in}{1.850264in}}%
\pgfpathlineto{\pgfqpoint{6.029648in}{1.849287in}}%
\pgfpathlineto{\pgfqpoint{6.030760in}{1.848308in}}%
\pgfpathlineto{\pgfqpoint{6.034794in}{1.847324in}}%
\pgfpathlineto{\pgfqpoint{6.035909in}{1.846338in}}%
\pgfpathlineto{\pgfqpoint{6.039997in}{1.845348in}}%
\pgfpathlineto{\pgfqpoint{6.041111in}{1.844354in}}%
\pgfpathlineto{\pgfqpoint{6.045391in}{1.843357in}}%
\pgfpathlineto{\pgfqpoint{6.046503in}{1.842356in}}%
\pgfpathlineto{\pgfqpoint{6.051135in}{1.841351in}}%
\pgfpathlineto{\pgfqpoint{6.052249in}{1.840343in}}%
\pgfpathlineto{\pgfqpoint{6.056755in}{1.839332in}}%
\pgfpathlineto{\pgfqpoint{6.057867in}{1.838316in}}%
\pgfpathlineto{\pgfqpoint{6.062537in}{1.837297in}}%
\pgfpathlineto{\pgfqpoint{6.063650in}{1.836274in}}%
\pgfpathlineto{\pgfqpoint{6.068493in}{1.835248in}}%
\pgfpathlineto{\pgfqpoint{6.069606in}{1.834217in}}%
\pgfpathlineto{\pgfqpoint{6.074424in}{1.833183in}}%
\pgfpathlineto{\pgfqpoint{6.075537in}{1.832145in}}%
\pgfpathlineto{\pgfqpoint{6.080818in}{1.831103in}}%
\pgfpathlineto{\pgfqpoint{6.081930in}{1.830057in}}%
\pgfpathlineto{\pgfqpoint{6.087277in}{1.829007in}}%
\pgfpathlineto{\pgfqpoint{6.088391in}{1.827953in}}%
\pgfpathlineto{\pgfqpoint{6.093251in}{1.826896in}}%
\pgfpathlineto{\pgfqpoint{6.094363in}{1.825834in}}%
\pgfpathlineto{\pgfqpoint{6.099621in}{1.824768in}}%
\pgfpathlineto{\pgfqpoint{6.100735in}{1.823698in}}%
\pgfpathlineto{\pgfqpoint{6.105939in}{1.822624in}}%
\pgfpathlineto{\pgfqpoint{6.107051in}{1.821546in}}%
\pgfpathlineto{\pgfqpoint{6.112719in}{1.820463in}}%
\pgfpathlineto{\pgfqpoint{6.113831in}{1.819377in}}%
\pgfpathlineto{\pgfqpoint{6.119554in}{1.818286in}}%
\pgfpathlineto{\pgfqpoint{6.120666in}{1.817191in}}%
\pgfpathlineto{\pgfqpoint{6.126415in}{1.816091in}}%
\pgfpathlineto{\pgfqpoint{6.127529in}{1.814987in}}%
\pgfpathlineto{\pgfqpoint{6.133504in}{1.813879in}}%
\pgfpathlineto{\pgfqpoint{6.134617in}{1.812766in}}%
\pgfpathlineto{\pgfqpoint{6.136936in}{1.812766in}}%
\pgfpathlineto{\pgfqpoint{6.138049in}{1.811649in}}%
\pgfpathlineto{\pgfqpoint{6.140393in}{1.811649in}}%
\pgfpathlineto{\pgfqpoint{6.141504in}{1.810527in}}%
\pgfpathlineto{\pgfqpoint{6.143990in}{1.810527in}}%
\pgfpathlineto{\pgfqpoint{6.145102in}{1.809401in}}%
\pgfpathlineto{\pgfqpoint{6.147585in}{1.809401in}}%
\pgfpathlineto{\pgfqpoint{6.148698in}{1.808270in}}%
\pgfpathlineto{\pgfqpoint{6.151305in}{1.808270in}}%
\pgfpathlineto{\pgfqpoint{6.152417in}{1.807135in}}%
\pgfpathlineto{\pgfqpoint{6.155213in}{1.807135in}}%
\pgfpathlineto{\pgfqpoint{6.156325in}{1.805995in}}%
\pgfpathlineto{\pgfqpoint{6.158814in}{1.805995in}}%
\pgfpathlineto{\pgfqpoint{6.159925in}{1.804850in}}%
\pgfpathlineto{\pgfqpoint{6.162390in}{1.804850in}}%
\pgfpathlineto{\pgfqpoint{6.163501in}{1.803701in}}%
\pgfpathlineto{\pgfqpoint{6.166211in}{1.803701in}}%
\pgfpathlineto{\pgfqpoint{6.167325in}{1.802547in}}%
\pgfpathlineto{\pgfqpoint{6.169929in}{1.802547in}}%
\pgfpathlineto{\pgfqpoint{6.171042in}{1.801388in}}%
\pgfpathlineto{\pgfqpoint{6.173691in}{1.801388in}}%
\pgfpathlineto{\pgfqpoint{6.174804in}{1.800224in}}%
\pgfpathlineto{\pgfqpoint{6.177569in}{1.800224in}}%
\pgfpathlineto{\pgfqpoint{6.178680in}{1.799056in}}%
\pgfpathlineto{\pgfqpoint{6.181266in}{1.799056in}}%
\pgfpathlineto{\pgfqpoint{6.182380in}{1.797882in}}%
\pgfpathlineto{\pgfqpoint{6.185277in}{1.797882in}}%
\pgfpathlineto{\pgfqpoint{6.186389in}{1.796703in}}%
\pgfpathlineto{\pgfqpoint{6.189138in}{1.796703in}}%
\pgfpathlineto{\pgfqpoint{6.190249in}{1.795520in}}%
\pgfpathlineto{\pgfqpoint{6.193053in}{1.795520in}}%
\pgfpathlineto{\pgfqpoint{6.194165in}{1.794331in}}%
\pgfpathlineto{\pgfqpoint{6.197052in}{1.794331in}}%
\pgfpathlineto{\pgfqpoint{6.198166in}{1.793138in}}%
\pgfpathlineto{\pgfqpoint{6.201083in}{1.793138in}}%
\pgfpathlineto{\pgfqpoint{6.202195in}{1.791939in}}%
\pgfpathlineto{\pgfqpoint{6.205399in}{1.791939in}}%
\pgfpathlineto{\pgfqpoint{6.206511in}{1.790735in}}%
\pgfpathlineto{\pgfqpoint{6.209448in}{1.790735in}}%
\pgfpathlineto{\pgfqpoint{6.210560in}{1.789525in}}%
\pgfpathlineto{\pgfqpoint{6.213480in}{1.789525in}}%
\pgfpathlineto{\pgfqpoint{6.214593in}{1.788311in}}%
\pgfpathlineto{\pgfqpoint{6.217703in}{1.788311in}}%
\pgfpathlineto{\pgfqpoint{6.218816in}{1.787091in}}%
\pgfpathlineto{\pgfqpoint{6.221862in}{1.787091in}}%
\pgfpathlineto{\pgfqpoint{6.222975in}{1.785866in}}%
\pgfpathlineto{\pgfqpoint{6.226057in}{1.785866in}}%
\pgfpathlineto{\pgfqpoint{6.227169in}{1.784635in}}%
\pgfpathlineto{\pgfqpoint{6.230228in}{1.784635in}}%
\pgfpathlineto{\pgfqpoint{6.231340in}{1.783399in}}%
\pgfpathlineto{\pgfqpoint{6.234401in}{1.783399in}}%
\pgfpathlineto{\pgfqpoint{6.235513in}{1.782158in}}%
\pgfpathlineto{\pgfqpoint{6.238826in}{1.782158in}}%
\pgfpathlineto{\pgfqpoint{6.239940in}{1.780911in}}%
\pgfpathlineto{\pgfqpoint{6.243168in}{1.780911in}}%
\pgfpathlineto{\pgfqpoint{6.244280in}{1.779658in}}%
\pgfpathlineto{\pgfqpoint{6.247579in}{1.779658in}}%
\pgfpathlineto{\pgfqpoint{6.248690in}{1.778400in}}%
\pgfpathlineto{\pgfqpoint{6.252078in}{1.778400in}}%
\pgfpathlineto{\pgfqpoint{6.253190in}{1.777136in}}%
\pgfpathlineto{\pgfqpoint{6.256404in}{1.777136in}}%
\pgfpathlineto{\pgfqpoint{6.257517in}{1.775866in}}%
\pgfpathlineto{\pgfqpoint{6.260829in}{1.775866in}}%
\pgfpathlineto{\pgfqpoint{6.261941in}{1.774590in}}%
\pgfpathlineto{\pgfqpoint{6.265545in}{1.774590in}}%
\pgfpathlineto{\pgfqpoint{6.266657in}{1.773309in}}%
\pgfpathlineto{\pgfqpoint{6.270023in}{1.773309in}}%
\pgfpathlineto{\pgfqpoint{6.271135in}{1.772021in}}%
\pgfpathlineto{\pgfqpoint{6.274598in}{1.772021in}}%
\pgfpathlineto{\pgfqpoint{6.275709in}{1.770728in}}%
\pgfpathlineto{\pgfqpoint{6.278969in}{1.770728in}}%
\pgfpathlineto{\pgfqpoint{6.280082in}{1.769428in}}%
\pgfpathlineto{\pgfqpoint{6.283583in}{1.769428in}}%
\pgfpathlineto{\pgfqpoint{6.284695in}{1.768123in}}%
\pgfpathlineto{\pgfqpoint{6.288115in}{1.768123in}}%
\pgfpathlineto{\pgfqpoint{6.289227in}{1.766811in}}%
\pgfpathlineto{\pgfqpoint{6.292678in}{1.766811in}}%
\pgfpathlineto{\pgfqpoint{6.293789in}{1.765493in}}%
\pgfpathlineto{\pgfqpoint{6.297344in}{1.765493in}}%
\pgfpathlineto{\pgfqpoint{6.298457in}{1.764169in}}%
\pgfpathlineto{\pgfqpoint{6.301991in}{1.764169in}}%
\pgfpathlineto{\pgfqpoint{6.303102in}{1.762839in}}%
\pgfpathlineto{\pgfqpoint{6.306739in}{1.762839in}}%
\pgfpathlineto{\pgfqpoint{6.307852in}{1.761502in}}%
\pgfpathlineto{\pgfqpoint{6.311410in}{1.761502in}}%
\pgfpathlineto{\pgfqpoint{6.312521in}{1.760159in}}%
\pgfpathlineto{\pgfqpoint{6.315988in}{1.760159in}}%
\pgfpathlineto{\pgfqpoint{6.317101in}{1.758809in}}%
\pgfpathlineto{\pgfqpoint{6.320852in}{1.758809in}}%
\pgfpathlineto{\pgfqpoint{6.321964in}{1.757453in}}%
\pgfpathlineto{\pgfqpoint{6.325776in}{1.757453in}}%
\pgfpathlineto{\pgfqpoint{6.326888in}{1.756090in}}%
\pgfpathlineto{\pgfqpoint{6.330420in}{1.756090in}}%
\pgfpathlineto{\pgfqpoint{6.331532in}{1.754720in}}%
\pgfpathlineto{\pgfqpoint{6.335253in}{1.754720in}}%
\pgfpathlineto{\pgfqpoint{6.336366in}{1.753344in}}%
\pgfpathlineto{\pgfqpoint{6.340008in}{1.753344in}}%
\pgfpathlineto{\pgfqpoint{6.341120in}{1.751960in}}%
\pgfpathlineto{\pgfqpoint{6.344830in}{1.751960in}}%
\pgfpathlineto{\pgfqpoint{6.345942in}{1.750570in}}%
\pgfpathlineto{\pgfqpoint{6.349663in}{1.750570in}}%
\pgfpathlineto{\pgfqpoint{6.350775in}{1.749173in}}%
\pgfpathlineto{\pgfqpoint{6.354648in}{1.749173in}}%
\pgfpathlineto{\pgfqpoint{6.355760in}{1.747769in}}%
\pgfpathlineto{\pgfqpoint{6.359751in}{1.747769in}}%
\pgfpathlineto{\pgfqpoint{6.360864in}{1.746358in}}%
\pgfpathlineto{\pgfqpoint{6.364733in}{1.746358in}}%
\pgfpathlineto{\pgfqpoint{6.365846in}{1.744939in}}%
\pgfpathlineto{\pgfqpoint{6.369800in}{1.744939in}}%
\pgfpathlineto{\pgfqpoint{6.370911in}{1.743513in}}%
\pgfpathlineto{\pgfqpoint{6.374842in}{1.743513in}}%
\pgfpathlineto{\pgfqpoint{6.375953in}{1.742080in}}%
\pgfpathlineto{\pgfqpoint{6.379717in}{1.742080in}}%
\pgfpathlineto{\pgfqpoint{6.380828in}{1.740640in}}%
\pgfpathlineto{\pgfqpoint{6.384746in}{1.740640in}}%
\pgfpathlineto{\pgfqpoint{6.385857in}{1.739192in}}%
\pgfpathlineto{\pgfqpoint{6.389944in}{1.739192in}}%
\pgfpathlineto{\pgfqpoint{6.391056in}{1.737736in}}%
\pgfpathlineto{\pgfqpoint{6.394981in}{1.737736in}}%
\pgfpathlineto{\pgfqpoint{6.396092in}{1.736273in}}%
\pgfpathlineto{\pgfqpoint{6.400203in}{1.736273in}}%
\pgfpathlineto{\pgfqpoint{6.401315in}{1.734802in}}%
\pgfpathlineto{\pgfqpoint{6.405316in}{1.734802in}}%
\pgfpathlineto{\pgfqpoint{6.406429in}{1.733324in}}%
\pgfpathlineto{\pgfqpoint{6.410522in}{1.733324in}}%
\pgfpathlineto{\pgfqpoint{6.411635in}{1.731837in}}%
\pgfpathlineto{\pgfqpoint{6.415741in}{1.731837in}}%
\pgfpathlineto{\pgfqpoint{6.416854in}{1.730343in}}%
\pgfpathlineto{\pgfqpoint{6.420881in}{1.730343in}}%
\pgfpathlineto{\pgfqpoint{6.421992in}{1.728840in}}%
\pgfpathlineto{\pgfqpoint{6.426239in}{1.728840in}}%
\pgfpathlineto{\pgfqpoint{6.427351in}{1.727329in}}%
\pgfpathlineto{\pgfqpoint{6.431349in}{1.727329in}}%
\pgfpathlineto{\pgfqpoint{6.432461in}{1.725810in}}%
\pgfpathlineto{\pgfqpoint{6.436540in}{1.725810in}}%
\pgfpathlineto{\pgfqpoint{6.437651in}{1.724283in}}%
\pgfpathlineto{\pgfqpoint{6.441635in}{1.724283in}}%
\pgfpathlineto{\pgfqpoint{6.442748in}{1.722748in}}%
\pgfpathlineto{\pgfqpoint{6.446701in}{1.722748in}}%
\pgfpathlineto{\pgfqpoint{6.447813in}{1.721203in}}%
\pgfpathlineto{\pgfqpoint{6.451907in}{1.721203in}}%
\pgfpathlineto{\pgfqpoint{6.453019in}{1.719651in}}%
\pgfpathlineto{\pgfqpoint{6.456993in}{1.719651in}}%
\pgfpathlineto{\pgfqpoint{6.458105in}{1.718089in}}%
\pgfpathlineto{\pgfqpoint{6.462184in}{1.718089in}}%
\pgfpathlineto{\pgfqpoint{6.463296in}{1.716519in}}%
\pgfpathlineto{\pgfqpoint{6.467483in}{1.716519in}}%
\pgfpathlineto{\pgfqpoint{6.468595in}{1.714940in}}%
\pgfpathlineto{\pgfqpoint{6.472767in}{1.714940in}}%
\pgfpathlineto{\pgfqpoint{6.473879in}{1.713351in}}%
\pgfpathlineto{\pgfqpoint{6.477941in}{1.713351in}}%
\pgfpathlineto{\pgfqpoint{6.479053in}{1.711754in}}%
\pgfpathlineto{\pgfqpoint{6.483127in}{1.711754in}}%
\pgfpathlineto{\pgfqpoint{6.484238in}{1.710148in}}%
\pgfpathlineto{\pgfqpoint{6.488264in}{1.710148in}}%
\pgfpathlineto{\pgfqpoint{6.489377in}{1.708532in}}%
\pgfpathlineto{\pgfqpoint{6.493616in}{1.708532in}}%
\pgfpathlineto{\pgfqpoint{6.494729in}{1.706906in}}%
\pgfpathlineto{\pgfqpoint{6.498888in}{1.706906in}}%
\pgfpathlineto{\pgfqpoint{6.500000in}{1.705272in}}%
\pgfpathlineto{\pgfqpoint{6.503998in}{1.705272in}}%
\pgfpathlineto{\pgfqpoint{6.505109in}{1.703627in}}%
\pgfpathlineto{\pgfqpoint{6.509340in}{1.703627in}}%
\pgfpathlineto{\pgfqpoint{6.510452in}{1.701973in}}%
\pgfpathlineto{\pgfqpoint{6.514684in}{1.701973in}}%
\pgfpathlineto{\pgfqpoint{6.515796in}{1.700309in}}%
\pgfpathlineto{\pgfqpoint{6.519949in}{1.700309in}}%
\pgfpathlineto{\pgfqpoint{6.521060in}{1.698635in}}%
\pgfpathlineto{\pgfqpoint{6.525293in}{1.698635in}}%
\pgfpathlineto{\pgfqpoint{6.526405in}{1.696951in}}%
\pgfpathlineto{\pgfqpoint{6.530502in}{1.696951in}}%
\pgfpathlineto{\pgfqpoint{6.531614in}{1.695257in}}%
\pgfpathlineto{\pgfqpoint{6.535747in}{1.695257in}}%
\pgfpathlineto{\pgfqpoint{6.536859in}{1.693552in}}%
\pgfpathlineto{\pgfqpoint{6.540934in}{1.693552in}}%
\pgfpathlineto{\pgfqpoint{6.542046in}{1.691837in}}%
\pgfpathlineto{\pgfqpoint{6.546218in}{1.691837in}}%
\pgfpathlineto{\pgfqpoint{6.547330in}{1.690111in}}%
\pgfpathlineto{\pgfqpoint{6.551361in}{1.690111in}}%
\pgfpathlineto{\pgfqpoint{6.552473in}{1.688374in}}%
\pgfpathlineto{\pgfqpoint{6.556648in}{1.688374in}}%
\pgfpathlineto{\pgfqpoint{6.557759in}{1.686627in}}%
\pgfpathlineto{\pgfqpoint{6.561966in}{1.686627in}}%
\pgfpathlineto{\pgfqpoint{6.563078in}{1.684868in}}%
\pgfpathlineto{\pgfqpoint{6.567217in}{1.684868in}}%
\pgfpathlineto{\pgfqpoint{6.568328in}{1.683098in}}%
\pgfpathlineto{\pgfqpoint{6.572295in}{1.683098in}}%
\pgfpathlineto{\pgfqpoint{6.573408in}{1.681317in}}%
\pgfpathlineto{\pgfqpoint{6.577457in}{1.681317in}}%
\pgfpathlineto{\pgfqpoint{6.578568in}{1.679525in}}%
\pgfpathlineto{\pgfqpoint{6.582671in}{1.679525in}}%
\pgfpathlineto{\pgfqpoint{6.583782in}{1.677721in}}%
\pgfpathlineto{\pgfqpoint{6.587813in}{1.677721in}}%
\pgfpathlineto{\pgfqpoint{6.588925in}{1.675905in}}%
\pgfpathlineto{\pgfqpoint{6.593020in}{1.675905in}}%
\pgfpathlineto{\pgfqpoint{6.594131in}{1.674077in}}%
\pgfpathlineto{\pgfqpoint{6.598172in}{1.674077in}}%
\pgfpathlineto{\pgfqpoint{6.599284in}{1.672237in}}%
\pgfpathlineto{\pgfqpoint{6.603297in}{1.672237in}}%
\pgfpathlineto{\pgfqpoint{6.604409in}{1.670385in}}%
\pgfpathlineto{\pgfqpoint{6.608373in}{1.670385in}}%
\pgfpathlineto{\pgfqpoint{6.609484in}{1.668521in}}%
\pgfpathlineto{\pgfqpoint{6.613349in}{1.668521in}}%
\pgfpathlineto{\pgfqpoint{6.614461in}{1.666644in}}%
\pgfpathlineto{\pgfqpoint{6.618590in}{1.666644in}}%
\pgfpathlineto{\pgfqpoint{6.619702in}{1.664754in}}%
\pgfpathlineto{\pgfqpoint{6.623698in}{1.664754in}}%
\pgfpathlineto{\pgfqpoint{6.624810in}{1.662852in}}%
\pgfpathlineto{\pgfqpoint{6.628828in}{1.662852in}}%
\pgfpathlineto{\pgfqpoint{6.629940in}{1.660936in}}%
\pgfpathlineto{\pgfqpoint{6.634029in}{1.660936in}}%
\pgfpathlineto{\pgfqpoint{6.635141in}{1.659007in}}%
\pgfpathlineto{\pgfqpoint{6.639070in}{1.659007in}}%
\pgfpathlineto{\pgfqpoint{6.640181in}{1.657065in}}%
\pgfpathlineto{\pgfqpoint{6.644142in}{1.657065in}}%
\pgfpathlineto{\pgfqpoint{6.645254in}{1.655109in}}%
\pgfpathlineto{\pgfqpoint{6.649279in}{1.655109in}}%
\pgfpathlineto{\pgfqpoint{6.650391in}{1.653139in}}%
\pgfpathlineto{\pgfqpoint{6.654206in}{1.653139in}}%
\pgfpathlineto{\pgfqpoint{6.655317in}{1.651155in}}%
\pgfpathlineto{\pgfqpoint{6.659265in}{1.651155in}}%
\pgfpathlineto{\pgfqpoint{6.660377in}{1.649157in}}%
\pgfpathlineto{\pgfqpoint{6.664295in}{1.649157in}}%
\pgfpathlineto{\pgfqpoint{6.665407in}{1.647145in}}%
\pgfpathlineto{\pgfqpoint{6.669198in}{1.647145in}}%
\pgfpathlineto{\pgfqpoint{6.670309in}{1.645118in}}%
\pgfpathlineto{\pgfqpoint{6.674096in}{1.645118in}}%
\pgfpathlineto{\pgfqpoint{6.675208in}{1.643076in}}%
\pgfpathlineto{\pgfqpoint{6.679042in}{1.643076in}}%
\pgfpathlineto{\pgfqpoint{6.680153in}{1.641019in}}%
\pgfpathlineto{\pgfqpoint{6.683901in}{1.641019in}}%
\pgfpathlineto{\pgfqpoint{6.685012in}{1.638946in}}%
\pgfpathlineto{\pgfqpoint{6.688695in}{1.638946in}}%
\pgfpathlineto{\pgfqpoint{6.689807in}{1.636858in}}%
\pgfpathlineto{\pgfqpoint{6.693517in}{1.636858in}}%
\pgfpathlineto{\pgfqpoint{6.694628in}{1.634755in}}%
\pgfpathlineto{\pgfqpoint{6.698255in}{1.634755in}}%
\pgfpathlineto{\pgfqpoint{6.699366in}{1.632635in}}%
\pgfpathlineto{\pgfqpoint{6.702929in}{1.632635in}}%
\pgfpathlineto{\pgfqpoint{6.704041in}{1.630499in}}%
\pgfpathlineto{\pgfqpoint{6.707676in}{1.630499in}}%
\pgfpathlineto{\pgfqpoint{6.708787in}{1.628347in}}%
\pgfpathlineto{\pgfqpoint{6.712358in}{1.628347in}}%
\pgfpathlineto{\pgfqpoint{6.713469in}{1.626178in}}%
\pgfpathlineto{\pgfqpoint{6.716980in}{1.626178in}}%
\pgfpathlineto{\pgfqpoint{6.718092in}{1.623992in}}%
\pgfpathlineto{\pgfqpoint{6.721463in}{1.623992in}}%
\pgfpathlineto{\pgfqpoint{6.722574in}{1.621788in}}%
\pgfpathlineto{\pgfqpoint{6.726031in}{1.621788in}}%
\pgfpathlineto{\pgfqpoint{6.727143in}{1.619568in}}%
\pgfpathlineto{\pgfqpoint{6.730731in}{1.619568in}}%
\pgfpathlineto{\pgfqpoint{6.731842in}{1.617329in}}%
\pgfpathlineto{\pgfqpoint{6.735303in}{1.617329in}}%
\pgfpathlineto{\pgfqpoint{6.736414in}{1.615072in}}%
\pgfpathlineto{\pgfqpoint{6.739827in}{1.615072in}}%
\pgfpathlineto{\pgfqpoint{6.740939in}{1.612797in}}%
\pgfpathlineto{\pgfqpoint{6.744326in}{1.612797in}}%
\pgfpathlineto{\pgfqpoint{6.745438in}{1.610502in}}%
\pgfpathlineto{\pgfqpoint{6.748609in}{1.610502in}}%
\pgfpathlineto{\pgfqpoint{6.749721in}{1.608189in}}%
\pgfpathlineto{\pgfqpoint{6.753028in}{1.608189in}}%
\pgfpathlineto{\pgfqpoint{6.754140in}{1.605857in}}%
\pgfpathlineto{\pgfqpoint{6.757293in}{1.605857in}}%
\pgfpathlineto{\pgfqpoint{6.758405in}{1.603505in}}%
\pgfpathlineto{\pgfqpoint{6.761485in}{1.603505in}}%
\pgfpathlineto{\pgfqpoint{6.762597in}{1.601133in}}%
\pgfpathlineto{\pgfqpoint{6.765845in}{1.601133in}}%
\pgfpathlineto{\pgfqpoint{6.766956in}{1.598740in}}%
\pgfpathlineto{\pgfqpoint{6.770082in}{1.598740in}}%
\pgfpathlineto{\pgfqpoint{6.771193in}{1.596327in}}%
\pgfpathlineto{\pgfqpoint{6.774186in}{1.596327in}}%
\pgfpathlineto{\pgfqpoint{6.775298in}{1.593893in}}%
\pgfpathlineto{\pgfqpoint{6.778357in}{1.593893in}}%
\pgfpathlineto{\pgfqpoint{6.779469in}{1.591437in}}%
\pgfpathlineto{\pgfqpoint{6.782384in}{1.591437in}}%
\pgfpathlineto{\pgfqpoint{6.783495in}{1.588959in}}%
\pgfpathlineto{\pgfqpoint{6.786350in}{1.588959in}}%
\pgfpathlineto{\pgfqpoint{6.787461in}{1.586459in}}%
\pgfpathlineto{\pgfqpoint{6.790261in}{1.586459in}}%
\pgfpathlineto{\pgfqpoint{6.791372in}{1.583937in}}%
\pgfpathlineto{\pgfqpoint{6.794182in}{1.583937in}}%
\pgfpathlineto{\pgfqpoint{6.795293in}{1.581391in}}%
\pgfpathlineto{\pgfqpoint{6.798064in}{1.581391in}}%
\pgfpathlineto{\pgfqpoint{6.799176in}{1.578823in}}%
\pgfpathlineto{\pgfqpoint{6.801754in}{1.578823in}}%
\pgfpathlineto{\pgfqpoint{6.802866in}{1.576230in}}%
\pgfpathlineto{\pgfqpoint{6.805398in}{1.576230in}}%
\pgfpathlineto{\pgfqpoint{6.806509in}{1.573613in}}%
\pgfpathlineto{\pgfqpoint{6.809036in}{1.573613in}}%
\pgfpathlineto{\pgfqpoint{6.810148in}{1.570971in}}%
\pgfpathlineto{\pgfqpoint{6.812689in}{1.570971in}}%
\pgfpathlineto{\pgfqpoint{6.813801in}{1.568303in}}%
\pgfpathlineto{\pgfqpoint{6.816271in}{1.568303in}}%
\pgfpathlineto{\pgfqpoint{6.817383in}{1.565610in}}%
\pgfpathlineto{\pgfqpoint{6.819736in}{1.565610in}}%
\pgfpathlineto{\pgfqpoint{6.820847in}{1.562891in}}%
\pgfpathlineto{\pgfqpoint{6.823103in}{1.562891in}}%
\pgfpathlineto{\pgfqpoint{6.824215in}{1.560145in}}%
\pgfpathlineto{\pgfqpoint{6.826460in}{1.560145in}}%
\pgfpathlineto{\pgfqpoint{6.827571in}{1.557371in}}%
\pgfpathlineto{\pgfqpoint{6.829805in}{1.557371in}}%
\pgfpathlineto{\pgfqpoint{6.830917in}{1.554570in}}%
\pgfpathlineto{\pgfqpoint{6.833048in}{1.554570in}}%
\pgfpathlineto{\pgfqpoint{6.834159in}{1.551740in}}%
\pgfpathlineto{\pgfqpoint{6.836302in}{1.551740in}}%
\pgfpathlineto{\pgfqpoint{6.837413in}{1.548882in}}%
\pgfpathlineto{\pgfqpoint{6.839497in}{1.548882in}}%
\pgfpathlineto{\pgfqpoint{6.840608in}{1.545993in}}%
\pgfpathlineto{\pgfqpoint{6.842591in}{1.545993in}}%
\pgfpathlineto{\pgfqpoint{6.843702in}{1.543075in}}%
\pgfpathlineto{\pgfqpoint{6.845631in}{1.543075in}}%
\pgfpathlineto{\pgfqpoint{6.846742in}{1.540125in}}%
\pgfpathlineto{\pgfqpoint{6.848609in}{1.540125in}}%
\pgfpathlineto{\pgfqpoint{6.849721in}{1.537144in}}%
\pgfpathlineto{\pgfqpoint{6.851480in}{1.537144in}}%
\pgfpathlineto{\pgfqpoint{6.852592in}{1.534131in}}%
\pgfpathlineto{\pgfqpoint{6.854310in}{1.534131in}}%
\pgfpathlineto{\pgfqpoint{6.855422in}{1.531085in}}%
\pgfpathlineto{\pgfqpoint{6.857044in}{1.531085in}}%
\pgfpathlineto{\pgfqpoint{6.858155in}{1.528005in}}%
\pgfpathlineto{\pgfqpoint{6.859758in}{1.528005in}}%
\pgfpathlineto{\pgfqpoint{6.860870in}{1.524890in}}%
\pgfpathlineto{\pgfqpoint{6.862450in}{1.524890in}}%
\pgfpathlineto{\pgfqpoint{6.863562in}{1.521741in}}%
\pgfpathlineto{\pgfqpoint{6.865077in}{1.521741in}}%
\pgfpathlineto{\pgfqpoint{6.866188in}{1.518555in}}%
\pgfpathlineto{\pgfqpoint{6.867628in}{1.518555in}}%
\pgfpathlineto{\pgfqpoint{6.868739in}{1.515333in}}%
\pgfpathlineto{\pgfqpoint{6.870077in}{1.515333in}}%
\pgfpathlineto{\pgfqpoint{6.871189in}{1.512073in}}%
\pgfpathlineto{\pgfqpoint{6.872457in}{1.512073in}}%
\pgfpathlineto{\pgfqpoint{6.873569in}{1.508774in}}%
\pgfpathlineto{\pgfqpoint{6.874734in}{1.508774in}}%
\pgfpathlineto{\pgfqpoint{6.875845in}{1.505436in}}%
\pgfpathlineto{\pgfqpoint{6.876968in}{1.505436in}}%
\pgfpathlineto{\pgfqpoint{6.878080in}{1.502058in}}%
\pgfpathlineto{\pgfqpoint{6.879118in}{1.502058in}}%
\pgfpathlineto{\pgfqpoint{6.880229in}{1.498638in}}%
\pgfpathlineto{\pgfqpoint{6.881176in}{1.498638in}}%
\pgfpathlineto{\pgfqpoint{6.882288in}{1.495175in}}%
\pgfpathlineto{\pgfqpoint{6.883136in}{1.495175in}}%
\pgfpathlineto{\pgfqpoint{6.884248in}{1.491669in}}%
\pgfpathlineto{\pgfqpoint{6.885094in}{1.491669in}}%
\pgfpathlineto{\pgfqpoint{6.886205in}{1.488119in}}%
\pgfpathlineto{\pgfqpoint{6.886964in}{1.488119in}}%
\pgfpathlineto{\pgfqpoint{6.888075in}{1.484522in}}%
\pgfpathlineto{\pgfqpoint{6.888845in}{1.484522in}}%
\pgfpathlineto{\pgfqpoint{6.889956in}{1.480879in}}%
\pgfpathlineto{\pgfqpoint{6.890599in}{1.480879in}}%
\pgfpathlineto{\pgfqpoint{6.891711in}{1.477187in}}%
\pgfpathlineto{\pgfqpoint{6.892310in}{1.477187in}}%
\pgfpathlineto{\pgfqpoint{6.893421in}{1.473445in}}%
\pgfpathlineto{\pgfqpoint{6.893950in}{1.473445in}}%
\pgfpathlineto{\pgfqpoint{6.895061in}{1.469653in}}%
\pgfpathlineto{\pgfqpoint{6.895503in}{1.469653in}}%
\pgfpathlineto{\pgfqpoint{6.896615in}{1.465809in}}%
\pgfpathlineto{\pgfqpoint{6.897000in}{1.465809in}}%
\pgfpathlineto{\pgfqpoint{6.898111in}{1.461910in}}%
\pgfpathlineto{\pgfqpoint{6.898481in}{1.461910in}}%
\pgfpathlineto{\pgfqpoint{6.899592in}{1.457957in}}%
\pgfpathlineto{\pgfqpoint{6.899852in}{1.457957in}}%
\pgfpathlineto{\pgfqpoint{6.900963in}{1.453946in}}%
\pgfpathlineto{\pgfqpoint{6.901182in}{1.453946in}}%
\pgfpathlineto{\pgfqpoint{6.902294in}{1.449877in}}%
\pgfpathlineto{\pgfqpoint{6.902386in}{1.449877in}}%
\pgfpathlineto{\pgfqpoint{6.903497in}{1.445748in}}%
\pgfpathlineto{\pgfqpoint{6.903566in}{1.445748in}}%
\pgfpathlineto{\pgfqpoint{6.904678in}{1.441556in}}%
\pgfpathlineto{\pgfqpoint{6.904721in}{1.441556in}}%
\pgfpathlineto{\pgfqpoint{6.905832in}{1.432979in}}%
\pgfpathlineto{\pgfqpoint{6.906734in}{1.432979in}}%
\pgfpathlineto{\pgfqpoint{6.907846in}{1.424130in}}%
\pgfpathlineto{\pgfqpoint{6.908526in}{1.424130in}}%
\pgfpathlineto{\pgfqpoint{6.909638in}{1.414991in}}%
\pgfpathlineto{\pgfqpoint{6.910169in}{1.414991in}}%
\pgfpathlineto{\pgfqpoint{6.911281in}{1.405541in}}%
\pgfpathlineto{\pgfqpoint{6.911579in}{1.405541in}}%
\pgfpathlineto{\pgfqpoint{6.912691in}{1.395761in}}%
\pgfpathlineto{\pgfqpoint{6.912765in}{1.395761in}}%
\pgfpathlineto{\pgfqpoint{6.913877in}{1.380414in}}%
\pgfpathlineto{\pgfqpoint{6.914223in}{1.380414in}}%
\pgfpathlineto{\pgfqpoint{6.915335in}{1.364173in}}%
\pgfpathlineto{\pgfqpoint{6.915337in}{1.364173in}}%
\pgfpathlineto{\pgfqpoint{6.916451in}{1.334806in}}%
\pgfpathlineto{\pgfqpoint{6.916583in}{1.334806in}}%
\pgfpathlineto{\pgfqpoint{6.917701in}{1.256678in}}%
\pgfpathlineto{\pgfqpoint{6.917753in}{1.256678in}}%
\pgfpathlineto{\pgfqpoint{6.918182in}{0.580000in}}%
\pgfpathlineto{\pgfqpoint{6.918182in}{0.580000in}}%
\pgfusepath{stroke}%
\end{pgfscope}%
\begin{pgfscope}%
\pgfsetrectcap%
\pgfsetmiterjoin%
\pgfsetlinewidth{0.803000pt}%
\definecolor{currentstroke}{rgb}{0.000000,0.000000,0.000000}%
\pgfsetstrokecolor{currentstroke}%
\pgfsetdash{}{0pt}%
\pgfpathmoveto{\pgfqpoint{1.000000in}{0.440000in}}%
\pgfpathlineto{\pgfqpoint{1.000000in}{3.520000in}}%
\pgfusepath{stroke}%
\end{pgfscope}%
\begin{pgfscope}%
\pgfsetrectcap%
\pgfsetmiterjoin%
\pgfsetlinewidth{0.803000pt}%
\definecolor{currentstroke}{rgb}{0.000000,0.000000,0.000000}%
\pgfsetstrokecolor{currentstroke}%
\pgfsetdash{}{0pt}%
\pgfpathmoveto{\pgfqpoint{1.000000in}{0.440000in}}%
\pgfpathlineto{\pgfqpoint{7.200000in}{0.440000in}}%
\pgfusepath{stroke}%
\end{pgfscope}%
\begin{pgfscope}%
\pgfsetbuttcap%
\pgfsetmiterjoin%
\definecolor{currentfill}{rgb}{1.000000,1.000000,1.000000}%
\pgfsetfillcolor{currentfill}%
\pgfsetfillopacity{0.800000}%
\pgfsetlinewidth{1.003750pt}%
\definecolor{currentstroke}{rgb}{0.800000,0.800000,0.800000}%
\pgfsetstrokecolor{currentstroke}%
\pgfsetstrokeopacity{0.800000}%
\pgfsetdash{}{0pt}%
\pgfpathmoveto{\pgfqpoint{5.462858in}{2.593460in}}%
\pgfpathlineto{\pgfqpoint{7.102778in}{2.593460in}}%
\pgfpathquadraticcurveto{\pgfqpoint{7.130556in}{2.593460in}}{\pgfqpoint{7.130556in}{2.621238in}}%
\pgfpathlineto{\pgfqpoint{7.130556in}{3.422778in}}%
\pgfpathquadraticcurveto{\pgfqpoint{7.130556in}{3.450556in}}{\pgfqpoint{7.102778in}{3.450556in}}%
\pgfpathlineto{\pgfqpoint{5.462858in}{3.450556in}}%
\pgfpathquadraticcurveto{\pgfqpoint{5.435080in}{3.450556in}}{\pgfqpoint{5.435080in}{3.422778in}}%
\pgfpathlineto{\pgfqpoint{5.435080in}{2.621238in}}%
\pgfpathquadraticcurveto{\pgfqpoint{5.435080in}{2.593460in}}{\pgfqpoint{5.462858in}{2.593460in}}%
\pgfpathlineto{\pgfqpoint{5.462858in}{2.593460in}}%
\pgfpathclose%
\pgfusepath{stroke,fill}%
\end{pgfscope}%
\begin{pgfscope}%
\pgfsetrectcap%
\pgfsetroundjoin%
\pgfsetlinewidth{1.204500pt}%
\definecolor{currentstroke}{rgb}{0.121569,0.466667,0.705882}%
\pgfsetstrokecolor{currentstroke}%
\pgfsetdash{}{0pt}%
\pgfpathmoveto{\pgfqpoint{5.490636in}{3.338088in}}%
\pgfpathlineto{\pgfqpoint{5.629525in}{3.338088in}}%
\pgfpathlineto{\pgfqpoint{5.768414in}{3.338088in}}%
\pgfusepath{stroke}%
\end{pgfscope}%
\begin{pgfscope}%
\definecolor{textcolor}{rgb}{0.000000,0.000000,0.000000}%
\pgfsetstrokecolor{textcolor}%
\pgfsetfillcolor{textcolor}%
\pgftext[x=5.879525in,y=3.289477in,left,base]{\color{textcolor}{\rmfamily\fontsize{10.000000}{12.000000}\selectfont\catcode`\^=\active\def^{\ifmmode\sp\else\^{}\fi}\catcode`\%=\active\def%{\%}Ideal signal}}%
\end{pgfscope}%
\begin{pgfscope}%
\pgfsetrectcap%
\pgfsetroundjoin%
\pgfsetlinewidth{1.204500pt}%
\definecolor{currentstroke}{rgb}{1.000000,0.498039,0.054902}%
\pgfsetstrokecolor{currentstroke}%
\pgfsetdash{}{0pt}%
\pgfpathmoveto{\pgfqpoint{5.490636in}{3.134231in}}%
\pgfpathlineto{\pgfqpoint{5.629525in}{3.134231in}}%
\pgfpathlineto{\pgfqpoint{5.768414in}{3.134231in}}%
\pgfusepath{stroke}%
\end{pgfscope}%
\begin{pgfscope}%
\definecolor{textcolor}{rgb}{0.000000,0.000000,0.000000}%
\pgfsetstrokecolor{textcolor}%
\pgfsetfillcolor{textcolor}%
\pgftext[x=5.879525in,y=3.085620in,left,base]{\color{textcolor}{\rmfamily\fontsize{10.000000}{12.000000}\selectfont\catcode`\^=\active\def^{\ifmmode\sp\else\^{}\fi}\catcode`\%=\active\def%{\%}Ideal noise}}%
\end{pgfscope}%
\begin{pgfscope}%
\pgfsetrectcap%
\pgfsetroundjoin%
\pgfsetlinewidth{1.204500pt}%
\definecolor{currentstroke}{rgb}{0.172549,0.627451,0.172549}%
\pgfsetstrokecolor{currentstroke}%
\pgfsetdash{}{0pt}%
\pgfpathmoveto{\pgfqpoint{5.490636in}{2.930374in}}%
\pgfpathlineto{\pgfqpoint{5.629525in}{2.930374in}}%
\pgfpathlineto{\pgfqpoint{5.768414in}{2.930374in}}%
\pgfusepath{stroke}%
\end{pgfscope}%
\begin{pgfscope}%
\definecolor{textcolor}{rgb}{0.000000,0.000000,0.000000}%
\pgfsetstrokecolor{textcolor}%
\pgfsetfillcolor{textcolor}%
\pgftext[x=5.879525in,y=2.881762in,left,base]{\color{textcolor}{\rmfamily\fontsize{10.000000}{12.000000}\selectfont\catcode`\^=\active\def^{\ifmmode\sp\else\^{}\fi}\catcode`\%=\active\def%{\%}High quality data}}%
\end{pgfscope}%
\begin{pgfscope}%
\pgfsetrectcap%
\pgfsetroundjoin%
\pgfsetlinewidth{1.204500pt}%
\definecolor{currentstroke}{rgb}{0.839216,0.152941,0.156863}%
\pgfsetstrokecolor{currentstroke}%
\pgfsetdash{}{0pt}%
\pgfpathmoveto{\pgfqpoint{5.490636in}{2.726516in}}%
\pgfpathlineto{\pgfqpoint{5.629525in}{2.726516in}}%
\pgfpathlineto{\pgfqpoint{5.768414in}{2.726516in}}%
\pgfusepath{stroke}%
\end{pgfscope}%
\begin{pgfscope}%
\definecolor{textcolor}{rgb}{0.000000,0.000000,0.000000}%
\pgfsetstrokecolor{textcolor}%
\pgfsetfillcolor{textcolor}%
\pgftext[x=5.879525in,y=2.677905in,left,base]{\color{textcolor}{\rmfamily\fontsize{10.000000}{12.000000}\selectfont\catcode`\^=\active\def^{\ifmmode\sp\else\^{}\fi}\catcode`\%=\active\def%{\%}Low quality data}}%
\end{pgfscope}%
\end{pgfpicture}%
\makeatother%
\endgroup%
}
	\captionsetup{width=0.85\linewidth}
	\caption[Barcode rank simulation]{The barcode rank plot. In this simulation, 200 million reads are distributed across \num{10000} cells, and the whitelist is \num{800000} barcodes in length. The ceiling for mean number of reads per barcode is thus 200M reads/10k cell = 20k reads/cell (perfect signal), with a floor of 200M reads/800k barcodes = 250 reads/barcode (perfect noise).}
	\label{fig:kneeplot}
\end{figure}

In the "ideal signal" case, all reads are evenly distributed across a small number of barcodes that matches the number of input cells (e.g. \num{10000} in the figure). This means that there are no quality differences between real cells, and there are no barcodes with low read counts. The opposite scenario ("ideal noise") has an even distribution of all reads across \textit{all} barcodes, and we cannot distinguish true cell barcodes from noise barcodes. Real datasets fall between these two extremes. A high quality dataset exhibits a clear cut-off, where the number of reads per barcode sharply drops off in the transition between real cells and background noise. A dataset of lower quality will exhibit a more gradual transition between real cells and background noise. Based on this simple plot, we can estimate at a glance the rough quality characteristics of any single-cell dataset:

\begin{itemize}[noitemsep]
	\item The number of cells captured in the dataset is given by the barcode rank where true cells are separated from background noise ("knee"). If the dataset is sequenced to near-saturation, the total number of barcodes can even exceed the whitelist, as sequencing/PCR/synthesis errors introduce non-whitelisted species.
	\item The approximate number of reads associated with each true cell can be estimated from the y-axis.
	\item Overall data quality can be estimated by comparing the area under real cells and area under noise, or simply the "sharpness" of the cutoff.
	\item Total noise can be inferred by looking at the number of reads present after the true cell cutoff, taking into account that an equal fraction of noise is also present in the true cells, for example stemming from cross-contamination or bead impurities\sjcitep[.]{lareau2020inference}
\end{itemize}

\section{Single-cell data post-processing}
% [It would actually be better to structure this paragraph as such: "compared to rna-seq data, scrna-seq data has two problems: one it is more non-linear because of the added cell states and two drop-outs additionally make it more difficult to do statistical test. this is solved by...:"]

Conventional bulk RNA-seq analysis tools modeled gene counts using Poissonian (\verb|edgeR|\sjcitep[]{robinsonEdgeRBioconductorPackage2010}) or negative binomial (\verb|DESeq|\sjcitep[]{andersDifferentialExpressionAnalysis2010}) distributions. Statistical tests can then be used to identify statistically significant differentially expressed genes. Directly visualising correlations or variance between many RNA-seq or microarray samples is difficult due to the high-dimensionality of the data. Traditionally, this was done by first reducing the dimensionality of the data using PCA, and then visualising the first 2 principal components in a 2D space. Alternatively, heatmaps with dendrograms were used to order correlated samples.

Compared to bulk RNA-seq and microarray data, single-cell sequencing data is highly zero-inflated: when sampling transcripts in a single cell, the probability of weakly-expressed genes not being detected is high, leading to large amounts of "dropouts" (false negative gene expression). This sparsity, couplied with the more fine-grained nature of cell states in single-cell data further introduces non-linearity, since gradients of expression in the population are no longer averaged out. This problem was partially solved by the introduction of t-SNE, a technique first designed to visualise the MNIST dataset of handwritten digits\sjcitep[.]{maatenVisualizingHighDimensionalData2008} The earliest implementations of t-SNE were used for single-cell mass cytometry data\sjcitep[]{amirViSNEEnablesVisualization2013} (45 protein levels per cell), then single-cell qPCR\sjcitep[]{saadatpourCharacterizingHeterogeneityLeukemic2014} (175 genes) and later single-cell whole transcriptomics\sjcitep[]{macoskohighlyparallelgenomewide2015} (after PCA on \textasciitilde{\num{20000}}genes). t-SNE preserves local structure (close neighborhood of cells) at the cost of some concessions to global structure. Clusters of single cells can appear closer in the t-SNE embedding than they are in high-dimensional gene space. UMAP is a slightly newer dimensionality reduction algorithm which aims to better preserve global structure\sjcitep[.]{bechtDimensionalityReductionVisualizing2018} Some common criticisms of both t-SNE and UMAP is that they still distort global topologies\sjcitep[.]{chariSpeciousArtSinglecell2023} In addition, both methods are sensitive to hyperparameters (perplexity, number of neighbours) and stochasticity (e.g. initialisation). These tools are thus only used for exploration and quick visualisation of the data.

% \verb|SCDE| corrects for dropouts by taking into account aggregate expression levels across all cells to calculate dropout probabilities\sjcitep[.]{kharchenkoBayesianApproachSinglecell2014} These dropout probabilities are then used as priors to model gene expression levels in individual cells. Dropouts are weighted based on their prior probability, and likely relative (fold-change) gene expression levels are inferred. This approach ensures that highly-expressed genes are not penalized, and low-expressed genes are not inflated.

% Monocle\sjcitep[]{trapnellDynamicsRegulatorsFate2014} Advent of single-cell RNA-sequencing --> we can now look at gene expression dynamics (not possible in bulk rna-seq due to simpson's paradox). Differentiate human myoblasts to [skeletal muscle], sample 270 cells over 4 timepoints between 0 and 72 hours using the Fluidigm C1, 4M reads/cell, compare also with bulk rna-seq of each timepoint. Then Monocle: TopHat and CuffLinks are used to generate a count matrix, independent component analysis is used for dimreduc (why?). finally, a minimum-spanning tree (graph that connects cells with minimal total edge length, with edges representing transcriptional similarity), then monocle identifies the longest path along the MST. a starting cell (highly proliferative progenitor) is chosen and cells are ordered along the path. branch points may be detected. expression profiles are also smoothed along the trajectory using generalized additive models. now genes can be correlated to pseudotime.

% Simpson's paradox prohibits accurate investigation of gene expression dynamics in bulk samples. Single-cell datasets, however, can be both larger and much more fine-grained than bulk RNA-seq datasets sampled across different conditions or timepoints. \verb|Monocle| leverages this property inferring differentiation state from single-cell sequencing data sampled from differentiating cells\sjcitep[.]{trapnellDynamicsRegulatorsFate2014} High-dimensional gene expression space is reduced to statistically significant dimensions, and a minimum spanning tree is calculated between cells in this space. The longest path within this tree can then be used to calculate a "pseudotime" property. Beyond being able to order cells according to their differentiation state, genes can also be correlated with pseudotime to identify significant differentiation drivers.

By 2019, several single-cell sequencing technologies had emerged. Each had their own strengths and weaknesses, sometimes leading to technical artifacts in datasets even when assessing the same biological sample. Even the same technology used on different days could lead to batch effects due to sample degradation or other technical artifacts between experiments. This translates to cells clustering by dataset batch rather than biological state. \verb|Seurat| aims to facilitate integration of multiple single-cell datasets across different technologies and conditions\sjcitep[.]{stuartcomprehensiveintegrationsinglecell2019} The package uses Canonical Correlation Analysis (CCA)\sjcitep[]{hotellingRelationsTwoSets1936} to calculate linear combinations of genes (canonical correlation vectors, in this case "metagenes") that capture common information between two datasets. Cells are then co-projected in this shared metagene space to create a new "joint embedding" where clustering and differential expression analysis can be performed. The canonical vectors can be used to re-transform differential metagenes back to the original gene set used. \verb|Seurat|'s CCA implementation could also be used for label transfer: if one dataset is labeled or annotated, proximity in the joint embedding space can be used to transfer labels from one dataset to another. Similarly to Seurat, \verb|Scanpy| unifies preprocessing, highly variable gene selection, PCA, graph-based clustering, t-SNE/UMAP embedding, differential expression in a python package\sjcitep[.]{wolf2018scanpy} The cells by gene count matrix is stored along unstructured metadata in the \verb|AnnData| format.

Whereas scRNA-seq is limited to counts in \textasciitilde{\num{22000}} genes (human, mouse), the accessible genome can span hundreds of thousands of open chromatin regions which can each be sampled a maximum of two times per cell. These effects lead to highly-dimensional, but sparse count matrices. Latent semantic indexing (LSI), an approach borrowed from text analysis, can be used to perform cell clustering and dimensionality reduction\sjcitep[.]{cusanovichsinglecellatlasvivo2018} Briefly, single-cell accessibility is first counted in either annotated regions or 5 kb genome bins, binarized, and scaled according to inverse document frequency (peaks that are not commonly accessible are weighted higher). Singular value decomposition (SVD), is then used to decompose the resulting weighted cell-accessibility matrix into latent features across cells:

\[
	M = U\Sigma V^T =
	\underset{(m \times n)}{\begin{bmatrix}
			\cdot  & \cdot  & \cdot  & \cdots & \cdot  \\
			\cdot  & \cdot  & \cdot  & \cdots & \cdot  \\
			\vdots & \vdots & \vdots & \ddots & \vdots \\
			\cdot  & \cdot  & \cdot  & \cdots & \cdot  \\
			\cdot  & \cdot  & \cdot  & \cdots & \cdot
		\end{bmatrix}}
	=
	\underset{(m \times r)}{\begin{bmatrix}
			\cdot  & \cdot  & \cdots & \cdot  \\
			\cdot  & \cdot  & \cdots & \cdot  \\
			\vdots & \vdots & \ddots & \vdots \\
			\cdot  & \cdot  & \cdots & \cdot  \\
			\cdot  & \cdot  & \cdots & \cdot
		\end{bmatrix}}
	\underset{(r \times r)}{\begin{bmatrix}
			\sigma_1 & 0        & \cdots & 0        \\
			0        & \sigma_2 & \cdots & 0        \\
			\vdots   & \vdots   & \ddots & \vdots   \\
			0        & 0        & \cdots & \sigma_r
		\end{bmatrix}}
	\underset{(r \times n)}{\begin{bmatrix}
			\cdot  & \cdot  & \cdot  & \cdots & \cdot  \\
			\cdot  & \cdot  & \cdot  & \cdots & \cdot  \\
			\vdots & \vdots & \vdots & \ddots & \vdots \\
			\cdot  & \cdot  & \cdot  & \cdots & \cdot
		\end{bmatrix}}
\]

Here, $m$ is the number of peaks, $n$ is the number of cells, and $r = \min(m,n)$ is the rank of the decomposition. U represents the peak-composition of the latent features, $\Sigma$ is diagonal and scales each latent feature, and V\textsuperscript{T} represents the contribution of each latent feature to each cell. Based on $\Sigma$, we can select the top \verb|n| most important features and cluster cells in this reduced space using t-SNE. Differentially accessible regions (DARs) between clusters can then be calculated on the raw binary accessibility matrix using logistic regression, which is suited to model the odds of a binary event like chromatin accessibility.

% \verb|chromVAR| works in a fundamentally different way\sjcitep[.]{schepchromVAR2017} Instead of using hundreds of thousands of accessible peak regions, \verb|chromVAR| first groups these peaks into cistromes based on common TF motifs or pre-set functional annotations such as known enhancer peak sets. Accessibility within each cistrome is normalized and converted to z-scores, which are then used for t-SNE and clustering. Instead of known TF motifs, k-mers can also be used to identify novel motifs important in a given biological process.

In this work, \verb|cisTopic| was used extensively to generate and post-process all scATAC-seq count matrices\sjcitep[.]{bravogonzalezblascistopic2019} \verb|cisTopic| also uses a probabilistic approach to model chromatin accessibility. LDA (rather than LSI) is used to calculate two sets of distributions: the region-topic distribution (regulatory regions grouped into "topics") and the topic-cell distribution (probability of a topic being associated with a cell) which should approximate the original data when multiplied\sjcitep[.]{bleiLatentDirichletAllocation2003}

\[
	\underset{m \times n}{
		\begin{bmatrix}
			x_{11} & x_{12} & \cdots & x_{1n} \\
			x_{21} & x_{22} & \cdots & x_{2n} \\
			\vdots & \vdots & \ddots & \vdots \\
			x_{m1} & x_{m2} & \cdots & x_{mn}
		\end{bmatrix}}
	\approx
	N \cdot
	\underset{m \times k}{
	\begin{bmatrix}
		p_{r_1t_1} & p_{r_1t_2} & \cdots & p_{r_1t_k} \\
		p_{r_2t_1} & p_{r_2t_2} & \cdots & p_{r_2t_k} \\
		\vdots     & \vdots     & \ddots & \vdots     \\
		p_{r_mt_1} & p_{r_mt_2} & \cdots & p_{r_mt_k}
	\end{bmatrix}_{}}
	\cdot
	\underset{k \times n}{
		\begin{bmatrix}
			p_{t_1c_1} & p_{t_1c_2} & \cdots & p_{t_1c_n} \\
			p_{t_2c_1} & p_{t_2c_2} & \cdots & p_{t_2c_n} \\
			\vdots     & \vdots     & \ddots & \vdots     \\
			p_{t_kc_1} & p_{t_kc_2} & \cdots & p_{t_kc_n}
		\end{bmatrix}}
\]

LDA parameters (probabilities) are estimated using a Markov chain monte carlo approach: regions are initially randomly assigned to topics, and topic-region assignment is then updated iteratively in each cell based on two factors: (1) the topic's occurrence in the cell and (2) the region's occurrence in each topic across all other cells. Sampling is only performed in this latent topic-assignment vector for each region in each cell (collapsed Gibbs sampling), which reduces total computational complexity compared to sampling in both region-topic and topic-cell distribution. This approach is more efficient than other Bayesian approaches (e.g. variational inference) due to the extreme sparsity of the data. tSNE or UMAP can then be used to visualise the resulting normalized topic-cell distributions. LDA's topics are also interpretable\sjcitep[,]{steyversProbabilisticTopicModels2007} a property that is exploited to extract transcription factor binding motifs from the cis-regulatory topics found. Motif enrichment analyses can be performed using \verb|cisTarget|\sjcitep[.]{imrichova2015icisTarget} Regions and combinations of regions can also be annotated using nearest genes (\verb|ChIPseeker|)\sjcitep[.]{yu2015chipseeker}

\verb|cisTopic| also introduces a region score which rewards regions that are more specific to a topic, and penalises regions that are ubiquitous across topics:

\[
	\text{Region score} = \beta_{r,t} \left( \log \beta_{r,t} - \frac{1}{T} \sum_{t=1}^{T} \log \beta_{r,t} \right)
\]

Consider Shannon entropy:

\[H(X) = -\sum_{i} p(x_i) \log p(x_i)\]

Then the term: \[\beta_{r,t} \log \beta_{r,t}\] can be considered the local information content of the region in topic context, while:

\[\frac{\beta_{r,t}}{T} \sum_{t=1}^{T} \log \beta_{r,t}\]

...is the average log-probability of the region acrosss all topics. The region score is thus an average-normalized measure of a given region's information content, indicating how specific it is to each topic. Topics can then be scored for TF motif enrichment by counting TF motif occurrence in the score-ranked regions and calculating the area under this rank curve (\verb|AUCell|)\sjcitep[.]{aibarscenicsinglecell2017}

\verb|cisTopic|'s LDA approach was benchmarked against other popular tools such as \verb|chromVAR|, \verb|scABC|, \verb|BROCKMAN|, \verb|SCRAT|, \verb|Cicero|, and LSI, using both simulated single-cell ATAC-seq datasets and real data from FAC-sorted hematopoietic cells. \verb|cisTopic| typically yielded higher scores for cell-type clustering, especially in low-coverage scenarios, confirming that the method can more robustly capture biologically meaningful patterns in sparse single-cell ATAC-seq data.

% [There was something about how cistopic reduces dimensionality for both cells and regions somewhere, in the part where they compare previous analysis methods. highlight that.]

% Cicero[probably good to look at to]\sjcitep[.]{plinerCiceroPredicts2018}

\section{Deep-learning models}
Recently, the availability of massive public databases and availability of compute resources have led to widespread application of deep learning in biology. Advances in algorithms trace back to 1959, when Hubel and Wiesel showed that specific neurons in a cat's visual cortex respond to common spatial patterns (e.g. edges, corners, orientations) in specific locations within the cat's visual field\sjcitep[.]{hubel1959receptive} These patterns can overlap when layers of neurons pass information onwards, forming broader structures which ultimately allow the cat to recognise complex objects. This idea was borrowed to implement convolutional neural networks (CNNs) for handwritten digit recognition\sjcitep[.]{fukushima1980neocognitron,lecun1998gradient} Rasterized (vectorized) images are scanned by a "filter", converting the original data into a new representation. This new representation is then passed on to the next layer, where it is again scanned by a new filter, and so on until the final layer. CNNs are thus well suited for detecting spatial patterns, and integrating these patterns into larger, more global structures. as a result, CNNs have been successfully used in enhancer prediction models. An early example was DeepSEA\sjcitep[,]{zhou2015predicting} a CNN trained on ENCODE and Roadmap Epigenomics data (690 TF ChIP-seq + 125 DHS + 104 H3K27Ac/H3K4Me3 datasets). This model could simultaneously predict these 919 chromatin features at 1 base resolution with higher accuracy than previous tools (\verb|CADD|, \verb|GWAVA|). Saturation mutagenesis analyses showed that DeepSEA could de novo predict known allelic imbalances and eQTLs (using chromatin features as a proxy).

% The same logic is applied in [ref], where a CNN is trained to detect local patterns in DNA (cell-type specific peaks). In order to train a CNN on DNA sequences (peaks), DNA sequences were one-hot encoded (each position is encoded by a binary vector of length 4, and the sequence is represented as a matrix of these vectors of size $N \times L \times 4$). where $N$ is the number of sequences, $L$ is the sequence length (500 bp), and the nucleotides are one-hot encoded as binary vectors.

% \[
% 	\mathbf{X} \in \{0, 1\}^{N \times L \times 4}, \text{ e.g. }
% 	\begin{bmatrix}
% 		\begin{bmatrix}
% 			1 & 0 & 0 & 0 \\
% 			0 & 1 & 0 & 0 \\
% 			0 & 0 & 1 & 0 \\
% 			0 & 0 & 0 & 1 \\
% 			1 & 0 & 0 & 0
% 		\end{bmatrix}
% 		,
% 		\begin{bmatrix}
% 			0 & 1 & 0 & 0 \\
% 			1 & 0 & 0 & 0 \\
% 			0 & 0 & 1 & 0 \\
% 			0 & 0 & 0 & 1 \\
% 			0 & 1 & 0 & 0
% 		\end{bmatrix}
% 	\end{bmatrix}
% 	\text{ encodes for [ACGTA, CAGTC].}
% \]

% The output labels are represented by a given by:
% \[
% 	\mathbf{Y} \in [0, 1]^{N \times C}
% \]
% where $C$ denotes the number of cell types (e.g., $C = 81$). Each $Y_{ij}$ represents the probability that sequence $i$ is an enhancer for cell type $j$. For example:
% \[
% 	\begin{bmatrix}
% 		0.9 & 0.1 & 0.0 & \cdots & 0.2 \\
% 		0.1 & 0.8 & 0.1 & \cdots & 0.0 \\
% 		0.0 & 0.1 & 0.9 & \cdots & 0.1
% 	\end{bmatrix}
% \]

% indicates that sequence 1 is likely an enhancer for cell type 1 (0.9), sequence 2 for cell type 2 (0.8), and sequence 3 for cell type 3 (0.9).

% The convolution operation for filter $k$ is:
% \[
% 	z_{i,k} = f \left( \sum_{j=1}^{h} \sum_{m=1}^{4} X_{i+j-1, m} W_{j, m}^{(k)} + b_k \right)
% \]
% where $f(x) = \max(0, x)$ is the ReLU activation function.

% After pooling, the flattened feature map $\mathbf{P}$ is mapped to outputs as:
% \[
% 	\mathbf{\hat{Y}} = \sigma \left( \mathbf{P} \mathbf{V} + \mathbf{b} \right)
% \]
% where $\sigma$ is the sigmoid function.

% The binary cross-entropy loss is computed as:
% \[
% 	\mathcal{L}_{ij} = - \left[ Y_{ij} \log(\hat{Y}_{ij}) + (1 - Y_{ij}) \log(1 - \hat{Y}_{ij}) \right]
% \]
% with the total loss:
% \[
% 	\mathcal{L} = \frac{1}{N \cdot C} \sum_{i=1}^{N} \sum_{j=1}^{C} \mathcal{L}_{ij}
% \]

% DeepSEA: used chip-seq data from ENCODE and Roadmap Epigenomics, is a CNN , saturation mutagenesis (effect of every possible single-nucleotide mutation in the sequence), 1 kb window around a location of interest is one-hot encoded, various types of chromatin features (chip-seq tf binding from encode of 690 tfs, dnase i hypersensitivity 125 datasets, 104x histone marks (h3k27ac, h3k4me3),) - and this is binarized (if 200 bp window is in a peak of the modality or not), then a CNN is trained with an NLL loss function to predict the modalities, output is a vector of size 919 (sum of all datasets, which are different by tf or by cell type), matches allelic imbalance measurements! predicts whether a given segment of DNA is likely to be bound by regulatory proteins, effects of mutations on that binding. able to de novo predict known eQTLs from negative control SNPs (from 1000 genomes project). better than simpler models such as CADD (logistic regression), GWAVA (random forest classifier).

Several track-predicting foundation models have also been developed. In 2018 Basenji2 was published, a CNN trained on \num{4229} human epigenomic (DNase-seq, histone-ChIP) and transcriptional (CAGE) tracks across more than 500 cell and tissue types. With a large receptive field of 131 kb\sjcitep[,]{kelley2018sequential} the model could predict 128-bp binned coverage across the genome: users feed any 131-kb reference or mutant sequence to the pretrained model (or fine-tune on new tracks) and obtain coverage predictions for each of the \num{4229} tissue-specific tracks it was trained on. In 2021, Enformer replaced dilated convolutions with global self-attention and pushed the receptive field to 100 kb with a 200 kb input window\sjcitep[.]{avsec2021effective} Trained now on \num{5313} human and \num{1643} mouse tracks (TF and histone-mark ChIP, DNaseI, ATAC and CAGE), it predicts these tracks at 128 bp resolution, including long-range interactions. Borzoi is another CNN-attention hybrid, trained on thousands of CAGE, DNase, ATAC, and ChIP-seq datasets similarly to Basenji2 and Enformer, but also adding 866 human and 279 mouse RNA-seq tracks\sjcitep[.]{linder2025predicting} The model predicts these tracks now at 32 bp-resolution. Kundaje lab introduced ChromBPNet, a CNN trained on ENCODE ATAC-seq and DNase-seq tracks for the five tier 1 ENCODE cell lines (GM12878, HepG2, K562, IMR90, H1-hESC)\sjcitep[.]{pampari2025} A smaller sub-net learned each enzyme's sequence biases based on non-peak regions signal, allowing ChromBPNet to isolate each motif's true footprint independent from assay used. This allows the model to predict, at single base resolution, the accessibility read profile (shape + total counts) that either ATAC-seq or DNase-seq would produce.

Recent work has also shown that single-cell datasets can be mined directly to learn enhancer grammar and design new regulatory elements. scRNA-seq of nine patient-derived cultures uncovered an intermediate state that co-expresses melanocytic factors (SOX10, MITF) and mesenchymal regulators (SOX6, EGR3).\sjcitep[]{wouters2020robust} Follow-up bulk ATAC-seq at four time points after SOX10 knock-down showed gradual chromatin shifts rather than abrupt switches. Expanding to 26 melanoma samples from six species, Minnoye et al.\sjcitep[]{minnoye2020cross} found SOX10/MITF/TFAP2A-rich enhancers in melanocytic cells and AP-1/TEAD motifs in mesenchymal cells, with motif usage conserved in dog, horse, pig and zebrafish. The DeepMEL model, CNN-RNN trained on \num{339099} human ATAC peaks, predicted state-specific enhancer activity from sequence alone and guided the synthesis of new MES and MEL enhancers. A similar approach has been used to decode \textit{Drosophila} brain development. Janssens and colleagues profiled \num{240919} scATAC-seq nuclei from the developing \textit{Drosophila} brain across nine time pointset\sjcitep[.]{janssens2022decoding} Peaks were clustered with cisTopic to obtain topic sets of co-accessible, cell-type-specific regions. DeepFlyBrain was trained to predict each topic's aggregated accessibility from 2 kb of sequence. SHAP DeepExplainer scores highlighted the nucleotides most responsible for each prediction, revealing, for example, that \textit{eyeless} and \textit{mef2} motifs act synergistically when placed in close proximity, whereas repressor motifs can override activators. In-silico saturation mutagenesis confirmed the minimal length of active enhancers (as short as 49 bp), and reporter studies showed that 41 of 63 synthesized candidates drove the expected cell-type expression in transgenic flies. Building on these workflows, the \verb|CREsted| pipeline trains compact CNNs (e.g. DeepBICCN2, 6M parameters) directly on scATAC-seq peaks, interprets motif logic and performs de novo enhancer design. Starting from topic-assigned peaks or pseudobulk heights, models are first pretrained on the full peak set and then fine-tuned for cell-type specificity. In mouse motor cortex, \verb|CREsted| reached Pearson r \textasciitilde 0.82 on held-out peaks and retrieved the motifs underlying 171 in-vivo-validated enhancers. A human PBMC version recovered every known TF footprint at the CD79A and TCR\textalpha enhancers and captured 34\% more informative bases than Borzoi. DeepZebrafish, trained on \num{793273} peaks across \num{639} zebrafish cell-type/time-point combinations, served as an oracle for greedy in-silico evolution: five of seven 2-kb designs drove reporter expression specifically in cardiac muscle, somatic muscle or endothelial cells at 48 hpf. Domain or sample-focused models like DeepMEL, DeepFlyBrain, DeepZebraFish complement large foundation models such as Basenji2, Enformer and scFoundation: they require far less compute, yet excel at highly specific questions like cell-type enhancer logic.

Introduced in 2017, the transformer model architecture abandoned recurrence and convolutions in favour of multi-head self-attention blocks\sjcitep[.]{vaswani2023attentionneed} Every token attends to every other, so gradients propagate more directly and training parallelises across all sequence positions. This paper triggered the field-wide shift to attention-only networks that now dominate NLP, vision and sequence biology. For example, meta's Evolutionary Scale Modeling (ESM) applies transformer architecture (specifically, BERT) to protein language to predict masked tokens in amino acid sequences. The log-likelihood of proposed sequences serves as an estimate for how disruptive given mutations will be, allowing researchers to triage synthetic designs. ESM-3 also introduced multimodal masked token modeling: joint tokenized amino-acid sequence, atomic structure, and region annotations can be passed partially masked, and the model will jointly reconstruct missing tokens. This led to the generation of esmGFP, a completely synthetically designed fluorescent protein.

In a similar vein, the original DNABERT was an 86M parameter transformer trained on the 3.2 Gbp human reference genome to classify promoters, TF-binding sites, and other regulatory features\sjcitep[.]{ji2021dnabert} The model was competitive with ClinVar for SNV scoring, but had a relatively small context window of 512 bp, so could not compete in long range tasks such as enhancer-promoter pairing. Nucleotide Transformer (NT) extended the BERT approach to multi-species data: adding both genomic and evolutionary variety by using \num{3300} individual human genomes and 850 other organisms, totaling 262B tokens. NT consistently outperformed BERT and DeepSEA on promoter, enhancer, chromatin, and histone mark prediction, and enabled zero-shot scoring of SNPs and indels - achieving better accuracy than previous models on ClinVar benchmarks\sjcitep[.]{dallatorre2025nucleotide} Recently, Evolutionary Scale Modeling for Genomes (Evo 2) pushed this scaling further with a 40B-parameter model trained on 9.3T tokens spanning all domains of life\sjcitep[.]{brixi2025genome} The model scores variant effects across taxa, and supports up to 1 million tokens of context, enabling the generation of entire synthetic mitochondrial or bacterial genomes. Loss decreased smoothly with more data and parameters, with no saturation observed so far.

AlphaFold 1 was a CNN trained on protein structures (PDB) and multiple sequence alignments (MSAs) to predict residue-residue distances\sjcitep[.]{senior2020improved} AlphaFold 2 switched architecture, passing MSAs through Evoformer, a model with transformer-based self-attention, to predict spatial contacts and ultimately complete 3D protein structure\sjcitep[.]{jumper2021highly} These changes led it to achieve a median backbone accuracy of \SI{0.96}{\angstrom} in CASP14 - surpassing the next best method by \textasciitilde\SI{2}{\angstrom}. AlphaFold 3 included a generative diffusion model to handle atomic coordinates directly rather than relying on MSAs. As a result, the model is capable of predicting binding interactions across protein, DNA/RNA, and small molecules, and performs better than AlphaFold 2 on other tasks.

Geneformer applies the transformer playbook to a collection of 30M human scRNA-seq cells across various tissues, with cells tokenized as a ranked list of their top 2048 genes sorted by expression (no count information)\sjcitep[.]{theodoris2023transferlearning} A transformer architecture is then trained to fill in masked gene positions, implicitly learning patterns of co-expression and regulation across tissues and conditions present in the data. The model is designed to fine-tuned for specific tasks such as dosage-sensitivity classification, regulatory range prediction, and chromatin-state inference. scFoundation, another masked-prediction transformer, was trained on 100M single cells\sjcitep[.]{hao2024large} Varying size from 3M to 100M parameters showed that its loss curves followed the same token vs FLOPs power laws seen in language models.
